% !TEX program = xelatex
\documentclass[12pt,a4paper]{ctexbook}

% ==================== 页面 ====================
\usepackage[top=2.5cm,bottom=2.5cm,left=2.5cm,right=2.5cm]{geometry}

% ==================== 字体 ====================
\usepackage{fontspec}
\usepackage{iffont}
\IfFontExistsTF{Consolas}{\setmonofont{Consolas}}{\setmonofont{DejaVu Sans Mono}}

% ==================== 颜色 ====================
\usepackage[dvipsnames,svgnames,x11names]{xcolor}
\definecolor{codebg}{HTML}{F5F5F5}
\definecolor{codeframe}{HTML}{E0E0E0}
\definecolor{cppkeyword}{HTML}{1A1AFF}
\definecolor{cppcomment}{HTML}{2E8B2E}
\definecolor{cppstring}{HTML}{B22222}
\definecolor{linenumgray}{HTML}{999999}
\definecolor{toolheadbg}{HTML}{2D3E50}
\definecolor{yamlkey}{HTML}{7A3E9D}
\definecolor{diffadd}{HTML}{1A7F37}
\definecolor{diffdel}{HTML}{B62324}

% ==================== listings ====================
\usepackage{listings}

% ---- 所有代码块共用的排版（先定义，各语言 style 组合它）----
\lstdefinestyle{codebase}{
  basicstyle=\small\ttfamily,
  keywordstyle=\color{cppkeyword}\bfseries,
  commentstyle=\color{cppcomment}\itshape,
  stringstyle=\color{cppstring},
  numberstyle=\tiny\color{linenumgray},
  numbers=left,
  frame=single,
  rulecolor=\color{codeframe},
  framesep=6pt,
  breaklines=true,
  tabsize=4,
  columns=flexible,
  keepspaces=true,
  backgroundcolor=\color{codebg},
  showstringspaces=false,
  escapeinside={(*@}{@*)},
}

\lstdefinestyle{cppstyle}{
  style=codebase,
  language=C++,
  morekeywords={constexpr,consteval,constinit,decltype,auto,
    concept,requires,co_await,co_yield,co_return,
    noexcept,override,final,nullptr,static_assert,
    template,typename,namespace,using,mutable,volatile,
    explicit,operator,friend,inline,virtual,
    thread_local,alignas,alignof},
  deletekeywords={int,char,bool,float,double,long,short,unsigned,signed,void,wchar_t},
  emph={size_t,int32_t,uint32_t,int64_t,uint64_t,ssize_t,
    string,vector,map,set,unique_ptr,shared_ptr,optional,variant,any,
    span,string_view,FILE,wchar_t},
  emphstyle=\color{teal!80!black},
}
\lstset{style=cppstyle}

% ---- YAML：.clang-format / .clang-tidy / .clangd / pre-commit / Actions ----
\lstdefinestyle{yamlstyle}{
  style=codebase,
  morecomment=[l]{\#},
  morestring=[b]",
  morestring=[d]',
  alsoletter={-,:},
  morekeywords=[1]{Version,Checks,CheckOptions,HeaderFilterRegex,FormatStyle,
    User,WarningsAsErrors,CompilationDatabase,ExtraArgs,Add,Remove,
    Diagnostics,ClangTidy,UnusedIncludes,MissingIncludes,index,background-index,
    clang-tidy,InlayHints,Style,BasedOnStyle,IndentWidth,TabWidth,UseTab,
    ColumnLimit,BreakBeforeBraces,AlignAfterOpenBracket,AllowShortIfStatementsOnASingleLine,
    IncludeBlocks,SortIncludes,SeparateDefinitionBlocks,SpacesInAngles,
    Language,If,Else,Disable,Key,Name,files,exclude,repo,rev,hooks,id,entry,
    language_types,on,push_pull_request,steps,uses,runs-on,defaults},
  keywordstyle=[1]\color{yamlkey}\bfseries,
}

% ---- JSON：compile_commands.json / settings.json / CMakePresets.json ----
\lstdefinestyle{jsonstyle}{
  style=codebase,
  morestring=[b]",
  alsoletter={-},
  morekeywords={directory,file,output,arguments,command,version,
    cmakeMinimumRequired,name,generator,binaryDir,sourceDir,caches,
    toolchainFile,C_CXX_FLAGS,defineConstants,environments,
    editorMode,rulerColumn,tabSize,insertSpaces,files,settings},
  keywordstyle=\color{yamlkey}\bfseries,
}

% ---- Shell：命令行与工具调用 ----
\lstdefinelanguage{ShellLD}{
  morekeywords={if,then,else,elif,fi,for,in,do,done,while,case,esac,
    function,return,local,export,source,echo,cd,set,exec,trap,read,
    printf,test,exit,shift,eval,alias},
  sensitive=true,
  morecomment=[l]{\#},
  morestring=[b]",
  morestring=[d]',
}
\lstdefinestyle{shellstyle}{
  style=codebase,
  language=ShellLD,
  keywordstyle=\color{cppkeyword}\bfseries,
  emph={clang-format,clang-tidy,clangd,clang,clang++,cppcheck,iwyu,
    compile-commands,grep,sed,awk,xargs,jq,python3,pip,pip3,pwsh,bash,git,
    cmake,ninja,make,ccache,lldb,gdb,code,nvim,vim,emacs,tar,mkdir,touch,ls,cat,
    pre-commit,actions,gh,yamllint,wc,diff,find,sort,uniq,tee,rm},
  emphstyle=\color{teal!80!black}\bfseries,
}

% ---- CMake ----
\lstdefinelanguage{CMakeLD}{
  morekeywords={cmake_minimum_required,project,add_executable,
    add_library,target_link_libraries,target_include_directories,
    target_compile_definitions,target_compile_options,set,if,elseif,
    else,endif,find_package,option,message,install,file,
    target_sources,add_subdirectory,enable_language,set_target_properties,
    FetchContent_Declare,FetchContent_MakeAvailable,list,string},
  sensitive=false,
  morecomment=[l]{\#},
  morestring=[b]",
}
\lstdefinestyle{cmakestyle}{
  style=codebase,
  language=CMakeLD,
  keywordstyle=\color{cppkeyword}\bfseries,
}

% ---- 统一 diff：clang-format 与 clang-tidy 的改动 ----
\lstdefinestyle{diffstyle}{
  style=codebase,
  numbers=none,
  frame=none,
  backgroundcolor=\color{white},
  morecomment=[l][\color{diffdel}]{-},
  morecomment=[l][\color{diffadd}]{+},
  morecomment=[l][\color{blue!70!black}\bfseries]{@@},
}

% ==================== tcolorbox ====================
\usepackage[most]{tcolorbox}

\newtcolorbox{guideline}[1][]{
  enhanced,breakable,
  colback=blue!5!white,colframe=blue!75!black,
  fonttitle=\bfseries,title=要点,#1,
  boxed title style={colback=blue!75!black},
  attach boxed title to top left={yshift=-2mm,xshift=2mm},
  arc=2mm,boxrule=0.8mm,
}
\newtcolorbox{compare}[1][]{
  enhanced,breakable,
  colback=green!3!white,colframe=green!50!black,
  fonttitle=\bfseries,title=对比,#1,
  boxed title style={colback=green!50!black},
  attach boxed title to top left={yshift=-2mm,xshift=2mm},
  arc=2mm,boxrule=0.8mm,
}
\newtcolorbox{warning}[1][]{
  enhanced,breakable,
  colback=red!5!white,colframe=red!75!black,
  fonttitle=\bfseries,title=常见陷阱,#1,
  boxed title style={colback=red!75!black},
  attach boxed title to top left={yshift=-2mm,xshift=2mm},
  arc=2mm,boxrule=0.8mm,
}
\newtcolorbox{note}[1][]{
  enhanced,breakable,
  colback=orange!5!white,colframe=orange!80!black,
  fonttitle=\bfseries,title=注意,#1,
  boxed title style={colback=orange!80!black},
  attach boxed title to top left={yshift=-2mm,xshift=2mm},
  arc=2mm,boxrule=0.8mm,
}
\newtcolorbox{bestpractice}[1][]{
  enhanced,breakable,
  colback=purple!5!white,colframe=purple!60!black,
  fonttitle=\bfseries,title=最佳实践,#1,
  boxed title style={colback=purple!60!black},
  attach boxed title to top left={yshift=-2mm,xshift=2mm},
  arc=2mm,boxrule=0.8mm,
}

% ==================== 章节标题 ====================
\usepackage{titlesec}
\usepackage{fancyhdr}
\usepackage{graphicx}
\usepackage{amssymb}
\usepackage{amsmath}
\usepackage{tabularx}
\usepackage{booktabs}
\usepackage{longtable}
\usepackage{array}
\usepackage{multirow}
\usepackage{enumitem}
\usepackage{seqsplit}
\usepackage{float}

\titleformat{\chapter}[display]
  {\normalfont\huge\bfseries\color{toolheadbg}}
  {\chaptertitlename\ \thechapter}{20pt}{\Huge}
\titleformat{\section}
  {\normalfont\Large\bfseries\color{blue!70!black}}
  {\thesection}{1em}{}
\titleformat{\subsection}
  {\normalfont\large\bfseries\color{blue!60!black}}
  {\thesubsection}{1em}{}

\pagestyle{fancy}
\fancyhf{}
\fancyhead[LE,RO]{\thepage}
\fancyhead[RE]{\leftmark}
\fancyhead[LO]{\rightmark}
\setlength{\headheight}{14.5pt}

\setlist[itemize]{leftmargin=2em,itemsep=2pt}
\setlist[enumerate]{leftmargin=2em,itemsep=2pt}
\emergencystretch=6em

% ==================== 超链接 ====================
\usepackage{hyperref}
\hypersetup{
  colorlinks=true,
  linkcolor=blue!70!black,
  urlcolor=blue!60!black,
  bookmarks=true,
  pdfauthor={C++ Reference},
  pdftitle={C/C++ 语言服务器与代码质量工具参考手册：LSP、clangd、clang-format、clang-tidy},
}

% ==================== 自定义命令 ====================
\newcommand{\cpp}[1]{C++#1}
\newcommand{\Fn}[1]{\texttt{#1()}}
\newcommand{\Tp}[1]{\texttt{#1}}
\newcommand{\Opt}[1]{\texttt{-\/-#1}}
\newcommand{\Bk}{\discretionary{}{}{}}
\newcommand{\Lib}[1]{\textbf{\texttt{#1}}}
\newcommand{\Lang}[1]{\textsf{#1}}
\newcommand{\CL}{\Lang{C}}
\newcommand{\CPP}{\Lang{C++}}
\newcommand{\PY}{\Lang{Python}}
% 本书反复点名的三个工具，用固定宏指代，避免同一工具多种写法
\newcommand{\Tidy}{\Lib{clang-tidy}}
\newcommand{\Fmt}{\Lib{clang-format}}
\newcommand{\CD}{\Lib{clangd}}

% ====================================================================
\begin{document}

\frontmatter

\begin{titlepage}
\centering
\vspace*{4cm}
{\Huge\bfseries C/C++ 语言服务器与代码质量工具\par}
\vspace{1cm}
{\LARGE LSP、clangd、clang-format、clang-tidy\par}
\vspace{0.4cm}
{\Large 编译数据库、编辑器集成、静态分析与 CI 质量门\par}
\vfill
{\large \today\par}
\end{titlepage}

\tableofcontents

\mainmatter

\cleardoublepage

% ====================================================================
% 章节正文由 frag_chNN.tex 合并插入到下面这行之前（merge_frags.py）
% ---- frag_ch01.tex ----
\part{基础设施：工具链与编译数据库}
\chapter{工具链解剖：clang、libclang 与 clang-tools-extra}
\label{ch:toolchain}

在把 clang-format、clang-tidy、clangd 装进工作流之前，先要搞清楚一件事：
它们是同一棵树上长出来的三个分支，但这三个分支各自站在这棵树的不同高度上。
站得高（只看见 token）的工具快、宽松、不需要项目上下文；
站得低（看见完整语义）的工具准、慢、离不开编译数据库。
本章把 clang 前端拆开，说明每一层暴露什么接口、
哪些工具用了哪一层，最后解决最实际的麻烦：同一台机器上装了五份
clang-format，你的编辑器到底在用哪一份。

\section{clang 前端的三段流水线}
\label{sec:ch01-s1}

\CL\ 与 \CPP\ 编译器驱动 \texttt{clang++} 在真正生成机器码之前，
要经过三个界限分明的阶段。本册讨论的所有工具都只和前三个阶段打交道，
代码生成部分（优化、LLVM IR、汇编）几乎不影响编辑器体验，可以暂时忽略。

\begin{enumerate}
  \item \textbf{预处理}（\texttt{Preprocessor}）。展开 \texttt{\#include}、
        展开宏、求值 \texttt{\#if}、把注释剥掉，产出一条\textbf{token 流}。
        这一阶段的结果是扁平的、没有嵌套关系的字符序列。
  \item \textbf{解析}（\texttt{Parser}）。只按文法把 token 流组装成
        \textbf{AST}（抽象语法树）的骨架：声明、语句、表达式的嵌套结构。
        此时不判断类型是否匹配，也不做名字查找。
  \item \textbf{语义分析}（\texttt{Sema}）。名字查找、重载决议、隐式转换、
        模板实例化、\texttt{constexpr} 求值，绝大多数诊断消息在这里产生。
\end{enumerate}

用一个具体文件来看差异。清单 \ref{lst:ch01-demo} 里那段代码故意做了三件
不利于工具处理的事：把大括号藏进宏、\texttt{std::vector} 的头文件可能
找不到、留了一个未使用的变量。

\begin{lstlisting}[style=cppstyle,caption={带宏与不完整包含的示例文件 \texttt{demo.cpp}},label={lst:ch01-demo}]
// demo.cpp -- 三处"对工具不友好"的写法
#define BEGIN_SCOPE {
#define END_SCOPE   }

#include <vector>        // 假设搜索路径没配好，这一行会失败

template <typename T>
class Box {
public:
    explicit Box(T v) : v_(v) {}
    T get() const { return v_; }
private:
    T v_;
};

int main()
BEGIN_SCOPE
    Box<int> b(42);
    std::vector<int> xs = {1, 2, 3};
    int n = xs.size();   // 未使用：只有 Sema 阶段能发现
    return b.get();
END_SCOPE
\end{lstlisting}

token 流和 AST 的数量级差异，用两条命令就能亲眼看到
（清单 \ref{lst:ch01-tokenstream}）。

\begin{lstlisting}[style=shellstyle,caption={分别取出 token 流、预处理结果与 AST},label={lst:ch01-tokenstream}]
# 1) 只要 token：每个 token 一行，带来源与行列
clang++ -Xclang -dump-tokens -fsyntax-only demo.cpp | wc -l

# 2) 只要预处理后的扁平源码：宏在这里已经被替换掉
clang++ -std=c++17 -E demo.cpp | wc -l

# 3) 完整 AST：包含所有从系统头里拉进来的声明
clang++ -fsyntax-only -Xclang -ast-dump demo.cpp | wc -l
\end{lstlisting}

典型结果是第一、第二项几千到几万行，第三项几十万行——差别不在于
\texttt{main} 有多长，而在于第 5 行那个 \texttt{\#include <vector>}
被展开后带进来的整条依赖链（第 3 章会把这个倍率量化）。
把 \texttt{-dump-tokens} 的输出截一段看形状：

\begin{lstlisting}[style=shellstyle,caption={\texttt{-dump-tokens} 输出的形状（截断）},label={lst:ch01-tokenout}]
Token 'l' Identifier   demo.cpp:11:7   Length   8   Escaped 0
Token '(' l_paren      demo.cpp:11:12  Length   1   Escaped 0
Token 'T' Identifier   demo.cpp:11:16  Length   1   Escaped 0
Token 'v' Identifier   demo.cpp:11:18  Length   1   Escaped 0
Token ')' r_paren      demo.cpp:11:19  Length   1   Escaped 0
...
\end{lstlisting}

每个 token 都带着\textbf{文件名、行、列}，这正是 \Fmt\ 能工作的全部前提：
它只需要位置信息，加一个对 token 角色的粗略猜测，不需要同位置的 \Tp{T}
在这里到底是模板参数还是别的东西。

\subsection*{为什么只有 \Fmt\ 不吃完整 AST}

\Fmt\ 内部是一条独立的流水线：词法分析 → token 注解
（\texttt{Token\allowbreak Annotation}，判断这是左花括号还是 lambda 引导符？
这是模板的 \texttt{<} 还是小于号？）→ 排版决策（按 Style 字段换行与缩进）。
它不调用 \texttt{Sema}，不做名字查找，因此：

\begin{itemize}
  \item 能吃\textbf{编译不过}的代码：包含缺失、类型未知都无所谓；
  \item 能吃\textbf{片段}：选中一段代码只排这一段是可行的；
  \item 能吃宏包裹的整块代码：它有一层宏体预处理，会把形如
        清单 \ref{lst:ch01-demo} 第 2、3 行那种整块宏临时展开后再排；
  \item 不需要编译数据库——这是它能在 CI 里以秒级完成全仓检查的根本原因。
\end{itemize}

代价是：当注解猜错角色时，输出会莫名其妙。最常见的两类是位移运算符连写的
模板（\texttt{a << b<{}c>}）与函数式宏 cast。此时正确做法不是换格式化工具，
而是给那一行加上关闭格式化的注释（第 9 章展开）。

\Fmt\ 也不是完全拒绝 AST：当它拿得到编译数据库时，会做一次轻量解析来
辅助消歧，例如判断某个标识符是不是类型名，用于排序 \texttt{using} 声明。
拿不到就退化为纯 token 模式，格式化结果基本一致，只是少了几处依赖语义的决策。

\begin{guideline}[title={哪些能力属于哪一层}]
词法层能做的：缩进、换行、空格、include 排序。
AST 层多出来的：识别声明与定义的关系、把符号解析到具体 decl、
重命名与移动代码。Sema 层多出来的：类型诊断、模板实例化结果、
\texttt{constexpr} 的求值，以及绝大多数 \Tidy\ 检查需要的
「这个表达式有没有未定义行为」这类判断。层次越深，越依赖正确的编译命令，
也越慢。
\end{guideline}

\section{一个文件在工具手里经过的对象}
\label{sec:ch01-s2}

把 clang 前端当作对象图看，四个类撑起了工具需要的全部输入。
表 \ref{tab:ch01-objects} 按「工具实际会向它要什么」来组织。

\begin{table}[htbp]
\centering
\caption{clang 前端四个核心对象与工具的实际用法}
\label{tab:ch01-objects}
\begin{tabular}{@{}>{\raggedright\arraybackslash\small}p{0.26\textwidth}
                 >{\raggedright\arraybackslash\small}p{0.28\textwidth}
                 >{\raggedright\arraybackslash\small}p{0.28\textwidth}@{}}
\toprule
对象 & 它掌管什么 & 工具真正向它要的东西 \\
\midrule
\texttt{Preprocessor} &
  宏表、条件编译的求值结果、已加载文件列表、token lexer &
  \Fmt\ 只要它的 token 流；\CD\ 用它的宏表把光标处的宏展开显示出来
  （hover 时给出「展开后是……」）；include 分析类工具看它记录了
  哪些文件被谁包含 \\
\texttt{SourceManager} &
  所有已加载文件的行表、宏展开调用栈、\texttt{Source\allowbreak Location}
  到行列的换算 &
  任何「从 AST 里的一个位置回到编辑器里的某一行」都要过它；
  宏展开产生的诊断必须靠它才能指回真正的书写处 \\
\texttt{ASTContext} &
  AST 的内存池、内建类型与规范类型、目标信息（triple、字长、对齐） &
  \Tidy\ 的类型判断、\CD\ 的补全候选与 inlay hint
  （推导出来的 \Tp{auto} 实际类型）都从这里取 \\
\texttt{Diagnostic\allowbreak sEngine} &
  严重级别、\texttt{-Wxxx} 诊断组、抑制规则、输出格式化器 &
  所有工具的诊断出口；\Tidy\ 把自己的检查伪装成一个 warning 组塞进这里，
  好复用 clang 的呈现逻辑 \\
\bottomrule
\end{tabular}
\end{table}

\subsection*{诊断消息的格式怎么读}

所有诊断共用一个模板：\texttt{file:line:col: severity: text [-Wxxx]}。
清单 \ref{lst:ch01-diag} 是一段真实的 \texttt{clang++ -fsyntax-only} 输出，
包含三种最需要看懂的情形。

\begin{lstlisting}[style=shellstyle,caption={一段诊断输出：范围标记、宏展开链与目标文件},label={lst:ch01-diag}]
demo.cpp:23:9: warning: unused variable 'n' [-Wunused-variable]
   23 |     int n = xs.size();
      |         ^
demo.cpp:24:12: error: no matching function for call to 'add'
   24 |     return add(b, 1);
      |            ^~~
hdr/ops.h:7:5: note: expanded from macro 'TRY'
    7 |     add(__VA_ARGS__)
      |     ^
demo.cpp:12:1: error: expected '}'
   12 | BEGIN_SCOPE
      | ^
hdr/scope.h:3:19: note: expanded from macro 'BEGIN_SCOPE'
    3 | #define BEGIN_SCOPE {
      |                   ^
\end{lstlisting}

读法分五条：

\begin{itemize}
  \item 第一行的 \texttt{warning}、\texttt{error}、\texttt{note}、
        \texttt{fatal error} 是严重级别。\texttt{note} 永远依附于前一条
        主诊断，单独看没有意义。
  \item 尾部方括号里的 \texttt{-Wunused-variable} 是\textbf{诊断组名}，
        它就是开关的名字：\texttt{-Wno-unused-variable} 关掉它，
        \texttt{-Werror=unused-variable} 单独把它升级成错误。
        不知道组名就无法精确控制，所以不要用
        \texttt{-fno-caret-diagnostics} 之类把这一栏关掉。
  \item 源码行下方 \texttt{|} 里的 \texttt{\^{}} 与波浪线标出诊断覆盖的
        \textbf{源范围}（换算由 \texttt{SourceManager} 完成）。
        LSP 里诊断对象的 \texttt{range} 字段就来自这里，因此编辑器中
        波浪线的起止应当与终端里标记的起止一致；若明显不一致，
        多半是中间有宏被展开过。
  \item \texttt{expanded from macro} 表示真实位置在宏体里，此时行列指向
        宏定义所在文件；要回到调用处，看紧邻上一条诊断的行号。
  \item 部分诊断带 fix-it 提示，例如 \texttt{-Wformat} 建议把
        \texttt{\%d} 改成 \texttt{\%ld}。\Tidy\ 会把这类提示转成可应用的
        替换，所以它修过的某些编译器警告并非来自某个 check，
        而是来自 clang 自己的诊断。
\end{itemize}

\section{三层 API：libclang、Clang C++ API 与现成工具}
\label{sec:ch01-s3}

同一份前端对外有三层接口，选错层的代价是重写。

\begin{table}[htbp]
\centering
\caption{三层接口的定位对比}
\label{tab:ch01-api}
\begin{tabular}{@{}>{\raggedright\arraybackslash\small}p{0.16\textwidth}
                 >{\raggedright\arraybackslash\small}p{0.26\textwidth}
                 >{\raggedright\arraybackslash\small}p{0.33\textwidth}@{}}
\toprule
维度 & libclang（C API） & Clang C++ API（内部） \\
\midrule
接口形态 &
  以 \texttt{CX} 为前缀的不透明句柄加 C 函数，
  头文件 \texttt{clang-c/Index.h} &
  \texttt{clang::} 命名空间下的真实类：\texttt{ASTConsumer}、
  \texttt{RecursiveASTVisitor}、\texttt{Sema} \\
ABI 稳定性 &
  \textbf{稳定}，官方承诺向后兼容，适合做绑定层 &
  \textbf{不承诺}，随大版本甚至小版本重构，内部接口会改名 \\
能力 &
  遍历 decl、取位置与拼写、诊断、少量文本替换、补全与符号索引 &
  全部：改写类型、控制实例化、使用 \texttt{ASTMatcher} 匹配结构、
  接 \texttt{Rewriter} 做编辑、跑数据流分析 \\
典型用户 &
  Python、C\#、Rust 的语言绑定，IDE 插件，轻量代码检索 &
  \CD、\Tidy、clang-rename 全部用它；\Fmt\ 的语义辅助部分也用它 \\
代价 &
  细粒度控制少，性能一般，很多概念（模板依赖表达式、部分实例化）
  表达不出来 &
  必须跟着 LLVM 源码或对应发行版走，链接整个 clang 库 \\
\bottomrule
\end{tabular}
\end{table}

清单 \ref{lst:ch01-libclang} 是 libclang 的最小可用样例：打开一个文件，
把所有函数声明及其行列打印出来。注意它不需要编译数据库，
搜索参数靠手工传入第三个实参。

\begin{lstlisting}[style=cppstyle,caption={libclang：遍历 cursor 打印所有函数声明},label={lst:ch01-libclang}]
// 依赖：clang-c/Index.h 与 libclang 动态库
#include <clang-c/Index.h>
#include <cstdio>

static CXChildVisitResult onCursor(CXCursor cur, CXCursor parent,
                                   void *unused)
{
    (void)parent;
    (void)unused;
    if (clang_getCursorKind(cur) == CXCursor_FunctionDecl) {
        CXString name = clang_getCursorSpelling(cur);
        unsigned line = 0;
        clang_getSpellingLocation(clang_getCursorLocation(cur),
                                  nullptr, &line, nullptr, nullptr);
        std::printf("%s  line %u\n", clang_getCString(name), line);
        clang_disposeString(name);
    }
    return CXChildVisit_Recurse;   // 继续深入命名空间与类
}

int main(int argc, char **argv)
{
    if (argc < 2) {
        std::printf("usage: dumpfn <file.cpp>\n");
        return 2;
    }
    const char *args[] = { "-std=c++17", "-Iinclude" };
    CXIndex index = clang_createIndex(0, 1);
    CXTranslationUnit tu = clang_parseTranslationUnit(
        index, argv[1], args, 2, nullptr, 0, CXTranslationUnit_None);
    if (!tu) {
        std::printf("parse failed\n");
        clang_disposeIndex(index);
        return 1;
    }
    clang_visitChildren(clang_getTranslationUnitCursor(tu),
                        onCursor, nullptr);
    clang_disposeTranslationUnit(tu);
    clang_disposeIndex(index);
    return 0;
}
\end{lstlisting}

\begin{lstlisting}[style=shellstyle,caption={编译并运行 libclang 程序},label={lst:ch01-libclang-build}]
# Ubuntu，apt 装的 llvm-18
clang++ dumpfn.cpp -o dumpfn \
    -I/usr/lib/llvm-18/include -L/usr/lib/llvm-18/lib -lclang

# Windows，releases.llvm.org 默认安装
clang++ dumpfn.cpp -o dumpfn.exe ^
    -I"C:/Program Files/LLVM/include" ^
    -L"C:/Program Files/LLVM/lib" -lclang

./dumpfn demo.cpp
\end{lstlisting}

三个值得记住的实现细节：\texttt{clang\_createIndex} 的第二个参数为 1
表示把诊断打印到标准错误；遍历回调返回
\texttt{CXChildVisit\_\allowbreak Recurse} 才会进入函数体内部；
所有 \texttt{CXString} 必须配对 \texttt{clang\_disposeString}，
否则常驻进程会稳定泄漏内存。用 libclang 写跨文件索引是可行的，
但它给不出「这行代码在当前编译命令下究竟实例化成了什么」，
那需要 Clang C++ API。

\begin{warning}[title={libclang 不是 clang 的轻量版，而是另一套语义视图}]
同一个类型在 libclang 里可能只给回 \texttt{CXType\_Unexposed}：
模板依赖类型、部分实例化的结果、以及许多重载决议的中间态都不暴露。
用 libclang 做代码质量检查的项目，常见结局是判定逻辑里塞满针对
\texttt{Unexposed} 的特例，正确性靠运气。若你要的是「精确」，
从 \Tidy\ 的自定义 check（第 12、14 章）起步，比从 libclang 起步省时间。
\end{warning}

\section{clang-tools-extra 工具全家表}
\label{sec:ch01-s4}

\texttt{clang-tools-extra} 是 LLVM 单仓里存放 clang 衍生工具的目录，
和 clang 在同一次 release 里一起构建。表 \ref{tab:ch01-tools} 列出本册
真正会用到的那些，按依赖层次与是否需要编译数据库分类。
最后一栏是本册展开它的章节。

\begin{longtable}{@{}>{\raggedright\arraybackslash\small}p{0.19\textwidth}
                 >{\raggedright\arraybackslash\small}p{0.21\textwidth}
                 >{\raggedright\arraybackslash\small}p{0.13\textwidth}
                 >{\raggedright\arraybackslash\small}p{0.13\textwidth}
                 >{\raggedright\arraybackslash\small}p{0.11\textwidth}@{}}
\caption{clang-tools-extra 相关工具一览}\label{tab:ch01-tools} \\
\toprule
工具 & 依赖层次 & 需要数据库 & 自动修复 & 详见 \\
\midrule
\endfirsthead
\toprule
工具 & 依赖层次 & 需要数据库 & 自动修复 & 详见 \\
\midrule
\endhead
\Fmt &
  词法与 token 注解 & 否（有则更准） & 是，原地改写 & 第 9 章 \\
\CD &
  完整 AST 加后台索引 & \textbf{是}，缺失则退化 & 通过 code action & 第 5 章 \\
\Tidy &
  AST、ASTMatcher、数据流 & \textbf{是} & \Opt{fix} & 第 12 章 \\
clang-rename &
  AST 与符号表 & \textbf{是} & 原地 & 第 14 章 \\
clang-move &
  AST 与 include 关系 & \textbf{是} & 原地 & 第 14 章 \\
clang-include-fixer &
  预先生成的符号索引 & 是，或用索引文件 & 插入 include & 第 14 章 \\
clang-apply-replacements &
  只消费 YAML 替换文件 & 否 & 是 & 第 14 章 \\
clang-query &
  AST，交互式匹配器 & \textbf{是} & 否 & 第 14 章 \\
clang-scan-deps &
  预处理的 include 与 module 扫描 & 是，逐条命令 & 否 & 第 3、16 章 \\
include-cleaner &
  AST 记录的 use 信息 & \textbf{是} & 是，删或加 include & 第 15 章 \\
clang-check &
  AST 与 Sema，clang 自带 & \textbf{是} & 是，见下方说明 & 第 2 章 \\
\bottomrule
\end{longtable}

关于最后两行：\texttt{clang-check} 不在 extra 目录里，它是 clang 自带的
诊断驱动器，靠 \texttt{-fixit-errors} 与 \texttt{-fixit-warnings}
两个开关决定是否把 fix-it 提示真正写回文件；
而 \texttt{clang-query} 是交互式 REPL，适合回答
「全仓有多少个类继承自 \Tp{Base}」这类结构性问题。
两者都要求能编译通过，因此都需要第 2 章的编译数据库。

\section{版本、安装与 PATH 打架}
\label{sec:ch01-s5}

\subsection*{三者必须来自同一次 LLVM release}

\Fmt、\Tidy、\CD\ 和编译器 \texttt{clang} 在同一个仓库里发布，
\Opt{version} 输出里的版本号应当完全一致，如清单 \ref{lst:ch01-version}。

\begin{lstlisting}[style=shellstyle,caption={版本自检：三行必须落在同一 release},label={lst:ch01-version}]
clang --version        | head -1
clang-format --version
clangd   --version     | head -1

# 期望形如（小版本可以不同，主版本必须一致）：
# clang version 18.1.8
# clang-format version 18.1.8
# clangd version 18.1.8
\end{lstlisting}

版本不一致会带来三类具体故障：\texttt{.clang-tidy} 里写的 check 名在
另一个版本不存在，\Tidy\ 报 \texttt{unknown check} 直接退出；
\texttt{.clangd} 里的配置键被当作未知键静默忽略，你以为是配置没生效；
\Fmt\ 与 \CD\ 内嵌格式化器对同一文件给出不同结果，于是出现
「编辑器里排过了，CI 说不合格」。第三类是本册第 10 章整章要解决的问题。

\subsection*{五个常见安装源}

\begin{itemize}
  \item Ubuntu 与 Debian：\texttt{apt install clang-tools-18}，
        可执行文件在 \texttt{/usr/lib/\allowbreak llvm-18/\allowbreak bin/}，
        另有 \texttt{/usr/bin/clang-format-18} 这类带后缀的名字。
        发行版默认的无后缀名字可能指向别的版本，
        用 \texttt{update-alternatives} 显式绑定。
  \item macOS Homebrew：\texttt{brew install llvm@18}。这个 formula 是
        keg-only 的，\textbf{不会自动进 PATH}，需要把
        \texttt{/opt/homebrew/opt/llvm@18/bin} 追加到 \texttt{PATH} 里
        （Intel 机型前缀是 \texttt{/usr/local/opt}）。
        与此同时 \texttt{brew install clang-format} 会装一个独立 formula，
        两个来源的版本必然漂移。
  \item releases.llvm.org 官方安装包：Windows 的 MSI 与
        Linux 的 tar.xz，默认落在 \texttt{/usr/local/} 或
        \texttt{C:\textbackslash Program Files\textbackslash LLVM}。
  \item Visual Studio 自带：可选组件「C++ Clang tools for Windows」，
        目录在 \texttt{VC\textbackslash Tools\textbackslash Llvm\textbackslash x64\textbackslash bin}。
        它的版本跟着 VS 走，常常落后于 LLVM 主线若干个小版本。
  \item choco 与 scoop：\texttt{choco install llvm}、
        \texttt{scoop install llvm}，都是滚动更新，
        升级时全家桶一起跳版本，无法单独钉住某个工具。
\end{itemize}

\begin{lstlisting}[style=shellstyle,caption={查出机器上到底有几份 clangd},label={lst:ch01-which}]
# Windows cmd
where clangd

# Windows PowerShell：一次看全，并显示每个命中的来源目录
Get-Command clangd -All | Select-Object -ExpandProperty Source

# POSIX shell
which -a clangd
readlink -f "$(which clangd)"

# 确认某个 clang-format 属于哪次 LLVM release
"/usr/local/bin/clang-format" --version
\end{lstlisting}

Windows 上标准 LLVM 目录的形状值得记住，因为它直接决定内置头能否被找到，
交叉编译那一节（第 2 章的 \Opt{query-driver}）会回到这里：

\begin{lstlisting}[style=shellstyle,caption={Windows 上 LLVM 安装的目录结构},label={lst:ch01-win-tree}]
C:\Program Files\LLVM
|-- bin\    clang.exe clang++.exe clang-format.exe clang-tidy.exe
|           clangd.exe clang-scan-deps.exe clang-rename.exe
|           libclang.dll llvm-config.exe
|-- include\   clang-c\Index.h          <-- libclang 的 C API 头
|-- lib\clang\18\include\   stdint.h    <-- 编译器内置头
|                           immintrin.h 等一组 intrinsic 头
\end{lstlisting}

三个目录各有用途：\texttt{bin} 进 PATH；\texttt{include} 是编译
libclang 程序时要 \texttt{-I} 的；
\texttt{lib\textbackslash clang\textbackslash 主版本号\textbackslash include}
是编译器自带的内置头，路径里带主版本号，所以升级 LLVM 之后旧路径会失效。
把这个绝对路径写死在脚本里是常见的坏味道，正确做法是用
\texttt{clang -print-resource-dir} 现场查询。

\begin{warning}[title={VS 自带 clang 与独立 LLVM 混装的典型症状}]
症状：\CD\ 在每个文件上都报 \texttt{'cstdint' file not found}，
或者报 \texttt{unknown argument}，但命令行 \texttt{clang++} 编译完全正常。
原因：编辑器启动的 \texttt{clangd.exe} 来自 VS 组件目录，
而项目真实使用的编译器是独立安装的 LLVM。这一份 \CD\ 找不到与自己配套的
内置头目录，也读不到 MSVC 需要的 \texttt{INCLUDE} 与 \texttt{LIB}
环境变量——它们只在 \texttt{vcvars64.bat} 打开的 shell 里存在。
对策：收敛 PATH，让 \texttt{where clangd} 只剩一条；
或者在 VS Developer Command Prompt 里启动编辑器；
或者给 \texttt{.clangd} 配 \Opt{query-driver} 指向真实编译器（第 2 章）。
\end{warning}

\begin{bestpractice}[title={把工具版本钉死的一次性投入}]
\begin{enumerate}
  \item 选一个主版本（比如 18），团队内只允许这一个主版本。
  \item CI 里显式安装\textbf{带版本号}的包，例如
        \texttt{apt install clang-tools-18}，不用 \texttt{llvm} 这种滚动包名。
  \item 本地也把 PATH 收敛到单一目录。VSCode 用户在
        \texttt{settings.json} 里写死
        \texttt{"clangd.path"} 指向绝对路径，不要让编辑器自己找。
  \item 仓库根放一个记录版本号的文件，CI 第一步用 \Opt{version}
        与之比对，不符则整条流水线失败。
\end{enumerate}
这条纪律能消掉本册第 10 章与第 17 章里将近一半的疑难杂症。
\end{bestpractice}

% ---- frag_ch02.tex ----
\chapter{compile\_commands.json：编译数据库的格式与生成}
\label{ch:compdb}

第 1 章说过：凡是要看懂语义的工具，都需要知道\textbf{这个文件究竟是怎么被编译的}。
这份信息的载体就是一个 JSON 数组，业界统一叫\textbf{编译数据库}
（compilation database）。它不是可选项：数据库缺失时 \CD\ 依然会启动，
只是所有语义功能悄悄降级，你会以为「clangd 不准」。本章把格式逐字段拆开，
覆盖 CMake、Ninja、Make、Meson、Bazel 五条生成途径与交叉编译的坑，
最后给一条可以照抄的完整实战链。

\section{格式：四个字段决定一切}
\label{sec:ch02-s1}

文件是一个数组，每个元素描述\textbf{一个翻译单元}。必须的键只有三个：
\texttt{directory}、\texttt{file}，以及 \texttt{command} 或 \texttt{arguments}
二者取其一；\texttt{output} 可选。清单 \ref{lst:ch02-min} 是一个包含两个
翻译单元的最小完整数据库。

\begin{lstlisting}[style=jsonstyle,caption={最小完整数据库：两个翻译单元，分别用 command 与 arguments},label={lst:ch02-min}]
[
  {
    "directory": "/home/me/proj/build",
    "command": "clang++ -std=c++17 -Iinc -c src/net.cpp",
    "file": "/home/me/proj/src/net.cpp"
  },
  {
    "directory": "/home/me/proj/build",
    "arguments": [
      "clang++", "-c", "-std=c++17", "-DDEBUG=1",
      "-I/home/me/proj/include", "-o", "main.o",
      "/home/me/proj/src/main.cpp"
    ],
    "file": "/home/me/proj/src/main.cpp",
    "output": "main.o"
  }
]
\end{lstlisting}

\begin{itemize}
  \item \texttt{directory}：执行编译命令时的工作目录。\textbf{必须是绝对路径，
        并且真实存在}。所有相对路径（包括 \texttt{file} 里写的相对路径）
        都以它为基准解释。
  \item \texttt{file}：翻译单元的源文件。可以是相对路径（相对
        \texttt{directory}），也可以绝对；但工具用它做键来查找条目，
        路径解析不一致会导致「数据库里明明有这个文件却说不在数据库里」。
  \item \texttt{command}：一条完整的 shell 命令字符串，包含编译器可执行文件、
        所有选项、以及 \texttt{-c} 与 \texttt{-o}。
  \item \texttt{arguments}：同一条命令\textbf{已经拆好}的参数数组，
        第一个元素是编译器。
  \item \texttt{output}：目标文件路径，供少数工具推断依赖，可省。
\end{itemize}

\subsection*{为什么优先用 arguments}

工具的读法决定了这一点：\CD、\Tidy\ 这类基于 Clang 工具库的程序
\textbf{从不启动 shell}。拿到 \texttt{command} 字符串后，它们用自己的
词法器按 POSIX 规则切词；拿到 \texttt{arguments} 数组则直接跳过这一步。
差别在两种情况下会暴露：

\begin{enumerate}
  \item 路径里有空格。Windows 上 \texttt{-I"C:\textbackslash Program Files\textbackslash inc"}
        这种写法，切词器对引号与反斜杠的处理和真实 shell 并不完全一致，
        拆出来的参数会错位。
  \item 值里带 shell 元字符。\texttt{-DMSG="hello world"}、
        \texttt{-DVER=\$ revision}、反引号、单引号嵌套，都是重灾区。
  \item 混入 shell 结构。有人手工编辑数据库，加了
        \texttt{mkdir -p out \&\& clang++ \dots} 或
        \texttt{CC=g++ clang++ \dots}，切词器会把这些当成编译器参数。
\end{enumerate}

因此：能生成 \texttt{arguments} 的形式就优先用它；只能得到
\texttt{command} 时，务必确认那串内容里没有 shell 语法。
表 \ref{tab:ch02-field} 给出字段的严格程度。

\begin{table}[htbp]
\centering
\caption{工具对数据库字段的实际要求}
\label{tab:ch02-field}
\begin{tabular}{@{}>{\raggedright\arraybackslash\small}p{0.17\textwidth}
                 >{\raggedright\arraybackslash\small}p{0.22\textwidth}
                 >{\raggedright\arraybackslash\small}p{0.36\textwidth}@{}}
\toprule
键 & 是否必需 & 常见错误写法与后果 \\
\midrule
\texttt{directory} & 必需 &
  写成相对路径；指向已被删除的构建目录；换了机器后旧绝对路径残留
  ——后果是文件根本读不到，\CD\ 报 \texttt{not found} \\
\texttt{file} & 必需 &
  拼成 \texttt{filename} 或 \texttt{source}；与 \texttt{directory}
  拼接后指向不存在的文件——该 TU 被静默跳过 \\
\texttt{command} & 二选一 &
  含 shell 变量或未加引号的空格路径；缺少 \texttt{-c}
  导致被当成链接命令 \\
\texttt{arguments} & 二选一 &
  把整条命令塞进一个元素；混入 \texttt{echo} 之类非编译器 token \\
\texttt{output} & 可选 &
  一般无需关心；少数工具用它定位对象文件 \\
\bottomrule
\end{tabular}
\end{table}

\section{生成途径一：CMake 与 CMakePresets}
\label{sec:ch02-s2}

CMake 是绝大多数 \CPP\ 项目的答案，一条开关即可：

\begin{lstlisting}[style=shellstyle,caption={用 CMake 生成数据库（只 configure 就够）},label={lst:ch02-cmake}]
cmake -B build -G Ninja \
      -DCMAKE(*@\allowbreak@*)_EXPORT(*@\allowbreak@*)_COMPILE(*@\allowbreak@*)_COMMANDS=ON

ninja -C build            # 可选：configure 已经写出数据库
ls build/compile_commands.json
\end{lstlisting}

关键事实与限制：

\begin{itemize}
  \item 这个开关\textbf{只对 Makefile 与 Ninja 系生成器有效}。
        Visual Studio 与 Xcode 生成器不会写出
        \texttt{compile\_\allowbreak commands\_\allowbreak json}，
        文档层面就是如此，不是 bug。
  \item 数据库写在\textbf{二进制目录}（\texttt{build/}）根下。
        \CD\ 的搜索规则是：从被打开文件的目录向上走，
        在每一层找 \texttt{compile\_\allowbreak commands\_\allowbreak json}，
        同时\textbf{额外检查该层下的 \texttt{build/} 子目录}。
        所以只要用 \texttt{-B build}，就完全不需要手工做软链接。
  \item 新增或删除源文件后必须\textbf{重新 configure}，
        数据库不会自动更新，否则新文件不在表内。
  \item 想永久打开它，可以在项目 \texttt{CMakeLists.txt} 开头写
        \texttt{set(CMAKE\_\allowbreak EXPORT\_\allowbreak COMPILE\_\allowbreak COMMANDS\ ON)}，
        但更推荐只由 Preset 或命令行控制，避免污染 CI 的 Release 构建。
\end{itemize}

VS 与 Xcode 生成器的处理办法只有一个正确的方向：\textbf{为了工具链再配一个
Ninja 构建目录}，它和正式构建互不干扰。清单 \ref{lst:ch02-preset}
是 CMakePresets 的写法，把编译器与开关一起钉住，团队里每个人
\texttt{cmake --preset dev} 得到同一个数据库。

\begin{lstlisting}[style=jsonstyle,caption={CMakePresets.json：dev 预设固定 Ninja 并导出数据库},label={lst:ch02-preset}]
{
  "version": 3,
  "configurePresets": [
    {
      "name": "dev",
      "hidden": false,
      "generator": "Ninja",
      "binaryDir": "${sourceDir}/build",
      "cacheVariables": {
        "CMAKE_BUILD_TYPE": "Debug",
        "CMAKE(*@\allowbreak@*)_EXPORT(*@\allowbreak@*)_COMPILE(*@\allowbreak@*)_COMMANDS": "ON"
      }
    }
  ]
}
\end{lstlisting}

Windows 上让 Ninja 找到 MSVC 的标准做法（清单 \ref{lst:ch02-win-ninja}）：
先在 VS Developer Command Prompt 里执行，或用 \texttt{vswhere} 定位后
手工导入环境。这一步和数据库本身无关，但没有它
\texttt{cmake -G Ninja} 会连编译器都找不到。

\begin{lstlisting}[style=shellstyle,caption={Windows：导入 MSVC 环境后再生成 Ninja 数据库},label={lst:ch02-win-ninja}]
:: 在 cmd 中执行；VS 的安装路径按你的版本调整
set VSDIR=C:\Program Files\Microsoft Visual Studio\2022\Community
call "%VSDIR%\VC(*@\allowbreak@*)\Auxiliary(*@\allowbreak@*)\Build\vcvars64.bat"
cmake -B build -G Ninja \
      -DCMAKE(*@\allowbreak@*)_EXPORT(*@\allowbreak@*)_COMPILE(*@\allowbreak@*)_COMMANDS=ON
\end{lstlisting}

\begin{note}[title={数据库只需 configure，不需要真的编译成功}]
\ttfamily CMAKE\_\allowbreak EXPORT\_\allowbreak
COMPILE\_\allowbreak COMMANDS\normalfont\ 在生成构建系统的那一刻就写文件，
内容来自 CMake 的属性，不来自实际编译。也就是说：项目当前编译不过、
甚至一个对象都还没产出，数据库照样可用。反过来，
如果 \texttt{build/} 里的数据库和源码树不同步
（改了 \texttt{CMakeLists.txt} 却没重新 configure），
工具看到的仍然是旧内容。
\end{note}

\section{生成途径二：Ninja、Make、Meson、Bazel}
\label{sec:ch02-s3}

\subsection*{Ninja 自带导出}

已经用 Ninja 描述构建的项目（不经过 CMake 也行），可以直接问 Ninja 要：

\begin{lstlisting}[style=shellstyle,caption={从 build.ninja 导出编译数据库},label={lst:ch02-ninja}]
ninja -C build -t compdb > /tmp/compile_commands.json

# 只要某类目标的命令，缩小数据库体积
ninja -C build -t compdb objects > compile_commands.json
\end{lstlisting}

\texttt{-t compdb} 接受目标类型名作为参数（可选），输出的 JSON 里
\texttt{command} 是原始规则串，若你的 \texttt{build.ninja} 用了
\texttt{\$} 变量插值，导出结果已经是插值后的形式。

\subsection*{Make 项目：bear 与 \texttt{compiledb}}

裸 Makefile 不携带结构化信息，只能\textbf{观察真实执行}再还原。
两条路：

\begin{itemize}
  \item \texttt{bear}（Build EAR）：通过 \texttt{LD\_PRELOAD}
        注入一个拦截 \texttt{exec} 系列调用的库，把每次真实执行的
        进程与参数记录下来，构建结束后写成数据库。macOS 上用
        \texttt{DYLD\_INSERT\_LIBRARIES} 或 interpose 机制实现同样效果。
        用法是 \texttt{bear -I -- make -j8}，
        \texttt{-I} 表示忽略子目录里的内部命令。它只能用于
        真正会调用编译器的构建；\textbf{Windows 上没有对应机制}，
        nmake 与 MSBuild 都不能用 bear。
  \item \texttt{compiledb}：一个 \PY\ 脚本，
        \texttt{pip install compiledb}。它解析
        \texttt{make -n}（伪执行）或既有构建日志，不需要注入，
        因此跨平台。典型用法
        \texttt{compiledb -n make}；日志重放用
        \texttt{compiledb -s build.log}。
        注意伪执行拿不到变量展开后的内容时，产出的命令可能是
        \texttt{\$CC} 这种未展开形式，必须检查。
\end{itemize}

\subsection*{Meson 与 Bazel}

Meson 原生就写数据库：\texttt{meson setup build} 之后
\texttt{build/compile\_\allowbreak commands\_\allowbreak json} 已经存在，
可以直接把 \texttt{build/} 当作 \texttt{Compilation\allowbreak Database} 目录，
或者按 Meson 文档的建议在项目根做一个软链接。

Bazel 走的是完全不同的模型（构建图 + action cache），
没有内置导出。社区做法是写一个 \emph{aspect}（切面），
在分析阶段收集每个编译 action 的参数再拼成数据库，
常见实现的名字是 \texttt{bazel-compilation-database}。
本册不展开 aspect API；需要的话，把它当作「构建系统外挂一个收集器」
来理解即可。

\section{真实世界的长命令与交叉编译}
\label{sec:ch02-s4}

清单 \ref{lst:ch02-long} 是数据库里一条 \texttt{command} 的内容
（为了排版拆成多行，实际是\textbf{一行}）。它同时包含
\texttt{-I}、\texttt{-D}、\texttt{-std}、\texttt{-isystem} 与 \texttt{-o}，
这类命令是工具行为异常的主要来源。

\begin{lstlisting}[style=jsonstyle,caption={一条典型 command 的组成（实际为 JSON 里的一行）},label={lst:ch02-long}]
/usr/bin/clang++
  -DQT_NO_DEBUG -DVERSION=3 -DMSG=\"hello\"
  -I/home/me/proj/include
  -I/home/me/proj/build/gen
  -isystem /opt/sysroot/armhf/usr/include
  -std=gnu++17 -O0 -g -fno-omit-frame-pointer
  -Wall -Wextra
  -c -o CMakeFiles/app.dir/src/main.cpp.o
  /home/me/proj/src/main.cpp
\end{lstlisting}

逐项的意义：\texttt{-I} 决定项目头能否找到；\texttt{-isystem}
决定第三方头是否被当作系统头（\Tidy\ 默认不报告系统头里的诊断，
\CD\ 也不会为它们做诊断）；\texttt{-std=} 决定语言模式，
漏写就会让 \cpp{17} 的结构化绑定报错；\texttt{-D} 决定条件编译分支，
漏写会让整块代码「消失」；\texttt{-c}/\texttt{-o} 告诉工具这是编译而非链接。

\subsection*{交叉编译：内置头找不到}

数据库里的编译器是 \texttt{arm-none-eabi-g++} 之类的交叉工具链时，
基于 Clang 的工具会遇到两个问题：它不知道该工具链\textbf{自带的头在哪}，
也不知道目标平台的\textbf{数据模型}。典型症状是
\texttt{'stddef.h' file not found}，而真实交叉编译一切正常。

解法是 \CD\ 的 \Opt{query-driver}：允许它去问那个真实编译器，
用一次 \texttt{-E -v} 的伪编译把系统头搜索路径抄回来。

\begin{lstlisting}[style=shellstyle,caption={用 query-driver 让 clangd 学习真实工具链},label={lst:ch02-qdriver}]
# 允许 clangd 调用匹配这一模式的编译器（注意是绝对路径或 glob）
clangd --query-driver=/usr/bin/arm-none-eabi-g++

# 多个来源用逗号分隔
clangd --query-driver=/usr/bin/arm-*-,/opt/gcc-12/bin/*-g++

# 配合数据库目录与 sysroot
clangd --compile-commands-dir=build-arm \
       --query-driver=/usr/bin/arm-none-eabi-g++ \
       --check=src/main.cpp
\end{lstlisting}

\Opt{query-driver} 是命令行选项而非配置文件选项，
出于安全考虑不能写进 \texttt{.clangd}，必须在启动参数里给。
与之配合的三个手段：

\begin{itemize}
  \item \texttt{sysroot}：交叉工具链的根。数据库里应当带着
        \texttt{--sysroot=/opt/sysroot/armhf}，否则即使内置头找到了，
        libc 头仍然缺失。
  \item \texttt{-\/-target=}：目标三元组，例如
        \texttt{arm-none-eabi}、\texttt{x86\_64-w64-mingw32}。
        它决定内建宏（\texttt{\_\_ARM\_EL\_\_} 之类）与数据模型，
        缺失时条件编译会走错分支。
  \item \texttt{-isystem}：手工追加只读头目录，适合 vendor SDK。
        与 \texttt{-I} 的区别在于诊断归属，见上一段。
\end{itemize}

在 \texttt{.clangd} 里做补充也是常见做法，清单 \ref{lst:ch02-clangdyaml}
展示最常用的三个动作：指定数据库目录、追加参数、删除不认识的参数。

\begin{lstlisting}[style=yamlstyle,caption={.clangd：修数据库而不是修构建系统},label={lst:ch02-clangdyaml}]
# 放在项目根，或数据库同级目录
CompileFlags:
  CompilationDatabase: build-arm     # 数据库所在目录
  Add:
    - -std=gnu++17
    - --target=arm-none-eabi
    - -isystem/opt/vendor/sdk/inc
  Remove:
    - -fno-omit-frame-pointer        # clang 不认的旧工具链选项
    - -mcpu=*
\end{lstlisting}

\texttt{Remove} 的价值在 Windows 上尤其明显：MSVC 的
\texttt{/EHsc}、\texttt{/MD}、\texttt{/std:c++17} 这类以斜杠开头的选项
完全无法被 Clang 系工具解析，用 \texttt{Remove: [/EHsc, /MD]}
剥掉它们，比把整套构建换成 clang-cl 省事得多。

\section{数据库之外的文件}
\label{sec:ch02-s5}

数据库只列 \texttt{.cpp}（以及被显式编译的 \texttt{.c}、\texttt{.cc}），
头文件从来不在里面。于是打开一个 \texttt{.h} 时，\CD\ 必须\textbf{猜}一份命令：

\begin{enumerate}
  \item 先看有没有「同一个文件名的 \texttt{.cpp}」在数据库里，
        用它的命令加一组启发式调整；
  \item 否则从索引里找一个包含这个头的翻译单元，借用它的命令；
  \item 都不行就退到 \textbf{fallback command}（等价于
        \texttt{clang++ -fsyntax-only <文件>}，没有项目的 \texttt{-I}）。
\end{enumerate}

第 3 种情况下，你会看到「所有项目内的符号都跳到定义失败、
补全列表只剩语言关键字」，但没有任何错误提示——这是最难被注意到的
降级模式。判断方法：在打开该头文件时看状态栏，或跑
\texttt{clangd --check=include/mylib.h}，
输出里会写它用的是推断命令还是回退命令。

另一个相关字段是 \Tidy\ 的 \texttt{HeaderFilterRegex}：
数据库只决定「用哪条命令检查」，而这个字段决定「报告是否覆盖头文件」。
两者不要混淆，第 12 章展开。

\section{实战链：从三个文件到可读的输出}
\label{sec:ch02-s6}

工程定义如清单 \ref{lst:ch02-cmakelists}，目标是拿到数据库并验证它可用。

\begin{lstlisting}[style=cmakestyle,caption={CMakeLists.txt：三个源文件的最小工程},label={lst:ch02-cmakelists}]
cmake_minimum_required(VERSION 3.21)
project(demo LANGUAGES CXX)

set(CMAKE_CXX_STANDARD 17)
set(CMAKE_CXX_STANDARD_REQUIRED ON)

add_library(core STATIC src/util.cpp src/util.h)
target_include_directories(core PUBLIC include src)

add_executable(app src/main.cpp)
target_link_libraries(app PRIVATE core)
\end{lstlisting}

\begin{lstlisting}[style=cppstyle,caption={src/main.cpp：故意留一处会被诊断的问题},label={lst:ch02-main}]
#include "util.h"
#include <vector>

int main()
{
    std::vector<int> xs = {3, 1, 2};
    int total = sum(xs);
    int unused = total * 2;      // 触发 -Wunused-variable
    return total > 0 ? 0 : 1;
}
\end{lstlisting}

\begin{lstlisting}[style=shellstyle,caption={完整操作序列：生成、查看、验证数据库},label={lst:ch02-chain}]
mkdir -p src include && cd proj

# 1) configure，同时导出数据库
cmake -B build -G Ninja -DCMAKE_EXPORT_COMPILE_COMMANDS=ON

# 2) 看数据库里到底有几条、都是什么命令
python3 -c "import json;print(len(json.load(open('build/compile_commands.json'))))"
jq -r '.[].file' build/compile_commands.json

# 3) 让 clangd 只读一个文件，把整条流水线跑一遍并打印统计
clangd --compile-commands-dir=build --check=src/main.cpp
\end{lstlisting}

清单 \ref{lst:ch02-dbentry} 是 \texttt{build/compile\_commands.json}
中与 \texttt{main.cpp} 对应的那一条（拆行显示）。

\begin{lstlisting}[style=jsonstyle,caption={数据库里 main.cpp 对应的条目
（为排版拆行，实际是一行；JSON 字符串内部不允许换行）},label={lst:ch02-dbentry}]
{
  "directory": "/home/me/proj/build",
  "command": "/usr/bin/clang++ -DAPP_VERSION=1
    -I/home/me/proj/include -I/home/me/proj/src
    -std=gnu++17
    -o CMakeFiles/app.dir/src/main.cpp.o
    -c /home/me/proj/src/main.cpp",
  "file": "/home/me/proj/src/main.cpp"
}
\end{lstlisting}

\begin{lstlisting}[style=shellstyle,caption={clangd --check 的输出（截断，右侧为解读）},label={lst:ch02-check}]
I[..] clangd version 18.1.8
I[..] Features: linux
I[..] PID: 31742
I[..] Working directory: /home/me/proj
I[..] config: Loading .clangd files from user directory
I[..] config: Loading .clangd files from file directory
I[..] Loaded compilation database from /home/me/proj/build
I[..] ASTWorker building file /home/me/proj/src/main.cpp version=1
I[..] Building preamble...
I[..] Building AST...
I[..] Indexing AST...
I[..] Testing features at each token
I[..] All checks passed!
\end{lstlisting}

读这四行就够：\texttt{Loaded compilation database from \dots} 确认数据库被找到
且路径正确；\texttt{Building preamble} 到 \texttt{Building AST} 的耗时
是冷启动的主要部分；\texttt{Indexing AST} 说明它在为跨文件功能建索引；
最后一行如果变成 \texttt{not found in the compilation database}，
那么本章前面的字段规则就要逐条对照。

\begin{warning}[title={生成数据库时的四个高频错误}]
\begin{enumerate}
  \item 把数据库放到 \texttt{/tmp} 却指望工具找到它：\texttt{directory}
        字段里的绝对路径若不存在，条目会被整体忽略。
  \item 手工复制了别的机器的数据库：\texttt{directory} 与 \texttt{file}
        全是旧绝对路径，症状是「数据库找到了但每个文件都不在数据库里」。
  \item 在 \texttt{command} 里保留 generator expression（CMake 的
        \texttt{\$<TARGET\_PROPERTY:\dots>}）。这类内容只有构建系统能求值，
        工具会把它当成未知选项。
  \item 用响应文件（\texttt{clang @args.rsp}）传参。\texttt{@}
        由编译器驱动在当前工作目录展开，工具的工作目录和
        \texttt{directory} 未必一致，结果时有时无。
\end{enumerate}
\end{warning}

\section{排错表：症状、根因、对策}
\label{sec:ch02-s7}

\begin{longtable}{@{}>{\raggedright\arraybackslash\small}p{0.24\textwidth}
                 >{\raggedright\arraybackslash\small}p{0.26\textwidth}
                 >{\raggedright\arraybackslash\small}p{0.30\textwidth}@{}}
\caption{编译数据库相关故障速查}\label{tab:ch02-debug} \\
\toprule
症状 & 根因 & 对策 \\
\midrule
\endfirsthead
\toprule
症状 & 根因 & 对策 \\
\midrule
\endhead
\CD\ 状态栏显示「no compilation database」 &
  数据库不在文件目录的祖先里，或 \texttt{build} 目录名不是这个 &
  \Opt{compile-commands-dir} 显式指定；或在 \texttt{.clangd}
  写 \texttt{Compilation\allowbreak Database} \\
每个文件都报「not in the compilation database」 &
  \texttt{directory} 指向不存在的路径，或 \texttt{file} 键名拼错 &
  \texttt{jq} 检查 \texttt{directory} 与 \texttt{file} 是否存在；
  重新 configure \\
新增的文件没有语义 &
  configure 之后没再跑，数据库仍是旧列表 &
  重新执行 \texttt{cmake -B build}；CI 里把 configure 作为固定步骤 \\
头文件里补全只剩关键字 &
  走的是回退命令，没有项目的 \texttt{-I} &
  让同名 \texttt{.cpp} 进入数据库；或用 \texttt{Add} 手工补 \texttt{-I} \\
\texttt{stddef.h}/\texttt{cstdint} 找不到 &
  工具链内置头目录没被学习（交叉编译、VS 自带 clang） &
  \Opt{query-driver} 指向真实编译器；检查
  \texttt{clang -print-resource-dir} \\
未知选项 \texttt{/EHsc} 之类 &
  数据库里的编译器是 MSVC 的 \texttt{cl.exe} &
  \texttt{CompileFlags} 的 \texttt{Remove} 剥掉斜杠选项，
  或改用 clang-cl \\
\texttt{@file.rsp} 时而生效时而不生效 &
  响应文件按当前工作目录展开，与 \texttt{directory} 不一致 &
  生成数据库时禁用响应文件，或把路径写成绝对路径 \\
\Tidy\ 报 \texttt{unknown check}（与数据库无关但常被混淆） &
  check 名不属于当前版本的 \Tidy &
  先确认工具版本组合（第 1 章），再检查数据库 \\
\bottomrule
\end{longtable}

\begin{bestpractice}[title={把数据库当成构建产物对待}]
\begin{enumerate}
  \item 数据库永远由构建系统生成，不手工维护；要改参数就改构建脚本
        或用 \texttt{.clangd} 的 \texttt{Add}、\texttt{Remove}。
  \item 把 \texttt{build/} 整体加入 \texttt{.gitignore}，
        数据库与对象文件一样不进版本库。
  \item CI 里在跑质量工具之前先跑
        \texttt{cmake -B build -G Ninja
        -DCMAKE\_EXPORT\_COMPILE\_COMMANDS=ON}，
        并加一步存在性检查：文件不存在就失败，别把降级模式带进流水线。
  \item 多构建目录（Debug 与交叉编译）并存时，为每个目录配一份
        \texttt{.clangd}，用 \texttt{CompilationDatabase} 分别指向，
        比在编辑器里切启动参数可靠。
\end{enumerate}
\end{bestpractice}

% ---- frag_ch03.tex ----
\chapter{索引与解析的代价：C\allowbreak/\allowbreak C++ 工具为什么慢}
\label{ch:indexing-cost}

前两章把工具看懂一个文件所需的两块地基铺好了：第 1 章说明\textbf{语义视图}
从何处来（clang 前端的预处理、解析、语义分析三段），第 2 章说明\textbf{这个
文件究竟怎么被编译}由谁来告诉工具（\texttt{compile\_\allowbreak commands\_\allowbreak json}）。
本章要正面回答一个所有人在第一次把 \CD\ 装进大型项目时都会撞上的问题：
\textbf{为什么打开一个文件要转这么久，为什么 CPU 会打满，为什么内存要吃掉两
GB。}

答案不在 LSP 协议里，也不在 \CD\ 的实现里，而在 \CPP\ 这门语言的
\textbf{翻译单元模型}里。它是所有 \CPP\ 工具成本的共同来源；谁也没有绕过它，
只能在它之上选择不同的取舍。把这条成本线量出来，才知道每一分缓解手段花在了
哪里。

\section{翻译单元：打开一个文件等于重做整个工程的一个切片}
\label{sec:ch03-s1}

\subsection*{一个 .cpp 展开后到底有多大}

\CPP\ 没有模块边界意义上的「一个项目的类型系统」，编译与一切分析的基本单位是
\textbf{翻译单元}（translation unit，简称 TU）：一个源文件加上它经过
预处理器递归展开的\textbf{全部}头文件，拼成的那一条扁平 token 流。第 1 章
讲过，解析、语义分析、诊断都发生在这条流上。工具要看懂
\texttt{widget.cpp} 里一行 \texttt{v.push\_back(x)}，就必须先把
\texttt{widget.cpp} \textbf{include 到的每一个头}都读进来，否则
\texttt{std::vector} 是谁、\texttt{push\_back} 的重载集从何而来都无从谈起。

这个「include 全部头文件」的倍率有多夸张？用两个纯前端的动作就能实测，不依赖
任何编辑器。清单 \ref{lst:ch03-e-line} 用 \texttt{clang++ -E} 只做预处理，
把展开后的 token 流打到标准输出，再用 \texttt{wc} 数行数。

\begin{lstlisting}[style=shellstyle,caption={用 -E 的行数实测 include 的倍率},label={lst:ch03-e-line}]
# size of the source file itself
wc -l src/widget.cpp
# 412 src/widget.cpp

# preprocess only, then count the expanded lines
clang++ -std=c++20 -Iinclude -E src/widget.cpp | wc -l
# 187342

# drop the C++ standard library to isolate project headers
clang++ -std=c++20 -Iinclude -E -nostdinc++ src/widget.cpp | wc -l
# 60118
\end{lstlisting}

同一个文件，源码 \texttt{412} 行，展开后 \texttt{187342} 行，倍率约
\texttt{450}。这个倍率随项目里 STL、\texttt{<tuple>}、
\texttt{<functional>}、以及层层转包的公共头\textbf{只增不减}。编辑器每打开
一个这样的文件，就等于把这 \texttt{187342} 行\textbf{重新预处理 + 重新解析}
一遍；你切十个标签页，就是十次。这是所有 \CPP\ 语义工具贵的根因，不是某
一个实现的缺陷。

\subsection*{头文件树：-H 告诉你倍率从哪来}

想知道这 \texttt{450} 倍是被谁乘出来的，\Opt{H} 会打印出 include 的嵌套树，
点号个数就是嵌套深度。清单 \ref{lst:ch03-h} 是一段真实形态的输出。

\begin{lstlisting}[style=shellstyle,caption={clang -H 打印 include 树，点号是深度},label={lst:ch03-h}]
clang++ -std=c++20 -Iinclude -fsyntax-only src/widget.cpp -H
. ./include/widget.h
.. /usr/include/c++/12/vector
... /usr/include/c++/12/bits/stl_vector.h
.... /usr/include/c++/12/bits/alloc_traits.h
..... /usr/include/c++/12/bits/ptr_traits.h
.... /usr/include/c++/12/bits/stl_construct.h
... /usr/include/c++/12/bits/vector.tcc
.. ./include/net/socket.h
... ./include/net/error.h
# the leaf count of this tree is the real size of the TU
\end{lstlisting}

从输出可以直接读出两件事。其一，\texttt{widget.cpp} 只写了两个
\texttt{\#include}，但递归展开后落到十几个头，其中 STL 那条支路又各自转包了
更深的实现头。其二，很多深度是\textbf{传递依赖}：你 include \texttt{vector}，
就白送了 \texttt{stl\_vector.h}、\texttt{alloc\_traits.h}、
\texttt{ptr\_traits.h} 一整套。想压缩体积，就得从这些转包下手（前置声明与
PIMPL 的意义在此，第 5 节量化它的收益）。

\subsection*{阶段耗时：-ftime-report 把成本摊开}

倍率解释了「要读多少」，但没解释「时间花在哪」。清单
\ref{lst:ch03-ftime} 用 \Opt{ftime-report} 让 clang 把前端各阶段的墙钟时间列
出来。它只编译一个 TU，不链接，输出即工具的真实负担分布。

\begin{lstlisting}[style=shellstyle,caption={-ftime-report 拆解单个 TU 的各阶段耗时},label={lst:ch03-ftime}]
clang++ -std=c++20 -Iinclude -c src/widget.cpp -ftime-report -o /dev/null
===------------------------------===
                          Timer Details
===------------------------------===
  Total Execution Time: 1.44 seconds (1.46 wall clock)

   ---Wall Time---  --- Name ---
   0.98 ( 68.4%)  Parsing
   0.24 ( 16.8%)  Sema::TypeChecking
   0.10 (  7.1%)  Instantiate Decls
   0.06 (  4.2%)  Preprocessor
   0.03 (  2.1%)  ASTConsumer HandleTranslationUnit
   0.03 (  2.1%)  (pool) other
\end{lstlisting}

读这张表的关键结论：\textbf{解析（Parsing）几乎吃掉七成}。预处理把 18 万行
token 吐出来其实很快（这里是 \texttt{4.2\%}），真正贵的是把这 18 万行
token 组装成 AST、再把模板实例化出可见声明的过程。一个 \texttt{widget.cpp}
展开后可能有\textbf{十几万个声明}进入符号表，绝大多数是 STL 与传递依赖送
的、你这行代码根本用不到的。这也解释了 \CD\ 为什么冷启动慢、编辑时首个
补全延迟高——它每改一次就要重走这条流水线。

\section{各工具在同一个模型上的不同取舍}
\label{sec:ch03-s2}

既然成本来自「打开即重做整个 TU」，不同的工具与语言服务器选择的只是
\textbf{怎么少做一点}。表 \ref{tab:ch03-cost} 把本册会点名的几类机制放在同一
个尺度上比较：一次编辑（改一个函数的返回类型）之后，各机制到底要重做多少工
作。这里的对比只看「重做的范围」，具体到 \CD\ 的实现细节留给第 5 章，
跨实现对比留给第 6 章。

\begin{longtable}{>{\raggedright\arraybackslash}p{0.17\textwidth}
                  >{\raggedright\arraybackslash}p{0.14\textwidth}
                  >{\raggedright\arraybackslash}p{0.50\textwidth}}
\caption{同一 TU 模型下，各机制「一次编辑后要重做多少工作」的对比}\label{tab:ch03-cost}\\
\toprule
机制 & 需要语义? & 改一处后要重做的范围 \\
\midrule
\endfirsthead
\toprule
机制 & 需要语义? & 改一处后要重做的范围 \\
\midrule
\endhead
\CD（完整 AST + 后台索引） & 是 &
  当前 TU 从改动点起重新解析其依赖子树，模板再实例化一遍；
  跨 TU 的引用/定义靠\textbf{后台索引}补齐，索引本身按文件增量重写 \\
\midrule
ccls（preamble 共享） & 是 &
  把公共头前缀烘成可共享的 preamble，命中前缀就不重复解析；
  改动落在 preamble 之后时只重解析尾段（细节见第 6 章） \\
\midrule
rust-analyzer（自研增量） & 是，但无预处理器 &
  \Lang{Rust} 没有 \texttt{\#include} 与宏展开倍率，源码即结构；
  它在持久化的 syntax tree 上做增量重算，改一个函数只脏化该函数所在节点。
  本册只把它当\textbf{对照组}：它快，是因为语言没有 TU 展开这一层 \\
\midrule
Tree-sitter（纯语法） & 否 &
  只按文法把文本切成 CST，毫秒级、无需编译命令，改一处只重 parse 可见的
  那一小段。代价是\textbf{没有语义}：它不知道 \texttt{push\_back} 属于
  \texttt{std::vector}，因此\textbf{跨文件跳转与真实 find-references 做不了} \\
\bottomrule
\end{longtable}

这张表澄清了一个常见误解：「换更快的语言服务器就能摆脱这个成本」。不能。
只要你要的是\textbf{语义级}功能（跳到定义、找引用、按类型补全、诊断），
就必须把改动文件所在的整个 TU 重新解析；能优化的只是\textbf{共享}（preamble、
modules）与\textbf{范围}（增量、脏区）两件事。而 Tree-sitter 之所以「快」，
恰恰是它放弃了语义那一侧——这也是为什么纯 Tree-sitter 编辑器能做语法高亮和
括号匹配，却做不出可靠的跨文件重构。

\begin{note}[title={索引不是数据库，它只是跨 TU 的补充}]
后台索引\textbf{不能}替代编译数据库，也\textbf{不能}替代当前文件的完整 AST。
数据库告诉工具「这个文件怎么编译」，当前 TU 的 AST 让工具「看懂眼前这一屏」。
索引解决的只有第三件事：「\textbf{这个符号在别的翻译单元里还有谁引用、在哪里
定义}」。所以索引没建完时，当前文件的补全与诊断照常工作，只有跨文件的
find-references、workspace symbol 搜索会残缺。别指望索引能让单文件解析变快。
\end{note}

\section{后台索引：clangd 怎么回答跨文件的问题}
\label{sec:ch03-s3}

\subsection*{它产出什么，存在哪里}

\CD\ 的 \Opt{background-index} 打开后，会在\textbf{编译数据库里每个 TU 的
输出目录附近}建一个 \texttt{.cache/clangd/index/} 目录，把每个文件解析出的
符号（定义、声明、引用位置）序列化成若干 \texttt{.idx} 分片存盘。清单
\ref{lst:ch03-index-dir} 是它的真实形态。

\begin{lstlisting}[style=shellstyle,caption={.cache/clangd/index 的内容与体积量级},label={lst:ch03-index-dir}]
ls -lh .cache/clangd/index | head
# widget.cpp.7A3F1C20D5B8E904.idx   512K
# socket.h.0C1D2E3F4A5B6C7D.idx     180K
# vector.9F8E7D6C5B4A3210.idx       1.2M
# ...
du -sh .cache/clangd/index
# 142M            <- a mid-size repo; a big one can reach several GB

# name = basename of the source + a content hash + .idx
# renaming or moving files leaves orphan shards until the dir is wiped
\end{lstlisting}

体积与 TU 数量和头文件复杂度成正比：STL 这类被反复解析的头会以「每个引用它
的 TU 都贡献一份」的方式膨胀，所以上面 \texttt{142M} 里相当一部分是
\texttt{vector}、\texttt{string} 之类的系统头副本。冷启动时 \CD\ 会遍历数据库
里所有 TU，把索引建满，这个过程 CPU 打满、可能持续几十秒到几分钟，正是第 1
节那条流水线的批量版。

\subsection*{为什么跨 TU 的 find-references 要等索引建完}

在当前文件里找一个函数的引用，靠当前 TU 的 AST 就够。但「谁在别的 .cpp 里
调用了我刚改签名的函数」这个问题，答案散落在成百上千个 TU 里，每个都要各自
解析一遍。索引做的事，就是\textbf{把这些跨 TU 的关系预先摊平成一张可查询的
表}。所以：

\begin{itemize}
  \item 索引未建完时，跨文件 find-references 只会返回「已经索引到的那部分
        文件」，结果是\textbf{不全}而非报错，容易让人误判「没人用」。
  \item 索引是\textbf{尽力而为的补充}，可能陈旧。改了一个头之后，引用它的
        其它 TU 要等后台重新解析并覆盖对应分片才更新；在此之前旧数据仍在。
  \item 结果太多时，用 \Opt{limit-results}\Bk（默认几十到几百）截断
        workspace 搜索与 find-references 的返回条数，避免一次查询把 UI 卡死；
        只想少占 CPU，可以 \Opt{background-index}\Bk 与
        \Opt{background-index-priority=low}\Bk 组合，把索引线程降到最低优先级，
        让前台编辑永远优先。
\end{itemize}

\subsection*{大仓库：把索引搬到外面}

当工程大到本地重建一次索引要几十分钟，正确做法是\textbf{不索引全部}，或
\textbf{预先在 CI 上索引一次再分发}。两条路：

\begin{itemize}
  \item \textbf{外部静态索引}：CI 上跑一遍 \CD（或用
        \texttt{clangd-indexing-tool} 离线生成），把 \texttt{.idx} 汇总成一个
        \texttt{.index} 文件随仓库或产物发布；本地 \CD\ 用
        \Opt{index-file}\Bk 加载它，只补索引自己真正改动的部分，冷启动从
        「重建全部」降到「加载一个文件」。
  \item \textbf{索引服务}：monorepo 用
        \texttt{clangd\allowbreak-index\allowbreak-server}\Bk 常驻，本地
        \CD\ 通过 gRPC 把跨 TU 查询委托给它，本机只做当前文件的解析。
\end{itemize}

清单 \ref{lst:ch03-clangd-yaml} 是 \texttt{.clangd} 里等价的配置写法
（字段含义见第 5 章）。命令行开关与配置文件二选一，配置文件更适合团队固化。

\begin{lstlisting}[style=yamlstyle,caption={.clangd 里对索引范围与外部索引的配置},label={lst:ch03-clangd-yaml}]
# Background: Skip turns off the index; only the open TU is parsed
Index:
  Background: Build
  External:
    # load a shard set prebuilt in CI
    File: build/monorepo.clangd-index
    Server:
      # or query a shared index service
      Url: grpc://index.internal:8000
\end{lstlisting}

\begin{warning}[title={把 \texttt{.cache/clangd} 提交进版本库的后果}]
\texttt{.cache/clangd/index} 里是\textbf{带绝对路径与内容哈希的二进制分片}，
换一台机器、换一个人、甚至把仓库换个目录，哈希就对不上，旧分片全成孤儿。
把它们提交进版本库会带来：仓库体积暴涨（动辄数 GB）、diff 无意义冲突、
CI 缓存命中率归零、别人拿到的索引在你的机器上根本用不了。正确做法是把
\texttt{.cache/} 写进 \texttt{.gitignore}，需要预索引就用外部
\texttt{.index} 单文件（上一小节）随发布产物走，绝不要提交原始分片目录。
\end{warning}

\section{preamble 与模块：共享头前缀的两条路}
\label{sec:ch03-s4}

第 1 节的倍率说到底是「每个 TU 都从零把同一批公共头重新解析一遍」。既然
\texttt{<vector>}、\texttt{<string>}、项目公共头的展开结果对同项目 99\% 的
文件都一样，能不能\textbf{只解析一次，大家共用}？这就是 preamble 与 modules
两条思路的共同动机。

\textbf{preamble 共享}的做法：把源文件里到第一个「不稳定」内容为止的那段前缀
（通常就是一串 \texttt{\#include} 加上它们的展开）烘成一块可复用的解析结果，
后面真正的业务代码只解析这一小段。ccls 把这条做到了极致（第 6 章展开），
\CD\ 内部对同一会话里反复打开的文件也会复用已解析 AST 的一部分。它的局限很
直白：前缀里只要有一个宏、一行注释位置变化让哈希失效，整块 preamble 就要
重烘。

\textbf{Clang Modules} 想用更彻底的方式：把每个头预编译成一个二进制模块
（\texttt{.pcm}），用 \Opt{fmodules} 让 \texttt{import} 直接加载编译好的
语义而不是重新解析文本。听起来是终极答案，但它始终没成为 \CPP\ 工具链的
标配，原因是硬的：

\begin{itemize}
  \item \textbf{宏污染与可见性}。模块系统要求把宏严格隔离（\texttt{export
        macro} 等），而现实代码库里的头互相依赖宏、宏又随
        \texttt{-D}、平台、构建类型变化，任何一个宏集合不同就触发重编或
        「module not found / build failed」。
  \item \textbf{可复现性}。\texttt{.pcm} 的哈希绑定了编译器版本、命令行、
        甚至绝对路径，团队里一人升级 clang、一人改了 \texttt{-D}，模块就整体
        失效，缓存变成黑洞。
  \item \textbf{ABI 与分发}。模块没有稳定的跨版本 ABI，\texttt{.pcm}
        无法像 \texttt{.o}/\texttt{.a} 那样在机器间可靠分发，网络文件系统上的
        模块缓存一致性问题更是长期痛点。
\end{itemize}

\cpp{20} 的具名模块（named modules）试图从语言层面重做这件事，但工具链对它
的支持目前仍远不如对 header 的支持成熟，索引与编译数据库都还在追赶。本册把
它当趋势而非现成的提速手段：真正能立刻用的缓解，仍是下一节的范围控制与第 5
节的依赖扇出治理。

\section{增量与失效：改一个公共头会波及多少 TU}
\label{sec:ch03-s5}

编辑单个文件的成本由第 1 节决定；但\textbf{维护}成本由另一件事决定：一个
公共头被改，会连锁让多少个 TU 失效、进而让索引重建多少。这决定了「改一行
头文件，CI 和所有人的 clangd 卡多久」。

要估计这个「依赖扇出」，用 include 图就够了。清单
\ref{lst:ch03-scandeps} 用 \texttt{clang-scan-deps} 从编译数据库抽出每个 TU
真正依赖哪些头（工具细节与用法在第 14 章，这里只看它给出的依赖信息如何服务于
成本估计）。

\begin{lstlisting}[style=shellstyle,caption={用 clang-scan-deps 抽出每个 TU 的头依赖},label={lst:ch03-scandeps}]
clang-scan-deps -format=dump -j=8 \
  -compilation-database=build/compile_commands.json \
  > deps.txt

# 反查：某个公共头被多少 TU 直接依赖（扇出的下界）
grep -c "include/net/socket.h" deps.txt
# 1271        <- 改 socket.h，这 1271 个 TU 的语义都可能变

# 也可以只从 -H 的树粗估同一件事
clang++ -std=c++20 -Iinclude -fsyntax-only src/net.cpp -H 2>&1 \
  | grep -c "net/socket.h"
\end{lstlisting}

\texttt{1271} 这个数就是「改 socket.h 后，索引要为多少个 TU 重解析」的估计
下界（这还只是直接 include；传递 include 另算，但对判断量级足够）。它解释了
为什么一个被全工程 include 的巨型公共头如此危险：它的任何改动都会让后台索引
把上千个 TU 重走一遍第 1 节那条流水线。

\texttt{\#include} 的\textbf{顺序}也影响成本，但方式不同：\CD\ 在解析头文件
时，需要一个「代表性的 TU」来推断这个头该怎么编译。它的判定规则是——优先找与
头同名的源文件（\texttt{foo.h} 配 \texttt{foo.cpp}），这叫 main include；
找不到就从 include 图里挑一个把它当非 main include 的 TU。这条链能走通，取决
于你的 include 组织与编译数据库是否自洽，具体规则留给第 5 章。这里只需要知道：
\textbf{把公共头塞进几乎所有源文件的传递依赖里，会让这个判定更贵也更不稳}。

\section{预算与缓解：症状到手段的对照}
\label{sec:ch03-s6}

把前五节的机制落到运维层面，就是一张「症状 → 根因 → 缓解」的对照表。表
\ref{tab:ch03-budget} 按你实际会看到的症状组织，每条缓解都能回溯到本章某节
的原理。缓解里涉及的 \texttt{.clangd} 字段与 \Tidy\ 联动，字段速查在附录 A，
IWYU 与依赖治理在第 15 章。

\begin{longtable}{>{\raggedright\arraybackslash}p{0.19\textwidth}
                  >{\raggedright\arraybackslash}p{0.26\textwidth}
                  >{\raggedright\arraybackslash}p{0.36\textwidth}}
\caption{C++ 语言服务器的成本症状、根因与缓解}\label{tab:ch03-budget}\\
\toprule
症状 & 根因 & 缓解手段 \\
\midrule
\endfirsthead
\toprule
症状 & 根因 & 缓解手段 \\
\midrule
\endhead
打开文件后 CPU 打满数十秒 &
  冷启动把数据库里所有 TU 走一遍解析+索引（第 1、3 节） &
  编辑期关掉后台索引：\Opt{background-index=false}\Bk
  或 \texttt{Index.Background} 设 \texttt{Skip}，只在 CI 建；
  大仓库改用预生成的外部 \texttt{.index} \\
\midrule
内存吃到数 GB 不退 &
  每个已打开 TU 的完整 AST 常驻，公共头被逐 TU 复制 &
  少开标签页；缩小 \texttt{-I} 范围；
  拆巨型 TU、减传递 include（第 5 节）；
  限制索引并发 \Opt{limit-results} \\
\midrule
补全首秒无结果 / 延迟高 &
  当前 TU 尚未解析完，模板未实例化 &
  等首解析完成再重度使用；把巨型头换成前置声明减少解析量；
  用 preamble 型服务器（ccls，第 6 章） \\
\midrule
find-references 结果不全 &
  后台索引还没建到那些 TU（第 3 节） &
  等索引建完；把无关目录从数据库剔除；
  用外部索引一次性覆盖全仓 \\
\midrule
改一个公共头后全工程卡 &
  依赖扇出巨大，索引整体失效重建（第 5 节） &
  用 \texttt{clang-\allowbreak scan-\allowbreak deps} 监控扇出并拆分公共头；
  把稳定接口抽成独立小头；
  \Opt{background-\allowbreak index-\allowbreak priority=low} 让前台优先 \\
\bottomrule
\end{longtable}

缓解里有一处需要特别小心的取舍：\textbf{前置声明}能显著缩小 TU 的解析量（不
include 那个头，就不解析它的整棵子树），听起来是免费的提速。但第 15 章的
IWYU（include-what-you-use）会从\textbf{另一个}角度建议你把缺失的直接 include
补回来。两者目标不同：前置声明为解析省钱，IWYU 为「每个符号来自哪」的可追溯性
和「改实现头不触发重编」的构建稳定性服务。别把它们当成互相矛盾的建议乱用；
先明确你此刻优化的是\textbf{工具解析成本}还是\textbf{构建成本}。

\begin{bestpractice}[title={大仓库先只索引自己的 src，第三方头按需解析}]
面对一个带巨型第三方依赖（Boost、protobuf、自研 SDK）的仓库，不要一上来就
全量索引。分三步压缩冷启动与常驻内存：其一，把编译数据库的范围限制到
\texttt{src/}，让 \CD\ 只对\textbf{你自己的} TU 建后台索引；其二，对第三方头
不要放进 include 图的常解析路径，需要时用外部 \texttt{.index} 或
\texttt{Index.External} 单独喂给它们，本地只解析当前真正打开的 TU；其三，把
被全工程 include 的稳定接口头拆小，用 \texttt{clang-scan-deps} 把扇出最高的
头排进重构队列。这样第一次打开工程等待的是「你的代码」，而不是「整个依赖
树」。
\end{bestpractice}

\section*{本章小结}
本章把「\CPP\ 工具为什么慢」还原到一个不可回避的事实：打开一个文件等于把
它的整个翻译单元重新预处理并解析一遍，而 include 倍率（本章实测约
\texttt{450} 倍）把这个动作放大得很贵。各工具只是在同一个模型上选不同的共享
与范围策略；索引解决的是跨 TU 的关系，不能替当前文件减负。preamble 与
modules 是共享头前缀的两条路，后者因宏、可复现性、ABI 分发而没能普及。
最后，把症状对到缓解，记住前置声明与 IWYU 优化的是不同成本。下一章开始进入
Part II：LSP 协议本身——这些昂贵的解析结果，究竟通过哪些请求、以什么时序
送达编辑器。

% ---- frag_ch04.tex ----
% ==================== 第 4 章：LSP 协议 ====================
\part{语言服务器}
\chapter{LSP 协议：生命周期、能力协商与消息模型}\label{ch:lsp-protocol}

Language Server Protocol（LSP）由微软在 2016 年随 VS Code 一并公布，
它解决的问题只有一个：把「编辑器对代码语义的提问」和「分析器对代码语义的回答」
拆成两个可以独立演进的进程。在此之前，每款编辑器都要为每门语言重写一遍
「跳转到定义」；在此之后，工作量变成 $M+N$ 而非 $M\times N$。
LSP 的全部内容都建立在 JSON-RPC 2.0 之上，但它在外面包了一层自己的帧格式，
并规定了严格的生命周期状态机与能力协商机制。
本章按「线格式 $\to$ 生命周期 $\to$ 文档同步 $\to$ 语义能力 $\to$ 位置模型 $\to$ 版本边界」
的顺序展开，目标是让你能直接读懂抓包日志、自己实现或调试一个客户端。
\CPP{} 语言服务器的具体实现与调优见第 5 章。

% --------------------------------------------------------------------
\section{线格式与消息三分类}\label{sec:ch04-s1}

LSP 消息通过 stdio（也可为 socket/pipe）传输，每一帧由三部分组成：
一条或多条头部、一个空行（\texttt{\textbackslash r\textbackslash n\textbackslash r\textbackslash n}
中的第二个 CRLF）、以及 UTF-8 编码的 JSON 正文。
头部至少包含 \texttt{Content-Length}，表示正文的\textbf{字节数}；
\texttt{Content-Type} 可省略，缺省即 \texttt{application/vscode-jsonrpc; charset=utf-8}，
但 \CD\ 总会把它发出来。清单~\ref{lst:ch04-wire} 是一段真实形态的抓包节选：
客户端发一个 \texttt{textDocument/hover}\Bk{} 请求，服务端回响应，
随后又主动推了一条通知。

\begin{lstlisting}[style=jsonstyle,caption={一帧一帧读：请求、响应与通知的原始报文},label={lst:ch04-wire}]
Content-Length: 208
Content-Type: application/vscode-jsonrpc; charset=utf-8

{
  "jsonrpc": "2.0",
  "id": 1,
  "method": "textDocument/hover",
  "params": {
    "textDocument": { "uri": "file:///d%3A/proj/src/main.cpp" },
    "position": { "line": 41, "character": 8 }
  }
}
Content-Length: 96
Content-Type: application/vscode-jsonrpc; charset=utf-8

{
  "jsonrpc": "2.0",
  "id": 1,
  "result": {
    "contents": {
      "kind": "markdown",
      "value": "int compute(int n)"
    }
  }
}
Content-Length: 124
Content-Type: application/vscode-jsonrpc; charset=utf-8

{
  "jsonrpc": "2.0",
  "method": "window/logMessage",
  "params": {
    "type": 3,
    "message": "preamble rebuilt in 340ms"
  }
}
\end{lstlisting}

逐段拆解清单~\ref{lst:ch04-wire}：
第一帧带 \texttt{id}: 1 与 \texttt{method}，是\textbf{请求}（Request）；
第二帧带同样的 \texttt{id} 与 \texttt{result}，是对它的\textbf{响应}（Response）；
第三帧有 \texttt{method} 却没有 \texttt{id}，是服务端主动发来的\textbf{通知}
（Notification），不需要也不允许回包。三种消息的判别规则汇总在表~\ref{tab:ch04-msg}。

\begin{longtable}{>{\raggedright\arraybackslash}p{0.13\textwidth}
                  >{\raggedright\arraybackslash}p{0.26\textwidth}
                  >{\raggedright\arraybackslash}p{0.19\textwidth}
                  >{\raggedright\arraybackslash}p{0.22\textwidth}}
\caption{LSP 的三类消息与判别规则}\label{tab:ch04-msg}\\
\toprule
种类 & 必含字段 & 期待回复? & 典型例子 \\
\midrule
\endfirsthead
\toprule
种类 & 必含字段 & 期待回复? & 典型例子 \\
\midrule
\endhead
请求（Request） & \texttt{id}、\texttt{method}、\texttt{params} & 是：同 \texttt{id} 的 \texttt{result} 或 \texttt{error} & \texttt{initialize}、\texttt{shutdown}、\texttt{hover} \\
\midrule
响应（Response） & \texttt{id}，加 \texttt{result} 或 \texttt{error} 之一 & ---（它本身就是回复） & \texttt{error.code}：\texttt{MethodNotFound}（$-32601$） \\
\midrule
通知（Notification） & \texttt{method}、\texttt{params}，无 \texttt{id} & 否 & \texttt{initialized}、\texttt{didOpen}、\texttt{logMessage} \\
\bottomrule
\end{longtable}

\texttt{id} 可以是整数或字符串，由发起方分配、响应方原样回填，允许多个请求在途并发；
客户端必须按 \texttt{id} 配对，而不能按到达顺序假设一一对应。
对不再需要的请求，用取消通知（\texttt{\$/cancelRequest}\Bk{}，参数
\texttt{\{id: 1\}}）告知服务端尽快放弃计算；服务端可能早已完成，此时忽略取消即可。
另一个容易踩的兼容性问题：JSON-RPC 2.0 允许把多个请求打包成数组批量发送，
而 LSP \textbf{明确禁止}批量数组——看到服务端对 \texttt{[{...}]} 直接断开连接，
多半是客户端实现了 JSON-RPC 却忽略了这条裁剪。

\begin{note}[title={Content-Length 数的是 UTF-8 字节，不是字符}]
帧头长度按 UTF-8 字节计数。若客户端用「字符串长度」算长度（常见于按 UTF-16
码元计数再乘错系数的实现），正文里一旦混入中文注释或多字节符号，字节数就会对不上：
服务端会多读或少读若干字节，轻则解析失败，重则后续每一帧的边界全部错位、
协议流彻底报废。正确做法是对序列化后的字节串求 \texttt{len(bytes)}。
\end{note}

% --------------------------------------------------------------------
\section{生命周期与能力协商}\label{sec:ch04-s2}

LSP 把一次会话分成五个阶段，状态机固定如下：
\begin{enumerate}
  \item \texttt{initialize}（请求）：客户端上报自身能力与项目上下文；
  \item 服务端回复 \texttt{capabilities}：声明自己支持什么、怎么支持；
  \item \texttt{initialized}（通知）：客户端确认协商结束，此后才允许正常流量；
  \item 正常工作阶段：文档同步、语义请求、诊断推送……
  \item \texttt{shutdown}（请求）$\to$ \texttt{exit}（通知）：前者让服务端释放
        状态并回 \texttt{null}，后者要求进程立即退出——退出码有约定，
        未调用 \texttt{shutdown} 就收到 \texttt{exit} 时应返回非零，方便外层
        看门狗判断这是异常收场。
\end{enumerate}

\texttt{initialize} 的参数里，\CPP{} 项目中最要紧的四项是：
\texttt{processId}（父进程号，服务端可据此在父进程意外死亡时自杀，避免僵尸）、
\texttt{rootUri}（工作区根目录，URI 形式；早年的 \texttt{rootPath} 已废弃）、
\texttt{capabilities}（客户端能力大清单）、
\texttt{workspaceFolders}（多根工作区，每项含 \texttt{uri} 与 \texttt{name}）。
清单~\ref{lst:ch04-init-req} 给一份精简但字段真实的请求。

\begin{lstlisting}[style=jsonstyle,caption={initialize 请求（客户端 $\to$ 服务端）节选},label={lst:ch04-init-req}]
{
  "jsonrpc": "2.0",
  "id": 0,
  "method": "initialize",
  "params": {
    "processId": 18420,
    "clientInfo": { "name": "my-editor", "version": "0.3.1" },
    "rootUri": "file:///d%3A/proj",
    "workspaceFolders": [
      { "uri": "file:///d%3A/proj", "name": "proj" }
    ],
    "capabilities": {
      "general": { "positionEncodings": ["utf16", "utf8"] },
      "workspace": {
        "workspaceFolders": true,
        "configuration": true,
        "didChangeWatchedFiles": { "dynamicRegistration": true },
        "didChangeConfiguration": {
          "dynamicRegistration": true
        }
      },
      "textDocument": {
        "publishDiagnostics": {
          "relatedInformation": true,
          "versionSupport": true
        },
        "completion": {
          "completionItem": {
            "snippetSupport": true,
            "documentationFormat": ["markdown", "plaintext"],
            "resolveSupport": {
              "properties": ["detail", "documentation"]
            }
          }
        },
        "semanticTokens": { "requests": { "full": true } },
        "hover": { "contentFormat": ["markdown", "plaintext"] }
      }
    }
  }
}
\end{lstlisting}

\begin{lstlisting}[style=jsonstyle,caption={initialize 响应（服务端 $\to$ 客户端）节选},label={lst:ch04-init-resp}]
{
  "jsonrpc": "2.0",
  "id": 0,
  "result": {
    "capabilities": {
      "textDocumentSync": {
        "openClose": true,
        "change": 2,
        "willSave": true,
        "willSaveWaitUntil": false,
        "save": { "includeText": false }
      },
      "completionProvider": {
        "triggerCharacters": [".", ":", ">"],
        "resolveProvider": true
      },
      "hoverProvider": true,
      "signatureHelpProvider": {
        "triggerCharacters": ["(", ","]
      },
      "definitionProvider": true,
      "declarationProvider": true,
      "typeDefinitionProvider": true,
      "implementationProvider": true,
      "referencesProvider": true,
      "renameProvider": { "prepareProvider": true },
      "documentSymbolProvider": true,
      "workspaceSymbolProvider": true,
      "codeActionProvider": true,
      "documentFormattingProvider": true,
      "inlayHintProvider": { "resolveProvider": false },
      "positionEncoding": "utf-16",
      "semanticTokensProvider": {
        "legend": {
          "tokenTypes": ["namespace", "type", "function", "macro"],
          "tokenModifiers": [
            "declaration", "static", "deprecated"
          ]
        },
        "full": true,
        "range": true
      },
      "workspace": {
        "workspaceFolders": {
          "supported": true,
          "changeNotifications": true
        }
      }
    }
  }
}
\end{lstlisting}

响应侧最重要的字段是 \texttt{textDocumentSync}（文档同步策略，
数字 1 表示 full、2 表示 incremental，也可以像清单这样写成对象展开）与
一系列 \texttt{*Provider} 能力。注意 \texttt{semanticTokensProvider.legend}：
\texttt{tokenTypes} 与 \texttt{tokenModifiers} 的\textbf{下标顺序}就是后面
token 流里引用的编号表，协商阶段定好后整份文档都用它解码（见第~\ref{sec:ch04-s4} 节）。

\begin{guideline}[title={能力协商决定后面哪些请求能发}]
\begin{itemize}
  \item 服务端没声明 \texttt{implementationProvider}，就不要发
        \texttt{textDocument/implementation}\Bk{}——得到的是
        \texttt{MethodNotFound}（$-32601$），而行为良好的客户端应直接把
        「查找实现」菜单项禁用，根本不发请求。
  \item 客户端没声明 \texttt{snippetSupport}，服务端回传的
        \texttt{\$\{1:xxx\}} 占位符会被当成字面文本插入，编辑器里
        看到一堆大括号。
  \item 没声明 \texttt{positionEncoding} 的客户端，默认被按 UTF-16 对待；
        一旦源码含非 ASCII 字符，偏移就会错（第~\ref{sec:ch04-s5} 节）。
  \item 一切「支持运行时再注册」的能力走
        \texttt{client/registerCapability} 通知动态协商，静态清单只是保底。
\end{itemize}
\end{guideline}

% --------------------------------------------------------------------
\section{文档同步：内存副本与版本号}\label{sec:ch04-s3}

服务端看不到你的磁盘，也看不到你的编辑器缓冲区——它只认协议里同步过来的文本。
文档在服务器内存中的生命周期由五个通知构成：
\texttt{textDocument/didOpen}（打开，带全文与 \texttt{languageId}，版本号从 1 起）、
\texttt{textDocument/didChange}（修改）、
\texttt{willSaveWaitUntil}（保存前，服务端可同步返回编辑，客户端等待并应用，
因此它是阻塞调用，重格式化类工作不要放这里）、
\texttt{didSave}（保存，声明了 \texttt{includeText} 时再带全文）、
\texttt{didClose}（关闭，通知服务端回收内存；\CD\ 会延迟回收以容忍切标签页）。

修改的同步方式由 \texttt{textDocumentSync.change} 决定：
\textbf{full}（1）模式每次把整篇文档重发，实现简单但大文件昂贵；
\textbf{incremental}（2）模式只发 \texttt{Range}（起止位置）加替换文本，
清单~\ref{lst:ch04-didchange} 是典型报文。注意 \texttt{version}：
每次 \texttt{didChange} 必须单调递增，服务端用它判断自己内存里的副本是哪一版。

\begin{lstlisting}[style=jsonstyle,caption={incremental 同步下的 didChange 通知},label={lst:ch04-didchange}]
{
  "jsonrpc": "2.0",
  "method": "textDocument/didChange",
  "params": {
    "textDocument": {
      "uri": "file:///d%3A/proj/src/main.cpp",
      "version": 12
    },
    "contentChanges": [
      {
        "range": {
          "start": { "line": 41, "character": 4 },
          "end":   { "line": 41, "character": 10 }
        },
        "text": "std::ranges::sort"
      }
    ]
  }
}
\end{lstlisting}

版本号不只是记账工具。语义类请求（补全、重命名、代码动作）返回的
\texttt{WorkspaceEdit} 是基于「回答那一刻」的 AST 计算出来的：
若客户端拿着它回放到一份已经被后续 \texttt{didChange} 改过的文档上，
插入点、区间跨度全部失效。规范要求在
\texttt{WorkspaceEdit.changes}/\texttt{documentChanges} 里带上目标文档的
\texttt{version}（\texttt{TextDocumentEdit} 的 \texttt{textDocument} 字段
为 \texttt{VersionedDocumentUri\{uri, version\}}），
客户端发现版本对不上必须\textbf{拒绝应用}。实战中「按了 code action
却报编辑失败」的根因几乎都是缓冲区版本落后于请求发出时的版本。

% --------------------------------------------------------------------
\section{语义能力清单}\label{sec:ch04-s4}

LSP 把「对代码的提问」切成一组标准方法。理解它们最快的方式是把它映射回
\CPP{} 编译器前端：一次完整的 \texttt{clang++ -fsyntax-only} 内部有预处理、
解析、\texttt{clang::Sema} 语义分析和 AST；语言服务器的工作就是把这些
一次性过程拆成可反复询问的持久服务。

\subsection{提问/回答类：光标换到哪就答到哪}
\texttt{completion}（光标处能写什么；服务端遍历作用域链加候选排序，
列表太长时用 \texttt{completionItem/resolve} 二次惰性补齐
\texttt{detail}/\texttt{documentation}）、
\texttt{hover}（这个标识符的声明长什么样）、
\texttt{signatureHelp}（正在调哪个重载，来自 \texttt{Sema} 的重载决议候选集）、
\texttt{definition} 与 \texttt{declaration}（\texttt{DeclRefExpr} 指向的
\texttt{Decl} 在哪）、\texttt{typeDefinition}（顺着 \texttt{TypedefType}
剥到本质类型）、\texttt{implementation}（虚函数在派生类里的全部覆写）、
\texttt{references}（反向引用表，可带 \texttt{context.includeDeclaration}）、
\texttt{documentSymbol}（当前文件层级化大纲，返回
\texttt{DocumentSymbol[]} 树）与 \texttt{workspaceSymbol}
（全项目按名字模糊搜，返回平铺 \texttt{SymbolInformation[]}，依赖索引）、
\texttt{prepareRename}（当前光标允不允许改名，可返回占位名）与
\texttt{rename}（跨文件替换，返回 \texttt{WorkspaceEdit}）、
\texttt{formatting}/\texttt{rangeFormatting}/\texttt{onTypeFormatting}
（整篇/选区/输入时的 \texttt{TextEdit[]}，与第 9 章 \Fmt{} 同一引擎）。

\subsection{诊断：push 与 pull 两种模型}
经典模型是 push：服务端在任何时刻发现问题就推
\texttt{textDocument/}\hspace{0pt}\texttt{publish}\hspace{0pt}\texttt{Diagnostics}，
参数含 \texttt{uri}、可选 \texttt{version}
和诊断数组；每条 \texttt{Diagnostic} 有 \texttt{range}、
\texttt{severity}（1 Error、2 Warning、3 Info、4 Hint）、
\texttt{source}（如 \texttt{clang} 或 \texttt{clang-tidy}）、\texttt{code}
与 \texttt{message}，以及可选 \texttt{relatedInformation}
（\texttt{\{location, message\}} 列表，用来表达「此处未初始化是因为第 12 行
构造函数没写它」这类跨位置理由）。
新模型是 pull：客户端显式调 \texttt{textDocument/diagnostic}\Bk{}，
服务端回 \texttt{DocumentDiagnosticReport}（\texttt{kind} 为
\texttt{full}/\texttt{unchanged}，附 \texttt{resultId} 供增量复用），
另有 \texttt{workspace/diagnostic} 和 \texttt{diagnostic/refresh}。
两种模型的取舍见下方对比框。

\begin{compare}[title={push 与 pull 诊断的取舍}]
\begin{itemize}
  \item push：实现简单、编辑器标红即时；但后台线程随时发包，大项目里
        与用户输入抢资源，客户端很难做「只在保存后检查」这类节流。
  \item pull：请求发起权在客户端，可按可见文件/前台优先级调度，
        \texttt{resultId} 支持无变化复用；代价是需要缓存与刷新通知配合，
        实现复杂度更高。
  \item 现实是过渡期并存：\CD\ 目前仍以 push 为主，pull 路径在 3.17 起
        作为可选能力存在，第 15 章讨论 sanitizer/分析器时再划一次边界。
\end{itemize}
\end{compare}

\subsection{token 流与内联提示}
\texttt{semanticTokens} 解决「纯语法高亮不够」的问题：语法着色只能看到关键字
和字符串，分不清 \texttt{Foo f;} 里 \texttt{Foo} 是类型还是变量
（需要知道有没有 \texttt{typedef int Foo}），也分不清宏展开后的实体。
语义 token 来自 AST：服务端先按文档切出 token，标上类型
（对照 legend 的下标）与修饰位（\texttt{declaration}、\texttt{static} 等）。
返回的 \texttt{data} 是扁平整数数组，五个一组：
\texttt{deltaLine}、\texttt{deltaStartChar}（相对\textbf{上一个} token 的行偏移与
列偏移，列按 UTF-16 码元）、\texttt{length}、\texttt{tokenType}、
\texttt{tokenModifiers} 位掩码。相对编码让「前面插入一行」的增量
\texttt{semanticTokens/full/delta} 只需重发受影响部分。
\texttt{inlayHint} 则返回要「画在源码空隙里」的灰色提示（推断类型、
具名实参、循环边界），\texttt{label} 可拆多段并带
\texttt{paddingLeft}/\texttt{paddingRight}。

\subsection{方法到实现：一张对照表}
表~\ref{tab:ch04-methods} 把每个方法对应到 \CD\ 内部入口与等价的「手工 \CPP{} 操作」，
后者指如果你不用语言服务器、直接拿 \Lib{clang} 库手写会做什么。

\begin{longtable}{>{\raggedright\arraybackslash}p{0.26\textwidth}
                  >{\raggedright\arraybackslash}p{0.27\textwidth}
                  >{\raggedright\arraybackslash}p{0.27\textwidth}}
\caption{LSP 方法 $\to$ clangd 内部入口 $\to$ 等价的手工 \CPP{} 操作}\label{tab:ch04-methods}\\
\toprule
LSP 方法 & \CD\ 内部入口 & 等价的手工操作 \\
\midrule
\endfirsthead
\toprule
LSP 方法 & \CD\ 内部入口 & 等价的手工操作 \\
\midrule
\endhead
\texttt{textDocument/}\hspace{0pt}\texttt{completion} & \texttt{CodeComplete}，走 \texttt{clang::Sema} 补全钩子 & 在光标处插入伪代码后重新解析，遍历作用域链收集候选 \\
\midrule
\texttt{hover}/\texttt{definition} & \texttt{ASTReader} 查 \texttt{DeclRefExpr} $\to$ \texttt{Decl} & 拿到 AST 后查引用指向的声明并格式化类型 \\
\midrule
\texttt{references} & \texttt{XRef} 表加索引（\texttt{SymbolIndex}） & \texttt{RecursiveAST\hspace{0pt}Visitor} 全树扫引用，跨文件需项目索引 \\
\midrule
\texttt{workspaceSymbol} & \texttt{BackgroundIndex} 的 slab 查询 & 对全项目符号表做名字检索 \\
\midrule
\texttt{publishDiag-}\hspace{0pt}\texttt{nostics} & \texttt{Diagnostic}\hspace{0pt}\texttt{Consumer} 收集 \texttt{StoredDiagnostic} & 编译一遍 \texttt{-fsyntax-only} 抄警告 \\
\midrule
\texttt{rename} & \texttt{collect}\hspace{0pt}\texttt{Replacements}（引用集驱动） & 对每个引用位置做 \texttt{Lexer} 级替换 \\
\midrule
\texttt{formatting} & \Fmt{} 的 \texttt{format::reformat} & 直接调用格式化库，见第 9 章 \\
\midrule
\texttt{semanticTokens/}\hspace{0pt}\texttt{full} & \texttt{SemanticTokens} 模块按 AST 分类 token & \texttt{Tokenizer} 切词后用 \texttt{Sema} 结果逐个定性 \\
\midrule
\texttt{codeAction} & 按诊断 \texttt{code} 匹配 quick fix（含 include 整理） & 依据编译器提示手工改写源码 \\
\bottomrule
\end{longtable}

\subsection{工作区级能力}
多根工作区由 \texttt{workspaceFolders} 支持：客户端在
\texttt{initialize} 里给出初始列表，运行中通过
\texttt{workspace/}\Bk\texttt{didChangeWorkspaceFolders} 通知增删；
服务端为每个文件夹分别定位编译数据库（第 2 章）。
长时间操作（\CD\ 的后台索引）用
\texttt{window/}\Bk\texttt{workDoneProgress} 系列请求建进度条，
\texttt{\$/progress} 通知推进度。

% --------------------------------------------------------------------
\section{位置模型：UTF-16 码元计数的陷阱}\label{sec:ch04-s5}

LSP 的 \texttt{Position} 由 \texttt{line} 与 \texttt{character} 组成，
都从 0 起算；\texttt{Range} 是 \texttt{\{start, end\}} 两个位置。
陷阱在 \texttt{character}：规范默认按 \textbf{UTF-16 码元}计数——这是历史包袱，
因为最初的客户端是 JavaScript 写的，其字符串天然按 UTF-16 存。
而 \CPP{} 源码在磁盘上是 UTF-8 字节流、在 \Lib{clang} 内部按字节定位
（\texttt{SourceLocation}），三者互不相等。看清单~\ref{lst:ch04-utf16} 的例子。

\begin{lstlisting}[style=cppstyle,caption={含中文的行里，三种计数法给出的列号各不相同},label={lst:ch04-utf16}]
// 源文件为 UTF-8；第 2 行内容（含前导 4 空格）如下：
//     int 年龄 = 30;
// 求 '=' 所在列：
//   UTF-8 字节: 4 + 4 + 6 + 1 = 15  (每个汉字 3 字节)
//   UTF-16 码元: 4 + 4 + 2 + 1 = 11  (BMP 汉字 1 码元)
//   码点:       与 UTF-16 相同，除非遇到 emoji（2 码元）
int main() {
    int 年龄 = 30;   // 光标停在 '=' 上
    return 年龄;
}
\end{lstlisting}

同一行同一位置，UTF-16 计数是 \texttt{character}: 11，按 UTF-8 字节计数则是 15，
差了 $2\times2$（两个汉字各多算 2 字节）。若客户端按字节上报位置，
\texttt{hover} 会问在「龄」字里、\texttt{rename} 会改错位置；
反过来，服务端按 UTF-16 回 \texttt{Range}，客户端按码点解释，
含 emoji（BMP 外，占 2 个 UTF-16 码元）的注释之后所有位置集体左移。
3.17 的解法是把它变成可协商项：客户端在
\texttt{initialize} 里用 \texttt{general.positionEncoding}
声明支持的编码（\texttt{utf-8}、\texttt{utf-16}、\texttt{utf-32}），
服务端在响应里\textbf{选定并回填}一个（\CD\ 目前回 \texttt{utf-16}，
见清单~\ref{lst:ch04-init-resp}）。三者的语义差异见表~\ref{tab:ch04-enc}。

\begin{longtable}{>{\raggedright\arraybackslash}p{0.14\textwidth}
                  >{\raggedright\arraybackslash}p{0.24\textwidth}
                  >{\raggedright\arraybackslash}p{0.21\textwidth}
                  >{\raggedright\arraybackslash}p{0.20\textwidth}}
\caption{三种位置编码对同一串字符的计数}\label{tab:ch04-enc}\\
\toprule
编码 & 计数单位 & 「中文emoji」的长度 & 换算成本 \\
\midrule
\endfirsthead
\toprule
编码 & 计数单位 & 「中文emoji」的长度 & 换算成本 \\
\midrule
\endhead
\texttt{utf-8} & 字节 & 4 汉字 12 字节加 emoji 4 字节，共 16 & \Lib{clang} 内部原生，直接可用 \\
\midrule
\texttt{utf-16} & UTF-16 码元 & 4 加 2，共 6 & LSP 默认；JS/Java 客户端零成本，其余要转换 \\
\midrule
\texttt{utf-32} & 码点 & 5 & 与 Python 直觉一致，转换代价最高 \\
\bottomrule
\end{longtable}

\begin{warning}[title={中文注释导致补全、悬停、改名集体错位}]
\begin{itemize}
  \item \textbf{症状}：纯英文文件一切正常；加了中文注释或中文标识符字符串后，
        hover 显示隔壁符号、\texttt{rename} 报错甚至改坏源码、高亮右移。
  \item \textbf{原因}：客户端按 UTF-8 字节或码点上报 \texttt{character}，
        而服务端按默认 UTF-16 解释；每个非 ASCII 字符都在后面埋了偏移差。
  \item \textbf{对策}：三选一，全项目统一——
        （a）在 \texttt{initialize.capabilities} 里声明并尊重
        \texttt{general.positionEncoding}；
        （b）客户端在边界处做字节 $\to$ 码元的完整转换（注意 BMP 外字符）；
        （c）万不得已才把含非 ASCII 的行临时换成占位符验证是不是编码问题。
        排查时可用下一章的 \Opt{check} 模式对照 \CD\ 眼中的位置。
\end{itemize}
\end{warning}

% --------------------------------------------------------------------
\section{版本演进与协议边界}\label{sec:ch04-s6}

3.17（2023 年春冻结）是目前的事实基线：pull 诊断（\texttt{textDocument/diagnostic}
与 \texttt{workspace/diagnostic}）、\texttt{semanticTokens/full/delta}、
\texttt{positionEncoding}、\texttt{inlayHint} 都已进入规范，
\CD\ 与各主流客户端在此之上互通。
4.0 演进方向上最重要的两件事，一是把 pull diagnostics 从可选提升为
主模型（push 退居兼容层），二是让 \texttt{semanticTokens/full/delta}
这类增量传输成为 token 流的标准消费方式，同时清理各方法里标了
deprecated 的旧字段。写新客户端时建议：3.17 全量实现，4.0 增量按能力开关接入。

最后划清边界——LSP \textbf{不管}以下三件事：
\texttt{initialize} 之后，「这个翻译单元该用什么命令行编译」是编译数据库的事（第 2 章）；
「源码怎么变成可执行文件」是构建系统的事；
「在断点上停下来看变量」是调试的事（DAP，第 7 章）；
CMake 目标的查询与执行则要靠客户端另接 CMake Tools 之类的通道（第 16 章）。
理解这条边界，才能解释为什么 \CD\ 明明「认识」你的代码，却对
「怎么把项目编出来」一无所知。下一章进入 \CD\ 的进程内部，
看它如何消费编译数据库、以及 \texttt{.clangd} 里每一个配置键到底改了什么。

\begin{bestpractice}[title={实现客户端时守住三条纪律}]
\begin{itemize}
  \item \textbf{版本纪律}：每次 \texttt{didChange} 严格 \texttt{version}+1，
        应用 \texttt{WorkspaceEdit} 前核对版本，不匹配就丢回重算，绝不硬套。
  \item \textbf{取消纪律}：补全、\texttt{workspaceSymbol}、诊断这类慢请求，
        用户一敲新字符就发 \texttt{\$/cancelRequest}；否则队列里积压的
        旧请求会让新结果晚于旧结果显示。
  \item \textbf{编码纪律}：启动时记录协商结果
        （\texttt{positionEncoding}），在唯一一处边界完成转换，
        内部一律用自己语言的索引，不让 UTF-16 概念泄漏到业务代码里。
\end{itemize}
\end{bestpractice}

% ---- frag_ch05.tex ----
\chapter{\CD\ 架构与 \texttt{.clangd} 配置}\label{ch:clangd}

\CD\ 是把 \Lib{clang} 前端包装成语言服务器的官方实现：它不是「另一个编译器」，
而是让编译器前端常驻、可被反复询问的那层胶水。也正因为如此，
它的很多怪脾气（要等 preamble、索引没建完跳转不全、配置改了要重启）
都能从架构里推出因果。本章先给出进程内部的地图，再逐键讲清
\texttt{.clangd} 配置，最后用一张排错表收尾。

% --------------------------------------------------------------------
\section{进程架构：一条串行队列加一座后台书库}\label{sec:ch05-s1}

一个 \CD\ 进程内部分成两大部分。前台是请求处理链：
\texttt{ClangdLSPServer}（协议对账：\texttt{initialize} 握手、参数解析、
回包编码）$\to$ \texttt{ClangdServer}（功能入口与并发管理）$\to$
\texttt{ASTWorker}（真正持有当前翻译单元的线程）。
后台是索引系统：\texttt{BackgroundIndex} 在空闲算力里解析项目里的其他文件，
把符号与交叉引用写进 \texttt{.clangd/}\hspace{0pt}\texttt{index/} 下的 slab 文件。

为什么几乎所有请求都要排队进 \texttt{ASTWorker}？
因为\textbf{翻译单元是单点可变状态}：一份 AST、一张宏表、一个已展开的
预处理上下文，\Lib{clang} 前端并不支持并发读改。用户边打字，\CD\ 边做三件事：
重建 preamble（全部 \texttt{\#/include} 部分，只在 include 列表变化时重做）、
重解析当前文件（有约两秒的去抖）、把结果喂给请求。
跨文件的能力（\texttt{workspaceSymbol}、全项目 \texttt{references}）则不碰 AST，
直接查后台索引。

\begin{guideline}[title={请求排队带来的三个可观察现象}]
\begin{itemize}
  \item 刚 \texttt{didChange} 完立刻按补全，回复慢——它在等 AST 重建；
        这不是死锁，是队列。
  \item \texttt{\$/cancelRequest}\Bk{} 只能取消「还没开始」的请求；
        已经进到 ASTWorker 手里的活儿会做完再丢弃结果。
  \item 后台索引期间 CPU 高、跳转结果不全，日志出现
        \texttt{building index} 行属正常，索引落盘后消失。
\end{itemize}
\end{guideline}

排查「为什么这个文件没有补全」不必启动编辑器：\Opt{check} 模式
把 \CD\ 变成命令行自检器，对单个文件走一遍完整流水线并打印每一步；
\Opt{log}=\texttt{verbose} 则输出平时丢掉的细节（清单~\ref{lst:ch05-check}）。

\begin{lstlisting}[style=shellstyle,caption={用命令行触发一次完整自检},label={lst:ch05-check}]
:: 对单个文件走一遍流水线，不启动编辑器
clangd --check=src/render.cpp --compile-commands-dir=build

:: 排错时再加上详细日志，落盘慢慢看
clangd --check=src/render.cpp --log=verbose > clangd.log 2>&1
\end{lstlisting}

一份典型的 \Opt{check} 输出会把流水线逐步摊开（节选，见清单~\ref{lst:ch05-checkout}）。

\begin{lstlisting}[style=shellstyle,caption={check 与 verbose 的真实输出节选},label={lst:ch05-checkout}]
I[09:47:03.112] clangd version 17.0.6
I[09:47:03.141] pid: 21484
I[09:47:03.203] Working directory: D:/proj
I[09:47:03.412] Loaded compilation database from
    D:/proj/build/compile_commands.json
I[09:47:03.455] Compile command from CDB is: [D:/proj/src]
    g++ -std=c++17 -Iinclude -c render.cpp
V[09:47:03.910] Building preamble...
I[09:47:05.023] Building AST...
V[09:47:05.771] Indexing symbols...
I[09:47:06.314] Testing all features...
I[09:47:06.889] All checks completed, 0 errors
\end{lstlisting}

日志前缀 \texttt{I}/\texttt{V}/\texttt{E} 分别是 info/verbose/error。
读 \Opt{check} 输出先看三处：\texttt{Loaded compilation database} 路径对不对、
\texttt{Compile command} 展开后的 \texttt{-std} 与 include 目录对不对、
最后一行 \texttt{All checks completed, N errors}。
若卡在 \texttt{Building preamble} 之后没有下文，多半是某个头文件里
藏了病态模板展开，用 \Opt{log}=\texttt{verbose} 找出正在处理的 include。

% --------------------------------------------------------------------
\section{\texttt{.clangd}：逐键拆解}\label{sec:ch05-s2}

\texttt{.clangd} 是 YAML，作用于「\CD\ 如何解析、如何表现」，
不参与真实构建。表~\ref{tab:ch05-keys} 按组列出关键键位；
每组之下值得单独说明的，按顺序讲。

\textbf{CompileFlags} 管编译命令的来路与补丁。\Opt{CompilationDatabase}
给出编译数据库所在目录（相对配置文件所在目录解析）；\Opt{Add}
把若干参数追加到展开后命令的末尾；\Opt{Remove} 按键名前缀删除参数，
支持 \texttt{*} 通配；\Opt{DetectMissingExtensions} 打开后，
头文件可以反推同名源文件的编译命令来解析自己。

\textbf{Diagnostics} 管诊断的产出面。\Opt{ClangTidy} 是二级映射：
\texttt{Add}、\texttt{Remove} 增删启用规则（支持前缀通配），
\texttt{CheckOptions} 给单条规则传参。\Opt{UnusedIncludes}
取 \texttt{None}、\texttt{Missing}、\texttt{Strict} 三档强度，
控制「没用的 include」要不要提示、能否一键删除；
\Opt{MissingIncludes} 对称地管「用了却没包含」。\Opt{Suppress}
按诊断名静音，见~\ref{sec:ch05-s4}。

\textbf{Index} 管索引。\Opt{Background} 取 \texttt{Build} 或 \texttt{Skip}；
\texttt{External} 之下 \Opt{File} 指向一个预生成的 \texttt{.idx}、
\Opt{Index} 指向 slab 目录；\Opt{StandardLibrary} 决定是否索引系统头。

\textbf{InlayHints} 的 \Opt{Enabled} 是总开关，其余八键各管一类提示：
参数名、参数类型、返回类型、限定符、指定初始化器、冒号、
数组大小、模板依赖。\textbf{Hover} 的 \Opt{ShowColor} 给颜色字面量画色块，
\Opt{StoreProfile} 控制悬停里是否展示性能剖面。\textbf{Completion} 的
\Opt{ArgumentSnippets} 让函数补全带占位符，\Opt{IncludePatterns}
限定自动补 include 的候选集，\Opt{AllScopes} 决定是否补出非当前
作用域的符号。顶层的 \Opt{Style} 取 \texttt{file} 时格式化跟随
\texttt{.clang-format}，找不到就落到 \Opt{FallbackStyle}。

\begin{longtable}{>{\raggedright\arraybackslash}p{0.26\textwidth}
                  >{\raggedright\arraybackslash}p{0.14\textwidth}
                  >{\raggedright\arraybackslash}p{0.40\textwidth}}
\caption{\texttt{.clangd} 主要键：类型与作用}\label{tab:ch05-keys}\\
\toprule
键 & 取值 & 作用 \\
\midrule
\endfirsthead
\toprule
键 & 取值 & 作用 \\
\midrule
\endhead
\texttt{CompileFlags./}\hspace{0pt}\texttt{Compilation}\hspace{0pt}\texttt{Database}
  & 路径 & 编译数据库所在目录或文件 \\
\midrule
\texttt{CompileFlags./}\hspace{0pt}\texttt{Add}
  & 列表 & 追加参数到展开命令的末尾 \\
\midrule
\texttt{CompileFlags./}\hspace{0pt}\texttt{Remove}
  & 列表 & 按键名前缀删参数，支持 \texttt{*} 通配 \\
\midrule
\texttt{CompileFlags./}\hspace{0pt}\texttt{Detect}\hspace{0pt}\texttt{Missing}\hspace{0pt}\texttt{Extensions}
  & 布尔 & 头文件反推同名源文件的编译命令 \\
\midrule
\texttt{Diagnostics./}\hspace{0pt}\texttt{ClangTidy}
  & 映射 & 增删 tidy 规则、传 \texttt{CheckOptions} \\
\midrule
\texttt{Diagnostics./}\hspace{0pt}\texttt{Unused}\hspace{0pt}\texttt{Includes}
  & 三档 & 无用 include 的提示与可修复性 \\
\midrule
\texttt{Diagnostics./}\hspace{0pt}\texttt{Suppress}
  & 列表 & 按诊断名静音（通配见正文） \\
\midrule
\texttt{Index./}\hspace{0pt}\texttt{Background}
  & 枚举 & \texttt{Build} 建，\texttt{Skip} 不建 \\
\midrule
\texttt{Index./}\hspace{0pt}\texttt{External./}\hspace{0pt}\texttt{File}
  & 路径 & 挂载预生成的外部索引 \\
\midrule
\texttt{InlayHints.*}
  & 布尔 & 总开关加八类行内提示 \\
\midrule
\texttt{Hover./}\hspace{0pt}\texttt{ShowColor}
  & 布尔 & 颜色字面量渲染色块 \\
\midrule
\texttt{Completion./}\hspace{0pt}\texttt{Argument}\hspace{0pt}\texttt{Snippets}
  & 布尔 & 函数补全带参数占位符 \\
\midrule
\texttt{Style}、\texttt{FallbackStyle}
  & 枚举 & 格式化风格与找不到配置时的兜底 \\
\bottomrule
\end{longtable}

清单~\ref{lst:ch05-config} 给一份覆盖面最大的示例（多数键保持默认即可，
这里是为了把键位摆全）。缩进即层级，布尔写 \texttt{Yes}/\texttt{No}
或 \texttt{true}/\texttt{false} 均可。

\begin{lstlisting}[style=yamlstyle,caption={\texttt{.clangd} 全键位示例（覆盖面优先，非推荐组合）},label={lst:ch05-config}]
# .clangd —— 项目根配置，入版本库
CompileFlags:
  CompilationDatabase: build      # 编译数据库所在目录
  Add: [-std=c++20, -Wall]        # 追加到命令末尾
  Remove: [-mwindows, -Wa,*]      # 删掉不合拍的参数
  DetectMissingExtensions: true   # 头文件反推源文件
Diagnostics:
  ClangTidy:
    Add: [modernize-use-nullptr, performance-*]
    Remove: [modernize-use-trailing-return-type]
    CheckOptions:
      readability-identifier-naming.VariableCase: lower_case
  UnusedIncludes: Strict          # 没用的 include 直接可见
  MissingIncludes: Strict
  Suppress: [misc-unused-parameters]
Index:
  Background: Build               # 改 Skip 可关闭后台索引
  External:
    File: D:/proj/idx/base.idx    # 预生成的外部索引
  StandardLibrary: Yes
InlayHints:
  Enabled: Yes
  ParameterNames: Yes
  ParameterTypes: No
  ReturnType: No
  Qualifiers: No
  Designators: No
  Colon: No
  ArraySizes: Yes
  Dependent: No
Hover:
  ShowColor: Yes                  # 画出颜色样本
  StoreProfile: No
Completion:
  ArgumentSnippets: Yes           # 补全带占位参数
  IncludePatterns: ['*.hpp']      # 自动 include 候选集
  AllScopes: Yes
Style: file                       # 格式化跟随 .clang-format
FallbackStyle: LLVM               # 找不到配置文件时的兜底
\end{lstlisting}

\begin{note}[title={Add 追加在末尾，Remove 支持通配}]
\begin{itemize}
  \item \texttt{Add} 的参数追加在展开后编译命令行的\textbf{末尾}：
        对「后者覆盖前者」的成对开关（\texttt{-std} 家族），
        \texttt{Add} 能盖过数据库里的旧值；反过来，需要排在输入文件
        之前的选项（如 \texttt{-include}）放 \texttt{Add} 并不保险。
  \item \texttt{Remove} 按键名前缀匹配，尾部 \texttt{*} 是通配：
        \texttt{-Wa,*} 会把汇编透传选项连同其参数一并移除。
        删除发生在展开后的命令行上，与数据库里的书写方式无关。
\end{itemize}
\end{note}

% --------------------------------------------------------------------
\section{条件配置与三文件分工}\label{sec:ch05-s3}

同一份 \texttt{.clangd} 里可以放多个 YAML 片段，用 \texttt{---} 分隔，
每段以 \Opt{If} 开头声明生效条件：\Opt{PathMatch} 接正则列表，
匹配请求文件路径；\Opt{True} 表示无条件；后段的 \Opt{Else}
兜住前面所有片段都没命中的文件（清单~\ref{lst:ch05-if}）。
片段之间未冲突的键按普通片段合并。

\begin{lstlisting}[style=yamlstyle,caption={按路径分区：If 条件片段与 Else 兜底},label={lst:ch05-if}]
# 片段一：只作用于项目源文件
If:
  PathMatch: [src/.*\.cpp, src/.*\.hpp]
CompileFlags:
  Add: [-DCORE=1]
---
# 片段二：True 表示无条件，对所有文件生效
If:
  True:
Diagnostics:
  UnusedIncludes: Strict
---
# 片段三：Else 兜住前面没命中的文件
If:
  Else:
Index:
  Background: Skip
\end{lstlisting}

片段多了难维护时，用顶层 \Opt{Import} 把专区配置拆成子文件
（清单~\ref{lst:ch05-import}）；被导入文件里的条件片段照常生效，
路径正则始终以请求文件为基准。

\begin{lstlisting}[style=yamlstyle,caption={Import 拆分专区配置},label={lst:ch05-import}]
# 项目根 .clangd：只留基线，专区拆分
Import:
  - configs/tests.clangd          # 测试树的放宽项
  - configs/third_party.clangd    # 三方头文件静音
CompileFlags:
  CompilationDatabase: build
\end{lstlisting}

最后是三个配置文件的分工：\texttt{.clangd}、\texttt{.clang-format}、
\texttt{.clang-tidy}（表~\ref{tab:ch05-files}）。查找次序都是
「从当前文件所在目录逐级向上，直到项目根」；用户级配置在
\texttt{\textasciitilde/.config/clangd/config.yaml}，需以
\Opt{enable-config} 开启，优先级低于任何项目内配置。

\begin{longtable}{>{\raggedright\arraybackslash}p{0.18\textwidth}
                  >{\raggedright\arraybackslash}p{0.27\textwidth}
                  >{\raggedright\arraybackslash}p{0.34\textwidth}}
\caption{三文件分工：谁读它、从哪找}\label{tab:ch05-files}\\
\toprule
文件 & 读取方 & 查找次序与注意 \\
\midrule
\endfirsthead
\toprule
文件 & 读取方 & 查找次序与注意 \\
\midrule
\endhead
\texttt{.clangd} & 只有 \CD\ 读 & 从当前文件目录向上到项目根；改完需重启服务器才全量生效 \\
\midrule
\texttt{.clang-}\hspace{0pt}\texttt{format}、\texttt{.clang-tidy} & \Fmt、\Tidy 与 \CD\ 都读 & 同样逐级向上；\CD\ 只借它们做格式化与诊断，不会替你跑构建 \\
\midrule
用户级 \texttt{config.yaml} & \CD\ 读 & 位于 \texttt{\textasciitilde/.config/}\hspace{0pt}\texttt{clangd/}；须 \Opt{enable-config}；被项目配置覆盖 \\
\bottomrule
\end{longtable}

\begin{bestpractice}[title={团队基线只放 CompileFlags 与 Diagnostics}]
\texttt{.clangd} 是工具行为配置，不是构建配置。基线文件只应写
「怎么解析」和「怎么诊断」（\Opt{CompilationDatabase}、\Opt{Add}、
\texttt{ClangTidy}、\Opt{UnusedIncludes}），把 InlayHints、Hover、
Completion 这类个人偏好留给成员在用户级配置里自定义——
这样行为差异可解释，代码评审里的配置抖动也会显著变少。
\end{bestpractice}

% --------------------------------------------------------------------
\section{诊断噪声治理}\label{sec:ch05-s4}

\CD\ 的诊断来自三层：编译器本身的 \texttt{-W} 家族警告、
clang-tidy 检查（用 \Opt{clang-tidy} 命令行开关或
\texttt{Diagnostics./}\hspace{0pt}\texttt{ClangTidy} 片段启用）、
以及 \CD\ 自产的 \texttt{unused-includes} 一类提示。
治理手段有四件，按侵入性从低到高排。

第一件是 \Opt{Suppress}：按「诊断来源名」静音，支持三种写法——
编译器诊断用 \texttt{clang-diagnostic-} 前缀加警告名，
\texttt{clang-diagnostic-*} 可一键关掉全部编译器诊断；
tidy 检查直接用检查名（可带前缀通配）；\CD\ 自产的提示用
\texttt{unused-includes} 这类来源名（清单~\ref{lst:ch05-suppress}）。

\begin{lstlisting}[style=yamlstyle,caption={Suppress 与 WarningsAsErrors 配合},label={lst:ch05-suppress}]
Diagnostics:
  Suppress:
    - clang-diagnostic-*            # 编译器警告全静默
    - misc-unused-parameters        # tidy 检查名
    - unused-includes               # clangd 自产提示
  WarningsAsErrors:                 # 白名单升格回错误
    - bugprone-use-after-move
    - clang-diagnostic-uninitialized
\end{lstlisting}

第二件是 \Opt{WarningsAsErrors}：把名单内的警告在编辑器里按错误渲染，
与 \Opt{Suppress} 互补——先大面积静音，再把真正危险的几条点亮。
第三件回到命令行层：在 \texttt{CompileFlags./}\hspace{0pt}\texttt{Add} 里加
\texttt{-Wno-xxx}，让警告在解析阶段就不产生；它与 \Opt{Suppress}
效果相近，但会同时影响该 TU 里一切基于诊断的功能。
第四件是针对 include 卫生的三档开关：\Opt{UnusedIncludes} 与
\Opt{MissingIncludes} 各取 \texttt{None}、\texttt{Missing}、
\texttt{Strict}，决定提示是否存在、是否可一键修复。

\begin{warning}[title={别把 \texttt{.clangd} 当编译器配置乱加 \texttt{-Werror}}]
症状：编辑器里到处红线，命令行构建却零警告。原因：在 \texttt{Add} 里
塞了 \texttt{-Werror}，\CD\ 解析时把一切警告升格成错误，
而真实构建的命令行并没有这个开关——两边看到的「错误」从此对不上。
对策：编辑器侧的严格度交给 \Opt{WarningsAsErrors} 白名单精确控制，
构建层的警告策略留在构建系统里，同步给 \CD{}。
\end{warning}

% --------------------------------------------------------------------
\section{交叉编译与外部索引}\label{sec:ch05-s5}

交叉工程里 \CD\ 默认拿宿主头文件解析目标代码，系统头一飘红，
全项目跟着遭殃。正解是把「目标工具链的信息」喂给它：
系统根用 \texttt{-isysroot}（编译数据库里有最好；没有就在
\texttt{CompileFlags./}\hspace{0pt}\texttt{Add} 里补，或退而求其次用
\Opt{sysroot} 命令行参数）；编译器内建头与宏用
\Opt{query-driver}——\CD\ 会执行白名单里的交叉编译器抽取系统头路径，
白名单必须是与编译数据库\textbf{同一支}编译器，否则抽出来的路径
对不上号。数据库不在源码根时，用 \Opt{compile-commands-dir} 指路。
另外建议 \Opt{header-insertion}=\texttt{never}：补全不应往交叉文件里
插宿主头文件。

索引侧则解决「首次打开没有跨文件能力」：大项目可在 CI 里离线建库
（\texttt{clangd-indexer} 产出 \texttt{.idx}，或以索引服务
\texttt{clangd-index-server} 的形态部署），本地 \CD\ 用
\texttt{Index./}\hspace{0pt}\texttt{External./}\hspace{0pt}\texttt{File}
或 \Opt{index-file} 挂载现成文件，再关掉后台索引
（清单~\ref{lst:ch05-cross}）。

\begin{lstlisting}[style=shellstyle,caption={交叉编译启动参数与外部索引挂载},label={lst:ch05-cross}]
:: 交叉工程：让 clangd 向目标工具链要系统头
D:/ide/bin/clangd.exe ^
  --compile-commands-dir=build/arm ^
  --query-driver="D:/x-tools/arm-none-eabi-g++.exe" ^
  --header-insertion=never

:: 外部索引：CI 离线生成，本地只挂载
clangd-indexer --executor=all-TUs ^
  build/arm/compile_commands.json > build/arm/all.idx
D:/ide/bin/clangd.exe --index-file=build/arm/all.idx ^
  --background-index=false
\end{lstlisting}

两点提醒：\Opt{query-driver} 等价于允许 \CD\ 执行外部程序，
白名单要写死具体可执行文件而不是宽泛的 \texttt{*}；
外部索引与源码不同步时跨文件结果会「跳旧位置」，
CI 里应随构建产物一起重新生成。

% --------------------------------------------------------------------
\section{排错手册}\label{sec:ch05-s6}

表~\ref{tab:ch05-debug} 覆盖六个高频故障。通用流程：先
\Opt{check}=\texttt{file.cpp} 看流水线在哪一步断了，再对照本表。

\begin{longtable}{>{\raggedright\arraybackslash}p{0.17\textwidth}
                  >{\raggedright\arraybackslash}p{0.33\textwidth}
                  >{\raggedright\arraybackslash}p{0.32\textwidth}}
\caption{\CD{} 高频故障：症状、根因与处置命令}\label{tab:ch05-debug}\\
\toprule
症状 & 根因 & 处置 \\
\midrule
\endfirsthead
\toprule
症状 & 根因 & 处置 \\
\midrule
\endhead
完全没有补全 & 没找到编译数据库，\CD\ 用兜底命令行（无 include 路径、默认标准）解析，AST 根本建不起来 & \Opt{check}=\texttt{main.cpp} 看 \texttt{Loaded compilation database}；修数据库见第 2 章，或 \texttt{CompileFlags./}\hspace{0pt}\texttt{Compilation}\hspace{0pt}\texttt{Database} 指路 \\
\midrule
到处飘红、标准头也报错 & 交叉编译工具链的系统头未识别：缺 \texttt{-isysroot} 与内置头路径 & 加 \Opt{query-driver}=\hspace{0pt}\texttt{"**/}\hspace{0pt}\texttt{arm-none-eabi-g++"} 并重启；确认路径含 \texttt{**} 通配 \\
\midrule
跳转到声明不跳定义 & 定义在另一个 \texttt{.cpp}，当前 TU 里没有；后台索引尚未建成或被关闭 & 等 \texttt{building index} 结束；确认 \texttt{Index./}\hspace{0pt}\texttt{Background} 未关、\texttt{.clangd/index} 未被清理 \\
\midrule
\texttt{references} 漏项目内文件 & 只查了当前 TU；或索引过期、被 \texttt{Remove} 掉的编译参数使部分文件从未入索引 & 重建索引（删 \texttt{.clangd/index} 后重启）；用 \Opt{check} 逐一验证问题文件能否拿到编译命令 \\
\midrule
宏体内补全/悬停失效 & 宏展开位置的请求落到预处理层，\CPP{} 语义在展开前无法回答（\CD\ 日志出现 \texttt{not a pointer to expression}） & 属已知限制：在宏展开后的等价写法上提问，或临时展开宏定位；复杂宏改写见第 14 章 \\
\midrule
启动即退出、反复重启 & \texttt{.clangd} YAML 解析失败、配置目录不可写，或旧版本索引 slab 与新版本不兼容 & 用 \Opt{check} 带 \Opt{disable-config} 二分定位；删除 \texttt{.clangd/index} 重生成；升级后避免混用不同版本的 \CD\ 与索引文件 \\
\bottomrule
\end{longtable}

两条通则可解一半问题：其一，\textbf{\CD\ 的输出完全由「编译数据库加
\texttt{.clangd} 补丁」决定}，怀疑行为异常先拿 \Opt{check} 看它展开出的
真实命令行；其二，\textbf{索引是异步的}，跨文件能力刚启动时永远不完整，
用日志里 \texttt{indexed} 计数判断是否收敛。行为仍诡异时，
把 \Opt{log}=\texttt{verbose} 的片段留着——向 \CD\ 提 issue 时它就是复现脚本。
配置键的完整清单在附录 A，命令行速查在附录 B；协议层的问题（能力没协商上、
消息错序）请回看第 4 章，索引的存储结构在第 7 章继续讲。

% ---- frag_ch06.tex ----
\chapter{LSP 实现对比：ccls、vscode-cpptools、Eclipse CDT 与 rust-analyzer}\label{ch:lsp-implementations}

第 4 章把 LSP 拆成了报文，第 5 章把 \CD\ 的内部结构摊开。本章换一个视角：
真正可用的 C/C++ 语言服务器有哪几类，它们在索引方式、内存占用与项目配置上各做了什么取舍，
以及为什么最后都会落到同一件事上——你必须把编译命令行交给它。

\CD\ 位于 \texttt{clang-tools-extra/clangd}，是 LLVM 项目的正式组成部分，自 LLVM 7 起随每个
release 打包发布。这个「官方地位」带来两个直接后果。

其一是演进节奏：\CD\ 的能力边界由 clang 前端决定。clang 支持的语言标准（\cpp{11} 到
\cpp{23}）、属性、模块化（\texttt{import} 与 BMI）什么时候能用，取决于前端本身走到哪一步；
\CD\ 的 LSP capabilities 也随大版本推进，一个 LLVM release 一次，节奏比任何第三方实现都稳定。
其二是配套工具的零成本对齐：\Fmt\ 与 \Tidy\ 与 \CD\ 共用同一份 clang AST 与同一条
\texttt{ParseAST} 路径，所以「编辑器里看到的语法树」和「\Tidy\ 报问题的语法树」是同一棵树。
这一点在第 10 章的 CI 门禁里会反复用到：本地与流水线用同一份 AST，就不会出现
「编辑器不报、CI 报」的分类分歧。

优点之外，三条短板必须写清楚。

\textbf{内存。}\CD\ 每个活动文件都要保留 preamble（已算完宏与头的 PCH）、AST 与索引查询缓存。
小项目几百 MB 量级；几万到十几万个 TU 的仓库里，后台索引可以吃满若干 GB 内存，
并同时占用同样数量级的核。缓解手段是调整 \Opt{pch-storage}、开或关 \Opt{background-index}、
用 \Opt{limit-results} 限制返回量，以及在 \texttt{.clangd} 里把不需要索引的目录排除掉。

\textbf{后台索引重。}第一次索引是全量的：\CD\ 会为编译数据库里的每个文件跑一遍前端。
这就是第 3 章讲的「索引的代价」在实现层面的体现——索引不是缓存，是一次重新编译。

\textbf{构建系统集成要靠外接。}\CD\ 自己不会读 \texttt{CMakeLists.txt}，也不会执行
\texttt{make}。它只认 \texttt{compile\_commands.json}，所以你要么让构建系统导出，要么拦截：

\begin{lstlisting}[style=cmakestyle,caption={让 CMake 自己产出编译数据库：clangd 最省事的接法},label={lst:ch06-cmake}]
set(CMAKE_EXPORT_COMPILE_COMMANDS ON)

add_library(core src/parser.cpp src/lexer.cpp)
target_compile_definitions(core PUBLIC CORE_ENABLE_TRACE=1)
target_compile_features(core PUBLIC cxx_std_20)
target_include_directories(core PUBLIC include)

add_executable(app src/main.cpp)
target_link_libraries(app PRIVATE core)
\end{lstlisting}

\begin{lstlisting}[style=shellstyle,caption={三种拿到编译数据库的方式，以及单 TU 自检},label={lst:ch06-compdb}]
# 1) CMake：直接读构建目录里的 compile_commands.json
cmake -S . -B build -DCMAKE_EXPORT_COMPILE_COMMANDS=ON

# 2) Makefile 或自研构建：拦截编译器调用
bear --output compile_commands.json -- make -j8

# 3) 已有构建目录：软链到仓库根，避免把 build 里那份手改
ln -sf build/compile_commands.json .

# 单个 TU 解析结果如何，比看编辑器波浪线更快
clangd --check=src/parser.cpp \
  --compile-commands-dir=build --log=verbose
\end{lstlisting}

\begin{warning}[title={把同一份编译数据库同时喂给多个引擎}]
如果你为了对比效果，在同一个窗口里同时启用 \CD\ 与另一个引擎（例如 \texttt{cpptools}），
同一个文件会收到两套诊断：两套波浪线、两份 completion 列表、两个名字相近的 quick fix。
症状通常表现为「补全里出现两次同样的候选」和「自动修复选哪个都一样但结果不同」。
正确做法是\textbf{每个工作区只留一个语义引擎}：另一个要么整窗口禁用，要么只用它做
交叉验证，并把它的诊断输出关掉（\texttt{C\_Cpp.errorSquiggles} 设为 \texttt{disabled}）。
\end{warning}

\section{ccls 与 cquery：省内存的代价}\label{sec:ch06-s2}

ccls 是同一条技术线上的另一条路：它同样基于 clang 前端，但索引策略不同。ccls 建的是
IBCM（Identifier Based Cross Module Call Graph，按标识符组织的跨文件调用图），
把符号与引用关系落进 \texttt{.ccls-cache} 目录；它的 preamble 共享更激进，
相同命令行前缀的文件复用同一份 preamble，因此常驻内存与索引体积通常低于 \CD\ 的
后台索引，在低配机器和巨型老仓库上是有真实吸引力的。

ccls 的接入方式和 \CD\ 一样：项目根放 \texttt{compile\_commands.json}，或者退一步写一个
\texttt{.ccls} 文件（一行一个 flag，首行是工程根）。客户端侧配置走初始化选项：

\begin{lstlisting}[style=jsonstyle,caption={ccls 的初始化选项片段（不同客户端的外层键名略有差别）},label={lst:ch06-ccls}]
{
  "initOptions": {
    "clang": { "extraArgs": ["-std=c++20", "-Wall"] },
    "cache": { "directory": ".ccls-cache" },
    "index": {
      "initial": "clang",
      "onChange": true,
      "onWrite": true,
      "threadsCount": 0
    }
  }
}
\end{lstlisting}

其中 \texttt{index.onWrite} 控制「保存文件时是否立刻重新索引」——把它打开能得到
接近 \CD\ 的实时性，代价是每次保存都会重算受影响子图。ccls 还有一些非标准的自定义请求，
如 \texttt{\$ccls/call}（谁调用我 / 我调用谁）、\texttt{\$ccls/member}、
\texttt{\$ccls/inheritance}，这些是 \CD\ 没有对等能力的地方（\CD\ 走标准的
\texttt{callHierarchy/*} 与 \texttt{typeHierarchy/*}）。

必须写清现状：\textbf{ccls 已归档，停止维护}；它的前身 cquery 更早停摆（cquery 用
\texttt{.cquery} 项目文件与独立的 JSON 索引，今天不需要再考虑）。选 ccls 意味着承担什么：

\begin{itemize}
  \item clang 前端版本被钉死在归档点，新语法（\cpp{20}/\cpp{23} 的较新特性、模块化）
        会在编辑器里报错，而你的编译器已经支持它们；
  \item LSP 规范的新能力（标准 semantic tokens、inlay hints、pull diagnostics、
        progress 通知）不会补上，客户端越新越容易踩到未实现的方法；
  \item 协议层 bug 与安全相关的崩溃只能自己打补丁、自己编译；
  \item 生态：主流量化配置片段、CI 集成脚本都是围绕 \CD\ 写的，迁移成本由你付。
\end{itemize}

结论：ccls 值得用在一类场景——旧标准、旧编译器、内存吃紧、且你愿意自己维护构建的仓库。
新项目一律用 \CD\ 或商业引擎，不要拿一个停止演进的前端当团队的地基。

\begin{note}[title={闭源引擎的误报无法自查根因}]
用 ccls 或开源 \CD\ 时，遇到一个可疑诊断，你可以 \texttt{clang -Xclang -ast-dump} 把 AST
打出来，或者直接读 \texttt{clangd/SemaRebuilds.cpp}、\texttt{Index/Background.cpp} 确认
它是重用了 preamble 还是重跑了 \Fn{ParseAST}。用商业或闭源引擎时这条路径不存在：
你能看到「它报了什么」，看不到「它为什么报」，也无法确认它是解析失败后的降级结果
还是真正的语义判断。可行的替代办法是做\textbf{双引擎对照实验}：同一份源码、同一份编译
数据库，让开源引擎也跑一遍，两边不一致处优先怀疑配置（宏、标准版本、头搜索路径）。
\end{note}

\section{cpptools：闭源引擎与其配置文件}\label{sec:ch06-s3}

VS Code 的 C/C++ 扩展（仓库名 \texttt{vscode-cpptools}）是微软自家的实现，核心 IntelliSense
引擎闭源。它和 \CD\ 的根本差别是\textbf{配置来源}：\CD\ 要求你提供每个 TU 的真实命令行，
而 cpptools 允许你在 \texttt{c\_cpp\_properties.json} 里手写一份「IntelliSense 用的编译环境」，
它自己拼出等价命令行。

\begin{lstlisting}[style=jsonstyle,caption={c\_cpp\_properties.json：把编译环境显式写给闭源引擎},label={lst:ch06-cpptools}]
{
  "version": 4,
  "configurations": [
    {
      "name": "linux-gcc",
      "includePath": [
        "${workspaceFolder}/include",
        "${workspaceFolder}/third_party/**"
      ],
      "defines": ["CORE_ENABLE_TRACE=1", "__linux__"],
      "compilerPath": "/usr/bin/g++-13",
      "compilerArgs": ["-std=c++20"],
      "cStandard": "c17",
      "cppStandard": "c++20",
      "intelliSenseMode": "linux-gcc-x64",
      "configurationProvider": "ms-vscode.cmake-tools"
    }
  ]
}
\end{lstlisting}

各字段的实际含义：\texttt{includePath} 是头搜索路径，写 \texttt{**} 表示递归；
\texttt{defines} 是强制前置的宏；\texttt{compilerPath} 让引擎去问编译器要系统头路径
与内建宏（这一步与 \CD\ 的 \Opt{query-driver} 是同一件事）；\texttt{intelliSenseMode}
决定 ABI 宏与平台假设，取值形如 \texttt{windows-msvc-x64} 或 \texttt{linux-clang-x64}。
\texttt{configurationProvider} 则声明「配置由另一个扩展供给」——设成
\texttt{ms-vscode.cmake-tools}（或 \texttt{vector-of-bool\allowbreak.cmake-tools}）之后，
CMake Tools 解析工程时得到的 include 路径、宏与标准版本会自动灌进 IntelliSense，
这是它比 \CD\ 少一道手续的地方。

引擎内部分层也要知道，因为它决定「为什么跳不准」。
\textbf{AST 引擎}（\texttt{intelli\Bk SenseEngine} 设为 \texttt{default}）按 TU 做真正的语义解析，
重载决议、模板实例化、条件编译都对。另一档是 tag parser，把同一项设为
\texttt{tagParser} 就生效：它只扫符号表，不解析包含关系与宏，速度很快，
却会把同名但不同物的候选一起给你，跨文件重载也可能选错。
它的价值在于大仓降级与 Windows 平台。cpptools 对微软的标准库支持得最好，
\texttt{\_\_declspec} 这类扩展关键字也解析得对，这是开源引擎在 Windows 上常年的短板。
代价写在配置里。手写 \texttt{includePath} 一旦漏了某个目录，引擎并不报错，
它只会安静地把那个头文件当成未解析，然后给你一堆「未知类型名」。

\begin{compare}[title={能跳转不等于语义正确}]
「按 F12 能跳到定义」这一条几乎筛不出东西：ctags、tag parser、只做正则匹配的插件都能跳。
语义正确要求的是另一件事：\textbf{编辑器里的语法树，必须和编译器在那条命令行下
构造出来的那一棵一致}。差别在这四类用例里立刻显形。
\begin{itemize}
  \item 条件编译：\texttt{\#ifdef} 分支里的代码，tag parser 当成有效代码；
  \item 重载与模板：跳到哪一个，需要真正的重载决议；
  \item 宏：\texttt{\#define} 展开出的类型与成员，非语义级引擎看不见；
  \item 同一头文件在不同 TU 下含义不同，宏定义一变，跳转目标就该跟着变。
\end{itemize}
选型时请用这四类用例做判定，而不是问「有没有补全」。
\end{compare}

\section{IDE 内置引擎与商业解析器}\label{sec:ch06-s4}

\textbf{Eclipse CDT 的语言服务器。}CDT 把自家解析器包成语言服务器，实现走 LSP4J，
索引由 CDT 的 \texttt{IndexManager} 维护在工作区元数据里，符号模型是 CDT 的 DOM 与 AST。
它可以导入 \texttt{compile\_commands.json}，但原生的配置单位仍是 CDT 的工程模型，
也就是 Managed Build 与 \texttt{.cproject}。因此在 CMake 工程上想保持一致的体验，
通常要把工程重新导入成 CDT 工程。Theia 里的 C 与 C++ 支持复用的也是同一套 LSM。
它的短板是 LSP 新特性跟进较慢，inlay hints 与 pull diagnostics 一类常常缺；
优势是 CDT 索引能持久化、能跨工程引用，以及 EPL-2.0 许可允许你自由嵌进自己的 IDE。

\textbf{CLion。}JetBrains 用的是自研 C/C++ 解析器（在 PSI 之上再叠一层语义模型），
配置来源是从 \texttt{CMakeLists.txt}、Gradle 或 Makefile 直接解析出来的编译参数——
所以用户不需要显式提供编译数据库，但代价是「它解析出的命令行是不是你构建时那条」
变成了一个需要核对的黑盒（Project Properties 窗口可以看每个 TU 的实际 flags）。
CLion 同时提供可切换的后端：较新版本里可以把解析/补全交给 \CD\ 跑（实验性开关），
\Tidy\ 与 \Fmt\ 则一直是外挂进程。适合「要一体化 IDE、要重构与本地历史」的团队，
不适合追求零依赖与命令行的场景，专有许可与 JVM 内存是明面上的成本。

\textbf{Qt Creator。}默认后端是 \texttt{clangcode} 插件，基于 libclang 的 C API
（\texttt{clang\_cursor\_} 一族），只拿诊断、代码操作与跳转，索引与跨文件引用能力弱，
编译参数从 \texttt{.pro}/\texttt{.qmake} 或 CMake 工程里推。新版本提供以 \CD\ 为代码模型的
可选路径（在 C++ 相关设置里切换，Qt Creator 负责生成编译数据库再把 \texttt{\_clangd.conf}
交给它）。缺口是：libclang C API 拿不到完整的 Sema 结果，模板重载与自动推导会退化；
收益是启动快、与 Qt 元对象系统（\texttt{Q\_OBJECT}/信号槽）的耦合处理得比裸 \CD\ 好。

\section{非 C++ 对照：为什么 C++ 离不开编译数据库}\label{sec:ch06-s5}

把视线挪出 C/C++，能反证出结论。这些语言的语言服务器都活得很好，但它们的项目配置形态完全不同。

\textbf{rust-analyzer} 是自己写的解析器（rowan 语法树）加自研的名字解析与类型推导，
它\textbf{不需要编译数据库}：配置来自 Cargo 元数据（\texttt{cargo metadata} 给出 crate 图、
依赖、feature 与 flag），每个 crate 的源码清单是确定的。宏用 mbe（macro-by-example）
自研展开，过程宏则要求把 \texttt{proc-macro} 编译成动态库后单独加载。

\textbf{gopls} 用 \texttt{go list}（底层是 \texttt{go/packages}）一次拿到构建配置，
而 Go 的语义单位天然是 package：单个文件的类型检查可以只靠它所属 package 的文件集完成，
不依赖宏、不依赖头搜索顺序、也不依赖编译 flag 的组合。这就是「单文件语义完整」的含义。

\textbf{pyright} 与 \textbf{tsserver} 更彻底：Python 与 TypeScript 没有「一条 TU 对应一条命令行」
这个概念。pyright 只需要工作区根与 \texttt{pyrightconfig.json}（外加从 site-packages 找 stub），
tsserver 需要 \texttt{tsconfig.json}——那是路径映射与语言版本，而不是预处理器状态。

{\small\setlength{\tabcolsep}{3pt}
\begin{longtable}{>{\raggedright\arraybackslash}p{0.15\textwidth} >{\raggedright\arraybackslash}p{0.27\textwidth} >{\raggedright\arraybackslash}p{0.13\textwidth} >{\raggedright\arraybackslash}p{0.25\textwidth}}
\caption{非 C/C++ 语言服务器的项目配置来源}\label{tab:ch06-noncpp}\\
\toprule
服务器 & 配置来源 & 每 TU 命令行 & 单文件语义是否完整 \\
\midrule
\endfirsthead
\toprule
服务器 & 配置来源 & 每 TU 命令行 & 单文件语义是否完整 \\
\midrule
\endhead
\bottomrule
\endfoot
rust-analyzer & cargo metadata 与自研 mbe 宏展开 & 无 & 部分，宏须先展开 \\
gopls & go list 得到的 package 级文件集 & 无 & 完整 \\
pyright & 工作区根与 pyrightconfig.json & 无 & 完整 \\
tsserver & tsconfig.json 的路径映射 & 无 & 完整 \\
clangd 一类 C/C++ 引擎 & compile\_commands.json & 有，且必须显式给 & 不完整，须带全部依赖头 \\
\end{longtable}}

对照之后，C++ 的那条要求就显出来了：\textbf{一个 TU 的语义由它的命令行决定}。
同一份 \texttt{parser.cpp}，配上 \texttt{-DCORE\_ENABLE\_TRACE=1} 与不配，
\texttt{\#if} 分支里的成员是否存在都不一样；\Opt{std=c++17} 与 \Opt{std=c++20}
会让同一个声明从合法变成非法；include 顺序改变会换掉被 \texttt{\#pragma once}
挡住的那一份头；交叉编译时 \texttt{-target}、\texttt{-mcpu} 与 sysroot 又改掉了
所有 \texttt{\_\_ARM\_NEON\_\_} 一类的内建宏。语言服务器不可能猜这些，
只能被告知。C++ 没有「package 就是语义单位」这种结构，也没有可信的默认 flag 集，
这就是编译数据库在这一行里是必需、而在其它四行里是「否」的原因。

顺带一句实现层结论：\CD\ 之所以在每个文件里都跑一次真前端、内存与 CPU 都高，
正是上面这些条件要求它不能走捷径。任何号称「省掉了编译数据库」的 C/C++ 引擎，
要么是在自己重建它（cpptools 的 \texttt{compilerPath}、CLion 的工程解析），
要么就是在降级成语法级的近似。

\section{横向对比与三档选型}\label{sec:ch06-s6}

{\small\setlength{\tabcolsep}{3pt}
\begin{longtable}{>{\raggedright\arraybackslash}p{0.13\textwidth} >{\raggedright\arraybackslash}p{0.19\textwidth} >{\raggedright\arraybackslash}p{0.16\textwidth} >{\raggedright\arraybackslash}p{0.24\textwidth}}
\caption{C/C++ 语言服务器能力矩阵（一）：协议能力与许可}\label{tab:ch06-cap}\\
\toprule
实现 & Semantic Tokens & Inlay Hints & 许可证 \\
\midrule
\endfirsthead
\toprule
实现 & Semantic Tokens & Inlay Hints & 许可证 \\
\midrule
\endhead
\bottomrule
\endfoot
clangd & 原生，支持 delta 增量 & 原生，参数名与推导类型 & Apache 2.0，含 LLVM 例外条款 \\
ccls & 未实现标准接口，靠客户端主题着色 & 无 & MIT \\
cpptools & 部分，与内置主题的作用域耦合 & 原生，可按提示类型开关 & 闭源，扩展与问题跟踪仓库开源 \\
CDT LSM & 部分 & 无 & EPL 2.0 \\
CLion & 不适用，走自家编辑器模型 & 部分 & 专有 \\
\end{longtable}}

{\small\setlength{\tabcolsep}{3pt}
\begin{longtable}{>{\raggedright\arraybackslash}p{0.13\textwidth} >{\raggedright\arraybackslash}p{0.15\textwidth} >{\raggedright\arraybackslash}p{0.20\textwidth} >{\raggedright\arraybackslash}p{0.11\textwidth} >{\raggedright\arraybackslash}p{0.23\textwidth}}
\caption{C/C++ 语言服务器能力矩阵（二）：编译数据库、索引方式与代价}\label{tab:ch06-cost}\\
\toprule
实现 & 编译数据库 & 索引方式 & 内存 & 适用场景 \\
\midrule
\endfirsthead
\toprule
实现 & 编译数据库 & 索引方式 & 内存 & 适用场景 \\
\midrule
\endhead
\bottomrule
\endfoot
clangd & 必需，缺失则退回 fallback flags & 后台索引，preamble 复用，缓存落盘 & 中高 & 要与 \Tidy\ 同 AST 的语义精度 \\
ccls & 必需，或根目录一个 .ccls 文件 & IBCM 按标识符建跨文件调用图 & 中 & 旧标准、内存吃紧、能自维护 \\
cpptools & 可选，由配置文件生成等价命令行 & AST 引擎为主，tag parser 兜底 & 中 & Windows 与 MSVC 头生态 \\
CDT LSM & 可选，可导入也可用 CDT 工程模型 & CDT 索引库驻留工作区 & 高 & Eclipse 与 Theia 产品线 \\
CLion & 不需用户提供，IDE 自行解析 & 自研模型，可切换 clangd 后端 & 高 & 一体化 IDE 与重构优先 \\
\end{longtable}}

三档选型，每档都说清理由与代价。

\textbf{个人最小摩擦：}\CD\ 加 CMake 的
\texttt{CMAKE\_EXPORT\_COMPILE\_COMMANDS} 设为 \texttt{ON}（非 CMake 用
\texttt{bear}），编辑器只装一个 clangd 扩展，配置尽量交给默认值。理由：一次配置覆盖
补全、跳转、重构、\Tidy\ 与格式化，认知负担最低。代价：首次索引要等，笔记本风扇要响，
Windows 上遇到 MSVC 头解析问题需要自己配 \Opt{query-driver}。

\textbf{团队统一：}\CD\ 加进仓库的 \texttt{.clangd} 与固定的编译数据库生成脚本，
外加 CI 里的 \texttt{clangd --check} 与 \Tidy\ 门禁。理由是「同一条命令行」这件事
必须由版本控制保证，否则十个人有十份 IntelliSense。代价：编译数据库变成需要维护的产物，
CMake 结构要能导出（见清单 \ref{lst:ch06-cmake}），跨平台还得处理 MSVC 的 flag 差异。

\begin{lstlisting}[style=yamlstyle,caption={提交进仓库的 .clangd：把命令行补齐与诊断策略固定下来},label={lst:ch06-clangd}]
CompileFlags:
  CompilationDatabase: build
  Add: [-std=c++20, -Wall, -Wextra]
  Remove: [-march=native, -fprofile-use]

Diagnostics:
  ClangTidy:
    Add: [bugprone-*, performance-*, readability-*]
    Remove: [modernize-use-trailing-return-type]
  UnusedIncludes: Strict
  MissingIncludes: Strict

Index:
  Background: Build

InlayHints:
  Enabled: Yes
  ParameterNames: Yes
  DeducedTypes: Yes
\end{lstlisting}

\textbf{大仓与交叉编译：}先做「索引范围控制」再谈引擎选择。做法是把编译数据库裁剪成
当前子系统的子集（脚本过滤 \texttt{compile\_commands.json}），把 \texttt{CompileFlags} 段的
\texttt{CompilationDatabase}
指向过滤后的目录，必要时关掉后台索引（\Opt{background-index} 与 \texttt{Index.Background: Skip}）
改用外部索引或只保留工作区符号。Yocto、Android Soong 这类构建不能用
\texttt{bear} 硬拦，要让构建系统本身导出（或从 \texttt{compile\_commands.json} 的中间产物里取）。
代价：引用查找的完整性下降，跨子系统的改名不再可靠，你必须接受「编辑器只对你正在改的
那一层负责」。

\begin{bestpractice}[title={用四类用例做选型}]
拿你自己仓库里最难的四段代码做判定：一个 \texttt{\#ifdef} 密集的头、一个模板重载集、
一个用宏生成的成员（\texttt{Q\_OBJECT} 一类）、一个跨三个目录的符号引用。
逐个检查跳转目标、悬停里的真实签名、rename 是否只改到该改的位置。
别问「有没有补全」——四类用例里有一项不稳，就把它列成已知缺陷写进团队文档。
\end{bestpractice}

% ---- frag_ch07.tex ----
\chapter{DAP 调试、构建任务与 editorconfig 的闭环}\label{ch:dap-and-build}

前两章把语言服务器讲完了：一个把源码变成语义结果的进程，靠 LSP 与编辑器对话。
但「能补全、能跳转」离「能交付」还差两段路——\textbf{把代码编译出来}和
\textbf{把编译出来的程序停下断点}。这一段就把这两件事接进同一个闭环里：
调试走 DAP，构建走 \texttt{tasks.json}，诊断落点走 problemMatcher，
而「人和机器各管哪一半排版」由 \texttt{.editorconfig} 与 \texttt{.clang-format} 分工。
本章仍不碰格式化字段本身（那是第 9 章），只讲它们如何被保存动作触发。

% --------------------------------------------------------------------
\section{DAP 的模型：一条会话式的、没有索引的连接}\label{sec:ch07-s1}

Debug Adapter Protocol（DAP）是 Microsoft 从 VS Code 的调试实现里抽出来的协议，
和 LSP 一样跑在 JSON 报文之上，一样用 \texttt{Content-Length} 帧头分帧。
差别全在「假设」上：

\begin{itemize}
  \item \textbf{会话式。}DAP 是一次调试会话的生命周期，从 \texttt{initialize} 到
        \texttt{terminated}；LSP 是编辑器打开期间常驻的服务。
  \item \textbf{没有索引。}adapter 不建符号表、不缓存 AST，它只是一个被调试进程的
        遥控器；所有状态都在目标进程和它的调试桩里。
  \item \textbf{单连接。}一个调试会话对应一个 adapter 进程，一对 stdio；
        LSP 服务器往往一个工作区一个实例，但内部并发服务多个文件。
  \item \textbf{没有 workspace 概念。}DAP 的 \texttt{initialize} 里没有
        \texttt{rootUri} 这一类能力协商，只有客户端能力（\texttt{pathFormat}、
        \texttt{linesStartAt1}、\texttt{columnsStartAt1}）。
\end{itemize}

还有一个容易忽略的对账差异：LSP 用请求里的 \texttt{id} 把响应配回去，
DAP 用外层信封的 \texttt{seq}（自增序号）与响应/事件里的 \texttt{request\_seq}
（回指触发它的那条请求）。响应里同时可能出现 \texttt{body} 与 \texttt{message}；
\texttt{success} 为 \texttt{false} 时看 \texttt{message}。

下面是一段真实会话的骨架（为省篇幅，删掉了无关事件；外层信封字段都是真的）：

\begin{lstlisting}[style=jsonstyle,caption={一次 cppdbg 断点会话的 DAP 报文骨架（请求/响应/事件三类）},label={lst:ch07-dap}]
{ "seq": 1, "type": "request",  "command": "initialize",
  "arguments": { "clientID": "vscode", "adapterID": "cppdbg",
                 "pathFormat": "path", "linesStartAt1": true,
                 "columnsStartAt1": true, "supportsVariableType": true,
                 "supportsRunToFunctionalBreakpointRequest": false } }
{ "seq": 2, "type": "event",    "event": "initialized" }
{ "seq": 3, "type": "response", "request_seq": 1, "success": true,
  "command": "initialize",
  "body": { "supportsConfigurationDoneRequest": true,
            "supportsConditionalBreakpoints": true,
            "supportsEvaluateForHovers": true,
            "supportsTerminateRequest": true,
            "exceptionBreakpointFilters": [
              { "filter": "all", "label": "Raised Exceptions",
                "supported": true } ] } }
{ "seq": 4, "type": "request",  "command": "setBreakpoints",
  "arguments": {
    "source": { "path": "/work/src/parser.cpp", "name": "parser.cpp" },
    "breakpoints": [ { "line": 42, "condition": "n > 1000" } ] } }
{ "seq": 5, "type": "request",  "command": "launch",
  "arguments": { "program": "/work/build/app", "cwd": "/work",
                 "stopAtEntry": false, "args": [ "-i", "sample.txt" ] } }
{ "seq": 6, "type": "request",  "command": "configurationDone" }
{ "seq": 7, "type": "event",    "event": "output",
  "body": { "category": "console", "output": "Reading...\n" } }
{ "seq": 8, "type": "event",    "event": "stopped",
  "body": { "reason": "breakpoint", "threadId": 1,
            "allThreadsStopped": true } }
{ "seq": 9, "type": "request",  "command": "stackTrace",
  "arguments": { "threadId": 1, "startFrame": 0, "levels": 20 } }
{ "seq": 10, "type": "request",  "command": "scopes",
  "arguments": { "frameId": 0 } }
{ "seq": 11, "type": "request",  "command": "variables",
  "arguments": { "variablesReference": 1001 } }
{ "seq": 12, "type": "request",  "command": "stepNext",
  "arguments": { "threadId": 1 } }
{ "seq": 13, "type": "event",    "event": "terminated",
  "body": { "restartSupported": false } }
\end{lstlisting}

逐段解释这些字段：

\begin{itemize}
  \item \texttt{initialize} 的 \texttt{adapterID} 决定「要启动哪个 adapter」，
        由客户端配置里的 \texttt{type} 提供；\texttt{pathFormat} 告诉 adapter
        路径是 \texttt{path} 风格还是 \texttt{uri} 风格；\texttt{linesStartAt1}
        与 \texttt{columnsStartAt1} 是行列基准声明。
  \item \texttt{initialized} 是\textbf{事件}，由 adapter 发给客户端，意思是
        「握手完成，你可以开始下断点、发启动参数了」。客户端在此之前发
        \texttt{setBreakpoints} 是被允许的，但约定俗成的做法是等它。
  \item \texttt{setBreakpoints} 每次带\textbf{该文件的全量断点数组}，不是增量；
        响应 \texttt{body.}\Bk\texttt{breakpoints} 里每项有 \texttt{verified}，
        false 表示这个位置根本没法停（常见于行号落在函数外、或没带调试信息）。
  \item \texttt{launch} 与 \texttt{attach} 是两种启动方式：前者由 adapter 拉起进程，
        后者接一个已在运行的 PID 或远程端口。
  \item \texttt{configurationDone} 是「本回合配置到此为止」的信号，adapter 收到它
        才会真正让程序跑起来；\texttt{stopAtEntry} 为 true 时停在入口发 \texttt{stopped}。
  \item \texttt{stopped} 事件的 \texttt{reason} 常见取值有 \texttt{breakpoint}、
        \texttt{step}、\texttt{pause}、\texttt{exception}、\texttt{entry}。
  \item 读变量要走三层：\texttt{stackTrace} 拿 \texttt{stackFrames}（每帧有
        \texttt{id}、\texttt{name}、\texttt{source}、\texttt{line}），
        \texttt{scopes} 按 \texttt{frameId} 拿作用域（\texttt{Local}、\texttt{Global}，
        每项带 \texttt{variables}\Bk\texttt{Reference} 与 \texttt{expensive}），
        \texttt{variables} 按 \texttt{variables}\Bk\texttt{Reference} 拿实际变量。
        \texttt{variables}\Bk\texttt{Reference} 为 0 表示这个值是叶子、不可再展开。
  \item 单步是三条独立请求：\texttt{stepNext}（跳过）、\texttt{stepIn}（进入）、
        \texttt{stepOut}（返回），都带 \texttt{threadId}；\texttt{continue} 恢复运行。
  \item \texttt{output} 事件的 \texttt{category} 有 \texttt{console}、\texttt{stdout}、
        \texttt{stderr}、\texttt{telemetry}；控制台面板与调试输出分流就是看这个字段。
  \item \texttt{terminated} 表示会话结束，之后只剩 \texttt{disconnect}；
        \texttt{restartSupported} 决定客户端是否显示「重启调试」。
\end{itemize}

\begin{note}[title={DAP 与 LSP 是两个进程，还是同一个扩展宿主}]
在 VS Code 里，语言服务器与调试 adapter 的进程模型是\textbf{不同}的：
\CD\ 是被扩展单独拉起的常驻子进程（一个工作区一个），而 DAP 的 adapter
常常就\textbf{活在扩展宿主进程里}——cppvsdbg、cppdbg 的控制端由扩展自己实现，
真正对着 gdb 说话的是它内嵌或另拉的 MI 客户端；CodeLLDB、\texttt{lldb-dap}
则是把 adapter 编成一个独立可执行文件、按需起停。所以「关掉编辑器 \CD\ 还活着」
和「调试会话一断 adapter 就退」同时成立，是两套生命周期，别按一个来推。
在 Neovim/Emacs 侧，DAP 客户端（如 nvim-dap）永远另起独立 adapter 进程，
与 LSP 服务器完全无关。
\end{note}

% --------------------------------------------------------------------
\section{C/C++ 侧的 adapter 生态}\label{sec:ch07-s2}

同一个「在 gdb/lldb 前面套一层 DAP」的活，生态里有几件不同的工具在做。

\textbf{cppdbg}：微软 C/C++ 扩展里的跨平台调试器，后端是 MI 模式
（\texttt{MIMode} 取 \texttt{gdb} 或 \texttt{lldb}），用 \texttt{miDebuggerPath}
指定调试器可执行文件。它的特色是 \texttt{setupCommands}：一串原样发给 MI 的
命令，用来开关漂亮打印、忽略某类错误、设置非停止信号。

\textbf{cppvsdbg}：只在 Windows 有效，走 VC 的 Native/Managed 调试引擎，
不需要 gdb/lldb，对 PDB 与 \texttt{\#pragma} 相关行为处理得最好；
配置字段与 cppdbg 有交集但不完全相同（它有 \texttt{consoleType}、\texttt{symbolOptions}）。

\textbf{CodeLLDB}（\texttt{vadimcn.vscode-lldb}）：把 LLDB 包成独立 adapter，
启动快、变量展开干净，macOS/Android 上比 cppdbg 省心；
配置字段是另一套（\texttt{terminal}、\texttt{env}、\texttt{preRunCommands}）。

\textbf{\texttt{lldb-dap}}：LLVM 官方的 LLDB 调试适配器（早期叫
\texttt{lldb-vscode}），与 \CD\ 同一个 release 里出货，纯 DAP、无扩展宿主，
命令行/脚本化驱动最方便。

\textbf{cdt-gdb-adapter}：Eclipse CDT 出的无头 gdb adapter，
给 Eclipse/Theia 产品线复用，也能被裸 DAP 客户端直接驱动。

{\small\setlength{\tabcolsep}{3pt}
\begin{longtable}{>{\raggedright\arraybackslash}p{0.14\textwidth} >{\raggedright\arraybackslash}p{0.20\textwidth} >{\raggedright\arraybackslash}p{0.18\textwidth} >{\raggedright\arraybackslash}p{0.24\textwidth}}
\caption{C/C++ 常用 debug adapter 的适用面}\label{tab:ch07-adapter}\\
\toprule
adapter & 后端 & 平台 & 定位与取舍 \\
\midrule
\endfirsthead
\toprule
adapter & 后端 & 平台 & 定位与取舍 \\
\midrule
\endhead
\bottomrule
\endfoot
cppdbg & gdb 或 lldb（MI 模式） & 跨平台 & 与 C/C++ 扩展一体，\texttt{setupCommands} 灵活；MI 解析偶尔卡大对象 \\
cppvsdbg & VC 调试引擎 & 仅 Windows & PDB、混合托管/原生体验最好；不可移植 \\
CodeLLDB & LLDB（自带构建） & mac/Linux/Win & 展开干净、启动快；配置字段自成一派 \\
\texttt{lldb-dap} & LLDB（LLVM 出货） & 跨平台 & 与 \CD\ 同版本、可裸脚本驱动 \\
cdt-gdb-adapter & gdb & 跨平台 & Eclipse/Theia 生态；独立进程 \\
\end{longtable}}

\texttt{launch.json} 里 cppdbg 的几个关键字段逐项（这是最容易记混的一层）：

\begin{itemize}
  \item \texttt{program}：被调试的可执行文件；\texttt{cwd}：进程工作目录，
        相对路径与「找得到数据文件」都靠它。
  \item \texttt{environment}：注入的环境变量数组，形如
        \texttt{\{ "name": "LD\_LIBRARY\_PATH", "value": "..." \}}；
        \texttt{args} 是命令行参数数组。
  \item \texttt{MIMode}：\texttt{gdb} 或 \texttt{lldb}；\texttt{miDebuggerPath}
        指对应调试器；\texttt{miDebugger}\Bk\texttt{ServerAddress} 则接远程 gdbserver。
  \item \texttt{setupCommands}：每项 \texttt{\{ description, text, ignoreFailures \}}，
        例如打开 libstdc++ 漂亮打印。
  \item \texttt{stopAtEntry}：进 main 就停；\texttt{externalConsole}：
        另开控制台窗口（Linux 下要 terminal x-terminal-emulator 一类）。
  \item \texttt{symbolLoadInfo}：控制符号加载策略（\texttt{mode} 取
        \texttt{loadAll} 或按模块过滤，带 \texttt{excludeList}），
        大型程序启动慢时用它裁剪。
  \item \texttt{preLaunchTask} 与 \texttt{postDebugTask}：调试前后各挂一个
        \texttt{tasks.json} 的 label，把「先构建再调试」接上（见 \ref{sec:ch07-s3}）。
\end{itemize}

% --------------------------------------------------------------------
\section{构建任务闭环：tasks.json 与 CMake 联动}\label{sec:ch07-s3}

编辑器里的「构建」在 VS Code 里就是 \texttt{tasks.json}。一个任务要么
\texttt{type} 为 \texttt{shell}（经 shell 解释，能用管道与通配），要么为
\texttt{process}（直接 exec，参数不再被 shell 改写）。
\texttt{group} 设为 \texttt{\{ "kind": "build", "isDefault": true \}} 后，
\texttt{Ctrl+Shift+B} 与 \texttt{preLaunchTask} 就能引用它。

下面三件套是能直接用的最小可用配置：构建任务、调试配置、工作区设置。

\begin{lstlisting}[style=jsonstyle,caption={.vscode/tasks.json：导出编译数据库并构建，带 problemMatcher},label={lst:ch07-tasks}]
{
  "version": "2.0.0",
  "tasks": [
    {
      "type": "shell",
      "label": "cmake configure",
      "command": "cmake",
      "args": ["-S", ".", "-B", "build",
               "-DCMAKE_EXPORT_COMPILE_COMMANDS=ON",
               "-DCMAKE_BUILD_TYPE=Debug"],
      "group": "build",
      "problemMatcher": ["$gcc"]
    },
    {
      "type": "shell",
      "label": "cmake build",
      "command": "cmake",
      "args": ["--build", "build", "--target", "app"],
      "group": { "kind": "build", "isDefault": true },
      "dependsOn": "cmake configure",
      "problemMatcher": ["$gcc"]
    }
  ]
}
\end{lstlisting}

\begin{lstlisting}[style=jsonstyle,caption={.vscode/launch.json：cppdbg 跑本地 gdb，调试前先构建},label={lst:ch07-launch}]
{
  "version": "0.2.0",
  "configurations": [
    {
      "name": "(gdb) launch app",
      "type": "cppdbg",
      "request": "launch",
      "program": "${workspaceFolder}/build/app",
      "args": ["-i", "sample.txt"],
      "stopAtEntry": false,
      "cwd": "${workspaceFolder}",
      "environment": [],
      "externalConsole": false,
      "MIMode": "gdb",
      "miDebuggerPath": "/usr/bin/gdb",
      "setupCommands": [
        { "description": "enable pretty printing",
          "text": "-enable-pretty-printing", "ignoreFailures": true },
        { "description": "ignore SIGPIPE",
          "text": "handle SIGPIPE nostop noprint", "ignoreFailures": true }
      ],
      "preLaunchTask": "cmake build"
    }
  ]
}
\end{lstlisting}

\begin{lstlisting}[style=jsonstyle,caption={.vscode/settings.json：把保存动作与 CMake Tools 固定下来},label={lst:ch07-settings}]
{
  "cmake.configureOnOpen": true,
  "cmake.buildDirectory": "${workspaceFolder}/build",
  "C_Cpp.errorSquiggles": "disabled",
  "clangd.path": "clangd",
  "clangd.arguments": [
    "--compile-commands-dir=build",
    "--background-index",
    "--header-insertion=never"
  ],
  "editor.formatOnSave": true,
  "editor.defaultFormatter": "ms-vscode.cpptools",
  "editor.rulers": [100],
  "files.candidateEndOfLine": { "sourceFile": "lf" }
}
\end{lstlisting}

CMake Tools 与 \texttt{CMakePresets.json}（第 16 章）是联动的：装了 presets 之后
\texttt{cmake.configureOnOpen} 会按 preset 初始化，状态栏上依次是
configure / build / test / launch / debug 五个命令，点 debug 会走
「preset 里的调试器 + tasks」而不是裸 \texttt{launch.json}。
两边都能配置时，\textbf{以 presets 为事实源}，\texttt{launch.json} 只留一个引用，
避免同一份程序路径维护两处。

\begin{warning}[title={preLaunchTask 失败时调试器静默不启动}]
最常见的困惑是「按 F5 没反应，也不报错」。十有八九是 \texttt{preLaunchTask}
那个任务失败或 label 拼错了：构建一停，调试就\textbf{根本不启动}，而错误只出现在
「终端 / 输出」面板里，调试面板一无所获。排查步骤：先在任务里手动跑一次
\texttt{cmake build}；确认 \texttt{launch.json} 的 \texttt{preLaunchTask}
字符串与 \texttt{tasks.json} 的 \texttt{label} 逐字相等（含空格）；
任务 \texttt{type} 是 \texttt{shell} 时留意 Windows 与 POSIX 对引号处理不同。
临时验证可先把 \texttt{preLaunchTask} 删掉，确认程序本身能起，再接回去。
\end{warning}

% --------------------------------------------------------------------
\section{problemMatcher：把编译器诊断映射进编辑器}\label{sec:ch07-s4}

problemMatcher 干的事很朴素：\textbf{拿正则扫任务输出，把一行诊断变成一个可点击的
标记}。字段构成：

\begin{itemize}
  \item \texttt{owner}：这条 matcher 的归属名，用来去重与并列展示
        （多个 matcher 只有 owner 相同时才互相覆盖）。
  \item \texttt{fileLocation}：\texttt{absolute} 或 \texttt{relative}
        （相对路径以某目录为基准，常配
        \texttt{["relative", "\$\{workspaceFolder\}"]}）。
  \item \texttt{pattern}：一组正则，每个捕获组通过命名约定映射到
        \texttt{file}、\texttt{line}、\texttt{column}、\texttt{message}、
        \texttt{code}、\texttt{kind}（kind 里 \texttt{warning}/\texttt{error} 决定颜色）。
        跨多行诊断用 pattern 数组，每行一条正则。
  \item \texttt{source}：诊断来源名，会显示在标记上（如 \texttt{gcc}、\texttt{clang-tidy}）。
  \item 内建 matcher：\texttt{\$gcc}、\texttt{\$clang}（由 C/C++ 扩展贡献，
        处理 clang 的 \texttt{file:line:col: warning: msg [-W...]}）、
        \texttt{\$msCompile}（MSVC 的 \texttt{file(line): error C1234: msg}）。
\end{itemize}

\Tidy\ 的输出格式和编译器几乎一样，所以也能做成 matcher——
这是把静态分析接进「一键任务」的关键（增量格式化与 CI 侧的做法见第 10、16 章）：

\begin{lstlisting}[style=jsonstyle,caption={自定义 problemMatcher：把 clang-tidy 诊断映射成可点击标记},label={lst:ch07-tidy-matcher}]
{
  "version": "2.0.0",
  "tasks": [
    {
      "type": "shell",
      "label": "clang-tidy current TU",
      "command": "clang-tidy",
      "args": [
        "-p", "build", "src/parser.cpp",
        "--checks=bugprone-*,performance-*"
      ],
      "problemMatcher": {
        "owner": "clang-tidy",
        "fileLocation": ["relative", "${workspaceFolder}"],
        "pattern": [
          {
            "regexp": "^(.+):(\\d+):(\\d+): (\\w+): (.+)\\[(.+?)\\]",
            "file": 1, "line": 2, "column": 3,
            "severity": 4, "message": 5, "code": 6
          }
        ]
      }
    }
  ]
}
\end{lstlisting}

注意 \texttt{clang-tidy} 会把「来自 \texttt{\#include} 的 note」也打出来，
一条诊断常占多行；上面的单行 pattern 只抓主行，够用但会丢掉附带说明。
要全抓就在 pattern 数组里补第二、第三条正则，并把非捕获字段留空。

% --------------------------------------------------------------------
\section{.editorconfig：人机分工的分界线}\label{sec:ch07-s5}

\texttt{.editorconfig} 是给\textbf{人用的编辑器}看的，字段极少，语义极窄：

\begin{itemize}
  \item \texttt{root = true}：声明查找到此为止，不再往上级目录找。
  \item \texttt{indent\_style}：\texttt{space} 或 \texttt{tab}；
        \texttt{indent\_size}：每级缩进空格数（或 \texttt{tab} 表示跟随 tab\_width）；
        \texttt{tab\_width}：一个 tab 显示成几格。
  \item \texttt{max\_line\_length}：行宽软上限。
  \item \texttt{end\_of\_line}：\texttt{lf} 或 \texttt{crlf}；
        \texttt{insert\_final\_newline}：文件末尾是否留换行；
        \texttt{trim\_trailing\_whitespace}：是否去行尾空格；
        \texttt{charset}：\texttt{utf-8}、\texttt{latin1} 等。
  \item 选择器用 glob：\texttt{[*]} 全仓生效，\texttt{[*.c]}、\texttt{[*.h]}、
        \texttt{[*.\{c,h\}]} 按类型覆盖。
\end{itemize}

\Fmt\ 认识 \texttt{.editorconfig}：它会按当前文件名去匹配 section，读取
\texttt{indent\_style}、\texttt{indent\_size}、\texttt{tab\_width}、
\texttt{max\_line\_length} 这几项，映射到 \texttt{UseTab}、\texttt{IndentWidth}、
\texttt{TabWidth}、\texttt{ColumnLimit}。想显式以它为准，可以把 \texttt{BasedOnStyle}
写成 \texttt{EditorConfig}（较新版本的 \Fmt\ 才支持）。\textbf{冲突时谁赢}：
在同一个 section 里，\texttt{.clang-format} 显式写的字段优先于 editorconfig 推出来的值；
editorconfig 只填那些 \texttt{.clang-format} 没明说的字段。

为什么这样分工：\texttt{.editorconfig} 表达的是「这个仓库人眼看舒服的样子」——
制表符还是空格、换行用哪种、行尾清不清，这些都是编辑器的行为，跟语法树无关。
而 \texttt{.clang-format} 表达的是「机器按 AST 该怎么排版」——断不断行、
参数对齐、include 分组，它需要对语法的理解。把「人写的规则」放 editorconfig、
把「机器排版」放 \texttt{.clang-format}，两边的关注点就不会互相打架。

{\small\setlength{\tabcolsep}{3pt}
\begin{longtable}{>{\raggedright\arraybackslash}p{0.28\textwidth} >{\raggedright\arraybackslash}p{0.20\textwidth} >{\raggedright\arraybackslash}p{0.32\textwidth}}
\caption{六份配置文件：谁读它、缺了会怎样}\label{tab:ch07-files}\\
\toprule
文件 & 职责 & 缺失后果 \\
\midrule
\endfirsthead
\toprule
文件 & 职责 & 缺失后果 \\
\midrule
\endhead
\bottomrule
\endfoot
\texttt{compile\_commands.json} & 每个 TU 的真实命令行；\CD\ 与 \Tidy\ 共读 & \CD\ 退回 fallback flags，头找不到、宏不对，跳转与诊断大面积失准 \\
\texttt{launch.json} & 告诉客户端怎么起 adapter、调试哪个程序 & 无法按 F5；每次手敲命令行调 \\
\texttt{tasks.json} & 构建/分析任务与其诊断 matcher & \texttt{preLaunchTask} 无处引用，改完不会自动构建 \\
\texttt{c\_cpp\_properties.json} & 给闭源 IntelliSense 手写编译环境 & cpptools 用默认探测，跨平台/交叉编译时误报 \\
\texttt{.clang-format} & 机器排版的逐字段规则 & \Fmt\ 走 \texttt{Style: LLVM} 等内置默认，团队风格散架 \\
\texttt{.clangd} & \CD\ 的编译 flag 增删、诊断策略、索引开关 & 命令行里的 flag 无法覆盖，\Tidy\ 集成与 include 插入策略回默认 \\
\end{longtable}}

\begin{lstlisting}[style=yamlstyle,caption={一份与 .clang-format 分工清晰的 .editorconfig},label={lst:ch07-editorconfig}]
root = true

[*]
charset = utf-8
end_of_line = lf
insert_final_newline = true
trim_trailing_whitespace = true
indent_style = space
indent_size = 2

[*.{c,cc,cpp,h,hpp}]
indent_size = 4
tab_width = 4
max_line_length = 100

[*.md]
trim_trailing_whitespace = false
\end{lstlisting}

% --------------------------------------------------------------------
\section{保存时动作链}\label{sec:ch07-s6}

一次「Ctrl+S」在装了这套工具之后，会串行触发好几件事，顺序决定了你会看到什么现象：

\begin{enumerate}
  \item 编辑器先发 \texttt{textDocument/}\Bk\texttt{willSaveWaitUntil}
        （若服务器声明了该能力），
        \textbf{同步等}服务器返回要改的位置，客户端把改动应用后再落盘。
  \item 落盘前/后依客户端设置跑格式化：\texttt{editor.formatOnSave}
        触发 \texttt{textDocument/}\Bk\texttt{formatting}，默认格式化器由
        \texttt{editor.}\Bk\texttt{defaultFormatter} 指定。
  \item \texttt{editor.codeActionsOnSave} 里的 \texttt{source.fixAll}
        触发「一次性应用所有 code action」；想只让 \CD\ 的 \Tidy\ 修复参与，
        用更窄的键（如 \texttt{source.fixAll.clangd}），值可以是
        \texttt{explicit} 或 \texttt{always}。
  \item 落盘后发 \texttt{didSave}，服务器重新解析并推诊断。
\end{enumerate}

风险在第 1 步：\texttt{willSaveWaitUntil}\Bk\ 是\textbf{阻塞}保存的，客户端有超时
（VS Code 默认约三秒）。如果服务器在这一路径里做了重活——重新解析整个 TU、
跑一遍慢的检查——那么超时一到，客户端就\textbf{放弃服务器给的改动}直接保存，
你会得到「修复有时生效有时不生效」的薛定谔体验；更糟的是网络/进程卡顿时
每次保存都等满超时，手感极差。原则：\texttt{willSaveWaitUntil}\Bk\ 只做便宜的、
局部的改动，重活留给显式的 code action。

其余两端各一句：Neovim 用 \texttt{BufWritePre} 自动命令手动调
\texttt{vim.lsp.buf.}\Bk\texttt{formatting\_sync()}（这就是同步的
\texttt{willSaveWaitUntil}），格式化则靠 LSP 客户端插件或 \texttt{:Format};
Emacs 的 \texttt{eglot} 提供 \texttt{eglot-}\Bk\texttt{format-on-save} 一类的 hook，
\texttt{M-x} 侧走 \texttt{before-save-hook}。机制名字不同，串的都是同一条链。

\begin{guideline}[title={把保存链固定成一条可复现的路径}]
团队里最好只留\textbf{一条}保存链：一个默认格式化器、一组
\texttt{codeActionsOnSave}、一个 \CD\ 实例。把「谁格式化」写死在
\texttt{settings.json} 并提交，否则两个人会得出两种结果。格式化与
\texttt{fixAll} 的先后要固定——先修复再排版，或反过来都行，但必须一致。
\end{guideline}

\begin{bestpractice}[title={调试配置只放仓库根，个人路径放 user 级}]
把 \texttt{launch.json}、\texttt{tasks.json} 提交到仓库的
\texttt{.vscode/}，因为它们对\textbf{所有人}都该一样：目标名、构建命令、
matcher 规则。凡是带个人绝对路径的东西（\texttt{miDebuggerPath} 指向你机器上的
某个 gdb、临时加宽的 \texttt{args}、你自己的环境变量）放进
\textbf{user 级 settings} 或 \texttt{launch.json} 的
\texttt{compounds}/user 配置，别写进仓库那份。
需要按机器区分时，用 \texttt{\$\{workspaceFolder\}}、
\texttt{\$\{env:VAR\}} 这类变量代替硬路径；真要分平台就在
\texttt{launch.json} 里配多个 \texttt{name}，而不是维护两份文件。
\end{bestpractice}

\begin{compare}[title={DAP 与 LSP 的握手对照}]
\begin{itemize}
  \item LSP：客户端发 \texttt{initialize}（带 \texttt{capabilities}）
        $\to$ 服务器回 \texttt{capabilities} $\to$ 客户端发 \texttt{initialized}
        通知。之后才有 \texttt{didOpen} 与各种请求。
  \item DAP：客户端发 \texttt{initialize}（带客户端能力）$\to$ adapter 回
        自身能力 $\to$ adapter 发 \texttt{initialized} \emph{事件}。
        之后客户端下断点、发 \texttt{launch}/\texttt{attach}、再 \texttt{configurationDone}。
  \item 方向差异很关键：LSP 的 \texttt{initialized} 是\textbf{客户端发出去}的通知，
        DAP 的 \texttt{initialized} 是\textbf{服务器发出来}的事件。把它俩记反，
        写客户端时序时必错（第 8 章会真的手写一遍）。
\end{itemize}
\end{compare}

闭环到这里就完整了：编辑（LSP）$\to$ 保存动作链（格式化与修复）$\to$
构建（tasks/CMake）$\to$ 诊断回填（problemMatcher 与 \CD\ 推送）$\to$
调试（DAP）。下一站，第 8 章把这套协议拆开——自己写一个客户端与服务器，
把上面所有「约定」变成会跑的行。

% ---- frag_ch08.tex ----
\chapter{自己动手：一个最小 LSP 客户端与服务器}\label{ch:write-your-own}

前两章把「用现成服务器」讲透了。本章反过来：亲手把 LSP 的两端各写一遍。
不是为了替代 \CD，而是因为\textbf{只有写过一次帧编解码、发过一次
\texttt{initialize} 握手、被一个 UTF-8 字节数坑过}，你才真正知道协议里哪根线
牵着哪盏灯。全章 \CPP\ 只依赖标准库与 POSIX/Win32 系统调用，
服务器侧才碰 clang 的头。

% --------------------------------------------------------------------
\section{客户端骨架：进程、帧、id 与时序}\label{sec:ch08-s1}

一个能跑的最小客户端由五块拼成，完整程序约一百五十行，划分如下：

{\small\setlength{\tabcolsep}{3pt}
\begin{longtable}{>{\raggedright\arraybackslash}p{0.24\textwidth} >{\raggedright\arraybackslash}p{0.58\textwidth}}
\caption{最小 LSP 客户端的模块划分与大致行数}\label{tab:ch08-mod}\\
\toprule
模块 & 职责（约行数） \\
\midrule
\endfirsthead
\toprule
模块 & 职责（约行数） \\
\midrule
\endhead
\bottomrule
\endfoot
\texttt{spawn\_server()} & 起子进程、接两根管道：POSIX \texttt{fork}/\texttt{execvp}，Windows \texttt{CreateProcess}（约 40 行） \\
\texttt{frame\_encode()} / \texttt{frame\_read()} & \texttt{Content-Length} 帧的编解码，长度按 UTF-8 字节数（约 30 行） \\
\texttt{request()} / \texttt{notify()} & 分配自增 id、拼 JSON、写帧；通知不带 id（约 25 行） \\
读循环线程 & 收帧、按 id 落进 \texttt{pending\_}、回调通知与服务器发起的请求（约 35 行） \\
握手与 \texttt{main} & \texttt{initialize} 到 \texttt{initialized} 到 \texttt{didOpen} 的时序（约 20 行） \\
\end{longtable}}

\subsection{起子进程}\label{sec:ch08-s1a}

\CD\ 默认只讲 stdio：它从标准输入读帧、往标准输出写帧。所以「连上服务器」
等价于「把一个子进程的 stdin/stdout 接到你手上」。两条平台的差别用注释写清：

\begin{lstlisting}[style=cppstyle,caption={起 \CD\ 子进程并把 stdio 接到管道（POSIX 实现，Windows 差异见注释）},label={lst:ch08-spawn}]
#include <fcntl.h>
#include <unistd.h>
#include <sys/wait.h>
#include <cstdlib>

struct ServerPipe { int to_stdin; int from_stdout; };

ServerPipe spawn_server(const char* const argv[]) {
  int in[2], out[2];
  if (pipe(in) != 0 || pipe(out) != 0) { /* 失败处理 */ }
  pid_t pid = fork();
  if (pid == 0) {                       // 子进程：clangd 自己
    dup2(in[0],  STDIN_FILENO);         // 它的 stdin = 管道读端
    dup2(out[1], STDOUT_FILENO);        // 它的 stdout = 管道写端
    close(in[0]); close(in[1]);
    close(out[0]); close(out[1]);
    execvp(argv[0], const_cast<char* const*>(argv));
    _exit(127);                          // exec 失败才会走到这里
  }
  close(in[0]); close(out[1]);
  return ServerPipe{in[1], out[0]};      // 父进程：往 in[1] 写、从 out[0] 读
  // Windows 差异：没有 fork/execvp。改用 CreateProcessW，
  // 用 SECURITY_ATTRIBUTES 建匿名管道，令 bInheritHandles=TRUE，
  // 并把 STARTUPINFOW 的 hStdInput / hStdOutput 指到管道两端。
}
\end{lstlisting}

\subsection{帧编解码}\label{sec:ch08-s1b}

LSP 的 base 协议只有一层封装：\texttt{Content-Length: <字节数>}\texttt{\textbackslash r\textbackslash n}
一个头，紧跟一个空行 \texttt{\textbackslash r\textbackslash n}，然后是 JSON 体。
\textbf{长度是 UTF-8 字节数，不是字符数}——这是本章反复要强调的第一号坑。

\begin{lstlisting}[style=cppstyle,caption={帧编解码核心：读满 Content-Length 个字节，按字节而非字符},label={lst:ch08-frame}]
#include <string>
#include <istream>
#include <cstddef>

std::string frame_encode(const std::string& utf8_body) {
  std::string out;
  out += "Content-Length: ";
  out += std::to_string(utf8_body.size());   // size() 已是字节数
  out += "\r\n\r\n";                          // 头与体之间那个空行
  out += utf8_body;
  return out;
}

bool frame_read(std::istream& in, std::string& body) {
  std::size_t len = 0;
  std::string line;
  while (std::getline(in, line)) {
    if (line.empty()) break;                  // 空行 = 头部结束
    if (!line.empty() && line.back() == '\r') line.pop_back();
    const std::string key = "Content-Length:";
    if (line.rfind(key, 0) == 0)
      len = static_cast<std::size_t>(std::stoul(line.substr(key.size())));
  }
  body.assign(len, '\0');
  in.read(&body[0], static_cast<std::streamsize>(len));  // 读恰好 len 字节
  return static_cast<std::size_t>(in.gcount()) == len;
}
\end{lstlisting}

\subsection{请求—响应循环}\label{sec:ch08-s1c}

请求带 \texttt{id}，响应靠同一个 \texttt{id} 配回去；通知不带 \texttt{id}，
也\textbf{不进待匹配队列}。响应到达顺序不等于发出顺序，所以匹配必须靠 id、
不能靠先进先出。

\begin{lstlisting}[style=cppstyle,caption={请求分配自增 id、写帧；通知不入队；读循环另起线程},label={lst:ch08-loop}]
#include <atomic>
#include <map>
#include <mutex>
#include <string>

class LspClient {
public:
  int request(const std::string& method, const std::string& params) {
    const int id = next_id_.fetch_add(1);
    std::string body = "{\"jsonrpc\":\"2.0\",\"id\":" + std::to_string(id)
                     + ",\"method\":\"" + method + "\",\"params\":" + params + "}";
    write_frame(body);
    return id;
  }
  void notify(const std::string& method, const std::string& params) {
    std::string body = "{\"jsonrpc\":\"2.0\",\"method\":\"" + method
                     + "\",\"params\":" + params + "}";
    write_frame(body);            // 无 id，因此绝不登记到 pending_
  }
  void read_loop() {
    std::string body;
    while (frame_read(stdin_stream_, body)) {
      const int id = parse_id(body);                 // 有 id => 响应
      if (id != 0) {
        std::lock_guard<std::mutex> g(mu_);
        pending_.emplace(id, std::move(body));       // 按 id 落位，不按到达序
      } else {
        dispatch_notification(body);                 // 通知 / 服务器发起的请求
      }
    }
  }
private:
  void write_frame(const std::string& body);
  int parse_id(const std::string& body);
  void dispatch_notification(const std::string& body);
  std::atomic<int> next_id_{1};
  std::mutex mu_;
  std::map<int, std::string> pending_;
  std::istream& stdin_stream_;
};
\end{lstlisting}

握手时序要背下来，写错必挂：客户端先\textbf{请求} \texttt{initialize}
（带 \texttt{capabilities}），服务器\textbf{响应}它并回自己的 capabilities；
随后客户端发 \texttt{initialized} 这条\textbf{通知}（注意方向：LSP 的
\texttt{initialized} 是客户端发出的，和第 7 章 DAP 里那条同名\emph{事件}正好相反）；
此后才有 \texttt{didOpen}、才有各种 \texttt{textDocument/} 请求。
关闭走 \texttt{shutdown} 请求再 \texttt{exit} 通知。

\begin{guideline}[title={先写客户端，再接真服务器}]
顺序很关键：先只写\textbf{客户端}，对着一个你能完全控制的假服务器（几行脚本、
或干脆一个把收到的帧原样回显的进程）跑通帧编解码、id 匹配和握手时序，
\emph{然后}才把 \texttt{execvp} 换成真 \CD。这样每一次协议假设
（长度到底按字节还是字符、\texttt{initialized} 谁发）都能被你单独证伪，
不会被服务器的语义复杂度淹没。
\end{guideline}

% --------------------------------------------------------------------
\section{不写代码也能验证协议：手工喂报文}\label{sec:ch08-s2}

\CD\ 提供两条不写客户端就能观察协议的路。一条是 \Opt{check}：
它把单个 TU 从头解析一遍并把诊断打到标准输出，不走 LSP，适合快速确认
「这个文件的编译命令对不对、头找没找到」。另一条是 \Opt{log}=verbose，
把内部每个请求、每次 AST 重建都写进日志。而想看\textbf{真实报文交互}，
最省事的办法是用管道手工喂帧：

\begin{lstlisting}[style=shellstyle,caption={用 printf 造 Content-Length 头，从管道喂一条 initialize 给 clangd},label={lst:ch08-pipe}]
# 手工组一条 initialize，观察 clangd 的握手响应。
# 用 wc -c 取字节数（别用 ${#var} 那种含 # 的写法，会和 JSON 里的字符混淆）。
body='{"jsonrpc":"2.0","id":1,"method":"initialize","params":{"capabilities":{}}}'
{
  printf 'Content-Length: %d\r\n\r\n' "$(printf '%s' "$body" | wc -c)"
  printf '%s' "$body"
  sleep 1
} | clangd --log=verbose 2>clangd.log

# 只想看单文件解析结果，不必走 LSP：
clangd --check=src/parser.cpp --compile-commands-dir=build
\end{lstlisting}

要点：\texttt{printf} 负责\textbf{字节}，管道负责\textbf{送}；\texttt{sleep} 一下
防止你在 \CD\ 还没读完就把 stdin 关掉。诊断落在 \texttt{clangd.log} 里。

% --------------------------------------------------------------------
\section{服务器方向的两条路}\label{sec:ch08-s3}

要「自己产出语义结果」，取决于你是要复用 clangd 还是要自己算。

\textbf{路 (a)：复用 \CD\ 的服务器 API。}\Lib{clang} 的
\texttt{clang-tools-extra/clangd} 把服务层做成了可编程对象：核心类
\texttt{clang::clangd::ClangdServer} 负责持有工作区、调度请求；
行为由一个配置对象（\texttt{Config}）描述——后台索引开不开、
inlay hints、\Tidy\ 钩子等都在那里设；对外能力则通过一组回调注册，
比如自定义的 code action、动态注册的能力。思路是：你不重写解析，
只把 \texttt{ClangdServer} 包在自己的 LSP 传输层后面。
\textbf{这条路的代价是 ABI 与稳定性}：它是 \Lib{clang} 内部头，
不承诺跨版本接口不变，跟进 LLVM 升级时要改代码；措辞上应把它当
「可参考、可 fork，但别当稳定 API」。

\textbf{路 (b)：用 \texttt{clang::ast\_matchers} 自己产诊断，再包一层 LSP。}
这条路更可控：你写匹配规则、写回调，把回调命中直接翻译成一条诊断，
再由 LSP 层 \texttt{publish}\Bk\texttt{Diagnostics} 推给编辑器。核心是
\texttt{MatchFinder} 注册「模式\(\to\)回调」，回调继承
\texttt{MatchFinder::MatchCallback}：

\begin{lstlisting}[style=cppstyle,caption={MatchFinder 注册与回调：命中裸指针返回的函数就产一条诊断},label={lst:ch08-matcher}]
#include "clang/ASTMatchers/ASTMatchFinder.h"

using namespace clang;
using namespace clang::ast_matchers;

class RawPtrReturn : public MatchFinder::MatchCallback {
public:
  void run(const MatchFinder::MatchResult& Result) override {
    const FunctionDecl* Fn = Result.Nodes.getNodeAs<FunctionDecl>("fn");
    if (!Fn) return;
    DiagnosticsEngine& DE = *Result.Diags;
    DE.Report(Fn->getLocation(), DE.getCustomDiagID(
        DiagnosticsEngine::Warning, "function returns a raw pointer"))
        << Fn->getSourceRange();
  }
};

// 把回调挂到 MatchFinder：命中"返回类型为裸指针"的函数声明。
void register_matchers(MatchFinder& Finder, RawPtrReturn& Callback) {
  Finder.addMatcher(
      functionDecl(returns(pointerType())).bind("fn"),
      &Callback);
}
\end{lstlisting}

\texttt{bind("fn")} 把命中的节点存进 \texttt{Result.Nodes}，回调里再用同名键取出——
这是 matcher 与 callback 之间唯一的契约，名字写错就永远取到 \texttt{nullptr}。

\begin{warning}[title={不要在客户端里做语义计算}]
新手最容易犯的架构错误，是把「这段代码是什么意思」的判断写进客户端：
在编辑器侧用正则数括号、靠缩进推断作用域、用文本搜索冒充引用查找。
症状是同一份代码在不同客户端结论不同、改名改坏、补全忽多忽少。
\textbf{语义属于服务器，客户端只负责呈现与传帧}。客户端该做的只有：
编码 UTF-16 偏移、维护打开文档的文本、把用户动作翻成请求。越界的
文本推断，早晚会和服务器打架。
\end{warning}

% --------------------------------------------------------------------
\section{一定会踩的坑}\label{sec:ch08-s4}

下面每一个都是「写完第一版客户端、接上真 \CD\ 之后」才会暴露的坑，
配上最小复现思路，逐条对应到你上面的代码。

\textbf{一、UTF-16 偏移与中文注释错位。}LSP 的位置用的是 UTF-16 码元计数
（见第 4 章），不是字节、不是码点。一行里只要有非 BMP 之外的中文，
\texttt{strlen} 得到的字节数是 UTF-8 的，字符数和 UTF-16 单元数又各不相同；
用字节数当列偏移，光标就会往右漂。复现：在某行放三个汉字，让服务器报
「第 4 列」，你的客户端按第 4 字节定位，落点差一整块。修复：偏移一律换算成
UTF-16 单元数再比较。

\textbf{二、\texttt{id} 复用与并发响应乱序。}若把 \texttt{id} 设成 32 位、
或每次握手重置计数器，长跑后会撞上「两个在途请求同 id」；更常见的是
\texttt{id} 只增但响应按到达顺序处理。复现：连发五个 \texttt{hover}，
第一个最慢，若按 FIFO 配对，第二到第五的浮窗内容全错位。修复：响应只看
\texttt{id}、从 \texttt{pending\_} 里取，绝不假设有序。

\textbf{三、\texttt{didChange} 增量合并顺序。}\texttt{TextDocumentContentChangeEvent}
数组要\textbf{按序}应用，前一条改了字节区间，后一条的 \texttt{range} 是相对
前一条之后的新文档。复现：一次保存产生两条 range 编辑，先应用第二条再第一条，
文本就串行错位。修复：严格按数组顺序 \texttt{splice}，并处理
\texttt{contentChanges} 不带 range（表示整篇替换）的特例。

\textbf{四、push 与 pull 诊断混用。}经典是服务器用
\texttt{publish}\Bk\texttt{Diagnostics} 主动推；较新的规范又加了
\texttt{textDocument/diagnostic} 拉取。两边同时开，会出现「一条诊断被推两遍、
或者清不掉」。复现：打开诊断拉取又留着推送回调，删掉一个错误后旧标记仍在。
修复：二选一；若都用，以拉取为准、把推送的 \texttt{uri} 去重。

\textbf{五、服务器崩溃后客户端收不到 \texttt{exit}。}\CD\ 段错误或被 OOM 杀掉时，
管道 EOF 但没有任何 \texttt{exit} 通知。复现：\texttt{kill -9} 服务器，客户端
卡在 \texttt{read}，界面显示「仍在索引」永不返回。修复：读循环把
\texttt{frame\_read} 返回 false 视为连接已断，主动进入重启逻辑。

\textbf{六、\texttt{initialize} 里 capabilities 谎报。}客户端没实现
\texttt{semanticTokens} 却在握手时声明支持，服务器就会猛推
\texttt{semanticTokens/}\Bk\texttt{full}\Bk\texttt{/delta}，你把不认识的帧丢进
解析器直接崩。复现：把 \texttt{capabilities} 里 \texttt{dynamicRegistration}
设 true 但没注册处理器。修复：只声明你真的能处理的项，宁可少报。

{\small\setlength{\tabcolsep}{3pt}
\begin{longtable}{>{\raggedright\arraybackslash}p{0.19\textwidth} >{\raggedright\arraybackslash}p{0.24\textwidth} >{\raggedright\arraybackslash}p{0.19\textwidth} >{\raggedright\arraybackslash}p{0.20\textwidth}}
\caption{协议坑：症状、定位手段与修复}\label{tab:ch08-pitfall}\\
\toprule
坑 & 症状 & 定位手段 & 修复 \\
\midrule
\endfirsthead
\toprule
坑 & 症状 & 定位手段 & 修复 \\
\midrule
\endhead
\bottomrule
\endfoot
UTF-16 偏移 & 光标/标记往右漂 & 中文行里比对列号与字节数 & 偏移换成 UTF-16 单元 \\
id 复用乱序 & 浮窗内容串号 & 连发同类请求看是否错位 & 按 id 配对，不信到达序 \\
didChange 顺序 & 文本越改越乱 & 回放两条 range 编辑 & 按数组顺序合并、处理整替 \\
push 与 pull 混用 & 诊断重复或清不掉 & 看是否两套回调都在推 & 二选一，或拉取为准去重 \\
崩溃无 exit & 界面卡「索引中」 & kill -9 后看是否收 EOF & 读失败即判断连并重启 \\
capabilities 谎报 & 未知帧导致崩溃 & 关掉未实现的声明复测 & 只声明真能处理的项 \\
\end{longtable}}

\begin{note}[title={用回归测试把协议行为钉死}]
LLVM 里 \texttt{clang-tools-extra} 有实验性的 LSP 客户端，\CD\ 自身也带一套
测试装置（基于 gtest 的单测，加上「喂一段请求、比一段期望响应」的
\texttt{.expect} 式夹具）。思路值得抄：\textbf{把每一条协议约定固化成一个
回归用例}——一段真实输入报文，配一份期望输出报文，逐字节比。将来谁改了
握手、改了编码，测试先炸，而不是等用户报「补全又坏了」。你自写的客户端同理：
先写「输入帧序列 \(\to\) 期望请求序列」的黄金用例，再谈重构。
\end{note}

\begin{bestpractice}[title={给每个请求打点往返时间}]
在请求发出时记 \texttt{id} 对应的起始时间戳，收到响应时算差值，按方法名
分桶（\texttt{completion}、\texttt{hover}、\texttt{diagnostic}）。性能问题
几乎总先在这里露头：某方法 P99 突然从几毫秒涨到几百毫秒，往往不是网络，
而是服务器在重解析、或你在 \texttt{willSaveWaitUntil}\Bk\ 里做了重活（见第 7 章）。
打点本身要做成无锁或有界缓冲，别为了测时延把热路径阻塞。
\end{bestpractice}

写完这两端，你就把「第 4 章的报文」和「第 5、6 章的服务器」缝在了一起：
客户端不再是个黑盒扩展，服务器产出的诊断也不再是天上掉下来的波浪线。
下一章转入 Part III，开始逐字段拆 \Fmt——那是一台按 AST 排版的机器，
而你现在已经知道它靠什么拿到那棵 AST。

% ---- frag_ch09.tex ----
% ==================== 第 9 章 ====================
\part{格式化}

\chapter{\Fmt\ 机制与 \texttt{.clang-format} 逐字段}\label{ch:clangformat}

本章把 \Fmt\ 拆开讲：它按什么顺序找配置、每个字段到底改变什么、哪些字段值得团队协商、
哪些字段只在特定代码形状下才看得出差别。读完你应该能在不查文档的情况下，对一段别扭的代码
说出「是哪个字段在管它」，也应该清楚 \Fmt\ 的能力边界在哪里。
第 10 章讨论如何把这些配置变成机器可执行的门禁，第 11 章把它和别的格式化工具放在一起比较。

\section{处理流程：从字节流到最短路径}\label{sec:ch09-s1}

\Fmt\ 一次运行内部大致是四步。理解它们的顺序，就理解了它为什么又强又脆。

\begin{enumerate}
  \item \textbf{词法切分}。它复用 \CL/\CPP\ 的词法器把源文件切成
    \Tp{FormatToken} 序列，每个 token 记住自己的原始空白、是否带行尾注释、
    是否属于宏、是否跨行。这一步\textbf{不做语法分析}，也没有类型信息。
  \item \textbf{语境标注}（\Tp{TokenAnnotator}）。逐个 token 判定它落在什么语境里：
    这个 \texttt{*} 是指针声明还是乘法，这个 \texttt{>} 是模板右尖括号还是移位，
    连写的 \texttt{>>} 要不要拆，\texttt{\{} 是语句块、初始化列表还是 lambda 体，
    \texttt{:} 是位域、访问修饰符、初始化列表还是三目。
  \item \textbf{分段处理}（handler）。真正排版之前，几个 handler 先把结构性空白摆好：
    命名空间成员要不要整体缩进、类里 public/private 段的缩进偏移
    （\Tp{ClassLayout}）、构造函数的初始化列表在哪儿断
    （\Tp{InitListHandler}，\texttt{BreakConstructorInitializers} 在这里生效）。
  \item \textbf{按候选求最短路径}。对每一个必须断行的逻辑行，它枚举所有合法的换行位置，
    给每种组合算代价（penalty，较新版本里陆续改叫 cost），再用最短路求解器挑代价最小的
    一组。早期版本的入口叫 \Tp{ColumnLayout}，现在拆成逐行匹配加线性求解，思路不变：
    \textbf{排版是一次带权搜索，不是一堆 if-else}。
\end{enumerate}

因为第 4 步是搜索，所以「在这里多一个空格，后面每段就更短」这类改动会牵动整行的决策。
这也是 \Fmt\ 的输出看起来「有道理但难预测」的原因：它不是「遇到逗号换行」这种规则表，
而是在 \texttt{ColumnLimit} 约束下最小化惩罚。

\subsection*{三个关键推论}
\begin{itemize}
  \item \textbf{它不需要编译数据库}。第 2 章里 \CD\ 依赖
    \texttt{compile\_commands.json}，而 \Fmt\ 只要文件名和一段文本就能工作。
    未写完的代码、写坏的模板、缺 include 的文件都能格式化，
    所以它可以挂在保存钩子和 CI 上，而不需要先把整个项目编起来。
  \item \textbf{它不是 AST 格式化器}。它不「知道」这段是不是函数定义，只对 token 之间的
    上下文做猜测。因此凡依赖类型或语法完整性的构造——宏展开后的形状、模板参数换行、
    属性（attribute）与 lambda 捕获——它只能猜。猜错的表现不是崩溃，
    而是\textbf{每次都能格式化成功，但结果不合你意}。
  \item \textbf{宏是二等公民}。\texttt{Q\_OBJECT} 一类靠宏名识别的构造，需要用
    \texttt{ForEachMacros}、\texttt{StatementMacros} 这样的名单字段显式告诉它，
    否则宏调用会被当成普通表达式参与断行。
\end{itemize}

\begin{guideline}[title={把 \Fmt\ 当排版器，不要当重构器}]
  \Fmt\ 只改变空白与换行，绝不改变语义：它不会加括号、不会补分号、不会重命名。
  反过来说，也别指望它替你决定「这个长参数列表该拆成几个语义组」，那是接口设计问题。
  真需要一个字段时，先确认它是空白类字段，再考虑要不要为它开一次团队讨论。
\end{guideline}

\section{风格从哪来：\Opt{style}、\Opt{dump-config} 与查找顺序}\label{sec:ch09-s2}

\subsection{内置风格与内联风格}\label{sec:ch09-s2-inline}

\Opt{style} 的取值是固定的那几个名字：\texttt{LLVM}、\texttt{GNU}、\texttt{Chromium}、
\texttt{Microsoft}、\texttt{Mozilla}、\texttt{WebKit}、\texttt{Stroustrup}、
\texttt{File}、\texttt{None}，或者一段内联的 \texttt{\{...\}}。

\begin{lstlisting}[style=shellstyle,caption={查看内置风格、导出现场生效的配置},label={lst:ch09-dump}]
# dump every key with the value that style uses
clang-format --style=Chromium --dump-config > chromium.def

# base on one style, override two keys inline
clang-format --style="{BasedOnStyle: microsoft, IndentWidth: 2}" a.cpp

# keep only the keys that actually differ
diff -u llvm.def my.def | grep -E '^[+-][A-Za-z]'
\end{lstlisting}

内联风格里的 \texttt{Inherits} 允许叠好几层，适合把「公司基础风格」做成一行参数下发。
但\textbf{内联只适合两三个字段}：一旦超过这个量就写成仓库里的 \texttt{.clang-format}，
因为内联字符串无法评审，也无法被编辑器侧的 \CD\ 复用。

\subsection{\texttt{BasedOnStyle} 的继承与覆盖}\label{sec:ch09-s2-inherit}

\texttt{.clang-format} 是 YAML，它和内置风格之间是\textbf{字段级覆盖}：基风格提供所有
字段的值，你的文件里写了哪个字段，就整体替换那一个字段的值。两个容易踩的点：

\begin{itemize}
  \item 覆盖是整值替换，不是局部修补。像 \texttt{BraceWrapping} 这种带子键的结构，
    早期版本只写其中两三个子键时，未写出的子键会退回\textbf{默认值}而不是基风格的值，
    后来的版本调整过这里的语义。稳妥写法是配 \texttt{BreakBeforeBraces: Custom}
    时把整个 \texttt{BraceWrapping} 块写全（含 \texttt{Custom: true}），不依赖继承。
  \item \texttt{Language} 决定用哪套词法与关键词表（\texttt{None}、\texttt{C}、
    \texttt{Cpp}、\texttt{Java}、\texttt{JavaScript}、\texttt{Proto}、
    \texttt{TableGen}），它和 \texttt{BasedOnStyle} 是两个正交维度。
    \CL\ 文件与 \CPP\ 文件想区别对待，就按 \texttt{Language: C} 再配一份文件，
    而不是指望 \texttt{Standard} 字段自动区分。
\end{itemize}

\subsection{配置文件的查找顺序}\label{sec:ch09-s2-lookup}

给定一个待格式化文件的路径，\Fmt\ 从\textbf{该文件所在目录}开始逐级向上找
\texttt{.clang-format}，第一个命中就用；同目录下也接受带 YAML 头或扩展名的变体。
一路找到仓库根都没有时，退到 \Opt{fallback-style}（默认 \texttt{LLVM}）。补充几条：

\begin{itemize}
  \item \Opt{style=file} 的意思才是「按上面的顺序找」；\Opt{style} 给了具体风格名则
    \textbf{完全不找}。CI 里请显式写 \Opt{style=file}，否则你会得到一份没人同意过的
    LLVM 风格。
  \item \Opt{style} 可以写成 \texttt{file:.clang-format-alt}，用于在同一仓库里并存多套风格。
  \item 环境变量 \texttt{CLANG\_FORMAT\_STYLE} 会作为找不到配置文件时的风格名，
    优先级低于命令行 \Opt{style}，高于 \Opt{fallback-style}。
  \item 较新的版本会读 \texttt{.editorconfig} 里的少量字段（缩进宽度、制表宽度、
    行宽上限、是否补末尾换行），其余仍以 \texttt{.clang-format} 为准。
    两处都写了同一个字段时不要指望「谁新听谁的」，团队里只保留一处。
\end{itemize}

\begin{note}[title={\Opt{style=file} 与 \Opt{fallback-style} 的优先级}]
  「找不到配置文件」是\textbf{静默的}：它直接用 \Opt{fallback-style}，不报警，也不返回非零。
  症状是格式化结果与仓库风格明显不符，或者 CI 报出一堆无法解释的差异。
  对策是显式写 \Opt{fallback-style=none}：
  \texttt{clang-format --style=file --fallback-style=none}。
  此时找不到配置就不动文件，问题立刻暴露。CI 里务必这样写。
\end{note}

\section{缩进、宽度与大括号}\label{sec:ch09-s3}

\subsection{缩进族}\label{sec:ch09-s3-indent}

\texttt{IndentWidth} 是基本步长（\texttt{LLVM}/\texttt{Chromium} 为 2，\texttt{Microsoft}
为 4），\texttt{ContinuationIndentWidth} 是续行相对第一列多缩进多少（默认 4），
\texttt{TabWidth} 只决定制表符按几列折算，\texttt{AccessModifierOffset} 常留 \texttt{-1}
让 \texttt{public:} 比类体少缩进一格（缩进 4 的风格里通常是 \texttt{-4}）。
\texttt{UseTab} 的取值是 \texttt{Never}、\texttt{ForIndentation}、
\texttt{ForContinuationAndIndentation}、\texttt{AlignWithSpaces}；
老配置里的 \texttt{true}/\texttt{false} 仍被接受，新仓库请写枚举名。

\texttt{IndentPPDirectives} 管预处理指令的缩进（\texttt{None}、\texttt{AfterHash}、
\texttt{BeforeHash}），\texttt{PPIndentWidth} 单独给它定步长；
\texttt{IndentCaseLabels} 决定 \texttt{switch} 里的 \texttt{case} 要不要再进一层；
\texttt{IndentGotoLabels} 较新版本默认 \texttt{true}；\texttt{NamespaceIndentation} 取
\texttt{None}、\texttt{Inner}、\texttt{All}；\texttt{LambdaBodyIndentation} 决定 lambda
体相对签名还是相对外层作用域缩进（\texttt{SignatureLocation}、\texttt{OuterScope}）。

\texttt{ColumnLimit} 是所有断行决策的靶心：\texttt{LLVM} 是 80，\texttt{Microsoft} 是 120，
\texttt{GNU} 是 0 表示不限宽（此时只做空白归一化）。改大它只影响断行位置；
改小则可能让一行了事的构造（\ref{sec:ch09-s3-short}）集体拆开。

\subsection{大括号：\texttt{BreakBeforeBraces} 与 \texttt{BraceWrapping}}\label{sec:ch09-s3-braces}

七个常用取值加 \texttt{Custom} 的语义差别见表 \ref{tab:ch09-braces}。
选 \texttt{Custom} 时逐条写 \texttt{BraceWrapping} 子键：\texttt{AfterFunction}、
\texttt{AfterClass}、\texttt{AfterStruct}、\texttt{AfterEnum}、\texttt{AfterUnion}、
\texttt{AfterNamespace}、\texttt{AfterControlStatement}（取值有 \texttt{Never}、
\texttt{Always}、\texttt{MultiLine}、\texttt{Macro}，更细的还有
\texttt{TrueMacros}、\texttt{FalseMacros}）、\texttt{BeforeElse}、\texttt{BeforeCatch}、
\texttt{BeforeWhile}、\texttt{BeforeLambdaBody}、\texttt{AfterCaseLabel}、
\texttt{AfterExternBlock}、\texttt{IndentBraces}，以及控制「空的 record 要不要拆成三行」的
\texttt{SplitEmptyRecord}、\texttt{SplitEmptyFunction}、\texttt{SplitEmptyNamespace}。

{\small
\begin{longtable}{%
  >{\raggedright\arraybackslash}p{0.16\textwidth}%
  >{\raggedright\arraybackslash}p{0.22\textwidth}%
  >{\raggedright\arraybackslash}p{0.22\textwidth}%
  >{\raggedright\arraybackslash}p{0.18\textwidth}%
  }
  \toprule
  取值 & 函数与类型的左括号 & 控制语句的左括号 & 常见出处 \\
  \midrule
  \endhead
  \texttt{Attach} & 同行 & 同行 & \texttt{LLVM}、\texttt{Chromium} \\
  \texttt{Linux} & 函数独占一行，类型同行 & 同行 & 内核风格，\texttt{GNU} 接近它 \\
  \texttt{MS} & 一律独占一行 & 同行 & \texttt{Microsoft} \\
  \texttt{Allman} & 独占一行 & 独占一行 & .NET、部分嵌入式团队 \\
  \texttt{Stroustrup} & 独占一行 & 同行；\texttt{else}、\texttt{catch} 与前面的闭合大括号同行 & 教材风格 \\
  \texttt{WebKit} & 独占一行 & 同行 & \texttt{WebKit} \\
  \texttt{GNU} & 缩进两格的独立行 & 独立行且缩进 & \texttt{GNU} 内置 \\
  \texttt{Custom} & 由 \texttt{Brace\-Wrapping} 子键逐条决定 & 同左 & 有历史代码库的团队 \\
  \bottomrule
  \caption{\texttt{BreakBeforeBraces} 各取值的语义对照}
  \label{tab:ch09-braces}
\end{longtable}
}

注意 \texttt{BreakAfterAttributes}（\texttt{Never}、\texttt{Always}、\texttt{Leave}）：
\cpp{11} 之后的 \texttt{[[nodiscard]]}、\texttt{\_\_attribute\_\_} 与 MIDL 属性写在返回类型之前，
它决定属性与返回类型之间要不要断行。这是典型的「新版本才有」的字段，
老二进制可能报未知键或直接忽略它，用 \Opt{dump-config} 确认你手上版本的键名。

\subsection{允许单行的那几个开关}\label{sec:ch09-s3-short}

\texttt{Allow\-Short\-If\-State\-ments\-On\-ASingleLine} 的取值是 \texttt{Never}、
\texttt{WithoutElse}、\texttt{AllIfElseAndElse}、\texttt{All}；布尔值 \texttt{true} 与
\texttt{false} 是旧写法，大致等价于 \texttt{All} 与 \texttt{Never}。
这个键的带后缀变体（\texttt{WithoutAssertions} 一类）在不同 LLVM 版本上拼法不同，
\textbf{一律以 \Opt{dump-config} 的输出为准}。
\texttt{AllowShortFunctionsOnASingleLine} 取 \texttt{None}、\texttt{Empty}、\texttt{Inline}、
\texttt{All}，其中 \texttt{Empty} 只把空函数体收成一行，是最不容易引起争议的档位。
\texttt{AllowShortLambdasOnASingleLine} 取 \texttt{None}、\texttt{Empty}、\texttt{Inline}、
\texttt{All}。同一族还有 \texttt{AllowShortBlocksOnASingleLine}、
\texttt{AllowShortCaseLabelsOnASingleLine}、\texttt{AllowShortLoopsOnASingleLine}、
\texttt{AllowShortEnumsOnASingleLine}，语义都是「够短就别拆」。

这一族是格式化争议的最大来源：同一段 \texttt{if (x) return;} 在四种配置下有四种长相，
而且每一种都「合法」。团队协商时只需要问一个问题——\textbf{提前返回要不要独占一行}。
定了就写进配置，别让个人编辑器的默认值漂移进来。

\section{对齐、断行与惩罚}\label{sec:ch09-s4}

\subsection{指针与引用的星号位置}\label{sec:ch09-s4-pointer}

\texttt{PointerAlignment} 取 \texttt{Left}（\texttt{int* p}）、\texttt{Right}
（\texttt{int *p}）、\texttt{Middle}（多星号声明里做视觉居中）。较新版本把引用拆出来
单独管：\texttt{ReferenceAlignment} 取 \texttt{Pointer}、\texttt{Left}、\texttt{Right}、
\texttt{Middle}，默认 \texttt{Pointer} 表示跟随 \texttt{PointerAlignment}。

\begin{warning}[title={\texttt{DerivePointerAlignment} 让同一仓库出现两种星号}]
  \texttt{DerivePointerAlignment: true}（\texttt{LLVM} 内置风格就是开的）会让 \Fmt\
  观察这个\textbf{文件}里已有的写法然后跟着用。症状：新人拷一段旧代码进来，
  于是同一文件里同时出现 \texttt{int* p} 与 \texttt{char *q}，而且各自都能
  「格式化通过」。原因：它按文件推断，而不是按仓库约定。
  对策：团队配置里显式写 \texttt{DerivePointerAlignment: false} 并给出
  \texttt{PointerAlignment}，再做一次全量格式化（做法见第 10 章）。
\end{warning}

\subsection{对齐族：视觉整齐与可维护性的交换}\label{sec:ch09-s4-align}

\texttt{AlignAfterOpenBracket}（\texttt{Align}、\texttt{DontAlign}、\texttt{BreakOnClose}）
决定断行后的参数是否对齐到左括号后一列。\texttt{AlignOperands}（\texttt{DontAlign}、
\texttt{Align}、\texttt{AlignQuietly}）决定 \texttt{||}、\texttt{+} 这类操作数要不要跟第一行
对齐。\texttt{AlignConsecutive\Bk Assignments}、\texttt{AlignConsecutive\Bk Declarations}、
\texttt{AlignConsecutive\Bk Macros}、\texttt{AlignConsecutive\Bk BitFields}
取值同为 \texttt{None}、\texttt{Consecutive}、\texttt{LeadingWhitespace}、
\texttt{EmptyLines}，管的是「相邻若干行的等号竖着排齐」。
\texttt{AlignTrailingComments} 取 \texttt{Left}、\texttt{Right}、\texttt{DontAlign}；
\texttt{AlignEscapedNewlines} 决定反斜杠续行的行尾反斜杠靠哪边；
\texttt{AlignArrayOfStructures}（\texttt{None}、\texttt{Left}、\texttt{Right}）
把结构体数组的字面量排成表格。

\begin{lstlisting}[style=diffstyle,caption={\texttt{AlignAfterOpenBracket} 与 \texttt{AlignOperands} 的可见差别},label={lst:ch09-align}]
--- a/render.cpp  (DontAlign + BreakBeforeBinaryOperators: None)
+++ b/render.cpp  (Align + AlignOperands: Align)
@@ draw @@
-  auto hits = scene.raycast(
-    origin,
-    direction,
-    maxDistance);
+  auto hits =
+    scene.raycast(origin, direction, maxDistance);
@@ filter @@
-  bool keep = visible
-    && !culled
-    && alpha > kEps;
+  bool keep =
+      visible &&
+      !culled &&
+      alpha > kEps;
\end{lstlisting}

列表 \ref{lst:ch09-align} 说明一件事：对齐类字段互相纠缠，改一个会把别处的断行决策一起改掉。
所以\textbf{不要逐个字段试}，正确流程是：拿一份有代表性的代码，\Opt{dump-config} 导出基线，
改一个字段，再 diff 两份配置与两次格式化结果。

\subsection{要不要把参数全塞回一行：BinPack 族}\label{sec:ch09-s4-binpack}

\texttt{BinPackArguments} 与 \texttt{BinPackParameters} 控制「一行塞得下就塞」。
关掉之后每个参数各占一行，diff 更稳定（新增一个参数只改一行），代价是长签名很占地方。
\texttt{AllowAllArgumentsOnNextLine} 与 \texttt{AllowAllParametersOf\-Declara\-tion\-On\-NextLine}
是配套的「一行放不下时，允不允许整体挪到下一行」，默认都是 \texttt{true}。
想要「左括号后立刻换行、右括号独占一行」的视觉效果，就把它们关掉并配
\texttt{AlignAfterOpenBracket: BreakOnClose}。

断行方向相关的还有：\texttt{BreakBefore\Bk BinaryOperators}（\texttt{None}、
\texttt{NonAssignment}、\texttt{All}）、\texttt{BreakBefore\Bk TernaryOperators}（布尔，
决定 \texttt{?:} 落在上一行尾还是下一行首）、\texttt{AlwaysBreak\Bk AfterReturnType}
（\texttt{None}、\texttt{All}、\texttt{TopLevel}，让返回类型单独一行；
同名作用更窄的 \texttt{AlwaysBreak\Bk AfterDefinition\Bk ReturnType} 已废弃）、
\texttt{AlwaysBreak\Bk BeforeMultilineStrings}、\texttt{BreakStringLiterals}
（超长字符串字面量是否按词断开；关掉它会容忍超列）。

构造器与继承列表方向上则是三个字段：

\begin{itemize}
  \item \texttt{BreakConstructorInitializers} 决定断在冒号前还是逗号前，
    取值有 \texttt{BeforeColon}、\texttt{AfterColon}、\texttt{BeforeComma}、
    \texttt{AfterComma}。
  \item \texttt{ConstructorInitializerIndentWidth} 决定初始化列表续行的缩进量。
  \item \texttt{ConstructorInitializerAllOnOneLineOrOnePerLine}
    决定「一行放不下时，要不要一行一个冒号项」。
\end{itemize}

继承列表的对应字段是 \texttt{BreakInheritanceList}，它对继承列表做的事和上面第三条
对初始化列表做的事等价；不要再去找早已删除的 \texttt{BreakBeforeInheritanceComma}。

\subsection{penalty：只在明显不对时才动}\label{sec:ch09-s4-penalty}

最短路径求解的代价函数由一组「惩罚」字段构成，思路是：超出一列一个字符收多少代价
（\texttt{PenaltyExcessCharacter}）、返回类型单独占一行收多少
（\texttt{PenaltyReturnTypeOnItsOwnLine}）、在第一个调用参数处断收多少
（\texttt{PenaltyBreakBeforeFirstCallParameter}），以及在赋值号后断、在注释里断、
在 \texttt{<} 模板参数处断各收多少。这一族的键名在版本之间有增删与改名，
所以\textbf{唯一可靠的做法是先 \Opt{dump-config} 再改}。

实践建议：penalty 族整体保持默认，最多动两个——把 \texttt{PenaltyReturnType\-OnItsOwnLine}
调高可以抑制「返回类型孤零零一行」，把 \texttt{PenaltyExcessCharacter} 调高可以让它
宁肯多断几行也不超列。动完之后必须全仓库重跑一次格式化并统计 diff 规模，
否则你不知道自己顺手改了多少文件。

\subsection{注释、空行与结构整理}\label{sec:ch09-s4-comment}

\texttt{ReflowComments} 会把过长的普通注释按词重排（\texttt{LLVM} 默认开）。
它和 \texttt{MaxEmptyLinesToKeep}（默认 1）、
\texttt{KeepEmptyLinesAtTheStartOfBlocks}（\texttt{false} 表示块首空行删掉）、
\texttt{SeparateDefinitionBlocks}（\texttt{Leave}、\texttt{Skip}、\texttt{Always}，
定义之间强制留空行）、\texttt{EmptyLineBeforeAccessModifier}
（\texttt{LogicalBlock}、\texttt{Never}、\texttt{Preserve}、\texttt{Always}）
共同决定纵向密度。

\texttt{CommentPragmas} 是一个正则，声明「以这些词开头的注释是有意标记，别重排」，
例如文档注释的标记、自己团队的 \texttt{NOLINT} 前缀。这是让 \Fmt\ 尊重注释语义的
正规入口，比把注释整段包进 \texttt{clang-format off} 好得多（见 \ref{sec:ch09-s6}）。

\texttt{FixNamespaceComments} 会补全或删掉命名空间结尾的
\texttt{// namespace foo} 注释，配合 \texttt{ShortNamespaceLines} 决定多少行之内的
短命名空间可以不写这条注释；\texttt{CompactNamespaces} 决定
\texttt{namespace a \{ namespace b \}} 能否并成一行；\texttt{Cpp11BracedListStyle}
让花括号初始化列表按 \cpp{11} 之后的习惯排版（内部不加空格、逗号后加空格）；
\texttt{Standard}（例如 \texttt{c++17}、\texttt{Latest}）影响若干字段的默认组合，
团队里请写同一个值。表 \ref{tab:ch09-fields} 把本章字段按「要不要协商」重新排了一次。

{\footnotesize
\begin{longtable}{%
  >{\raggedright\arraybackslash}p{0.28\textwidth}%
  >{\raggedright\arraybackslash}p{0.20\textwidth}%
  >{\raggedright\arraybackslash}p{0.13\textwidth}%
  >{\raggedright\arraybackslash}p{0.21\textwidth}%
  }
  \toprule
  字段 & 影响面 & 协商成本 & 建议 \\
  \midrule
  \endhead
  \texttt{ColumnLimit} & 几乎每一次断行 & 高：全仓库 diff & 一次定死，随仓库文档发布 \\
  \texttt{IndentWidth} & 每一行的前导空白 & 高 & 跟随 \texttt{BasedOnStyle}，别混搭 \\
  \texttt{UseTab} & 前导空白形态、review 工具显示 & 中 & 统一 \texttt{Never} 最省事 \\
  \texttt{BreakBeforeBraces} & 每个函数与控制的阅读节奏 & 高：口味之争 & 选 \texttt{Custom} 写全子键 \\
  \texttt{Allow\-Short\-If\-State\-ments\-On\-ASingleLine} & 提前返回的形态 & 中 & 定 \texttt{WithoutElse} 或 \texttt{Never} \\
  \texttt{AllowShortFunctions\-OnASingleLine} & 头文件里的访问器密度 & 中 & \texttt{Empty} 争议最小 \\
  \texttt{PointerAlignment} & 声明与引用的可读性 & 高 & 关掉 \texttt{DerivePointer\-Alignment} \\
  \texttt{BinPackParameters} & 函数签名的 diff 稳定性 & 中 & 想要稳定 diff 就关掉 \\
  \texttt{Align\-Consecutive\-Assignments} & 常量表的视觉 & 低 & 多数仓库设 \texttt{None} \\
  \texttt{AlignTrailingComments} & 行尾注释列位置 & 低 & 设 \texttt{Left}，改动少时更稳 \\
  \texttt{ReflowComments} & 注释重排，可能与文档工具冲突 & 中 & 用 \texttt{CommentPragmas} 兜住标记 \\
  \texttt{IncludeBlocks}、\texttt{SortIncludes} & include 区 & 中 & 与手写分组习惯对齐 \\
  \texttt{NamespaceIndentation} & 整块代码的缩进层级 & 高 & \texttt{None} 让行宽更值钱 \\
  \texttt{penalty/cost 族} & 少数边界形状 & 高 & 保持默认，最多动两个 \\
  \bottomrule
  \caption{字段、影响面与团队协商成本}
  \label{tab:ch09-fields}
\end{longtable}
}

\section{include 的分组与排序}\label{sec:ch09-s5}

\Fmt\ 对 \texttt{\#include} 只做两件事：\textbf{重排}与\textbf{空白归一}，
不增、不减、不补。补 include 是 IWYU 一类工具的活（第 15 章）。

\begin{itemize}
  \item \texttt{SortIncludes} 取 \texttt{Never}、\texttt{LCase}、\texttt{CaseSensitive}，
    较新版本另有 \texttt{Dependency}（尝试按依赖关系排）。它管「每组内部怎么排」，
    不管「组怎么分」。
  \item \texttt{IncludeBlocks} 取 \texttt{Preserve}（保留你手写的分组，只做组内排序）、
    \texttt{Regroup}（按 \texttt{IncludeCategories} 重新分组，默认）、
    \texttt{Merge}（把所有块拉平重排，会吃掉你故意留的空行）。
    症状「一保存我排的组就散架」，原因通常是 \texttt{Regroup} 配了过宽的正则，
    对策是 \texttt{IncludeBlocks: Preserve}。
  \item \texttt{IncludeCategories} 是条目列表，每条有 \texttt{Regex} 与 \texttt{Priority}，
    数字小的组排前面；\texttt{SortPriority} 可以只影响排序而不影响分组。
    正则匹配的是\textbf{引号或尖括号里的路径文本}，所以 \texttt{"foo.h"} 与
    \texttt{<foo.h>} 是两条不同的通道，写规则时两边都要覆盖。
  \item \texttt{MainIncludeStyle} 决定「与本翻译单元同名的那个 include」落在哪一组：
    \texttt{RegexBased}（默认，用 \texttt{Include\Bk IsMainRegex} 判定）或
    \texttt{CaseSensitive\Bk RegexBased}；\texttt{Include\Bk IsMainRegex} 默认容忍
    \texttt{-inl}、\texttt{\_test} 一类的后缀，\texttt{Include\Bk IsMainSourceRegex}
    管源文件侧的判定。
  \item 一条实用规则：把「本项目自己的头」设成最小 \texttt{Priority}，
    系统头与第三方头后置，这样违反依赖方向的引用会出现在最显眼的位置。
\end{itemize}

\begin{lstlisting}[style=yamlstyle,caption={一份四组的 include 排序配置},label={lst:ch09-include}]
IncludeBlocks: Regroup
SortIncludes: CaseSensitive
MainIncludeStyle: RegexBased
IncludeIsMainRegex: '([-_](test|unittest))?$'
IncludeCategories:
  - Regex: '^".*"$'         # project headers first
    Priority: 1
  - Regex: '^<project/'     # public headers, bracket form
    Priority: 2
  - Regex: '^<(gsl|fmt|spdlog)/'
    Priority: 4
  - Regex: '^<.*\.h>'       # C compatibility headers
    Priority: 6
  - Regex: '^<.*>'          # standard library
    Priority: 5
  - Regex: '.*'             # everything else
    Priority: 7
\end{lstlisting}

列表 \ref{lst:ch09-include} 前后的对照见 \ref{lst:ch09-incl-diff}。
注意只有 \texttt{Priority} 相同的条目才会被同一次排序打散，不同组之间必然留一个空行，
所以 \texttt{Priority: 5} 与 \texttt{Priority: 6} 的顺序由数字决定，
而不是由「先写哪个」决定（列表里 \texttt{.h} 那条写在标准库之前，实际却排在后面）。

\begin{lstlisting}[style=diffstyle,caption={同一个翻译单元应用 include 排序配置后的差异},label={lst:ch09-incl-diff}]
--- a/widget.cpp   (before: hand written order)
+++ b/widget.cpp   (after : IncludeBlocks: Regroup)
@@ includes @@
-#include <vector>
-#include "widget.h"
-#include <spdlog/spdlog.h>
-#include <algorithm>
-#include <cmath>
+#include "widget.h"
+
+#include <algorithm>
+#include <cmath>
+#include <vector>
+
+#include <spdlog/spdlog.h>
\end{lstlisting}

\section{局部控制与一份可直接用的团队配置}\label{sec:ch09-s6}

\subsection{四个逃生口}\label{sec:ch09-s6-local}

\begin{itemize}
  \item \texttt{// clang-format off} 到 \texttt{// clang-format on} 之间原样保留。
    适合位域表、手工对齐的常量矩阵、依赖列位置的 ASCII 图。
  \item \texttt{// clang-format skip} 挂在下一行之上，只跳过紧跟着的那一行，
    泄漏面比成对的 off/on 小得多，优先用它。
  \item 命令行侧：\Opt{lines=3,7} 与 \Opt{offsets=120,480} 只格式化指定范围，
    \Opt{cursor=N} 保证光标原地不动（编辑器集成用），
    \Opt{assume-filename=x.cpp} 让从标准输入读入的内容按某个文件名的规则处理。
  \item 编辑器侧的 \texttt{@formatter:off} 是 Eclipse/CDT 一类编辑器的私有标记，
    \Fmt\ 本身\textbf{不认识}它。如果团队里既有人写 \texttt{@formatter:off}
    又有人写 \texttt{clang-format off}，就会出现「一边被格式化一边没有」。
    请统一成 \texttt{clang-format} 注释，见列表 \ref{lst:ch09-local}。
\end{itemize}

\begin{lstlisting}[style=cppstyle,caption={三种局部控制写法的落位},label={lst:ch09-local}]
namespace gfx {
// clang-format off
  // Bit layout table: hand aligned, do not touch.
  //  [31:24] kind   [23:16] lane   [15:8] flags  [7:0] idx
  enum Kind : unsigned { kNone = 0, kQuad = 1 };
// clang-format on

constexpr int kLanes = 4;  // clang-format skip

void store(Pixel *dst) {
  dst->kind = kQuad;  // aligned by hand below
  // clang-format off
  dst->lane  = 3;
  dst->flags = 0;
  // clang-format on
}
}  // namespace gfx
\end{lstlisting}

\begin{warning}[title={不要用文件级 \texttt{Disable: true} 掩盖配置不一致}]
  \texttt{.clang-format} 里写 \texttt{Disable: true} 会让整个文件原样不动，
  看起来「解决了」风格争论，实际把问题永久留在仓库里：这个目录的代码从此不受任何字段约束，
  新二进制版本发布时它的 diff 会一次性爆发。
  正确做法是把有争议的那\textbf{一小段}用 \texttt{clang-format off}/\texttt{skip} 圈起来，
  并在代码评审里要求写明理由；文件级的 \texttt{Disable} 只留给生成代码，
  而且更好的方式是直接在门禁里排除生成目录（第 10 章）。
\end{warning}

\subsection{一份团队配置}\label{sec:ch09-s6-team}

列表 \ref{lst:ch09-team} 是可以直接抄走的起点：以 \texttt{Microsoft} 为基，
把最有争议的字段全部显式写出来，这样\textbf{换 LLVM 版本时默认值漂移不会偷偷改掉你的代码}。
每加一个字段都应该能回答「它改变哪一类 diff」。

\begin{lstlisting}[style=yamlstyle,caption={仓库根 \texttt{.clang-format} 示例},label={lst:ch09-team}]
---
Language: Cpp
BasedOnStyle: Microsoft
Standard: c++17
# --- indentation and width ---
IndentWidth: 4
TabWidth: 4
UseTab: Never
ContinuationIndentWidth: 4
ColumnLimit: 110
AccessModifierOffset: -4
IndentCaseLabels: true
IndentGotoLabels: true
IndentPPDirectives: AfterHash
NamespaceIndentation: None
# --- braces ---
BreakBeforeBraces: Custom
BraceWrapping:
  Custom: true
  AfterFunction: true
  AfterClass: true
  AfterControlStatement: Never
  AfterEnum: false
  AfterStruct: false
  BeforeElse: false
  IndentBraces: false
  SplitEmptyFunction: false
  SplitEmptyRecord: false
# --- short statements and breaking ---
AllowShortIfStatementsOnASingleLine: WithoutElse
AllowShortFunctionsOnASingleLine: Empty
AllowShortLambdasOnASingleLine: All
PointerAlignment: Left
DerivePointerAlignment: false
SpaceBeforeParens: ControlStatements
BinPackParameters: false
BinPackArguments: false
AlignAfterOpenBracket: Align
AlignOperands: Align
AlignTrailingComments: Left
ReflowComments: false
# --- blank lines and comments ---
MaxEmptyLinesToKeep: 1
KeepEmptyLinesAtTheStartOfBlocks: false
FixNamespaceComments: true
Cpp11BracedListStyle: true
# --- includes, extend as in the earlier listing ---
IncludeBlocks: Regroup
SortIncludes: CaseSensitive
\end{lstlisting}

\subsection{跑一遍看效果}\label{sec:ch09-s6-effect}

把列表 \ref{lst:ch09-team} 应用到一段按 \texttt{LLVM} 风格写就的旧代码上，
节选差异见列表 \ref{lst:ch09-team-diff}（原文缩进 2 格、引用符号靠右）：

\begin{lstlisting}[style=diffstyle,caption={团队配置作用在既有代码上的差异节选},label={lst:ch09-team-diff}]
--- a/cache.cpp   (before: LLVM, IndentWidth 2)
+++ b/cache.cpp   (after : team style, lst:ch09-team)
@@ insert @@
-void Cache::insert(std::string_view key,
-  Entry value) {
-  if (full() && !evictOne()) return;
-  map_.emplace(std::move(key), std::move(value));
+void Cache::insert(
+    std::string_view key,
+    Entry value) {
+    if (full() && !evictOne()) {
+        return;
+    }
+    map_.emplace(std::move(key), std::move(value));
@@ find @@
-Entry* Cache::find(const std::string& key) {
+Entry *Cache::find(const std::string &key) {
   auto it = map_.find(key);
-  return it == map_.end() ? nullptr : &it->second;
+  return it == map_.end() ? nullptr
+                          : &it->second;
@@ dtor @@
-Cache::~Cache() = default;
+Cache::~Cache() {
+
+}
\end{lstlisting}

四处变化分别来自：续行缩进（\texttt{Continuation\hspace{0pt}IndentWidth}）、
\texttt{BinPackParameters} 设为 \texttt{false}、
\texttt{AllowShortIfStatements\hspace{0pt}OnASingleLine}
管的是 \texttt{if} 本体而大括号仍由 \texttt{BraceWrapping} 决定、
\texttt{AlignOperands} 与三元的对齐。最后一处空析构被拆成三行是
\texttt{SplitEmptyFunction} 的效果——如果你的仓库里有大量 \texttt{= default} 的析构，
记得把它设成 \texttt{false}，否则一次性格式化会造出一堆空块。

\begin{bestpractice}[title={配置只放仓库根，子目录不放覆盖文件}]
  \Fmt\ 的查找是「从文件所在目录向上，取第一个命中」，所以子目录里一份
  \texttt{.clang-format} 会静默接管整个子树。症状：某个模块的风格莫名其妙地和别处不同，
  而 CI 还是绿的（因为 CI 也用了那份子配置）。原因多为从别的仓库拷配置时留下的残留。
  对策：只在仓库根放一份；确有必要覆盖（例如语言混合目录）时，
  把它当接口变更来评审，并在仓库文档里写清哪些目录允许覆盖。
  第三方代码更好的处理是从门禁里排除，而不是给它一份子配置。
\end{bestpractice}

\section*{本章要点}
\begin{itemize}
  \item \Fmt\ 是「词法 token 加语境猜测加带权最短路」，
    所以它不需要编译数据库，也解释不了所有输出。
  \item 配置优先级只需要记三档：命令行 \Opt{style} 覆盖一切；
    \texttt{.clang-format} 从文件所在目录向上取第一个命中；都没有则
    \Opt{fallback-style}。CI 里写 \Opt{fallback-style=none}。
  \item 真正需要团队协商的字段不到十个：\texttt{ColumnLimit}、\texttt{IndentWidth}、
    \texttt{UseTab}、\texttt{BreakBeforeBraces}、\texttt{AllowShort*} 一族、
    \texttt{PointerAlignment}、\texttt{BinPack*}、\texttt{IncludeBlocks} 与
    \texttt{SortIncludes}、\texttt{ReflowComments}。其余保持默认。
  \item 升级 LLVM 之前先 \Opt{dump-config} 存基线，比对键清单再决定何时全量重格式化；
    一次性全量格式化与屏蔽 blame 的做法见第 10 章。
\end{itemize}

% ---- frag_ch10.tex ----
% ==================== 第 10 章 ====================
\chapter{CI 门禁、增量格式化与 pre-commit}\label{ch:format-gate}

第 9 章解决「配置该长什么样」，本章解决「怎么让它在十个人、几百次提交之后还是这个样子」。
格式化想真正落地，需要三件事同时成立：本地改完就能自证通过、只影响改动过的行、
CI 用的二进制版本和本地一致。缺任何一件，仓库最后都会变成
「每隔几个月有人提交一个几万行的格式化 commit」。

\section{三种运行模式与退出码}\label{sec:ch10-s1}

\Fmt\ 的命令行只有三种用法，但 CI 里出问题的多半就在这三种的切换上。

\begin{enumerate}
  \item \textbf{回写}：\Opt{in-place}（短写 \texttt{-i}）直接改文件，
    \Opt{backup-file} 决定备份后缀，静默无输出，退出码只反映「有没有读到文件」。
    它适合本地与 pre-commit，不适合门禁。
  \item \textbf{只检查}：\Opt{dry-run}（短写 \texttt{-n}）不改文件，把每一处需要格式化的位置
    按 \texttt{file:line:col: warning: ... [-Wclang-format-violations]} 打到
    \textbf{标准错误}；配合 \Opt{Werror} 才返回非零。这是门禁的正确用法。
  \item \textbf{取差异}：\Opt{output-replacements-xml} 输出机器可读的替换列表
    （带 offset、长度与替换文本），可以喂给脚本自己算行列；
    人可读的 unified diff 由 \texttt{git clang-format --diff}（\ref{sec:ch10-s3}）或
    \texttt{clang-format-diff.py} 给出。
\end{enumerate}

\begin{lstlisting}[style=shellstyle,caption={三种模式的真实调用与输出形态},label={lst:ch10-modes}]
# 1) rewrite in place, keep a backup next to each file
clang-format -i --backup-file=.bak --style=file --fallback-style=none a.cpp

# 2) check only: messages go to stderr, exit code is 1 with -Werror
clang-format -n --Werror --style=file --fallback-style=none a.cpp
echo "exit=$?"

# 3a) machine readable replacements, one <replacement> per edit
clang-format --output-replacements-xml a.cpp > repl.xml
grep -c '<replacement ' repl.xml

# 3b) readable diff for one file, without touching it
git clang-format --diff HEAD -- a.cpp
\end{lstlisting}

\subsection{「检查失败了但没打印原因」}\label{sec:ch10-s1-exit}

这一类的排查顺序是固定的，按可能性从高到低：

\begin{itemize}
  \item \textbf{只写了 \Opt{Werror} 没写 \Opt{dry-run}}。\Opt{Werror} 的作用是把
    dry-run 产生的 violation 警告升级成错误；不带 \Opt{dry-run} 时它没有任何作用，
    于是「什么都没报，退出码 0」或者「直接改写了文件」。
  \item \textbf{只抓 stdout}。violation 走的是 stderr，CI 里做管道重定向时最容易丢掉，
    于是失败日志一片空白。
  \item \textbf{找不到配置文件}。\Opt{style=file} 找不到配置是静默的
    （第 9 章专门用一条 note 说明过），此时它按 \Opt{fallback-style} 判断，
    结果取决于你恰好装的是哪个默认风格：要么报出一堆无法解释的差异，要么永远绿。
    门禁里请固定写 \Opt{fallback-style=none}。
  \item \textbf{配置里有未知键或 YAML 写坏}。这会在 stderr 打出
    \texttt{Unknown key} 之类的信息，但部分版本的退出码仍是 0，
    看起来「绿了」。判据：\texttt{-n} 之后同时判断 stderr 非空。
  \item \textbf{shell 的文件名展开出错}。带空格的路径被拆成两个参数，
    \texttt{clang-format} 报「文件不存在」，而错误被 \texttt{2>/dev/null} 吃掉。
    对策见 \ref{sec:ch10-s5} 的 \texttt{-z} 与 \texttt{xargs -0}。
\end{itemize}

\begin{warning}[title={\Opt{Werror} 只对 \Opt{dry-run} 生效}]
  症状：CI 脚本里写了 \texttt{clang-format --Werror --style=file *.cpp}，
  门禁永远是绿的，而本地文件格式明显没对齐。原因：没有 \Opt{dry-run}，
  \Opt{Werror} 无的放矢，而 \texttt{*.cpp} 在部分 shell 里只展开当前目录。
  对策：固定成 \texttt{clang-format -n --Werror --style=file --fallback-style=none}
  加一份由 \texttt{git ls-files} 提供的文件清单，并且注意
  \Opt{Werror} 要写在 \Opt{dry-run} 之前或同侧，不同版本对顺序的容忍度不一致。
\end{warning}

\section{幂等性：两次运行必须给出同一份文件}\label{sec:ch10-s2}

格式化的第一性质是幂等：\texttt{format(format(x))} 必须等于 \texttt{format(x)}。
不幂等时症状非常隐蔽——本地和 CI 都「通过」（因为只检查一遍），
但每次保存都会产生一个新 diff，pre-commit 要跑两次才干净。
\textbf{非幂等几乎总是配置问题，不是代码问题}。

常见的四组起因：

\begin{itemize}
  \item \texttt{ColumnLimit} 卡在边界，再叠加 \texttt{AlignConsecutive\hspace{0pt}Assignments}
    与 \texttt{AlignOperands}。第一轮对齐把列撑宽，第二轮就换一种断行方案。
  \item \texttt{ReflowComments} 与带前缀的注释块（例如行首连续的 \texttt{-} 列表、
    ASCII 图）：重排后前缀形状变了，第二次再重排又变回去。
  \item \texttt{IncludeBlocks: Regroup} 与手写分组冲突：第一轮把组打散重排，
    第二轮的 \texttt{MainIncludeStyle} 判定结果和第一轮不同（因为第一个 include 变了）。
  \item \texttt{FixNamespaceComments} 配嵌套命名空间加 \texttt{NamespaceIndentation}
    的非 \texttt{None} 取值，注释补全与缩进互相牵引。
\end{itemize}

自查脚本骨架见列表 \ref{lst:ch10-idem}：对每个文件连续格式化两次，
比较第一次与第二次的产物。它比任何「跑一遍看有没有 diff」的写法都可靠。

\begin{lstlisting}[style=shellstyle,caption={幂等性自查：跑两次并比较},label={lst:ch10-idem}]
#!/usr/bin/env bash
set -euo pipefail
fail=0
while IFS= read -r -d '' f; do
  clang-format --style=file --fallback-style=none -i "$f"
  cp "$f" "$f.two"
  clang-format --style=file --fallback-style=none -i "$f"
  if ! cmp -s "$f.two" "$f"; then
    echo "NOT IDEMPOTENT: $f"
    diff -u "$f.two" "$f" || true
    fail=1
  fi
  rm -f "$f.two"
done < <(git ls-files -z -- '*.c' '*.cpp' '*.h' '*.hpp')
exit "$fail"
\end{lstlisting}

\begin{guideline}[title={把幂等性当成门禁的第一条测试}]
  顺序很重要：先让 \texttt{-n} 在格式化过的代码上返回 0（这叫「自洽」），
  再验证跑两次不变（这叫「幂等」）。前者不合格说明代码没格式化到位，
  后者不合格说明配置本身有矛盾。两者都通过，才值得去加 pre-commit 和 CI。
  排查时把 \texttt{ColumnLimit} 调大（例如 200）再试：如果非幂等消失，
  基本可以断定是断行与对齐的纠缠，而不是注释或 include 的规则冲突。
\end{guideline}

\section{只格式化改动行：\texttt{git-clang-format}}\label{sec:ch10-s3}

\texttt{git clang-format} 是随 clang-tools-extra 一起安装的子命令
（也可以把 LLVM 源码里那个单文件脚本放进 \texttt{PATH}），
它先用 \texttt{git diff} 求出改动的行范围，再只让这些范围内的行受 \Fmt\ 约束。

\begin{lstlisting}[style=shellstyle,caption={增量格式化的三种用法},label={lst:ch10-incremental}]
# format the staged version and refresh the index
git clang-format --staged

# print what would change, apply nothing
git clang-format --diff HEAD~ -- src/render.cpp

# format everything touched since a tag, all files
git clang-format --extensions h,hpp,cpp v1.2.0

# point at a specific binary so local and CI really use the same one
git clang-format --binary /usr/lib/llvm-18/bin/clang-format --staged
\end{lstlisting}

三个要点：

\begin{itemize}
  \item \Opt{staged} 处理的是暂存区内容而不是工作区，格式化后会回写工作区并
    重新 \texttt{git add}；不带 \Opt{staged} 时它处理工作区。
  \item \Opt{diff} 与 \Opt{staged} 互斥：前者只报告，后者只落地。
    把它放进 pre-commit 用 \Opt{staged}，放进 PR 检查用 \Opt{diff}。
  \item \Opt{extensions} 决定哪些后缀参与，默认集合不含 \texttt{.inl}、\texttt{.ipp}
    一类，需要显式列出，否则头文件的改动会绕过格式化。
\end{itemize}

\subsection{blame 污染与一次性全量格式化}\label{sec:ch10-s3-blame}

增量格式化解决了「日常改动的责任归属」，但每个仓库迟早要有一次全量格式化。
它的问题是污染 \texttt{git blame}：所有行都变成那个 commit 的作者。
标准做法是两份文件配合：

\begin{lstlisting}[style=shellstyle,caption={全量格式化提交与 blame 屏蔽},label={lst:ch10-blame}]
# one mechanical commit, nothing else in it
git ls-files -z -- '*.cpp' '*.hpp' '*.h' \
  | xargs -0 clang-format --style=file --fallback-style=none -i
git add -u && git commit -m "style: apply clang-format 18.1 across the repo"

# record its hash so blame can skip it
git rev-parse HEAD >> .git-blame-ignore-revs

# tell every clone to use that file
git config blame.ignoreRevsFile .git-blame-ignore-revs
git log --oneline -5
\end{lstlisting}

\begin{bestpractice}[title={格式化提交与逻辑提交分离}]
  一次提交里只允许一种变化：要么全是空白，要么全是语义。理由不只是 diff 干净——
  评审者无法在几万行空白改动里发现一个被顺手改错的条件；bisect 会被格式化提交
  拖成无效区间；而 cherry-pick 到发布分支时，混合提交几乎必然冲突。
  落地要求三条：commit message 用固定前缀（例如 \texttt{style:}），
  全量格式化单独成 PR 并在 CI 里跳过「同时改动逻辑」的检查，
  \texttt{.git-blame-ignore-revs} 在同一个 PR 里更新。
\end{bestpractice}

\section{pre-commit：把检查前移到提交动作}\label{sec:ch10-s4}

\texttt{pre-commit} 是一个 Python 写的钩子管理器，它在 \texttt{.git/hooks/pre-commit}
里装一个转发脚本，按 \texttt{.pre-commit-config.yaml} 拉取并缓存钩子环境。
\Fmt\ 的官方镜像仓库是 \texttt{mirrors-clang-format}。

\begin{lstlisting}[style=yamlstyle,caption={仓库根 \texttt{.pre-commit-config.yaml}},label={lst:ch10-precommit}]
repos:
  - repo: https://github.com/pre-commit/mirrors-clang-format
    rev: v18.1.8.0
    hooks:
      - id: clang-format
        name: clang-format (staged C/C++ files)
        args: [-i, --style=file, --fallback-style=none]
        files: ^(src|include|tests)/.*\.(c|h|cc|cpp|hpp)$
        exclude: ^(third_party/|build/|generated/|.*\.inl$)
  - repo: local
    hooks:
      - id: format-is-idempotent
        name: formatting is idempotent
        entry: ./tools/check-idempotent.sh
        language: system
        pass_filenames: false
        files: ^src/.*\.cpp$
\end{lstlisting}

逐字段说明：\texttt{repo} 用官方镜像而不是 \texttt{local}，好处是版本被 \texttt{rev} 钉死；
\texttt{rev} 必须是 \textbf{与项目实际使用的 LLVM 版本一致的 tag}（镜像的 tag 形如
\texttt{v18.1.8.0}，具体可用 tag 见该仓库 Releases），否则「pre-commit 修好的代码」
和「CI 检查的代码」用的不是同一个排版器；\texttt{args} 里的 \texttt{-i} 表示
就地改写后中止提交，让开发者重新 \texttt{git add}；\texttt{files} 与 \texttt{exclude}
都是正则，匹配的是相对仓库根的路径，排除 \texttt{third\_party/} 与生成目录是必须的，
否则你有一天会格式化别人的代码。

\begin{note}[title={\texttt{rev} 对齐：pre-commit、CI、编辑器三者同一个版本}]
  症状：本地 pre-commit 通过后，CI 仍报同一文件有 violation；
  或者 \CD\ 的「格式化」与命令行的结果不同。
  原因通常是三个版本各来一套：pre-commit 的 \texttt{rev} 是 18，CI 的镜像里
  \texttt{apt install clang-format} 装了 19，编辑器的 \CD\ 二进制自带 17。
  对策：把版本号写在一个地方（\texttt{.tool-versions}、CI 变量或一个
  \texttt{tools/format.sh}），其它三处都从它取值，并在 README 里写清。
  表 \ref{tab:ch10-version} 给出对齐矩阵。
\end{note}

\section{CI 里的门禁写法}\label{sec:ch10-s5}

\subsection{文件清单：\texttt{-z} 与 \texttt{xargs -0}}\label{sec:ch10-s5-list}

门禁的输入清单只能来自 \texttt{git ls-files}，不要来自 \texttt{find}（会把构建目录拖进来），
也不要来自 shell 通配（会漏目录、会被空格击溃）。

\begin{lstlisting}[style=shellstyle,caption={空格与中文文件名的正确处理},label={lst:ch10-xargs}]
# NUL separated stream, safe for names with spaces or UTF-8 bytes
git ls-files -z -- '*.cpp' '*.hpp' '*.h' '*.c' \
  | xargs -0 --no-run-if-empty \
      clang-format -n --Werror --style=file --fallback-style=none

# -0 consumes NUL records; --no-run-if-empty stops xargs from running
# clang-format with an empty argument list on a doc-only branch

# same check restricted to files touched by this branch
BASE=$(git merge-base HEAD origin/main)
git diff --name-only -z --diff-filter=ACMR "$BASE" -- '*.cpp' '*.h' \
  | xargs -0 --no-run-if-empty \
      clang-format -n --Werror --style=file
\end{lstlisting}

\begin{note}[title={文件名含空格时为什么必须 \texttt{-z} 配 \texttt{xargs -0}}]
  症状：\texttt{clang-format: error: no such file or directory: 'my'}，
  而 \texttt{my file.cpp} 明明在仓库里。原因：\texttt{xargs} 默认按空白切分，
  换行分隔的清单遇到空格就分家；\texttt{for f in \$(git ls-files)} 同理。
  对策：\texttt{git ls-files -z} 输出以 NUL 分隔，配 \texttt{xargs -0}；
  \texttt{bash} 里也可以先用 \texttt{mapfile} 按 NUL 读成数组，再整体加引号展开。
  PowerShell 侧管道传的是对象而不是文本，逐文件调用时参数不会被切开，
  代价是每个文件起一次进程。
\end{note}

\subsection{GitHub Actions 与 GitLab CI}\label{sec:ch10-s5-actions}

\begin{lstlisting}[style=yamlstyle,caption={矩阵化的格式化门禁（GitHub Actions）},label={lst:ch10-actions}]
name: quality
on:
  pull_request: {}
jobs:
  format:
    runs-on: ${{ matrix.os }}
    strategy:
      fail-fast: false
      matrix:
        os: [ubuntu-24.04, windows-2022]
        llvm: ["18"]
    defaults:
      run:
        shell: bash
    steps:
      - uses: actions/checkout@v4
      - name: install pinned clang-format
        run: python3 -m pip install "clang-format==18.1.8.*"
      - name: show the binary really used
        run: clang-format --version
      - name: check formatting
        run: |
          git ls-files -z -- '*.c' '*.cpp' '*.h' '*.hpp' \
            | xargs -0 --no-run-if-empty \
                clang-format -n --Werror --style=file \
                             --fallback-style=none
      - name: keep replacements as an artifact
        if: failure()
        run: |
          git ls-files -z -- '*.cpp' '*.h' \
            | xargs -0 clang-format --output-replacements-xml \
                > replacements.xml
      - uses: actions/upload-artifact@v4
        if: failure()
        with:
          name: replacements
          path: replacements.xml
\end{lstlisting}

列表 \ref{lst:ch10-actions} 里有三处刻意的设计。\Opt{version} 打印进日志，
出问题时先怀疑「CI 用的是不是我以为的那个二进制」；\texttt{pip install
clang-format==18.1.8.*} 是跨平台钉版本最省事的方式（Linux 也可用
\texttt{apt.llvm.org} 的固定包名，Windows 上它最稳）；
只在 \texttt{failure()} 时上传替换 XML，评审者下载后可以用脚本把替换精确落到行列。

把差异变成 PR 上的行内 annotation，有两条路：交给 reviewdog
（\texttt{git clang-format --diff} 的 unified diff 用
\texttt{reviewdog -f=diff} 消费，具体 flag 名随版本有变化，先在本地验证一次），
或者自己写十几行脚本把 diff 头转成 GitHub 的
\texttt{::error file=path,line=N::} 形式。两种做法都只报告、不修改。

\begin{lstlisting}[style=yamlstyle,caption={GitLab CI：只在 C/C++ 改动时触发},label={lst:ch10-gitlab}]
quality:format:
  stage: quality
  image: debian:trixie
  before_script:
    - apt-get update -qq
    - apt-get install -y -qq git clang-format-18
  rules:
    - changes:
        - "**/*.{c,h,hpp,cpp,cc}"
      when: always
    - when: never
  script:
    - git ls-files -z -- '*.c' '*.cpp' '*.h' '*.hpp'
      | xargs -0 --no-run-if-empty
          clang-format-18 -n --Werror --style=file --fallback-style=none
\end{lstlisting}

\texttt{rules: changes} 让格式化任务只在相关文件变更时出现，但它会让 MR 的必需检查
在「纯文档 MR」上显示为「未运行」，需要把策略配成允许跳过，
或者用第二条 \texttt{when: never} 之外的兜底规则显式放行。

\subsection{策略选择}\label{sec:ch10-s5-policy}

{\small
\begin{longtable}{%
  >{\raggedright\arraybackslash}p{0.18\textwidth}%
  >{\raggedright\arraybackslash}p{0.24\textwidth}%
  >{\raggedright\arraybackslash}p{0.16\textwidth}%
  >{\raggedright\arraybackslash}p{0.24\textwidth}%
  }
  \toprule
  策略 & 摩擦来源 & 适用团队规模 & 回滚成本 \\
  \midrule
  \endhead
  全量强制：仓库一次性格式化，之后 CI 拦下所有 violation &
    第一次 PR 巨大；老分支合并必冲突 & 10 人以上、有 code review 文化 &
    低：规则明确，逐文件回退即可 \\
  仅 diff：只检查本分支改动过的行 &
    老代码继续烂；同一文件两次改动风格不同 & 3--10 人的渐进治理 &
    中：要判断哪部分是历史遗留 \\
  仅建议：CI 只发 annotation，不阻断合并 &
    没人当真，问题长期累积 & 2 人以内或实验仓库 &
    高：等到必须强制时是全量重格式化 \\
  \bottomrule
  \caption{格式化门禁的三档策略}
  \label{tab:ch10-policy}
\end{longtable}
}

{\footnotesize
\begin{longtable}{%
  >{\raggedright\arraybackslash}p{0.20\textwidth}%
  >{\raggedright\arraybackslash}p{0.26\textwidth}%
  >{\raggedright\arraybackslash}p{0.18\textwidth}%
  >{\raggedright\arraybackslash}p{0.18\textwidth}%
  }
  \toprule
  组件 & 版本从哪来 & 谁说了算 & 不同步的后果 \\
  \midrule
  \endhead
  命令行 \Fmt & 系统包或 \texttt{pip} 安装的独立二进制 & CI 与钉版本文件 &
    排版规则整体漂移 \\
  \CD\ 内置格式化 & \CD\ 自带的 libclang-format & 编辑器安装的 \CD\ 版本 &
    保存即产生非预期改动 \\
  pre-commit 钩子 & \texttt{.pre-commit\-config.yaml} 的 \texttt{rev} &
    该 tag 对应的 LLVM 小版本 & 本地绿、CI 红 \\
  \texttt{git clang-format} & \Opt{binary} 或 \texttt{PATH} 第一个命中 &
    开发者环境 & 同一改动两人结果不同 \\
  \bottomrule
  \caption{格式化版本对齐矩阵}
  \label{tab:ch10-version}
\end{longtable}
}

\section{与 \CD\ 实时反馈的分工}\label{sec:ch10-s6}

正确的分工是「本地格式化负责体验，CI 负责权威」：保存即格式化，让作者第一时间看到
最终排版；CI 只做验证，不做修改。反过来（CI 自动 commit 格式化结果）会让分支上出现
机器提交，评审者更难追。

\subsection{客户端侧的三个配置差异}\label{sec:ch10-s6-client}

\begin{itemize}
  \item \textbf{配置解析路径不同}。\Fmt\ 从\textbf{文件所在目录}向上找
    \texttt{.clang-format}；部分编辑器的格式化插件从\textbf{工作区根}找，
    或者把 \texttt{.clang-format} 的内容读进内存再传给内置库。
    monorepo 里子目录有覆盖文件时，两者结果不同。
  \item \textbf{\CD\ 的 \texttt{FormatStyle} 是另一层}。\CD\ 自带格式化实现，
    它读 \texttt{.clangd} 里的 \texttt{FormatStyle}（默认 \texttt{file}，
    即回落到仓库配置）。如果有人在 \texttt{.clangd} 里写了
    \texttt{FormatStyle: "\{ IndentWidth: 4 \}"}，
    那么 \texttt{.clang-format} 就被绕过了——这是「\CD\ 格式化与命令行不一致」
    最常见的真实原因。
  \item \textbf{二进制版本不同}：\CD\ 用的是自带的 libclang-format，
    命令行用的是安装的 \Fmt；跨大版本时两者对同一个键的默认值可能不同。
\end{itemize}

\begin{lstlisting}[style=jsonstyle,caption={编辑器侧最小一致配置（VS Code 风格）},label={lst:ch10-editor}]
{
  "editor.formatOnSave": true,
  "C_Cpp.formatting": "clangFormat",
  "C_Cpp.clang_format_style": "file",
  "C_Cpp.clang_format_fallbackStyle": "none",
  "clang-format.style": "file",
  "clang-format.fallbackStyle": "none",
  "editor.rulers": [110],
  "[cpp]": { "editor.formatOnSave": true }
}
\end{lstlisting}

\begin{lstlisting}[style=yamlstyle,caption={\texttt{.clangd} 里与格式化有关的写法},label={lst:ch10-clangd}]
# keep clangd's formatter pointed at the repo config
FormatStyle: file
Diagnostics:
  ClangTidy:
    Remove: []
# do NOT put an inline style here unless you really want
# to bypass .clang-format for the whole subtree
\end{lstlisting}

\subsection{「CI 红、本地绿」的根因清单}\label{sec:ch10-s6-green}

\begin{enumerate}
  \item 版本不同：\ref{tab:ch10-version} 的四行里任意一行错位。
    判据：两边都打印 \Opt{version} 对比。
  \item 配置查找不同：CI 在仓库根跑、本地在子目录跑，
    或者 CI 用 \texttt{clang-format} 而本地 \CD\ 用被 \texttt{FormatStyle} 覆盖的那份。
  \item 清单不同：CI 只检查 \texttt{git ls-files}（已跟踪文件），
    本地格式化的是工作区里未 \texttt{add} 的新文件；反向也会发生
    （\texttt{.gitignore} 里的文件在 CI 中不存在）。
  \item 行尾与编码：\texttt{core.autocrlf} 在 Windows 检出为 CRLF，
    而配置里 \texttt{LineEnding} 或 \texttt{DeriveLineEnding} 让 \Fmt\ 认为应为 LF。
    CI 与本地用同一个 \texttt{.gitattributes} 是唯一稳妥的办法。
  \item 只格式化 diff 与格式化全文的差别：本地用 \texttt{git clang-format --staged}
    只动改动行，CI 检查全文——历史遗留的未格式化区域照样报错。
    对策是把门禁也限制在 diff 上（\ref{sec:ch10-s5-list} 的第二段），
    或者按表 \ref{tab:ch10-policy} 选全量策略。
\end{enumerate}

\section*{本章要点}
\begin{itemize}
  \item 门禁只用 \Opt{dry-run} 加 \Opt{Werror} 加 \Opt{fallback-style=none}；
    回写模式留给本地和 pre-commit。
  \item 先测自洽，再测幂等（列表 \ref{lst:ch10-idem}）；非幂等要改配置而不是改代码。
  \item 增量格式化靠 \texttt{git clang-format}，全量格式化必须单独成提交，
    并同步更新 \texttt{.git-blame-ignore-revs}。
  \item 把 LLVM 版本钉在唯一一处，让 \texttt{rev}、CI 镜像、\CD\ 与
    \Opt{binary} 四个来源都对齐（表 \ref{tab:ch10-version}）。
\end{itemize}

% ---- frag_ch11.tex ----
% ==================== 第 11 章 ====================
\chapter{\Lib{uncrustify}、\Lib{astyle} 与 \texttt{gofmt} 之问}\label{ch:other-formatters}

本章把视野拉开：先看 \Fmt\ 之外的两个仍在大量仓库里服役的工具 \Lib{astyle} 与
\Lib{uncrustify}，再回答一个经常被提问的问题——为什么 \CL/\CPP\ 世界到现在也没有
一个像 Go 那样唯一、无配置、人人都用的 \texttt{gofmt}。
这不是怀旧：仓库里往往同时存在三套排版器（老代码用 \Lib{astyle}、新代码用 \Fmt、
CI 里还有人挂着 \Lib{uncrustify}），它们的输出互相不可幂等，
最终演变成「一个 PR 改完，另一个 PR 又改回去」。

\section{\Lib{astyle}：四个风格、一份选项文件}\label{sec:ch11-s1}

\Lib{astyle} 是纯词法的：它读字符流、识别括号与运算符、按选项表补空白，
不需要编译数据库，也不建 AST。这个决定带来它的两个性格——极快、极稳，
以及对较新语法的支持滞后。

\subsection{选项模型}\label{sec:ch11-s1-model}

它的选项基本是「一个诉求一个开关」，没有继承关系，所以一份 \texttt{astyle.options}
文件就能完整描述风格。常用的形状见列表 \ref{lst:ch11-astyle}。

\begin{lstlisting}[style=shellstyle,caption={一份 \texttt{tools/astyle.options} 与它的调用},label={lst:ch11-astyle}]
# tools/astyle.options
--style=ansi
--indent=force
--indent-switches
--indent-col1-comments
--break-blocks
--pad-oper
--pad-header
--unpad-paren
--align-pointer=type
--align-reference=type
--attach-namespaces
--attach-extern-c
--convert-tabs
--max-code-length=110
--preserve-date
--suffix=none

# invoke it: quotes keep the glob inside astyle, not the shell
astyle --options=tools/astyle.options "src/*.cpp" "src/*.h"

# CI use: write only files that really changed, exit status tells you
astyle --options=tools/astyle.options --formatted "src/*.cpp"
\end{lstlisting}

几个字段的行为值得单独说：\Opt{style} 的 \texttt{ansi}、\texttt{linux}、
\texttt{otbs}（One True Brace）、\texttt{banner}（等价于 \texttt{allman}）
对应四种大括号习惯，另有 \texttt{java}、\texttt{gnu}、\texttt{stroustrup}、
\texttt{webkit} 等派生取值；\Opt{indent=force} 表示「所有非空行强制按规则缩进」，
不加它时 \Lib{astyle} 尊重你原有的缩进，这也是很多老仓库看起来「格式化了但没整齐」的原因；
\Opt{attach-namespaces} 决定命名块的左括号是否贴在同一行；
\Opt{max-code-length} 是「超过这个长度我才考虑断行」的阈值，
它不像 \texttt{ColumnLimit} 那样是硬约束，所以和 \Fmt\ 的结果不会逐行一致；
\Opt{formatted} 只在文件真的发生变化时写盘，并让退出码区分「无变化」与「有变化」，
这是 \Lib{astyle} 能接进 CI 的唯一正解。

\subsection{它擅长什么，会在哪里失手}\label{sec:ch11-s1-scope}

\Lib{astyle} 的可靠区间是「\CL{} 风格的缩进与空白」：改缩进宽度、统一大括号、
补运算符两侧空格、把指针星号靠到类型上。它的失手集中在依赖语法判断的构造上：
属性（\texttt{[[nodiscard]]}）、尾置返回类型、lambda 捕获表、
\texttt{requires} 子句、协程关键字这些写法，词法器只会把它们当成括号与运算符的组合，
于是可能出现「按普通减法断行」「把捕获表当数组下标对齐」这类结果。
列表 \ref{lst:ch11-modern} 是一段典型的现代写法，也是三家工具的试金石。

\begin{lstlisting}[style=cppstyle,caption={一段用于测格式化工具的现代写法},label={lst:ch11-modern}]
template <class T>
[[nodiscard]] auto lookup(const T &key) noexcept
    -> std::optional<Entry> {
  auto items = std::vector<Entry>{
      {key, 1},
      {key, 2},
  };
  auto pred = [&](const Entry &e) { return e.key == key; };
  switch (items.size()) {
    case 0: return std::nullopt;
    default:
      return std::ranges::find_if(items, pred) != items.end()
                 ? items.front()
                 : std::nullopt;
  }
}
\end{lstlisting}

适用场景可以总结成三条：老代码库只做「缩进与大括号」两件事的统一；
环境装不下 LLVM（嵌入式开发机、不能联网的构建镜像）；
需求是「有一个能离线跑、输出可预期、十年不换版本的工具」。
反过来，如果你需要 include 分组、模板与 lambda 的稳定排版、
或者想让格式化结果与 \CD\ 的实时格式化一致（第 5 章），
\Lib{astyle} 不是合适的权威。

\section{\Lib{uncrustify}：三类前缀与配置漂移}\label{sec:ch11-s2}

\Lib{uncrustify} 走的是另一条路：先做词法与自己的 token 类型标注（比 \Lib{astyle} 多一层
「这个 token 是什么」的判断），再套一张极大的逐构造规则表。
它的键族前缀就是它的文档：

\begin{itemize}
  \item \texttt{nl\_*}：换行。形如 \texttt{nl\_<token>\_<token>}，
    例如 \texttt{nl\_if\_brace}、\texttt{nl\_brace\_else}、
    \texttt{nl\_func\_type\_name}，取值多是 \texttt{ignore}/\texttt{add}/\texttt{force}
    或行数上下限（\texttt{nl\_max}）。
  \item \texttt{sp\_*}：空白。同样按 token 对命名，
    例如 \texttt{sp\_before\_ptr\_star}、\texttt{sp\_after\_comma}、
    \texttt{sp\_after\_semi}。
  \item \texttt{align\_*}：对齐，附带跨行数与列宽上限，
    例如 \texttt{align\_right\_cmt\_span}。
  \item 另有 \texttt{indent\_*}（缩进策略，\texttt{indent\_with\_tabs} 最常用）、
    \texttt{pos\_*}（把某个运算符放在行首还是行尾，例如 \texttt{pos\_bool}）、
    \texttt{ls\_*}（长行拆分）、\texttt{mod\_*}（改写代码本身，
    例如给长函数补结尾注释）、\texttt{code\_width}（目标宽度）。
\end{itemize}

\begin{lstlisting}[style=shellstyle,caption={\Lib{uncrustify} 的三个自查工具},label={lst:ch11-uncrustify}]
# language is chosen per run: c, cpp, cs, java, d, vala, ecma, pawn
uncrustify -l cpp -c repo.cfg -f old.cpp -o new.cpp

# dump every key with its default value for THIS binary
uncrustify -l cpp --update-config > defaults.cfg

# detect the style of a file and print all keys with detected values
uncrustify -l cpp -c base.cfg --detect --set-all -f old.cpp > detected.cfg

# compare: which keys did detection disagree with
diff -u defaults.cfg detected.cfg | grep -E '^[+-][a-z]' | head -40

# editor integration: format a fragment as if it came from that path
uncrustify -c repo.cfg -l cpp --frag --assume-filename old.cpp < part.cpp
\end{lstlisting}

\subsection{为什么它更啰嗦，也更可控}\label{sec:ch11-s2-why}

啰嗦是因为它把「构造 × 位置」笛卡尔积展开成键：每一种 token 对都要一条规则，
所以一份认真配过的 \texttt{uncrustify.cfg} 动辄上千行，键的数量是
\texttt{.clang-format} 的一个数量级以上。收益是控制权：
\texttt{pos\_*} 可以精确规定「逻辑运算符断在行尾还是行首」，
\texttt{mod\_*} 甚至能改写代码文本（补大括号、补命名空间结尾注释）。
\Fmt\ 的对应能力是「一个字段管一类形状 + 搜索决定断行位置」，
写起来短，但边界情况不由你说了算。

代价就是\textbf{配置漂移}：键名与语义在版本之间改过多次，
旧配置在新二进制上要么报「unknown option」，要么被静默忽略，
而另一个键的默认值恰好变了。

\begin{warning}[title={\texttt{nl\_*} 组合爆炸导致非幂等}]
  症状：同一份 cfg，跑第一次改 300 行，跑第二次又改 40 行；
  或者升级 \Lib{uncrustify} 之后，某个文件在两版之间来回被改。
  原因：\texttt{nl\_*}、\texttt{sp\_*}、\texttt{pos\_*}、\texttt{ls\_*} 是各自独立的
  局部规则，没有 \Fmt\ 那样的全局代价函数把它们收敛到一个稳定解；
  一条 \texttt{nl\_*} 要求某个左括号前必须换行，另一条 \texttt{pos\_*} 又要求相邻
  运算符回到同一行，两轮的净效果就可能不为零。
  对策三条：把 cfg 纳入版本控制并注释每个改过的键；钉死工具版本（同一个 tag 或同一个镜像）；
  每次升级先跑第 10 章的幂等自查脚本，再决定要不要接受新版本。
  真有冲突时，用 \texttt{uncrustify off}/\texttt{on} 注释圈出该段，
  别指望规则表自己收敛。
\end{warning}

\section{三种机制的能力边界}\label{sec:ch11-s3}

把「机制」讲清楚最重要：这三家没有一个基于完整 AST。
\Lib{astyle} 是词法；\Lib{uncrustify} 是词法加自己的 token 类型表；
\Fmt\ 是 token 加语境标注加带权搜索（四步流程已在第 9 章展开，此处不重复）。
真正的 AST 方案要求完整的语法树，也就要求编译数据库
（第 2 章），而且它无法处理未编译的分支——这是本质的取舍，不是实现好坏。

{\footnotesize
\begin{longtable}{%
  >{\raggedright\arraybackslash}p{0.14\textwidth}%
  >{\raggedright\arraybackslash}p{0.23\textwidth}%
  >{\raggedright\arraybackslash}p{0.22\textwidth}%
  >{\raggedright\arraybackslash}p{0.23\textwidth}%
  }
  \toprule
  维度 & \Lib{astyle} & \Lib{uncrustify} & \Lib{clang-format} \\
  \midrule
  \endhead
  机制 & 纯词法，选项即规则 & 词法加 token 类型表，逐构造规则 & token 加语境标注，断行是代价搜索 \\
  需要编译数据库 & 否 & 否 & 否（这也是它能对残缺代码工作的前提） \\
  模板、lambda、属性 & 支持滞后，按普通运算符处理 & 有专门键但需自己配全 & 较好，仍靠猜测与名单字段 \\
  宏感知 & 基本无，宏当成标识符 & 有名单与 \texttt{mod\_*} 补救 & 有 \texttt{ForEachMacros} 一类名单，仍按 token 猜 \\
  include 排序 & 无 & 无整体方案 & \texttt{IncludeBlocks} 与 \texttt{SortIncludes} \\
  断行控制粒度 & 粗（\Opt{max-code-length} 阈值） & 极细（键数量最多） & 中（靠 penalty 与对齐族） \\
  速度 & 最快 & 中 & 中（大文件仍远快于编译） \\
  幂等性 & 稳定，选项少故不易冲突 & 需自证，规则可互相撤回 & 稳定，但配置本身可能非幂等 \\
  可预测性 & 高，一行对应一个开关 & 中，需要读大量键 & 中低，输出由搜索给出 \\
  社区默认 & 无统一默认，随项目 & 无统一默认，随项目 & 有具名风格（LLVM、Chromium 等） \\
  许可证 & 宽松系（以官方发布页为准） & GPL 系，注意发布约束 & Apache-2.0 加 LLVM 例外 \\
  \bottomrule
  \caption{三种机制与三家工具的能力对照}
  \label{tab:ch11-mech}
\end{longtable}
}

\begin{compare}[title={能力边界：谁在什么情况下必须让位}]
  要「改动可控、输出十年不变」的嵌入式老仓库：\Lib{astyle} 够用且最省事。
  要「每一个构造都由我规定」并且愿意维护千行配置：\Lib{uncrustify}。
  要「和编辑器、\CD、代码评审工具共用同一套解析」，或者需要 include 排序、
  与 \cpp{20} 的新标准写法和平共处：\Fmt\ 是默认答案。
  三者的能力差别不在「支持多少字段」，而在\textbf{断行决策由谁给出}：
  前两家的换行是你写的规则，\Fmt\ 的换行是搜索算出来的。
这条边界决定了排错方式：规则型工具去查键，搜索型工具去查 penalty 与列限。
\end{compare}

\section{为什么 \CL/\CPP\ 没有 \texttt{gofmt}}\label{sec:ch11-s4}

这是本章的论证重点。先说清 \texttt{gofmt} 到底是什么：它无配置项、不可关闭、
由 Go 团队规定为唯一合法排版，编辑器与 \texttt{go fmt} 都转调它，
输出直接由语法树打印而来。它之所以能成立，是因为 Go 的设计替它扫清了障碍。

\subsection{\texttt{gofmt} 成立的四个前提}\label{sec:ch11-s4-go}

\begin{enumerate}
  \item \textbf{每个文件都能局部完整解析}。Go 没有头文件、没有包含展开、没有宏，
    一个 \texttt{.go} 文件的文本本身就是一棵完整语法树的输入；
    解析器不需要知道任何外部信息就能给出无歧义的 AST。
  \item \textbf{没有预处理器}。不存在「这一行在某个编译配置下根本不是代码」的情形，
    所以格式化器永远不会面对非法输入（除非文件真的写坏了，那时 \texttt{gofmt} 直接报错，
    而不是硬排版）。
  \item \textbf{语法本身刻意保持小}。运算符优先级固定、没有重载、没有模板参数列表，
    所以没有 \texttt{<} 与 \texttt{>} 的身份之争，也没有括号与类型转换的歧义。
  \item \textbf{委员会规定唯一排版，且输出与 AST 一一对应}。
    \texttt{gofmt} 先解析再打印，打印规则是规范的一部分，
    没有「同一棵树两种输出」的空间；不可配置也就没有配置漂移、没有幂等性问题、
    没有风格争论进入 code review。代价是个人偏好被彻底牺牲，收益是零讨论成本。
\end{enumerate}

\subsection{C/C++ 的四道墙}\label{sec:ch11-s4-walls}

\begin{enumerate}
  \item \textbf{预处理器使源文件在展开前不是合法 \CPP}。
    \texttt{\#if} 的分支里可能有一整段在当前配置下不成立的代码；
    \texttt{\#define} 可以把一行变成任意语法形状。
    想知道「这里到底是不是语句」，就得知道宏定义与激活分支，
    而这需要 include 路径与 \texttt{-D} 集合，也就是需要编译数据库。
    一个只吃文本的格式化器永远只能对「词法上像什么」做决定。
  \item \textbf{语法歧义无法靠局部信息消解}。\texttt{>>} 是移位还是两层模板闭合，
    \texttt{a * b} 是乘法还是声明，\texttt{\{} 是块、初始化列表、lambda 体还是
    语句表达式——这些在 \cpp{11} 放宽了 \texttt{>>} 规则之后仍然需要
    知道外层是否为模板、符号是否有类型含义。编译原理里的标准结论是：
    考虑到模板与宏，\CPP\ 的文本语法不属于可以用有限局部 lookahead
    可靠解析的类；能可靠解析的是它的某个子集。
    「一个保证正确解析所有合法 \CPP\ 文本的格式化器」在原则上不存在。
  \item \textbf{宏使「这段是什么语法」不可判定}。\texttt{LIST\_ADD(x, y)}
    是表达式、语句还是声明？三家工具的答案都是「看它像什么」，
    所以都必须靠名单（\Fmt\ 的 \texttt{ForEachMacros} 一族、
    \Lib{uncrustify} 的 \texttt{mod\_*} 与 token 类型）来打补丁。
    而 Go 没有这个概念。
  \item \textbf{方言与混编}。\CL、\CPP、Objective-C、CUDA、OpenCL、MSVC 扩展、
    GNU 扩展各自带了语法。举一个具体证据：\Fmt\ 的 \texttt{Language} 枚举里
    \textbf{根本没有} \texttt{ObjC} 这一项，Objective-C 只能借 \texttt{Cpp} 的词法去猜；
    \texttt{.cu} 里的 kernel 启动写法会被当成移位与比较运算符。
    任何「唯一排版」都必须先把语言边界谈清楚。
\end{enumerate}

还有一道非技术的墙：\textbf{没有权威}。ISO 标准只管空白用于分隔 token，
不管列与行；语言演化由提案驱动，没有任何机构有权规定「函数左括号必须同行」。
而「一次性改动全世界所有 \CPP\ 代码」在 blame、分支、跨仓库依赖上的成本无法支付。

\subsection{事实标准是可达成的近似}\label{sec:ch11-s4-fact}

结论不是「所以 C/C++ 排版无望」，而是「把 \texttt{gofmt} 拆成两半来实现」：
\textbf{工具负责执行，具名风格负责规定内容}。
\Fmt\ 的 \texttt{BasedOnStyle} 正是这个结构：LLVM 风格指南与内置 \texttt{LLVM} 风格同源、
Google 的内部规范对应 \texttt{Chromium}、\texttt{Microsoft} 对应微软的内部指南，
而 \texttt{Stroustrup} 与 C++ Core Guidelines 提供另一套「教材式」参照。
它们不是唯一，但每一套都\textbf{可执行、有 owner、能进 CI}（第 10 章）。
仓库里写清「我们以 \texttt{BasedOnStyle: LLVM} 为准，另外覆盖这七个字段」，
效果上就得到了一个 per-repo 的 \texttt{gofmt}：不可辩驳、机器判定、零讨论成本。

\begin{guideline}[title={选型先定「谁说了算」，再定工具}]
  三个候选答案，选一个写进仓库文档：
  （1）\textbf{委员会标准}——采用某个具名风格（LLVM、Chromium、Microsoft），
  本地差异全部视为 bug；（2）\textbf{仓库配置}——自定义 \texttt{.clang-format}
  加一份「为什么改这七个字段」的说明；（3）\textbf{历史惯例}——不强制，只做增量格式化。
  无论选哪个，都必须排除「两个排版器同时有权改写文件」这种状态，
  那才是真正的事故源（\ref{sec:ch11-s6}）。
\end{guideline}

\section{三个常被混为一谈的「格式」}\label{sec:ch11-s5}

\begin{itemize}
  \item \textbf{\Fmt\ 不检查正确性}。它做的是 pretty-print，不是校验：
    格式化通过不代表语法好，更不代表没有未定义行为。它也\textbf{不保证}在残缺输入上
    保持语义，所以对宏密集、条件编译复杂的文件，
    第一次引入格式化时要把「格式化后编译并跑测试」写进同一个 PR 的清单里
    （第 7 章的构建闭环正好承接这一步）。
    另外，原始字符串常量内部不会被改动，但它周围的换行与对齐会参与排版。
  \item \textbf{编译器 \texttt{-Wformat} 的「格式」是另一回事}。它检查的是
    \texttt{printf} 一类的\textbf{格式字符串}与实参类型是否匹配
    （\texttt{-Wformat=2} 还会管非常量格式串与可变参数调用），
    属于诊断与静态分析侧的事（第 12、13 章），和排版毫无关系。
    同理，\texttt{std::format} 与 \Tp{fmt} 里的「格式串」也和排版无关。
    搜索文档时请区分「代码格式化（formatting/pretty-print）」、
    「格式字符串（format string）」与「格式化输出（formatted output）」。
  \item \textbf{注释对齐是长期争议}。\texttt{AlignTrailingComments}
    与 \texttt{align\_*} 一族把行尾注释排到同一列。支持方的理由：
    位域表、状态机表、常量表里列对齐让信息密度骤增，
    手工维护这类表格本来就要对齐；反对方的理由：一次改一个变量名会让整块注释
    横向移动，制造与语义无关的 diff，而且对齐后的列在换行重排时最不稳定。
    可落地的折中是：注释默认 \texttt{DontAlign}，需要表格的地方用
    \texttt{// clang-format off}（或 \texttt{skip}）圈出一小块，
    并在评审时要求写明理由——这样两派的诉求都被显式记录下来。
\end{itemize}

\section{从 \Lib{astyle} 与 \Lib{uncrustify} 迁移}\label{sec:ch11-s6}

\subsection{选项映射}\label{sec:ch11-s6-map}

表 \ref{tab:ch11-map} 给出常见诉求的三家对应写法。
迁移不是把选项一一对抄，而是「先删掉能被 \texttt{BasedOnStyle} 覆盖的部分，
再只把 \Fmt\ 表达不了的诉求保留成显式字段」。

{\footnotesize
\begin{longtable}{%
  >{\raggedright\arraybackslash}p{0.13\textwidth}%
  >{\raggedright\arraybackslash}p{0.22\textwidth}%
  >{\raggedright\arraybackslash}p{0.19\textwidth}%
  >{\raggedright\arraybackslash}p{0.26\textwidth}%
  }
  \toprule
  诉求 & \Lib{astyle} & \Lib{uncrustify} & \Lib{clang-format} \\
  \midrule
  \endhead
  缩进宽度与制表符 & \Opt{indent=\-spaces=4}、\Opt{convert-tabs} &
    \texttt{indent\_with\-tabs} 与 \texttt{indent\_*} 族 &
    \texttt{IndentWidth}、\texttt{TabWidth}、\texttt{UseTab} \\
  大括号位置 & \Opt{style=ansi}、\texttt{linux}、\texttt{otbs}、\texttt{banner} &
    \texttt{nl\_*}，如 \texttt{nl\_if\-brace}、\texttt{nl\_brace\-else} &
    \texttt{BreakBefore\-Braces} 与 \texttt{BraceWrapping} \\
  运算符与参数空格 & \Opt{pad-oper}、\Opt{unpad-paren}、\Opt{pad-header} &
    \texttt{sp\_*}，如 \texttt{sp\_after\-comma} &
    \texttt{SpaceBeforeParens} 等一族布尔与枚举字段 \\
  指针与引用符号 & \Opt{align\-pointer=type}、\Opt{align\-reference=same} &
    \texttt{sp\_before\-ptr\-star}、\texttt{sp\_after\-ptr\-star} &
    \texttt{PointerAlignment} 与 \texttt{ReferenceAlignment} \\
  行宽上限 & \Opt{max-code-length}（阈值） & \texttt{code\_width} &
    \texttt{ColumnLimit}（硬约束，配 penalty） \\
  断行位置 & 由 \Opt{break\-after\-logical} 等零散开关决定 &
    \texttt{pos\_*} 与 \texttt{ls\_*} & 由代价搜索给出，可调 penalty 族 \\
  include 顺序 & 无 & 无整体方案 &
    \texttt{IncludeBlocks}、\texttt{SortIncludes}、\texttt{Include\-Categories} \\
  空行密度 & \Opt{break\-blocks}、\Opt{keep\-one\-line\-statements} &
    \texttt{nl\_max} 与 \texttt{nl\_blank\-line} 一族 &
    \texttt{MaxEmptyLines\-To\-Keep}、\texttt{Separate\-Definition\-Blocks} \\
  \bottomrule
  \caption{同一批排版诉求在三家工具里的对应写法}
  \label{tab:ch11-map}
\end{longtable}
}

\subsection{迁移七步}\label{sec:ch11-s6-steps}

\begin{lstlisting}[style=shellstyle,caption={把旧工具换成 \Fmt\ 的操作顺序},label={lst:ch11-steps}]
# 1) freeze the old config, keep it as documentation
git mv tools/astyle.options docs/formatters-2024/astyle.options

# 2) record the current formatting so later steps can measure drift
git ls-files -z -- '*.cpp' '*.h' > /tmp/list
xargs -0 < /tmp/list clang-format --style=file --dump-config > /tmp/cf.def

# 3) write the new .clang-format: BasedOnStyle plus the mapped keys
# 4) one mechanical commit, nothing else
xargs -0 < /tmp/list clang-format --style=file -i
git add -u && git commit -m "style: move from astyle to clang-format 18"

# 5) make blame skip it, in the same pull request
git rev-parse HEAD >> .git-blame-ignore-revs

# 6) remove the old formatter from every entry point
grep -rn "astyle\|uncrustify" .github .pre-commit-config.yaml tools/

# 7) prove the new gate is self-consistent and idempotent
./tools/check-idempotent.sh
\end{lstlisting}

第 6 步是全章最实际的一条：\textbf{只要旧工具还留在任何一个入口}
（CI 脚本、某个 Makefile 目标、某个 IDE 外部工具），
它和新工具就会互相把对方的输出改回去，形成无限往返。
判据很简单：连续两次「先跑新工具再跑旧工具」，若文件仍变化，说明两者不可共存。

列表 \ref{lst:ch11-roundtrip} 展示一个典型的往返：纯词法工具按「运算符两侧加空格」
处理三元与 \texttt{>}，而搜索式工具按列限与代价重新决定断行位置。

\begin{lstlisting}[style=diffstyle,caption={同一文件在两种机制之间往返的结果},label={lst:ch11-roundtrip}]
--- a/parse.cpp   (after: lexical formatter, pad all operators)
+++ b/parse.cpp   (after: clang-format on the same file)
@@ ternary @@
-      return ok
-              ? value
-              : fallback;
+      return ok ? value : fallback;
@@ brackets @@
-  auto v = map[
-    key
-  ];
+  auto v = map[key];
@@ decl @@
-void f(
-  int a,
-  int b) {}
+void f(int a, int b) {}
\end{lstlisting}

\subsection{反向迁移的代价}\label{sec:ch11-s6-reverse}

从 \Fmt\ 退回 \Lib{astyle} 或 \Lib{uncrustify} 的损失是结构性的，不只是「少几个字段」：
一是失去 include 分组与排序，依赖方向的问题不再一眼可见；
二是失去断行搜索，长签名与长三元会退化成「靠人工决定断在哪儿」，
而被 \Fmt\ 重排过的那些断行在纯词法工具里会被视为「你故意的」，从此不再修正；
三是 \texttt{.clang-format} 里靠名单字段撑起来的宏与 lambda 处理失效，
需要人工回滚大量注释与表格。
如果确实要退回去，先把 \texttt{ColumnLimit} 提到很高（例如 0 或 200）再迁移，
让 \Fmt\ 只做空白归一而几乎不断行，能显著减少回退后的手工修补量。

\section*{本章要点}
\begin{itemize}
  \item \Lib{astyle} 是词法工具，快、稳、可预测，但对模板、lambda、属性滞后；
    它适合「只统一缩进与大括号」的老代码库与装不下 LLVM 的环境。
  \item \Lib{uncrustify} 是规则表工具，\texttt{nl\_*}、\texttt{sp\_*}、\texttt{align\_*}
    三族加 \texttt{code\_width} 换来极细的控制权，代价是配置漂移与非幂等风险，
    升级前必须跑幂等检查。
  \item 三家都不是 AST 格式化器，所以都不需要编译数据库；
    真正的 AST 方案会立刻被预处理器与未编译分支卡住。
  \item 没有 \texttt{gofmt} 的原因是四道墙：预处理、语法歧义、宏、方言与无权威。
    可达成的近似是「工具执行 + 具名风格规定内容」，
    即一份写明 \texttt{BasedOnStyle} 的 \texttt{.clang-format}。
  \item 迁移的关键不是抄选项，而是让旧工具从所有入口消失，
    并把一次性全量格式化与 blame 屏蔽放在同一个 PR 里。
\end{itemize}

% ---- frag_ch12.tex ----
\part{静态分析}
\chapter{\Tidy\ 架构与 checks 选择}\label{ch:tidy-arch}
\section{\Tidy\ 的分层架构与自定义 check}\label{sec:ch12-s1}
\subsection{五层结构}\label{sec:ch12-s1a}
\begin{itemize}
\item \textbf{AST matcher 层}：\CPP\ 前端把源码建成抽象语法树后，\texttt{clang::ast\_matchers} 提供一套声明式 DSL——\texttt{decl(...)}、\texttt{callExpr(...)}、\texttt{hasAncestor(...)} 等谓词可以像组合正则表达式一样组合。所谓「对 AST 的正则」就是这个意思：模式声明一次，遍历引擎负责匹配整棵树，规则作者不写递归访问代码。
\item \textbf{规则实现层}：每个 check 一族一个子目录（\texttt{clang-tidy/}\allowbreak\texttt{bugprone/}、\texttt{clang-tidy/}\allowbreak\texttt{modernize/} 等），内部把 matcher 命中翻译为诊断信息与 fix-it。
\item \textbf{基类层}\Tp{clang::tidy::ClangTidyCheck}：约定三个必须理解的虚函数——\Fn{registerMatchers}（AST 模式类）、\Fn{check}（命中后动作）、\Fn{storeOptions}（把可配置项落到选项表；需要时还要重写 \Fn{registerOptions} 读回）。
\item \textbf{模块注册层}\Tp{ClangTidyModule}：一个模块注册一批 check，名字前缀（\texttt{bugprone-}、\texttt{modernize-}）就来自这里。
\item \textbf{driver 层}：对每个翻译单元（TU）建 AST、跑全部启用的 check，并把结果汇入统一的诊断列表；带 autofix 的结果由 \Tidy\ 内部的修复回写机制落盘，工具链细节见第 14 章。
\end{itemize}
\subsection{自定义 check 骨架}\label{sec:ch12-s1b}
下面是一段接口正确、但不追求可直接编译的骨架，展示「如果我要给公司写一条自定义 check」的最小工作量：
\begin{lstlisting}[style=cppstyle,caption={自定义 \Tidy\ check 骨架（示意）},label={lst:ch12-check}]
#include "clang-tidy/ClangTidyCheck.h"
#include "clang-tidy/ClangTidyModule.h"
#include "clang-tidy/ClangTidyModuleRegistry.h"

namespace clang::tidy::mycorp {

class NoSleepInLoopCheck : public ClangTidyCheck {
public:
  NoSleepInLoopCheck(llvm::StringRef Name, ClangTidyContext *Context)
      : ClangTidyCheck(Name, Context) {}

  void registerMatchers(ast_matchers::MatchFinder *Finder) override {
    Finder->addMatcher(
        ast_matchers::callExpr(
            ast_matchers::callee(
                ast_matchers::functionDecl(ast_matchers::hasName("sleep"))),
            ast_matchers::hasAncestor(ast_matchers::forStmt()))
            .bind("call"),
        this);
  }

  void check(const ast_matchers::MatchFinder::MatchResult &R) override {
    const auto *Call = R.Nodes.getNodeAs<clang::CallExpr>("call");
    diag(Call->getBeginLoc(),
         "do not busy-sleep inside a loop; use a timer or condition variable");
  }

  void storeOptions(ClangTidyOptions::OptionMap &Opts) override {
    Options.store(Opts, "Strict", Strict);
  }

private:
  bool Strict = false;
};

class MyCorpModule : public ClangTidyModule {
public:
  void addCheckFactories(ClangTidyCheckFactories &Factories) override {
    Factories.registerCheck<NoSleepInLoopCheck>("mycorp-no-sleep-in-loop");
  }
};

// 静态初始化对象：链接进来即完成注册
static ClangTidyModuleRegistry::Add<MyCorpModule> MyCorpReg(
    "mycorp-module", "Checks required by my company.");

} // namespace clang::tidy::mycorp
\end{lstlisting}
\subsection{外部模块与 \Opt{load}}\label{sec:ch12-s1c}
主干 \Tidy\ 只认识 \texttt{clang-tools-extra/clang-tidy/} 里编译进去的族。公司私有规则有两条挂接路径：其一，把上面的源码编进定制版 clang-tidy；其二，编译成独立动态库，用命令行加载——这是旧机制的写法，clang 17 之后还提供 \texttt{clang-tidy/plugin.h} 的插件入口，思路相同（把模块注册进 \Tp{ClangTidyModuleRegistry}）：
\begin{lstlisting}[style=shellstyle,caption={加载外部 check 模块并确认注册成功},label={lst:ch12-load}]
clang-tidy --load=./libMyCorpChecks.so main.cpp -p build \
  --checks='-*,mycorp-*' --list-checks
\end{lstlisting}
\begin{guideline}[title={分层决定排障方向}]
误报/漏报问题查 AST matcher 层与数据流层；「我的 check 怎么没被跑到」查注册层与 driver 层；「改了 \texttt{.clang-tidy} 没生效」几乎总是 driver 层的配置查找顺序问题，见第 12.4 节。
\end{guideline}
\section{数据流、CFG 与成本模型}\label{sec:ch12-s2}
\subsection{三类实现，三档代价}\label{sec:ch12-s2a}
第一类是\textbf{纯 AST 模式}：命中即报，不看不跨语句的状态。例如 readability-implicit-bool-conversion 只需判断「这个条件表达式是否把整型隐式转成了布尔」，函数签名是否 const 都能当场定，无需任何流分析。第二类是\textbf{数据流 check}：在 matcher 之外，对每个函数体用 \texttt{clang::analysis::cfg::buildCFG} 建控制流图（CFG），再跑活跃性（liveness）与污点/状态分析。bugprone-use-after-move 是典型：必须沿 CFG 的所有边传播「这个对象已被 move」的状态，并处理分支、循环与重新赋值——纯 AST 只能看到「用了」和「move 了」两个点，无法判断先后与可达性。代价是：更慢（每函数建图）、更容易误报（抽象有限）、更依赖完整 TU（模板实例化、头文件里的定义都要能解析）。第三类是\textbf{analyzer 桥接}：\Tp{clang::tidy::analyzer::Checker} 把 clang static analyzer 的引擎整个搬进 \Tidy ，路径敏感、带约束求解，能发现空指针解引用沿哪条路径成立，代价也最大（详见第 13.6 节）。
\begin{warning}[title={数据流 check 在边界处会「装死」}]
遇到复杂宏展开、协程帧、未解析的模板依赖时，\texttt{bugprone-*} 的数据流实现通常保守放弃而不是猜测。这意味着同一段代码，编译数据库配全了会报、配缺了不报——把「不报」误读成「没问题」是门禁设计的大忌。
\end{warning}
\subsection{成本对照表}\label{sec:ch12-s2b}
{\small
\begin{longtable}{>{\raggedright\arraybackslash}p{0.30\textwidth} >{\raggedright\arraybackslash}p{0.14\textwidth} >{\raggedright\arraybackslash}p{0.12\textwidth} >{\raggedright\arraybackslash}p{0.12\textwidth} >{\raggedright\arraybackslash}p{0.13\textwidth}}
\textbf{代表族/检查} & \textbf{需数据流} & \textbf{单 TU 成本} & \textbf{误报倾向} & \textbf{autofix}\\ \midrule
\endfirsthead
\textbf{代表族/检查} & \textbf{需数据流} & \textbf{单 TU 成本} & \textbf{误报倾向} & \textbf{autofix}\\ \midrule
\endhead
modernize-\allowbreak use-\allowbreak override 等模式类 & 否 & 低 & 低 & 是\\
readability-\Bk magic-numbers & 否 & 低 & 中（噪声） & 否\\bugprone-unnecessary-copy-initialization & 否 & 低 & 低 & 是\\
bugprone-use-after-move & 是（CFG 与活跃性） & 中 & 中 & 否\\
performance-unnecessary-value-param & 部分（类型大小） & 低--中 & 低 & 是（受限）\\
cppcoreguidelines-pro-bounds & 否 & 低 & 低 & 否\\
concurrency-mt-unsafe & 否 & 低 & 中 & 否\\
clang-analyzer-\Bk{*} & 是（路径敏感引擎） & 高 & 中--高 & 否\\
\end{longtable}}
\section{checks 命名族全景}\label{sec:ch12-s3}
check 名一律是 \texttt{族名-具体名}，族名决定它归属哪个模块注册表。下面按「定位」与「该不该默认开」过一遍全部官方族：
{\small
\begin{longtable}{>{\raggedright\arraybackslash}p{0.20\textwidth} >{\raggedright\arraybackslash}p{0.47\textwidth} >{\raggedright\arraybackslash}p{0.13\textwidth}}
\textbf{族} & \textbf{定位} & \textbf{默认开?}\\ \midrule
\endfirsthead
\textbf{族} & \textbf{定位} & \textbf{默认开?}\\ \midrule
\endhead
modernize & 把旧写法升级为 \cpp{11} 之后的惯用法：override、range-for、智能指针、\texttt{using}。语义基本等价，最适合自动化。 & 是\\
bugprone & 真 bug 高发区：use-after-move、悬垂、移位/交换错误、断言副作用。 & 是\\
performance & 无谓拷贝、多余 move、循环内拷贝、枚举过大等运行时开销。 & 是（个别 ABI 项除外）\\
readability & 风格与可维护性：magic numbers、大括号、函数长度、标识符长度。主观性强。 & 挑着开\\
cppcoreguidelines & C++ 核心指南条款映射，覆盖面大且部分激进（如 owning-memory）。 & 挑着开\\
clang-analyzer & clang static analyzer 的检查器借 \Tidy\ 通道输出，路径敏感，见第 13.6 节。 & 是（预算好算力）\\
misc & 杂项：const-correctness、no-recursion、include-cleaner。噪声集中营。 & 否\\
portability & 可移植性：restrict-system-includes、非标准类型/函数。 & 有跨平台需求时\\
concurrency & 线程安全：mt-unsafe（线程不安全 API）、无锁并发相关。 & 是\\
abseil & 强制/迁移 Google Abseil 库风格，项目不用 Abseil 毫无意义。 & 否\\
google & Google C++ 风格专属项（如 build-explicit-overrides）。 & 否\\
altera & FPGA 实时系统扩展（C 方言）。 & 否\\
fuchsia & Fuchsia 内核/系统规范专属项。 & 否\\
hicpp & AUTOSAR 式工业编码指南的 clang-tidy 表达，与 CERT 风格相近。 & 合规要求时\\
linuxkernel & Linux 内核源码树专属检查（内核编码风格与惯用法）。 & 否\\
llvmlibc & LLVM libc 项目专用约束。 & 否\\
mpi & MPI 调用用法检查。 & 用 MPI 时\\
objc & Objective-C 族。 & 否\\
openmp & OpenMP 用法检查。 & 用 OpenMP 时\\
zircon & Zircon 微内核规范。 & 否\\
\end{longtable}}
\section{配置：\texttt{Checks} 选择器与 \texttt{.clang-tidy} 结构}\label{sec:ch12-s4}
\subsection{\texttt{Checks} 的匹配语法}\label{sec:ch12-s4a}
选择器是逗号分隔的通配列表，\textbf{从前到后依次生效}，所以「先减后加、再加回个别」的顺序有意义。推荐姿势是 \texttt{-*} 打底、白名单逐族加，而不是一开始就 \texttt{*} 再往回关：
\begin{lstlisting}[style=shellstyle,caption={选择器与 \Opt{list-checks} 的正确用法},label={lst:ch12-checks}]
# 查看当前配置到底启用了哪些 check
clang-tidy --checks='-*,bugprone-*,-bugprone-easily-swappable-parameters' \
  --list-checks main.cpp -p build

# 命令行 --checks= 会【整体覆盖】.clang-tidy 里的 Checks 键
clang-tidy main.cpp -p build --checks='-*,modernize-use-override'
\end{lstlisting}
\subsection{\texttt{.clang-tidy} 的顶层键}\label{sec:ch12-s4b}
\begin{itemize}
\item \texttt{Checks}（小写别名 \texttt{check} 也可）：选择器列表。
\item \texttt{CheckOptions}：逐 check 的参数，见 12.4.3。
\item \texttt{HeaderFilterRegex}：正则决定哪些头文件内的诊断要显示；不写则默认只看被分析的 TU 自身。
\item \texttt{HeaderFileExtensions} / \texttt{ImplementationFileExtensions}：声明哪些后缀算头/实现文件（默认已含 \texttt{h,hpp,hh} 等，嵌入式项目常要加 \texttt{inc}）。
\item \texttt{CompilationDatabase}：个别版本的实验键，跨版本脚本请继续用 \texttt{-p} 或 \Opt{compile-commands-dir}，别指望写进配置文件。
\item \texttt{User}：署名，出现在部分诊断与建议文本里。
\item \texttt{FormatStyle}：autofix 后的重排版风格，\texttt{file} 表示读取 \texttt{.clang-format}（与第 9 章呼应）。
\item \texttt{WarningsAsErrors}：把指定 check 升级为 error 以便 CI 变红，见 12.5.3。
\item \texttt{ExtraArgs} / \texttt{ExtraArgsBefore}：向 clang 前端追加参数，如 \texttt{-std=c++20}；\texttt{Before} 排在原命令行之前。
\item \texttt{UseColor}：终端彩色输出，CI 里通常设 \texttt{false}。
\item \texttt{InheritParentConfig}：置 \texttt{true} 时与上层目录的配置合并而不是覆盖。
\end{itemize}
\subsection{查找顺序与继承}\label{sec:ch12-s4c}
\begin{lstlisting}[style=yamlstyle,caption={一份完整的 \texttt{.clang-tidy} 骨架},label={lst:ch12-full}]
Checks: >-
  -*,
  bugprone-*,
  -bugprone-easily-swappable-parameters,
  modernize-use-override,
  modernize-loop-convert,
  performance-unnecessary-value-param,
  readability-braces-around-statements
HeaderFilterRegex: '^(src|include)/.*'
FormatStyle: file
UseColor: false
WarningsAsErrors: >-
  bugprone-use-after-move,
  clang-analyzer-core.NullDereference
CheckOptions:
  readability-magic-numbers.IgnoredIntegerValues: '1;2;3;4;8;16'
  performance-for-range-copy.WarnOnAllAutoCopies: 'true'
ExtraArgs: ['-std=c++20']
\end{lstlisting}
\begin{note}[title={配置到底读的是哪一份}]
对某个源文件，\Tidy\ 从该文件所在目录开始逐级向上查找 \texttt{.clang-tidy}，找到第一份即停（除非该文件写了 \texttt{InheritParentConfig: true} 则继续向上合并）。命令行 \Opt{config=}（直接给 YAML 串）与 \Opt{config-file=}（给路径）优先级最高，完全绕过文件查找；\texttt{--config} 指定的内容不会与文件继承。调试时先跑 \texttt{clang-tidy --dump-config} 看它最终生效的配置。
\end{note}
\subsection{\texttt{CheckOptions} 的点分键名}\label{sec:ch12-s4d}
选项键 = \texttt{check全名.选项名}，必须点分、不能嵌套 YAML 树；值一律写成字符串（列表用分号分隔）。真实可用的例子：
\begin{lstlisting}[style=yamlstyle,caption={八个经确认存在的 \texttt{CheckOptions} 键},label={lst:ch12-opts}]
CheckOptions:
  readability-magic-numbers.IgnoredIntegerValues: '1;2;3;8'
  readability-magic-numbers.IgnoredFloatingPointValues: '1.0;0.0'
  readability-identifier-length.MinimumVariableNameLength: '2'
  readability-function-size.LineThreshold: '80'
  performance-for-range-copy.WarnOnAllAutoCopies: 'false'
  bugprone-argument-comment.StrictMode: 'true'
  cppcoreguidelines-special-member-functions.AllowSoleDefaultDtor: 'true'
  cppcoreguidelines-pro-type-vararg.IgnoreMacros: 'false'
\end{lstlisting}
键名随版本演进，动手前先跑 \texttt{clang-tidy --list-checks --checks=族名-*} 并查对应版本的文档页，不确定的键宁可不写——写错的键会被静默忽略。
\section{诊断抑制与严重级别}\label{sec:ch12-s5}
\subsection{\texttt{NOLINT} 三件套}\label{sec:ch12-s5a}
\begin{lstlisting}[style=cppstyle,caption={三种抑制写法与它们的生效范围},label={lst:ch12-nolint}]
void Handle(Request req) {
  auto body = std::move(req.body);
  // NOLINTNEXTLINE(bugprone-use-after-move)
  Log(body);                     // 只有下一行被抑制，且只抑制指定 check

  Send(body);                    // NOLINT 抑制「本行」报出的诊断
  Process();                     // NOLINT 后不带 check 名 = 该行全部 check 闭嘴

  // NOLINTBEGIN(readability-magic-numbers)
  int table[4] = {1024, 2048, 4096, 8192};  // 区间内只抑制这一族
  // NOLINTEND(readability-magic-numbers)
}
\end{lstlisting}
生效规则容易记错：\texttt{NOLINT} 只看它\textbf{所居那一行}产生的诊断；\texttt{NOLINTNEXTLINE} 只看下一行；带括号 check 名列表时只抑制列表内 check，不带则全部。诊断的行号取的是 check 报告的起始位置，宏展开后可能落在意想不到的行——抑制不干净先检查行号语义。
\subsection{从 \CD\ 视角与基线}\label{sec:ch12-s5b}
\CD\ 在编辑器侧只按 \texttt{.clangd} 中 \texttt{Diagnostics} 组的设置跑一个子集（第 5 章），并用同组的 \texttt{Suppress} 列表按 check 名（含 \texttt{clang-diagnostic-*}）过滤显示——这是显示层抑制，不影响 CI。存量代码的基线化没有官方开关，通行做法是：全量跑一遍 \Tidy ，把 \texttt{file:line:check} 三元组生成 suppress 清单，CI 里只对新出现的诊断计数——工具本身（如 \texttt{run-clang-tidy}）只提供原始输出，清单与比对要自己写脚本。
\subsection{严重级别与门禁}\label{sec:ch12-s5c}
\Tidy\ 的诊断默认都是 warning 级（另附 note 说明移动点/首次声明点）。让 CI 变红有三档：\texttt{WarningsAsErrors: '*'} 全部升级；指定列表只升级红线项（见清单 \ref{lst:ch12-full}）；命令行 \Opt{warnings-as-errors=} 临时覆盖。\Opt{errors-as-warnings} 反向操作——即便被升级为 error 也继续跑完并以 0 退出，适合迁移期观察。\Opt{quiet} 只保留 error 级输出，配合白名单红线使用。
\begin{warning}[title={autofix 会顺手重排整个 include 块}]
开着 \texttt{FormatStyle: file} 时，\Opt{fix} 在写回前按 \texttt{.clang-format} 重排被改动的区域；如果 \Fmt\ 配置与团队现用风格不一致，一次修复会制造几百行无关 diff。先对齐格式化配置，再开 autofix。
\end{warning}
\subsection{与 clang 诊断的分工}\label{sec:ch12-s5d}
clang 前端诊断（\texttt{-Wall}、\texttt{-Wextra} 一族）工作在语法/语义层：未用变量、隐式窄化、格式化串不匹配。\Tidy\ 工作在 AST 模式与数据流层：它能看到调用关系、跨语句状态，但看不到（也不该看）预处理前的词法问题。两者在 \Tidy\ 输出里共用一套通道——clang 诊断以 \texttt{clang-diagnostic-*} 名义出现，因此既能被 \texttt{Checks:\allowbreak '-clang-diagnostic-\allowbreak unused-parameter'} 关掉，也能被 \CD\ 的 \texttt{Diagnostics.\allowbreak Suppress} 过滤。
\begin{compare}[title={编译器警告与 \Tidy\ 检查}]
编译器警告随每次编译产生、零额外解析成本、覆盖窄而稳；\Tidy\ 需要编译数据库、需要完整解析 TU，覆盖宽但有算力与误报税。工程上两者都要：门禁里编译器警告全量常开，\Tidy\ 白名单逐步扩。
\end{compare}
\section{四套起步配置}\label{sec:ch12-s6}
\subsection{保守起步（第一天就能合入）}\label{sec:ch12-s6a}
\begin{lstlisting}[style=yamlstyle,caption={保守集：只挑几乎零误报且可自动修复的项},label={lst:ch12-cfg1}]
Checks: >-
  -*,
  bugprone-use-after-move,
  bugprone-uninitialized-object,
  modernize-use-override,
  modernize-use-equals-default,
  modernize-use-using,
  performance-unnecessary-value-param
HeaderFilterRegex: 'src/.*'
FormatStyle: file
\end{lstlisting}
取舍逻辑：全部是命中即真、autofix 安全的项，\texttt{-} 掉主观风格类，存量项目一晚上能清零。
\subsection{严格模式（内部产品，接受人工复核）}\label{sec:ch12-s6b}
\begin{lstlisting}[style=yamlstyle,caption={严格集：数据流与 analyzer 全开，红线显式化},label={lst:ch12-cfg2}]
Checks: >-
  -*,
  bugprone-*,
  -bugprone-easily-swappable-parameters,
  -bugprone-narrowing-conversions,
  performance-*,
  modernize-*,
  readability-identifier-length,
  cppcoreguidelines-init-variables,
  cppcoreguidelines-pro-bounds,
  concurrency-mt-unsafe,
  clang-analyzer-core.*,
  clang-analyzer-cplusplus.*
WarningsAsErrors: >-
  bugprone-use-after-move,
  clang-analyzer-core.NullDereference
HeaderFilterRegex: '.*'
FormatStyle: file
ExtraArgs: ['-std=c++20']
\end{lstlisting}
取舍逻辑：数据流与 analyzer 全开，但要接受单 TU 秒级成本；两条著名噪声项显式关掉，红线只有两条。
\subsection{开源库（面向公共 API，尊重使用者）}\label{sec:ch12-s6c}
\begin{lstlisting}[style=yamlstyle,caption={开源库集：只管 src，不管对外头文件的风格},label={lst:ch12-cfg3}]
Checks: >-
  -*,
  bugprone-*,
  -bugprone-easily-swappable-parameters,
  performance-*,
  modernize-avoid-c-arrays,
  modernize-use-override,
  readability-inconsistent-declaration-parameter-name
HeaderFilterRegex: '^(include|api)/.*'
ImplementationFileExtensions: ['cpp', 'cc']
WarningsAsErrors: '*'
ExtraArgsBefore: ['-std=c++17']
\end{lstlisting}
取舍逻辑：开源库的 ABI 与公共头是资产，\texttt{modernize-*} 全量会强推破坏兼容的改写（如 enum-size），因此只留白名单；对内的 src 目录才全量开。
\subsection{嵌入式（资源受限，\CL\ 混合）}\label{sec:ch12-s6d}
\begin{lstlisting}[style=yamlstyle,caption={嵌入式集：小 CPU 预算，突出 ISR/并发/边界},label={lst:ch12-cfg4}]
Checks: >-
  -*,
  bugprone-assert-side-effect,
  bugprone-infinite-loop,
  bugprone-misplaced-widening-cast,
  misc-no-recursion,
  portability-restrict-system-includes,
  readability-function-size,
  cppcoreguidelines-pro-bounds-constant-array-index
CheckOptions:
  readability-function-size.LineThreshold: '60'
  portability-restrict-system-includes.Includes: ['<stdint.h>', '<string.h>']
HeaderFileExtensions: ['h', 'inc']
ExtraArgs: ['-std=c17']
\end{lstlisting}
取舍逻辑：不跑 analyzer 与数据流族（交叉编译镜像里往往连完整头都凑不齐），聚焦递归、栈上大数组、边界与包含约束；\texttt{.inc} 后缀要显式声明，否则头文件里的诊断全部隐身。
\begin{bestpractice}[title={先小集合跑通，再扩}]
任何一套起步配置的落地顺序都一样：白名单三五个 check → 全库清零 → \texttt{WarningsAsErrors} 锁死这批 → 每两周扩一族。直接从 \texttt{*} 开始的项目没有活到第二周的。
\end{bestpractice}

% ---- frag_ch13.tex ----
\chapter{常用检查族实战}\label{ch:tidy-checks}
本章是按族组织的案例集。每个案例固定四段：\textbf{触发代码}（cppstyle 清单）→ \textbf{\Tidy\ 诊断原文}（shellstyle 清单）→ \textbf{修复后代码}（cppstyle 或 diffstyle）→ \textbf{一句备注}（能否自动修复、需要什么标准版本）。诊断行默认呈现为 \texttt{warning:}；一旦在 \texttt{WarningsAsErrors}（第 12 章）里点名该 check，同样文本会变成 \texttt{error:} 并使 \Tidy\ 以非零码退出。

\section{bugprone 族实战}\label{sec:ch13-s1}
bugprone 族是「命中即大概率是 bug」的族，建议整族开启后再逐条关噪声。
\subsection{bugprone-use-after-move}\label{sec:ch13-s1a}
\begin{lstlisting}[style=cppstyle]
std::string s = LoadPayload(id);
Send(std::move(s));
Audit(s.size());                 // s 已处于 moved-from 状态
\end{lstlisting}
\begin{lstlisting}[style=shellstyle]
session.cpp:3:8: warning: 's' was used after being moved [bugprone-use-after-move]
session.cpp:2:3: note: 's' has been moved here
\end{lstlisting}
\begin{lstlisting}[style=diffstyle]
-  Send(std::move(s));
-  Audit(s.size());
+  Audit(s.size());
+  Send(std::move(s));
\end{lstlisting}
不可自动修复（重排语句需要理解控制流）；基于 CFG 与活跃性分析，\cpp{11} 起有效。
\subsection{bugprone-unnecessary-copy-initialization}\label{sec:ch13-s1b}
\begin{lstlisting}[style=cppstyle]
void Render(const std::string &name) {
  std::string copy = name;       // 只是换个名字的冗余拷贝
  Draw(copy);
}
\end{lstlisting}
\begin{lstlisting}[style=shellstyle]
render.cpp:2:15: warning: variable 'copy' is initialized with a copy of
'name', but there is no use that requires the copy [bugprone-unnecessary-copy-initialization]
\end{lstlisting}
\begin{lstlisting}[style=diffstyle]
-  std::string copy = name;
-  Draw(copy);
+  Draw(name);
\end{lstlisting}
提供 autofix；纯 AST 检查，任何 \CPP\ 标准都有效。
\subsection{bugprone-dangling-handle}\label{sec:ch13-s1c}
\begin{lstlisting}[style=cppstyle]
std::string_view Name() {
  return std::string("temp");    // view 指向即将销毁的临时对象
}
\end{lstlisting}
\begin{lstlisting}[style=shellstyle]
name.cpp:2:10: warning: handle 'string_view' is dangling because
it points at a temporary owned by no one [bugprone-dangling-handle]
\end{lstlisting}
\begin{lstlisting}[style=diffstyle]
-std::string_view Name() {
-  return std::string("temp");
+std::string Name() {
+  return "temp";
\end{lstlisting}
不可自动修复（返回类型是接口的一部分）；针对 \cpp{17} 的 \Tp{string\_view} 类「把手」类型。
\subsection{bugprone-implicit-widening-of-multiplication-result}\label{sec:ch13-s1d}
\begin{lstlisting}[style=cppstyle]
int32_t w = 46341, h = 46341;
int64_t pixels = w * h;          // 乘法先在 32 位里溢出，再拓宽
\end{lstlisting}
\begin{lstlisting}[style=shellstyle]
img.cpp:2:18: warning: result of multiplication of 'w * h' computed
in 32-bit type then implicitly widened to 64-bit [bugprone-implicit-widening-of-multiplication-result]
\end{lstlisting}
\begin{lstlisting}[style=diffstyle]
-  int64_t pixels = w * h;
+  int64_t pixels = static_cast<int64_t>(w) * h;
\end{lstlisting}
提供 autofix（插入显式 cast）；图像尺寸、缓冲区字节数最容易踩。
\subsection{bugprone-assert-side-effect}\label{sec:ch13-s1e}
\begin{lstlisting}[style=cppstyle]
assert(Next(&cursor) && "advance failed");   // Next 会修改 cursor
\end{lstlisting}
\begin{lstlisting}[style=shellstyle]
parser.cpp:7:5: warning: assert condition has side effects;
function call 'Next' will be removed in NDEBUG builds [bugprone-assert-side-effect]
\end{lstlisting}
\begin{lstlisting}[style=diffstyle]
-  assert(Next(&cursor) && "advance failed");
+  bool ok = Next(&cursor);
+  assert(ok && "advance failed");
\end{lstlisting}
不可自动修复（是 bug 还是笔误需要人判断）；用选项 \texttt{AssertMacros} 声明项目自定义断言宏，否则默认只认 \texttt{assert}。
\subsection{bugprone-parent-virtual-call}\label{sec:ch13-s1f}
\begin{lstlisting}[style=cppstyle]
struct Base { virtual void render(); };
struct Derived : Base {
  void render() override { Base::Render(); }   // 大小写笔误：调到了重载别函数
};
\end{lstlisting}
\begin{lstlisting}[style=shellstyle]
derived.cpp:2:26: warning: parent call calls 'Render' which is not the
overridden virtual function 'render' [bugprone-parent-virtual-call]
\end{lstlisting}
\begin{lstlisting}[style=diffstyle]
-  void render() override { Base::Render(); }
+  void render() override { Base::render(); }
\end{lstlisting}
提供改名 fix-it；专治「override 里调基类却调错函数」。本族还值得开的：bugprone-argument-comment（用 \texttt{/*name=*/} 给布尔实参加注释，选项 \texttt{StrictMode}）、bugprone-swapped-operations（\texttt{a < b < c} 式误写）、bugprone-unused-return-value（较新版本提供，先 \texttt{--list-checks} 确认存在）。

\section{performance 族实战}\label{sec:ch13-s2}
\subsection{performance-unnecessary-value-param}\label{sec:ch13-s2a}
\begin{lstlisting}[style=cppstyle]
void Save(std::string path) {    // 从不转移所有权，却按值接收
  std::filesystem::path p(path);
  Write(p);
}
\end{lstlisting}
\begin{lstlisting}[style=shellstyle]
io.cpp:1:15: warning: the parameter 'path' is copied for each invocation
but only used as a const reference; consider making it a const reference [performance-unnecessary-value-param]
\end{lstlisting}
\begin{lstlisting}[style=diffstyle]
-void Save(std::string path) {
+void Save(const std::string &path) {
\end{lstlisting}
fix-it 存在但默认不自动应用——改签名要同步所有 TU 与声明处，跨 TU 的批量改写请走增量脚本逐处人审（第 10、14 章）。
\subsection{performance-move-const-arg}\label{sec:ch13-s2b}
\begin{lstlisting}[style=cppstyle]
void Clone(const std::string &src) {
  std::string dst = std::move(src);   // const 上的 move 只是拷贝
}
\end{lstlisting}
\begin{lstlisting}[style=shellstyle]
f.cpp:2:21: warning: std::move of the const variable 'src' has no effect;
remove std::move call or make 'src' non-const [performance-move-const-arg]
\end{lstlisting}
\begin{lstlisting}[style=diffstyle]
-  std::string dst = std::move(src);
+  std::string dst = src;
\end{lstlisting}
可自动修复；对依赖重载决议的模板代码偶有误报（编译器会选 const 重载，语义其实不同），模板重的项目先降级为 note。
\subsection{performance-for-range-copy}\label{sec:ch13-s2c}
\begin{lstlisting}[style=cppstyle]
for (const auto w : words) {     // 每个元素拷贝一次
  total += w.len();
}
\end{lstlisting}
\begin{lstlisting}[style=shellstyle]
sum.cpp:1:6: warning: the loop variable 'w' of type 'std::string' is
copied for each iteration; consider making it a reference [performance-for-range-copy]
\end{lstlisting}
\begin{lstlisting}[style=diffstyle]
-  for (const auto w : words) {
+  for (const auto &w : words) {
\end{lstlisting}
可自动修复；range-for 本身是 \cpp{11} 特性。选项 \texttt{WarnOnAllAutoCopies} 控制是否对 \texttt{auto} 也检查（见第 12 章 12.4.4）。
\subsection{performance-enum-size}\label{sec:ch13-s2d}
\begin{lstlisting}[style=cppstyle]
enum Color { Red, Green, Blue };   // 默认底层类型 unsigned（4 字节）
\end{lstlisting}
\begin{lstlisting}[style=shellstyle]
gfx.cpp:1:6: warning: the underlying type of 'Color' can be 'uint8_t'
which would reduce the size of the enum [performance-enum-size]
\end{lstlisting}
\begin{lstlisting}[style=diffstyle]
-enum Color { Red, Green, Blue };
+enum Color : uint8_t { Red, Green, Blue };
\end{lstlisting}
提供 fix-it；但改底层类型\textbf{改变 ABI 与内存布局}，跨模块共享的枚举、与外部交互的结构体禁用。该族还包括 performance-trivially-copyable（可平凡拷贝的小类型建议按值传参，注意与 13.2.1 是互补关系而非矛盾——前者管大对象，后者管小对象）。
\begin{warning}[title={autofix 与 ABI 的边界}]
enum-size、const-return-type、convert-member-functions-to-static 这类「fix-it 存在但改变接口」的检查，在导出库上必须整族降级为 note 或 NOLINT 区间管理；第 14 章的 apply-fixes 只应作用于私有实现文件。
\end{warning}

\section{modernize 族实战}\label{sec:ch13-s3}
\subsection{modernize-avoid-c-arrays}\label{sec:ch13-s3a}
\begin{lstlisting}[style=cppstyle]
void Init(char buf[64]);               // 形参早已退化为指针
int primes[] = {2, 3, 5, 7};
\end{lstlisting}
\begin{lstlisting}[style=shellstyle]
net.cpp:1:13: warning: do not declare C-style arrays, use 'std::array' instead [modernize-avoid-c-arrays]
net.cpp:2:5: warning: do not declare C-style arrays, use 'std::array' instead [modernize-avoid-c-arrays]
\end{lstlisting}
\begin{lstlisting}[style=cppstyle]
void Init(std::array<char, 64> &buf);
constexpr std::array<int, 4> primes = {2, 3, 5, 7};
\end{lstlisting}
局部数组可自动改写为 \Tp{array}，形参改写不可自动（牵连调用方）；与 C 库边界处直接 NOLINT。
\subsection{modernize-use-override}\label{sec:ch13-s3b}
\begin{lstlisting}[style=cppstyle]
struct Shape { virtual double Area() const; };
struct Circle : Shape { double Area() const; };
\end{lstlisting}
\begin{lstlisting}[style=shellstyle]
shape.h:2:17: warning: 'Circle::Area' can be marked [override|final] (fix available) [modernize-use-override]
\end{lstlisting}
\begin{lstlisting}[style=diffstyle]
-  struct Circle : Shape { double Area() const; };
+  struct Circle : Shape { double Area() const override; };
\end{lstlisting}
全自动修复，\cpp{11} 起；是 CI 门禁里性价比最高的一条——它同时防「签名漂移导致基类虚函数没被覆盖」。
\subsection{modernize-use-equals-default}\label{sec:ch13-s3c}
\begin{lstlisting}[style=cppstyle]
Widget() {}
~Widget() {}
\end{lstlisting}
\begin{lstlisting}[style=shellstyle]
widget.h:1:1: warning: use '= default' to define a trivial constructor [modernize-use-equals-default]
\end{lstlisting}
\begin{lstlisting}[style=diffstyle]
-Widget() {}
+Widget() = default;
\end{lstlisting}
可自动修复，\cpp{11}；姊妹检查 modernize-use-equals-delete 处理「想禁止但用旧声明技巧」的函数。
\subsection{modernize-loop-convert}\label{sec:ch13-s3d}
\begin{lstlisting}[style=cppstyle]
for (std::vector<int>::iterator it = v.begin(); it != v.end(); ++it)
  total += *it;
\end{lstlisting}
\begin{lstlisting}[style=shellstyle]
sum.cpp:1:1: warning: use range-based for loop instead (fix available) [modernize-loop-convert]
\end{lstlisting}
\begin{lstlisting}[style=diffstyle]
-  for (auto it = v.begin(); it != v.end(); ++it)
-    total += *it;
+  for (int elem : v)
+    total += elem;
\end{lstlisting}
可自动修复；迭代器有复杂算术、倒序、循环内 erase 时不改写——这些位置会留下旧循环，正好提示人工复核。
\subsection{modernize-use-emplace}\label{sec:ch13-s3e}
\begin{lstlisting}[style=cppstyle]
std::vector<std::string> names;
names.push_back(std::string("abc"));   // 先构造临时对象再移动
\end{lstlisting}
\begin{lstlisting}[style=shellstyle]
names.cpp:2:3: warning: use emplace_back instead of push_back (fix available) [modernize-use-emplace]
\end{lstlisting}
\begin{lstlisting}[style=diffstyle]
-  names.push_back(std::string("abc"));
+  names.emplace_back("abc");
\end{lstlisting}
可自动修复，\cpp{11}；注意 emplace 绕过 \texttt{explicit} 单参构造的检查、也可能在 \cpp{17} CTAD 下改选出不同重载——自动修复后必须保留编译这一步。该族还常用：modernize-make-unique（\texttt{reset(new T)} 转 \Fn{make\_unique}，可自动，需 \cpp{14}）、modernize-pass-by-value（构造函数值传参加移动成员，需 \cpp{11}）、modernize-use-using、modernize-deprecated-headers、modernize-replace-random-with-uniform（无 autofix，样本分布不同不能盲改）。

\section{readability 族实战}\label{sec:ch13-s4}
readability 族主观性强，务必白名单挑，不要整族开。
\subsection{readability-magic-numbers}\label{sec:ch13-s4a}
\begin{lstlisting}[style=cppstyle]
if (retries > 3) return Fail();
double timeout = 30.5;
\end{lstlisting}
\begin{lstlisting}[style=shellstyle]
cfg.cpp:1:15: warning: magic number '3' (integer) used, consider a named constant [readability-magic-numbers]
\end{lstlisting}
\begin{lstlisting}[style=diffstyle]
-  if (retries > 3) return Fail();
+  constexpr int kMaxRetries = 3;
+  if (retries > kMaxRetries) return Fail();
\end{lstlisting}
无 autofix（命名需要人）；用 \texttt{IgnoredIntegerValues} 白名单消化 0/1/2/8 这类常识值（第 12 章 12.4.4）。cppcoreguidelines-avoid-magic-numbers 是同一实现的别名，二者选一个写进配置即可。
\subsection{readability-identifier-length}\label{sec:ch13-s4b}
\begin{lstlisting}[style=cppstyle]
int d = ComputeDayOfYear(ts);
\end{lstlisting}
\begin{lstlisting}[style=shellstyle]
t.cpp:1:5: warning: variable name 'd' is too short, minimal length is 3 [readability-identifier-length]
\end{lstlisting}
\begin{lstlisting}[style=diffstyle]
-  int d = ComputeDayOfYear(ts);
+  int day = ComputeDayOfYear(ts);
\end{lstlisting}
无 autofix；clang-tidy 17 引入。数学循环里 \texttt{i}、\texttt{j} 是合法惯用法，先确认版本提供了循环变量豁免选项再启用，否则整片红。
\subsection{readability-braces-around-statements}\label{sec:ch13-s4c}
\begin{lstlisting}[style=cppstyle]
if (ok)
  Send();
\end{lstlisting}
\begin{lstlisting}[style=shellstyle]
tx.cpp:1:3: warning: statement should be wrapped in braces [readability-braces-around-statements]
\end{lstlisting}
\begin{lstlisting}[style=diffstyle]
   if (ok)
-    Send();
+  { Send(); }
\end{lstlisting}
可自动修复；价值不在美观，而在杜绝「后续插一行语句忘加括号」的悬挂 if。
\subsection{readability-function-size}\label{sec:ch13-s4d}
\begin{lstlisting}[style=cppstyle]
void ParsePacket(Buffer &b) {
  // ... 手写状态机，97 行 ...
}
\end{lstlisting}
\begin{lstlisting}[style=shellstyle]
parse.cpp:10:6: warning: function 'ParsePacket' exceeds configured line threshold: lines 97 > 80 [readability-function-size]
\end{lstlisting}
无 autofix；阈值键 \texttt{LineThreshold}、\texttt{StatementThreshold} 等（第 12 章 12.4.4）。作为「收税型」指标只配 warning，别进 \texttt{WarningsAsErrors}。
\subsection{readability-convert-member-functions-to-static}\label{sec:ch13-s4e}
\begin{lstlisting}[style=cppstyle]
class Point {
 public:
  double Scale() const { return 2.0; }   // 从不触碰 this
};
\end{lstlisting}
\begin{lstlisting}[style=shellstyle]
pt.cpp:2:10: warning: member function 'Scale' can be made static [readability-convert-member-functions-to-static]
\end{lstlisting}
\begin{lstlisting}[style=diffstyle]
-  double Scale() const { return 2.0; }
+  static double Scale() { return 2.0; }
\end{lstlisting}
可自动修复；公开类的头文件慎改——去掉 \texttt{this} 依赖是源码级接口变更。本族其余常用项：readability-const-return-type、readability-implicit-bool-conversion、readability-named-parameter、readability-else-after-return、readability-redundant-access-specifiers、readability-avoid-const-params-in-decls、readability-uppercase-literal-suffix（\texttt{3u} 转 \texttt{3U}，可自动修复）、readability-inconsistent-declaration-parameter-name——多为纯 AST、低开销，按口味逐个试开。

\section{cppcoreguidelines 族实战}\label{sec:ch13-s5}
\subsection{cppcoreguidelines-pro-bounds}\label{sec:ch13-s5a}
\begin{lstlisting}[style=cppstyle]
int buf[8];
std::size_t i = ReadIndex();
buf[i] = 1;                      // i 越界编译器不管
\end{lstlisting}
\begin{lstlisting}[style=shellstyle]
b.cpp:3:3: warning: do not use built-in subscript operator on raw array, consider using 'std::span' instead [cppcoreguidelines-pro-bounds]
\end{lstlisting}
\begin{lstlisting}[style=cppstyle]
int buf[8];
std::span<int, 8> view{buf};
if (i < view.size()) view[i] = 1;
\end{lstlisting}
无 autofix；clang-tidy 17 引入，配合 \cpp{20} 的 \Tp{span}。更细粒度的三个子检查——pro-bounds-array-to-pointer-decay、pro-bounds-pointer-arithmetic、pro-bounds-constant-array-index——可单独开关，嵌入式大数组项目常只留 pointer-arithmetic。
\subsection{cppcoreguidelines-pro-type-vararg}\label{sec:ch13-s5b}
\begin{lstlisting}[style=cppstyle]
printf("%s", user_input);        // 连合法用法也一起拦
\end{lstlisting}
\begin{lstlisting}[style=shellstyle]
f.cpp:1:1: warning: do not use variadic arguments [cppcoreguidelines-pro-type-vararg]
\end{lstlisting}
\begin{lstlisting}[style=cppstyle]
std::cout << user_input;
// C++23: std::print("{}", user_input);
\end{lstlisting}
无 autofix；它反对的是 vararg 本身而非某次误用，日志宏满天飞的项目用选项 \texttt{IgnoreMacros} 或干脆只在新增文件启用。
\subsection{cppcoreguidelines-init-variables}\label{sec:ch13-s5c}
\begin{lstlisting}[style=cppstyle]
void f(bool flag) {
  int count;                     // 未初始化就流向使用点
  if (flag) count = 1;
  Use(count);
}
\end{lstlisting}
\begin{lstlisting}[style=shellstyle]
f.cpp:2:7: warning: variable 'count' of type 'int' is not initialized (fix available) [cppcoreguidelines-init-variables]
\end{lstlisting}
\begin{lstlisting}[style=diffstyle]
-  int count;
+  int count{};
\end{lstlisting}
可自动修复（值初始化）；但它治标——真正的「未初始化读取」请交给 clang 的 \texttt{-Wsometimes-uninitialized} 与 MSan 定案（见 13.8）。
\subsection{cppcoreguidelines-special-member-functions}\label{sec:ch13-s5d}
\begin{lstlisting}[style=cppstyle]
class Conn {
 public:
  Conn(int fd) : fd_(fd) {}
  ~Conn() { close(fd_); }        // 写了析构，没写五法则其余四个
 private:
  int fd_;
};
\end{lstlisting}
\begin{lstlisting}[style=shellstyle]
conn.h:3:3: warning: class 'Conn' defines a user-defined destructor but
lacks copy/move constructors and assignment operators [cppcoreguidelines-special-member-functions]
\end{lstlisting}
\begin{lstlisting}[style=cppstyle]
class Conn {
 public:
  Conn(int fd) : fd_(fd) {}
  ~Conn() { close(fd_); }
  Conn(const Conn &) = delete;
  Conn &operator=(const Conn &) = delete;
 private:
  int fd_;
};
\end{lstlisting}
无 autofix（删除还是转移语义要人定）；选项 \texttt{AllowSoleDefaultDtor} 允许「只有 defaulted 析构」的类过关（第 12 章 12.4.4）。该族还值得分批开：cppcoreguidelines-owning-memory（裸 \texttt{new} 建议 \Tp{gsl::owner}，提供 fix-it，需 GSL 与 \cpp{11} 起）、cppcoreguidelines-avoid-goto、cppcoreguidelines-pro-type-static-cast-downcast（建议 \texttt{dynamic\_cast}，RTTI 开销需确认）、cppcoreguidelines-rvalue-reference-param-not-moved（形参声明为 \texttt{T\&\&} 却未在函数体内 move，无 autofix）、cppcoreguidelines-avoid-c-style-cast（可自动改写为具名 cast，\cpp{11}）。

\section{clang-analyzer 族实战}\label{sec:ch13-s6}
analyzer 族复用 clang static analyzer 引擎（第 12 章 12.2），路径敏感、无 autofix、成本最高。诊断文本通常带 note 行说明推断路径。
\subsection{clang-analyzer-core.NullDereference}\label{sec:ch13-s6a}
\begin{lstlisting}[style=cppstyle]
Node *n = Find(key);
int v = n->value;                // Find 允许返回 nullptr
\end{lstlisting}
\begin{lstlisting}[style=shellstyle]
graph.cpp:12:11: warning: attempt to dereference a null pointer [clang-analyzer-core.NullDereference]
graph.cpp:11:3: note: 'n' initialized to a null pointer value
graph.cpp:12:11: note: Dereference of null pointer
\end{lstlisting}
\begin{lstlisting}[style=diffstyle]
   Node *n = Find(key);
+  if (!n) return 0;
   int v = n->value;
\end{lstlisting}
修复就是加判空，但 analyzer 不会替你想清楚「返回 0 还是抛异常」；路径敏感意味着它只在「存在一条使 n 为空的可行路径」时报警——没有报不等于安全（13.8）。
\subsection{clang-analyzer-deadcode.DeadStores}\label{sec:ch13-s6b}
\begin{lstlisting}[style=cppstyle]
int result = -1;
result = Compute(cfg);           // -1 从未被读
\end{lstlisting}
\begin{lstlisting}[style=shellstyle]
calc.cpp:1:7: warning: Value stored to 'result' during its initialization is never read [clang-analyzer-deadcode.DeadStores]
\end{lstlisting}
\begin{lstlisting}[style=diffstyle]
-  int result = -1;
+  int result;
   result = Compute(cfg);
\end{lstlisting}
无 autofix；真正价值在于暴露「错误路径被覆盖」的残局——初值 \texttt{-1} 往往是忘了用的失败返回。同族的 clang-analyzer-cplusplus.NewDelete（new/delete 配对与泄漏）值得默认开，clang-analyzer-optin.performance.Padding（成员排布空洞）用于压缩嵌入式结构体。
\subsection{clang-analyzer-security.insecureAPI.strcpy}\label{sec:ch13-s6c}
\begin{lstlisting}[style=cppstyle]
char path[64];
strcpy(path, getenv("HOME"));
\end{lstlisting}
\begin{lstlisting}[style=shellstyle]
env.c:3:3: warning: Call to function 'strcpy' is insecure as it does not provide
bounding of the memory buffer; replace with a length-bounded function. CWE-119 [clang-analyzer-security.insecureAPI.strcpy]
\end{lstlisting}
\begin{lstlisting}[style=diffstyle]
-  strcpy(path, getenv("HOME"));
+  std::snprintf(path, sizeof(path), "%s", getenv("HOME") ? getenv("HOME") : ".");
\end{lstlisting}
无 autofix；\texttt{security.insecureAPI.*} 子族覆盖 gets、strcat、sprintf 等。注意 \texttt{optin.} 前缀的 checker 默认不随 \texttt{clang-analyzer-*} 全开——写选择器时点名 \texttt{clang-analyzer-optin.performance.Padding}。
\begin{note}[title={analyzer 的成本税}]
路径敏感引擎在解析循环、状态机、超长函数上可能触发路径爆炸并超时放弃。单 TU 秒级起步，CI 上务必只对增量行跑（第 10 章的门禁思路），或把 analyzer 族拆到夜间全量任务。
\end{note}

\section{concurrency 与 misc 族}\label{sec:ch13-s7}
\subsection{concurrency-mt-unsafe}\label{sec:ch13-s7a}
\begin{lstlisting}[style=cppstyle]
char *tok = strtok(buf, ",");
struct tm *t = localtime(&ts);
\end{lstlisting}
\begin{lstlisting}[style=shellstyle]
p.c:1:13: warning: 'strtok' is not thread safe [concurrency-mt-unsafe]
p.c:2:14: warning: 'localtime' is not thread safe [concurrency-mt-unsafe]
\end{lstlisting}
\begin{lstlisting}[style=diffstyle]
-  char *tok = strtok(buf, ",");
+  char *save = nullptr;
+  char *tok = strtok_r(buf, ",", &save);
\end{lstlisting}
无 autofix（POSIX 函数名各异）；它按「线程不安全 API 名单」命中，并不验证你真的多线程——单线程程序要么 NOLINT 要么白名单该文件。
\subsection{misc-no-recursion}\label{sec:ch13-s7b}
\begin{lstlisting}[style=cppstyle]
bool Validate(Node *n) { return Validate(n->next); }
\end{lstlisting}
\begin{lstlisting}[style=shellstyle]
v.cpp:1:6: warning: the 'Validate' function is recursive; consider an iterative rewrite [misc-no-recursion]
\end{lstlisting}
\begin{lstlisting}[style=diffstyle]
-  bool Validate(Node *n) { return Validate(n->next); }
+  bool Validate(Node *n) {
+    for (; n; n = n->next) if (!CheckOne(n)) return false;
+    return true;
+  }
\end{lstlisting}
无 autofix；嵌入式与内核态（栈深度受限）强烈推荐，见第 12 章 12.6.4 的嵌入式配置。
\subsection{misc-const-correctness 为何常被关}\label{sec:ch13-s7c}
\begin{lstlisting}[style=cppstyle]
void f() {
  int x = Compute();     // 全部局部变量都能加 const
  Print(x);
}
\end{lstlisting}
\begin{lstlisting}[style=shellstyle]
f.cpp:2:3: warning: variable 'x' of type 'int' can be referenced as const [misc-const-correctness]
\end{lstlisting}
\begin{lstlisting}[style=diffstyle]
-  int x = Compute();
+  const int x = Compute();
\end{lstlisting}
技术上可自动修复，但一个中型模块一次能报出上千条——信噪比灾难，故多数团队整条关闭，只在 code review 阶段人工挑刺。misc 族的 misc-include-cleaner 则建议单独跑：它按符号使用判断头文件冗余/缺失，是第 15 章 include-cleaner 工具的 \Tidy\ 前端，输出当清单消费而不是当门禁。

\section{与 sanitizer 的分界}\label{sec:ch13-s8}
静态检查回答「模式上是否可疑」，sanitizer 回答「这一次运行是否真的错了」。analyzer 虽路径敏感，但路径空间有限、跨过程建模有限，报出的永远是「可能」；反过来，越界写、释放后使用、数据竞争这些必须真实运行才能定案。
\begin{compare}[title={同一种坏味的两种执法}]
悬垂引用：\Tidy\ 的 bugprone-dangling-handle 能在编译期抓到明显形态，但跨函数逃逸要 ASan 的 quarantine 才能坐实；use-after-move 是\textbf{合法} \CPP\ 状态，任何 sanitizer 都不报，只能靠静态与约定。结论：静态管「未运行即可证伪」的，动态管「只有运行时可见」的，二者是互补的两道门。
\end{compare}
{\small
\begin{longtable}{>{\raggedright\arraybackslash}p{0.20\textwidth} >{\raggedright\arraybackslash}p{0.27\textwidth} >{\raggedright\arraybackslash}p{0.33\textwidth}}
\textbf{问题} & \textbf{\Tidy\ 能做到} & \textbf{必须交给}\\ \midrule
\endfirsthead
\textbf{问题} & \textbf{\Tidy\ 能做到} & \textbf{必须交给}\\ \midrule
\endhead
堆越界读写 & pro-bounds 拦模式，不拦运行时值 & ASan\\
释放后使用 & analyzer 可追短路径 NewDelete & ASan（定案）\\
数据竞争 & concurrency-mt-unsafe 只点名 API & TSan\\
未初始化读 & init-variables 补默认值 & MSan / 编译器警告\\
整数溢出语义 & implicit-widening-of-multiplication-result & UBSan\\
use-after-move & bugprone-use-after-move 独占（运行时无味） & 静态\\
\end{longtable}}

\section{启用速查表}\label{sec:ch13-s9}
汇总本章出现过的 check（版本以 clang-tidy 17--19 为基准，升级后先跑 \texttt{--list-checks} 复核）：
{\small
\begin{longtable}{>{\raggedright\arraybackslash}p{0.28\textwidth} >{\raggedright\arraybackslash}p{0.09\textwidth} >{\raggedright\arraybackslash}p{0.10\textwidth} >{\raggedright\arraybackslash}p{0.09\textwidth} >{\raggedright\arraybackslash}p{0.24\textwidth}}
\textbf{check} & \textbf{autofix} & \textbf{标准} & \textbf{完整 TU} & \textbf{典型误报场景}\\ \midrule
\endfirsthead
\textbf{check} & \textbf{autofix} & \textbf{标准} & \textbf{完整 TU} & \textbf{典型误报场景}\\ \midrule
\endhead
bugprone-use-after-move & 否 & \cpp{11} & 是 & move 后立刻重赋值再用的路径\\
bugprone-unnecessary-copy-initialization & 是 & 任意 & 否 & 拷贝是为了改副本\\
bugprone-dangling-handle & 否 & \cpp{17} & 否 & view 指向更长命的成员\\
bugprone-implicit-widening-of-multiplication-result & 是 & 任意 & 否 & 故意先窄算再截断\\
bugprone-assert-side-effect & 否 & 任意 & 否 & 断言宏被定义为始终求值\\
bugprone-parent-virtual-call & 是 & \cpp{11} & 否 & 有意调用基类重载\\
performance-unnecessary-value-param & 受限 & 任意 & 是 & 参数确实要按值转移\\
performance-move-const-arg & 是 & \cpp{11} & 否 & 模板里依赖重载决议的 move\\
performance-for-range-copy & 是 & \cpp{11} & 否 & 循环体要副本（存指针成员）\\
performance-enum-size & 是 & 任意 & 否 & 跨 ABI 共享枚举\\
modernize-avoid-c-arrays & 部分 & \cpp{11} & 否 & 与 C API 的边界缓冲区\\
modernize-use-override & 是 & \cpp{11} & 否 & 几乎无误报\\
modernize-loop-convert & 是 & \cpp{11} & 否 & 倒序/算术迭代器\\
modernize-use-emplace & 是 & \cpp{11} & 否 & explicit 构造与 CTAD\\
readability-magic-numbers & 否 & 任意 & 否 & 单位换算、协议常量\\
readability-identifier-length & 否 & 任意 & 否 & 数学惯用短名\\
readability-function-size & 否 & 任意 & 否 & 生成代码\\
cppcoreguidelines-pro-bounds & 否 & \cpp{20} & 是 & 手写边界检查的热点\\
cppcoreguidelines-pro-type-vararg & 否 & 任意 & 否 & 日志门面宏\\
cppcoreguidelines-init-variables & 是 & \cpp{11} & 否 & 全分支覆盖的延迟赋值\\
cppcoreguidelines-special-member-functions & 否 & \cpp{11} & 否 & 析构仅为 RAII 句柄\\
clang-analyzer-core.NullDereference & 否 & 任意 & 是 & 隐含前置条件的私有函数\\
clang-analyzer-deadcode.DeadStores & 否 & 任意 & 是 & 调试用的占位赋值\\
clang-analyzer-security.insecureAPI.* & 否 & 任意 & 是 & 尺寸受控的内部缓冲\\
concurrency-mt-unsafe & 否 & C 库 & 否 & 单线程程序\\
misc-no-recursion & 否 & 任意 & 否 & 深度天然有界的结构\\
misc-const-correctness & 是 & 任意 & 否 & 海量局部变量（噪声）\\
\end{longtable}}
\begin{bestpractice}[title={一族一族地加，一条一条地关}]
顺序永远是：白名单三条 check 全库清零 → 进 \texttt{WarningsAsErrors} → 加一族、删噪声项 → 重复。每删一项，把理由写进 \texttt{.clang-tidy} 的 YAML 注释；半年后接手的同事需要知道的是「为什么关」，而不是「关过」。
\end{bestpractice}

% ---- frag_ch14.tex ----
\chapter{autofix、apply-fixes 与重构工具}\label{ch:autofix}

第 12 章把 \Tidy\ 讲成一台「读编译数据库、吐诊断」的分析机，本章处理它的另一半身体：
每一条诊断除了人类可读的文字之外，还可以附带一份机器可执行的修复指令。
把这份指令就地写回源文件，就是 autofix；把它导出成文件、交由另一个工具批量应用，
就是 \Bk clang-apply-replacements\Bk 的工作。围绕这两条主线，LLVM 还散布着一批
「以替换文本为唯一输出」的重构工具：\texttt{clang-rename}、\texttt{clang-move}、
\texttt{clang-include-fixer}，以及一个用来探查 AST 的交互式外壳 \texttt{clang-query}。

它们共享同一个底层抽象：\Tp{tooling::Replacement}。理解这个抽象的表达能力与限制，
比记住十几个命令行开关更有价值——因为限制决定了「哪些族可以放心批量自动修，
哪些族必须人来过一遍」。本章末尾给出一条从 modernize 到 \Fmt、编译、\Fn{ctest}
的端到端流水线。

\section{修复的表示与 \texttt{-fix} 的边界}\label{sec:ch14-s1}

一条 fix-it hint 本质上是「一组互不重叠的区间替换」：每个区间由四个字段确定——
目标文件（\texttt{FilePath}）、起始字节偏移（\texttt{Offset}）、被覆盖的长度
（\texttt{Length}，插入操作为 0）、以及替换进去的新文本（\texttt{ReplacementText}）。
一条诊断可以携带多个区间：\texttt{readability-braces-around-statements}
要插入左花括号并补上右花括号，就是两个区间；跨文件的修复（例如把虚函数签名
从头文件改到实现文件）则在同一个诊断下挂两个不同 \texttt{FilePath} 的区间。
偏移的计数方式与它的坑，见第 \ref{sec:ch14-s4} 节。

\Tidy\ 侧的开关只有四个是日常会用到的：

\begin{itemize}
  \item \texttt{-fix}：把当前 check 提供的修复直接写回源文件。默认只在 TU \emph{没有}
        编译错误时应用，避免在残缺 AST 上做改写。
  \item \texttt{-fix-errors}：即使 clang 报了解释错误也照样应用修复。它的正确用法是
        「错误来自第三方头文件的噪声、而目标代码本身是完整的」，不要用它来绕过
        自己代码的语法错误。
  \item \Opt{fixes=\allowbreak modernize-use-nullptr,bugprone-*}：白名单，只有列出的
        check 才应用修复，其余照常只报不改。这是把「报告」与「改写」解耦的关键开关。
  \item \Opt{fix-notes}：决定挂在 note 级诊断上的修复是否参与应用。很多「建议改成另一种写法」
        的提示以 note 形式附在主诊断之后，默认不应用，需要时显式打开。
\end{itemize}

\subsection*{为什么修复保守到几乎不跨 TU}

\Tidy\ 的每个 worker 只看到一个翻译单元（TU）的 AST。这带来三条硬限制，它们不是实现偷懒，
而是模型本身决定的：

\begin{enumerate}
  \item \textbf{不会新增或调整 include}。若某个修复把 \texttt{NULL} 换成 \texttt{nullptr}、
        或把裸指针换成 \Tp{std::unique\_ptr}，被引用的名字是否已声明、头文件是否已包含，
        工具并不负责补齐。这就是 \texttt{modernize-deprecated-headers} 一类 check
        经常改出「看起来对、编译不过」的原因。
  \item \textbf{不会重排别的函数、不会移动代码}。修复只落在诊断自己圈定的区间内，
        不会为了风格统一去动相邻函数，也不会把成员顺序按访问说明符重排。
  \item \textbf{不保证幂等}。一次运行中，靠前的替换会使靠后的 \texttt{Offset} 失效，
        冲突的修复被丢弃而不是回退重写；于是「跑两次才稳定」是常见现象。
        判断标准很简单：第二次运行若还能报出同类诊断，说明这一族不该被自动化。
\end{enumerate}

\begin{guideline}[title={可以放心交给 -fix 的检查族}]
判据是「机械且可验证」：改动局部、不需要新依赖、编译器立刻能证实或证伪。
modernize 与 readability 里只改局部写法的族属于此类，例如
\texttt{modernize-use-nullptr} 与 \texttt{readability-braces-around-statements}；
会改签名、牵动头文件与调用点的族（例如 \texttt{modernize-pass-by-value}）
应当只报告、由人决定。至于 \texttt{clang-analyzer-*} 与 \texttt{bugprone-*}，
它们的大多数诊断根本没有修复文本，\texttt{-fix} 对它们是无操作。
\end{guideline}

\begin{warning}[title={-fix 直接改文件时的批量误改}]
\texttt{-fix} 与 \texttt{-fix-errors} 修改的是磁盘上的源文件，而不是编辑器缓冲区。
三条最容易踩的坑：一是把 \texttt{-checks=*} 与 \texttt{-fix} 同时打开，
一次运行改走上千处且混杂多个族，diff 无法审阅；二是在未纳入版本控制、
或本地已有未提交改动的目录上运行，工具会与你手工编辑的内容互相覆盖；
三是把生成代码与 \texttt{externals/} 之类的第三方目录一并卷入。
可行的 review 策略是：改用 \Opt{export-fixes} 只导出不写回（第 \ref{sec:ch14-s2} 节），
或按族分批、每次一个 check 一个提交，让每个提交都能独立回滚；
仓库保持干净，必要时用 \texttt{git stash} 兜底（本文只在流程上提到它，
片段里不替你执行任何版本控制命令）。
\end{warning}

\section{\texttt{--export-fixes} 与批量应用替换}\label{sec:ch14-s2}

比「直接改文件」更安全的路径是三步：导出、审阅、应用。\Opt{export-fixes=\allowbreak out.yaml}
让 \Tidy\ 把本次运行收集到的所有替换序列化进一个 YAML 文件，源文件保持不变。
它与 \texttt{-fix} 是互斥的两种输出通道，语义上更接近「写报告」，
因此可以在 CI 里对任意分支安全运行。

\begin{lstlisting}[style=yamlstyle,caption={\Tidy\ 导出的 fix 文件骨架：诊断驱动的格式},label={lst:ch14-fixyaml}]
---
MainSourceFile:  'src/parser.cpp'
Diagnostics:
  - DiagnosticName:  modernize-use-nullptr
    DiagnosticMessage:
      Message:  'Use nullptr'
      FilePath:  'src/parser.cpp'
      FileOffset:  812
      Replacements:
        - FilePath:  'src/parser.cpp'
          Offset:  812
          Length:  1
          ReplacementText:  'nullptr'
  - DiagnosticName:  readability-braces-around
    DiagnosticMessage:
      Message:  'Body should be inside braces'
      FilePath:  'src/parser.cpp'
      FileOffset:  1305
      Replacements:
        - FilePath:  'src/parser.cpp'
          Offset:  1305
          Length:  0
          ReplacementText:  ' {'
...
\end{lstlisting}

清单 \ref{lst:ch14-fixyaml} 里的字段名要逐字理解：\texttt{MainSourceFile} 是这次分析的入口文件；
\texttt{DiagnosticName} 给出 check 名，审阅时按它排序与切分；\texttt{FileOffset}
是诊断本身的位置，而真正执行改写的是 \texttt{Replacements} 列表里的
\texttt{Offset}/\texttt{Length}/\texttt{ReplacementText} 三元组。
注意 \texttt{Length: 0} 表示纯插入；文件末尾独立的一行 \texttt{...} 是 YAML 文档结束标记，
多份修复会拼成多个文档。

另一类工具（\texttt{clang-rename}、\texttt{clang-move}）导出的文件不含诊断信息，
只有扁平的替换清单，顶层键是 \texttt{FileEntries}：

\begin{lstlisting}[style=yamlstyle,caption={重构工具导出的替换清单：只有 FileEntries},label={lst:ch14-fileentries}]
---
MainSourceFile:  'src/widget.h'
FileEntries:
  - FilePath:  'src/widget.h'
    Offset:  1043
    Length:  6
    ReplacementText:  'Widget'
  - FilePath:  'src/widget.cpp'
    Offset:  221
    Length:  6
    ReplacementText:  'Widget'
...
\end{lstlisting}

应用端是 \Bk clang-apply-replacements\Bk，它接受一个或多个目录（递归扫描其中的
\texttt{.yaml} 与 \texttt{.fix} 文件），把清单里的替换全部落到磁盘：

\begin{itemize}
  \item 位置参数：包含 fix 文件的目录，通常是 \Opt{export-fixes} 指定的那个路径。
  \item \Opt{format}（默认为真）：应用完替换后立刻调用 \Fmt\ 重排；
        更稳的做法是同时给出 \Opt{style=file}，让这一步与第 9 章的仓库
        \texttt{.clang-format} 使用同一份风格来源。
  \item \Opt{ignore-insert-conflict}：插入型替换（\texttt{Length: 0}）在同一位置撞车时，
        默认整个文件放弃应用；打开后跳过冲突项、应用其余替换。
  \item \Opt{remove}：应用成功后删除对应的 fix 文件，便于把目录当作队列反复消费。
\end{itemize}

\begin{lstlisting}[style=shellstyle,caption={导出、审阅、应用的最小闭环},label={lst:ch14-apply}]
# 1) export only, sources untouched
clang-tidy src/parser.cpp -p build -- \
  -checks='modernize-use-nullptr' \
  --export-fixes=fixes/parser.yaml

# 2) review: how many fixes, which files
grep -c DiagnosticName fixes/parser.yaml
grep FilePath fixes/parser.yaml | sort -u

# 3) apply, reformat from .clang-format, then verify
clang-apply-replacements fixes --style=file
cmake --build build && ctest --test-dir build
\end{lstlisting}

改动的实际形态就是一串纯文本区间替换的结果，审阅 diff 时值得注意「格式还没跑」这一点：

\begin{lstlisting}[style=diffstyle,caption={应用之后、尚未运行 \Fmt\ 时的典型形态},label={lst:ch14-diff}]
--- a/src/parser.cpp
+++ b/src/parser.cpp
@@ -41,7 +41,8 @@
-  if (p == NULL)
-    return Err;
+  if (p == nullptr) {
+    return Err;
+  }
\end{lstlisting}

\begin{note}[title={apply-replacements 的冲突处理}]
同一区间被多个修复覆盖时，\Bk clang-apply-replacements\Bk 不会「合并」也不会「择优」：
同文件同区间的重叠替换即判为冲突，该文件本次应用整体放弃（或配合
\Opt{ignore-insert-conflict} 只丢弃冲突项）。所以「先按族分批导出」不只是 review 策略，
也是避免整文件白跑的实用手段。另一个常见现象是替换在导出与应用之间已经失效——
源文件被手工编辑或被上一批 apply 改动过，此时 \texttt{Offset} 与 \texttt{Length}
指向的内容与 \texttt{ReplacementText} 的预期不符，工具会拒绝应用。
稳妥的顺序始终是：导出与 apply 之间不碰源文件，apply 完再导出下一批。
\end{note}

\section{并行与增量：两个官方脚本}\label{sec:ch14-s3}

单文件调用 \Tidy\ 的瓶颈在 AST 构建，仓库越大越明显。LLVM 自带两个 Python 脚本，
一个负责全量并行，一个负责增量只跑改动面：

\begin{itemize}
  \item \textbf{run-clang-tidy}：读编译数据库，为每个 TU 起一个 \Tidy\ 进程。常用参数是
    \texttt{-p build}（数据库目录）、\texttt{-j 8}（并发度，与内存占用直接挂钩）、
    \texttt{-checks=}、\texttt{-config-file .clang-tidy}、\texttt{-header-filter=}、
    \texttt{-quiet}，以及 \texttt{-fix} 与 \texttt{-export-fixes}——后者指向\emph{目录}，
    脚本为每个 TU 写出一份独立 YAML，正好回避并发写同一文件的竞争，
    也因此能直接喂给 \Bk clang-apply-replacements\Bk。
  \item \textbf{clang-tidy-diff.py}：从标准输入读 unified diff，抽取被改动的文件名，
    只对这些文件跑 \Tidy，并把行号过滤到 diff 的 hunk 范围内。常用参数是
    \texttt{-p1}（diff 剥离层级，与 \texttt{git diff} 的输出对应）、
    \texttt{-path build}（数据库目录）、\texttt{-clang-tidy-binary}
    （指定可执行文件名，例如带版本后缀的 \texttt{clang-tidy-18}）、
    \texttt{-config-file}、\texttt{-quiet}。
\end{itemize}

\begin{lstlisting}[style=shellstyle,caption={CI 中增量门禁与夜间全量的组合骨架},label={lst:ch14-parallel}]
set -euo pipefail
BUILD=build-analyze
FIX=fixes
rm -rf "$FIX"; mkdir -p "$FIX"

# PR gate: only issues introduced by this change
git diff --unified=0 "origin/main...HEAD" \
  | clang-tidy-diff.py -p1 -path "$BUILD" -quiet \
      -clang-tidy-binary clang-tidy \
      -config-file .clang-tidy

# nightly: whole repo in parallel, export only
run-clang-tidy -p "$BUILD" -j 8 -quiet \
  -config-file .clang-tidy \
  -checks='modernize-use-nullptr,modernize-use-override' \
  -export-fixes "$FIX"

# the fix job is triggered by a human, not by the pipeline
clang-apply-replacements "$FIX" --style=file
\end{lstlisting}

两个脚本的定位不同，混用才划算：增量脚本负责「不再变坏」，成本与 diff 大小成正比，
可以放进每个 PR；全量脚本负责「逐步变好」，成本与仓库大小成正比，适合定时任务或
只在修复作业里跑。第 10 章讲增量格式化门禁，第 16 章把这两条腿并进
CMakePresets 与 Actions 的矩阵里。

\section{偏移的坑：字节偏移与 UTF-16 计数}\label{sec:ch14-s4}

同一处代码位置，在这套工具链里至少有三种数字表示，混用必然错位：

\begin{longtable}{>{\raggedright\arraybackslash}p{0.20\textwidth}>{\raggedright\arraybackslash}p{0.26\textwidth}>{\raggedright\arraybackslash}p{0.32\textwidth}}
\caption{三种位置坐标系}\label{tab:ch14-coord}\\
\toprule
坐标系 & 计数单位 & 出现在哪里 \\
\midrule
\endfirsthead
\toprule
坐标系 & 计数单位 & 出现在哪里 \\
\midrule
\endhead
clang 诊断与 fix & 自文件起始的\textbf{字节}偏移 & \texttt{-fix}、\Opt{export-fixes}、\texttt{clang-rename -offset} \\
LSP 位置 & \textbf{UTF-16} 码元（行内 character） & 第 4 章协议的 \Tp{Position}、\CD\ 的 codeAction 区间 \\
编辑器缓冲区 & 码点或列，实现相关 & Neovim、VS Code 内部游标 \\
\bottomrule
\end{longtable}

于是有两个实用结论。第一，含中文注释或字符串的文件里，\textbf{字节偏移}比「字符位置」大，
手工照行列号推算 \texttt{Offset} 一定错；要算就用脚本按 UTF-8 编码长度算，
或改用 \texttt{clang-rename} 的 \texttt{-line}/\texttt{-column}
（行列从 1 起，列同样按字节计）。

\begin{lstlisting}[style=shellstyle,caption={同一位置的字节偏移与 UTF-16 码元数},label={lst:ch14-offset}]
# count bytes, not characters
python3 - <<'PY'
src = open("src/parser.cpp", encoding="utf-8").read()
name = "ParseToken"
i = src.index(name)
print("byte offset :", len(src[:i].encode("utf-8")))
print("utf-16 units:", len(src[:i].encode("utf-16-le")) // 2)
PY
\end{lstlisting}

第二个坑是\emph{行尾与编码}：把 LF 文件按 CRLF 保存之后，后续所有 \texttt{Offset}
都要加上前面新增的字节数，旧 fix 文件与新内容不再匹配——这正是
\Bk clang-apply-replacements\Bk 拒绝应用的最常见原因之一，
与第 9 章 \Fmt\ 中 \texttt{LineEnding} 的约定相互呼应。

\subsection*{clang-rename：把偏移暴露到底的重命名}

\texttt{clang-rename} 的定位是「脚本化改名」，参数组合如下：

\begin{itemize}
  \item 定位旧名字：\texttt{-offset}（\textbf{字节偏移}），或 \texttt{-line} 与
        \texttt{-column}；同名多处定义时用 \texttt{-select} 或其短形式 \texttt{-U}
        挑选第几个候选。
  \item 指定新名字：\texttt{-name=}，可配合 \texttt{-prefix}/\texttt{-suffix}
        与 \texttt{-regex} 做模式化批量改名。
  \item 输出：默认原地改写；\Opt{export-replacements=\allowbreak dir}
        改为写出清单 \ref{lst:ch14-fileentries} 那种 \texttt{FileEntries} 格式，
        再交给 \Bk clang-apply-replacements\Bk 落地。
  \item 依赖：\texttt{--compile-commands-dir=build} 提供编译数据库；
        跨项目改名需要 \texttt{-M=} 指向一份预先建好的索引文件
        （\CD\ 的后台索引可直接复用其磁盘格式）。
\end{itemize}

为什么日常更推荐用 \CD\ 的 rename 而不是 \texttt{clang-rename}？三点理由。
\CD\ 走 LSP 的 \texttt{textDocument/rename}，位置单位是 UTF-16，编辑器算出的偏移可直接使用，
不需要本节这一套换算；改名结果以 \Tp{WorkspaceEdit} 返回，客户端可以先展示预览再落盘，
天然可撤销；跨 TU 的引用来自 \CD\ 已建好的后台索引，而 \texttt{clang-rename}
没有 \texttt{-M} 索引时只能覆盖当前 TU 及其可见的声明，容易漏改。
\texttt{clang-rename} 的价值在无编辑器场景：脚本、pre-commit 钩子、CI 里的机械改名。

\section{搬移、include 修补与 AST 探查工具}\label{sec:ch14-s5}

\subsection*{clang-move：把一类搬出去}

\texttt{clang-move} 的输入是旧头文件、新头文件与类名，输出是替换清单而不是文件：

\begin{lstlisting}[style=shellstyle,caption={用 clang-move 拆出头文件再统一应用},label={lst:ch14-move}]
clang-move src/widget.h --out-replacements-dir=fixes \
  --old-header=src/widget.h --new-header=src/widget2.h \
  --class=Widget --style=file

# the new header must exist and must compile on its own
clang-apply-replacements fixes --style=file
\end{lstlisting}

关键在于它只搬「声明与定义文本」：新头文件不存在时需要先建好；
搬走之后旧头文件里残留的 include 与被切断的引用，要靠 \texttt{clang-include-extractor}
之类的辅助工具补齐——后者从旧头文件抽出被搬走代码所需的最小依赖集合，
写成与 \texttt{clang-move} 输出同目录、同格式的清单，一次 apply 收口。

\subsection*{clang-include-fixer 与 include-cleaner：往回补与往外删}

\texttt{clang-include-fixer} 解决「缺 include」：给它符号名与源文件，
它去查一张「哪个头文件提供哪个符号」的表，把该 include 插进去。
常用参数：\texttt{-symbol=Widget}、\texttt{-find-all-matches}
（补齐所有未定义符号而非只补第一个）、\texttt{-insert-tidy}
（优先插入 \Tidy\ 建议的头文件），以及 \texttt{-query-driver=} 与 \texttt{-std=}
用来取得系统头文件搜索路径与语言标准。

\Bk include-cleaner\Bk 是 IWYU 思路在 \CD\ 生态里的落地，方向恰好相反：
给定编译数据库，它报告「应该删掉哪些 include、应该补上哪些 include」，
逐文件输出类似 IWYU 的下划线（未用）与加号（缺失）标记。
它有两种使用形态：作为独立工具在 CI 里批量跑，或作为 \CD\ 的诊断来源
（\texttt{.clangd} 中 \texttt{Diagnostics} 一节的未用 include 开关，第 5 章已述）。
\Tidy\ 侧的对应物是 \texttt{misc-include-cleaner}：同样的结论做成普通诊断，
并可附带 fix-it，从而进入本章前面的导出—应用通道。它有若干以
\texttt{misc-include-cleaner.} 为前缀的 CheckOptions，用于忽略头文件、
控制是否强制更新等；具体键名请以本地 \texttt{clang-tidy --list-checks}
与所装 LLVM 版本文档为准，不同版本增减相当频繁。
IWYU 本体，以及「\cpp{20} modules 之后头文件工具的价值曲线」放在第 15 章讨论。

\subsection*{clang-scan-deps：模块时代的依赖图}

\texttt{clang-scan-deps} 不产出替换。它从编译数据库出发扫描预处理与 import 语句，
输出每个 TU 需要哪些模块（module）与哪些头文件（header）的依赖记录，
供构建系统调度模块编译顺序。输出形态由 \texttt{-format=} 选择：
既可给出人类可读的摘要，也可给出供 \texttt{make} 消费的规则式形式，
或 JSON 记录式形式（每条含 TU、模块名、种类、文件路径等字段）。
它与 \texttt{clang --analyze} 那个静态分析器毫无关系，只是名字相近，
第 15 章会专门把这条分界线画清楚。

\subsection*{clang-query：写自定义 check 之前的探查工具}

\texttt{clang-query} 是一个 matcher 的 REPL，用来先验证「这个 AST 查询到底命中什么」，
再决定要不要为此写一个自定义 check（第 13 章）：

\begin{lstlisting}[style=shellstyle,caption={clang-query 交互会话},label={lst:ch14-query}]
$ clang-query src/widget.cpp -- -std=c++20 -Iinclude
clang-query> match cxxRecordDecl(hasName("Widget"))
clang-query> set print-ast-to-std-output
clang-query> match cxxMethodDecl(isVirtual())
clang-query> q
\end{lstlisting}

会话里的命令记不住不如写成脚本：\texttt{:help} 列出全部命令，\texttt{let}
给 matcher 命名以便复用，\texttt{enable}/\texttt{disable} 切换选项，\texttt{q} 退出。
批处理用 \texttt{-f=}，把查询文件与源文件、编译参数一起交给 CI，
就得到一份「本周还剩多少处旧写法」的计数报表；这也是把「重构意图」固化成可复查资产的最省事方式。

\section{端到端实战与工具速查表}\label{sec:ch14-s6}

下面这条流水线是本章所有片段的合成形态：一次只做一个族，改动可审阅、可回滚，
每一步都有明确的失败出口。

\begin{lstlisting}[style=shellstyle,caption={modernize 一批：导出、应用、格式化、编译、测试},label={lst:ch14-e2e}]
set -euo pipefail
FIX=fixes/modernize-$(date +%Y%m%d)
mkdir -p "$FIX"

# 1) whole-repo analysis, export only
run-clang-tidy -p build -j 8 -quiet \
  -config-file .clang-tidy \
  -checks='modernize-use-nullptr,modernize-use-bool-literals' \
  -export-fixes "$FIX"

# 2) scale check: stop if the blast radius is too large
grep -l FilePath "$FIX"/*.yaml | wc -l

# 3) apply replacements, formatting from .clang-format
clang-apply-replacements "$FIX" --style=file

# 4) run the formatter again to prove it is stable
clang-format --dry-run --Werror $(git ls-files '*.cpp' '*.h')

# 5) build and test; any failure means this batch is dropped
cmake --build build
ctest --test-dir build --output-on-failure
\end{lstlisting}

失败时的回退按步骤对应。第 1 步失败通常是数据库过期，重新生成
\texttt{compile\_commands.json}（第 2 章）；第 3 步报冲突，回到第 1 步把族拆细，
或加 \Opt{ignore-insert-conflict}；第 4 步 \texttt{--dry-run --Werror} 失败，
说明 \texttt{.clang-format} 与替换文本的缩进假设打架，先让 \Fmt\ 单独收敛一次；
第 5 步编译失败，说明这一族需要人工补 include 或前向声明（第 \ref{sec:ch14-s1} 节的第一条限制），
最经济的处理是丢弃本批改动、把该 check 从白名单里摘掉。
回退动作本身交给版本控制：一个族一个提交时，丢弃就是一行 revert。

\begin{longtable}{>{\raggedright\arraybackslash}p{0.22\textwidth}>{\raggedright\arraybackslash}p{0.13\textwidth}>{\raggedright\arraybackslash}p{0.16\textwidth}>{\raggedright\arraybackslash}p{0.12\textwidth}>{\raggedright\arraybackslash}p{0.16\textwidth}}
\caption{第 14 章工具族速查}\label{tab:ch14-tools}\\
\toprule
工具 & 输入 & 输出形态 & 数据库 & 回滚方式 \\
\midrule
\endfirsthead
\toprule
工具 & 输入 & 输出形态 & 数据库 & 回滚方式 \\
\midrule
\endhead
\Tidy\ \texttt{-fix} & 源文件 & 原地改文件 & 需要 & 版本控制丢弃 \\
\Tidy\ \texttt{--export-}\Bk\texttt{fixes} & 同上 & 写单个 YAML & 需要 & 无需回滚 \\
\texttt{clang-apply-}\Bk\texttt{replacements} & fix 目录 & 原地改文件 & 不需要 & 版本控制丢弃 \\
\texttt{run-clang-}\Bk\texttt{tidy} & 数据库 & 打印或写目录 & 需要 & 看是否 \texttt{-fix} \\
\texttt{clang-tidy-}\Bk\texttt{diff.py} & diff 与数据库 & 打印 & 需要 & 只读，无需 \\
\texttt{clang-rename} & 偏移或行列 & 原地或清单 & 需要 & 版本控制丢弃 \\
\texttt{clang-move} & 新旧头文件 & 写清单目录 & 需要 & 清单未应用则安全 \\
\texttt{clang-include-}\Bk\texttt{fixer} & 符号名 & 原地插 include & 建议提供 & 版本控制丢弃 \\
\texttt{include-}\Bk\texttt{cleaner} & 数据库 & 打印增删建议 & 需要 & 只读，无需 \\
\texttt{clang-scan-}\Bk\texttt{deps} & 数据库 & 打印依赖记录 & 需要 & 只读，无需 \\
\texttt{clang-query} & 源文件与查询 & 打印 AST 命中 & 需要 & 只读，无需 \\
\bottomrule
\end{longtable}

\begin{bestpractice}[title={把 autofix 当作一次一个族的提交}]
一个族一个提交、一个 PR，标题写明 check 名单，正文贴上命中文件数与编译测试结果，
这样 review 时能逐族判断「值不值得继续自动化」，出问题也能只回滚一个族。
配套的门禁约束是：CI 白名单里只允许列入清单 \ref{lst:ch14-e2e} 那种可编译验证的族，
\texttt{-fix} 只允许在专门的修复分支上运行，主线 PR 保持只读报告。
\end{bestpractice}

% ---- frag_ch15.tex ----
\chapter{cppcheck、IWYU、CodeQL 与 sanitizer 的分界}\label{ch:boundaries}

第 12 章到第 14 章把 \Tidy\ 从架构讲到批量应用，但它只是「静态分析」这一格里的\emph{一种}实现。
真实项目里同时活跃的引擎往往有四五套，它们之间的差别不在「谁更聪明」，而在三件更朴素的事：
分析发生在代码执行的哪个时刻、需要为它付出什么代价、它给出的结论有多强。
本章先把这条坐标轴画出来，再把 \Bk cppcheck\Bk、include-what-you-use（IWYU）、
clang 内置静态分析器、商业与云端扫描器（Coverity、Klocwork、PC-lint-plus、Parasoft、CodeQL）
和 sanitizer 家族逐个放上去，最后落到选型。

贯穿全章的一条主线值得先写下来：\textbf{静态手段能覆盖全部代码，但只能报「可能」；
动态手段只覆盖被执行到的路径，可一旦触发就是事实}。
所有工具组合、门禁分层与人力分派，都是从这条不等式推出来的。

\section{先画坐标轴：四层手段与各自的失效方式}\label{sec:ch15-s1}

第一层是编译器自己的诊断。它的成本几乎为零，因为分析与正常编译共享同一份 AST、同一笔时间预算；
代价是它只知道语言规则允许它知道的东西。值得默认打开的一组是
\texttt{-Wall}、\texttt{-Wextra}、\texttt{-Wshadow}、\texttt{-Wconversion}、
\texttt{-Wold-style-cast}、\texttt{-Wnon-virtual-dtor}，
再在既有基线清零之后用 \texttt{-Werror} 把警告提升为构建失败。

\begin{lstlisting}[style=shellstyle,caption={编译期诊断层：最小可用的一套警告},label={lst:ch15-warn}]
WARN="-Wall -Wextra -Wshadow -Wconversion"
WARN="$WARN -Wold-style-cast -Wnon-virtual-dtor"
WARN="$WARN -Wdocumentation -Wconditional-uninitialized"

# only once the existing baseline is clean
clang++ -std=c++20 $WARN -Werror -c src/parser.cpp
\end{lstlisting}

清单 \ref{lst:ch15-warn} 里的开关属于 clang 与 GCC 这一系，MSVC 的警告按 \texttt{/W} 等级
另有一套编号，两边不能照着抄。第二层是\emph{不执行代码}的静态分析：\Tidy、clang 内置分析器、
\Bk cppcheck\Bk、CodeQL 与 Coverity 都在这一层。它能报告永远跑不到的分支，
但结论只能停在「存在一条路径使之为真」。第三层是 sanitizer：只有被执行的路径有发言权，
可它一旦开口就无需人工判定。第四层是模糊测试，它把「这段代码究竟能不能被喂到」
改写成搜索问题，而判定用的 oracle（判定器）通常正是第三层。

\begin{longtable}{>{\raggedright\arraybackslash}p{0.13\textwidth}>{\raggedright\arraybackslash}p{0.16\textwidth}>{\raggedright\arraybackslash}p{0.19\textwidth}>{\raggedright\arraybackslash}p{0.16\textwidth}>{\raggedright\arraybackslash}p{0.14\textwidth}}
\caption{四层质量手段的发现窗口与误报方向}\label{tab:ch15-layers}\\
\toprule
层次 & 典型工具 & 发现窗口 & 误报方向 & 成本形态 \\
\midrule
\endfirsthead
\toprule
层次 & 典型工具 & 发现窗口 & 误报方向 & 成本形态 \\
\midrule
\endhead
编译期诊断 & 编译器警告 & 单个 TU，语言规则可见之处 & 几乎没有，报出来就该改 & 与构建同额，可忽略 \\
无执行静态分析 & \Tidy、内置分析器、\Bk cppcheck & 全部源码，含不可达分支 & 假阳性为主，需人判定 & 随代码规模与查询复杂度 \\
运行时动态 & ASan、UBSan、TSan & 只有被执行到的路径 & 假阴性为主，漏而不报 & 随测试覆盖率与运行时间 \\
模糊测试 & libFuzzer、AFL++ & 输入空间上的搜索 & 取决于判定 oracle & 随算力，可用并行摊薄 \\
\bottomrule
\end{longtable}

表 \ref{tab:ch15-layers} 里「误报方向」这一列是全章最有用的东西。静态层的失效方式是\emph{吵}：
它不懂宏展开后的真实类型、不懂模板实例化、不懂外部库的约定，于是把「可疑」当成「结论」。
动态层的失效方式是\emph{沉默}：覆盖率没到的地方它一句都不说，而沉默与健康的报表长得一样。
想清楚这一点，问题就不再是「该用哪个工具」，而是「哪一层的结果需要人来判定」。

\begin{guideline}[title={分层的顺序是依赖关系，不是偏好}]
顺序不能反。第一层不清零，\Tidy\ 的噪声会把人力吞掉：同一段代码编译器与检查族各报一遍，
先让 \texttt{-Werror} 生效，再谈 check 的选择。第三层必须有单元测试承载，
否则插桩后的二进制没有可执行路径，等于什么都没检查。第四层必须有第三层做判定，
否则 fuzzer 只是在随机地跑。把四层压成一层的诱惑（「我们只跑 CodeQL」）
通常同时放弃了编译期的零成本与运行时的确定性这两样东西。
\end{guideline}

\section{cppcheck：没有编译数据库也能跑}\label{sec:ch15-s2}

\Bk cppcheck\Bk 的前端是自己实现的词法切分加一棵简化 AST：它不做完整的预处理展开，
不做模板实例化，也不链接任何 clang 组件。这带来它最大的优点——
\textbf{不需要编译数据库也能工作}，在只有 Makefile、只有散装源码、
或者构建系统在本机根本跑不起来的仓库里仍然能出结果。
需要更准的搜索路径与宏定义时，也可以把数据库喂给它：
\Opt{project=build/\allowbreak compile\_commands.json}，或者照旧手工给 \texttt{-I} 与 \texttt{-D}。

\begin{lstlisting}[style=shellstyle,caption={\Bk cppcheck\Bk 的两种形态：门禁与裸扫},label={lst:ch15-cppcheck}]
# gate form: a real exit code is what makes it a gate
cppcheck --enable=warning,style,performance,portability \
  --std=c++20 --quiet --jobs=8 --inline-suppr \
  --suppressions-list=.cppcheck-suppress \
  --error-exitcode=2 --project=build/compile_commands.json

# no database at all: sweep a tree, every configuration
cppcheck --enable=all --inconclusive --force src/ include/
\end{lstlisting}

几个开关的含义要记住。\Opt{enable=} 决定启用哪些类别（\texttt{error}、\texttt{warning}、
\texttt{style}、\texttt{performance}、\texttt{portability}、\texttt{information}），
默认只有 \texttt{error}，「装了工具却没开类别」是最常见的空扫。
\Opt{inconclusive} 打开那些证据不足的规则，例如看不出构造函数是否初始化了成员时
给出的未初始化告警；它的误报明显更高，适合本地清理，不适合放进门禁。
\Opt{suppress=} 接受 \texttt{id:file:line} 形式的条目，也可以只写 \texttt{id}
或 \texttt{id:file} 来放大作用面；成批压制应该放进 \Opt{suppressions-list=}
指向的清单文件，与源码一起进版本控制、一起 review。
\Opt{inline-suppr} 允许在源码里用一行注释就地压制，粒度最细，也最容易被滥用。
\Opt{error-exitcode=2} 是把它变成质量门的关键：命中即以指定码退出，
CI 的步骤失败即可阻断合并。\Opt{check-config} 只报告配置层面的问题
（找不到的头文件、缺失的宏定义），是接入老仓库时\emph{第一步}该跑的命令——
先确认它真的看到了代码，再讨论它说了什么。输出排版由 \Opt{template} 控制，
机器可读的格式则用来把结果并进统一报表。

\begin{lstlisting}[style=shellstyle,caption={抑制清单：每条压制都带理由与归属},label={lst:ch15-suppress}]
# one entry per line: id[:file[:line]]
missingIncludeSystem:*

# reviewed by the parser team, kept because the
# pointer is validated inside a macro expansion
useStlAlgorithm:src/parser.cpp:128
nullPointer:src/legacy/arena.cpp:77
\end{lstlisting}

addon 是 \Bk cppcheck\Bk 的第二增长曲线：它把外部 \PY\ 脚本挂进同一趟分析，
用同一套结论格式输出。合规相关的 addon（用 \Opt{addon=misra} 请求的 MISRA 那一系列是其中最
确定的一个）适合在交付审计时给出「逐条规则编号加违规位置」的报表；
统计型 addon 用来数老式强转、数魔数。因为它们复用同一个输出通道，
\Opt{suppress} 与 \Opt{error-exitcode} 的语义一并继承，不需要另建一套门禁。

\begin{note}[title={cppcheck 的误报是有形状的}]
它不认识宏展开之后的真实类型，条件编译里的两条分支可能被同时当成成立；
模板只做到很浅的实例化，越依赖 \Tp{auto}、概念与偏特化的代码，
它越容易给出「看不懂的即疑似问题」；它不知道外部构造函数的行为，
于是未初始化成员与空指针解引用这一族会带上证据不足的标签。
反过来，它对 \CL\ 风格的老代码、对裸指针与定长数组、对越界与泄漏的模式识别廉价而有效，
这正是它作为「无构建环境兜底层」的价值。把它当第二个作者来读，而不是当法官来信，
噪声就可管理。
\end{note}

\section{include-what-you-use：用到的不等于包含的}\label{sec:ch15-s3}

IWYU 的公理只有一句：一个翻译单元应当直接包含它\emph{直接使用}的每个实体所在的头文件，
并且不包含它不直接使用的头文件。由此派生两个实践结论：一是传递包含
（靠别人头文件里的 \texttt{\#include} 拿到名字）是债，它会在依赖变更的那天集中偿付；
二是「只需要指针或引用」的实体应当用前向声明把编译依赖切断，让重编范围缩到最小——
这笔账与第 3 章「索引与解析的代价」是同一笔。
IWYU 需要一个\emph{真实}的编译数据库：它本质上是把 clang 前端换成带 include
用途追踪的版本，所以第 2 章的 \texttt{compile\_commands.json} 是它的地基。
它逐文件给出「该删的行」与「该加的行」，默认还会往源码里写建议注释。

\begin{lstlisting}[style=shellstyle,caption={IWYU 的并行跑法与单文件调试},label={lst:ch15-iwyu}]
# whole database, one IWYU process per translation unit
iwyu_tool.py -p build -- -Xiwyu --no_comments \
  -Xiwyu --max_line_length=100 -Xiwyu --cxx17ns

# a single file, with project knowledge injected
iwyu --mapping_file=project.imp src/widget.cpp \
  -std=c++20 -Iinclude
\end{lstlisting}

\texttt{-Xiwyu} 前缀后面跟的是 IWYU 自己的开关而不是编译器的：
\Opt{no\_comments} 只报告不写注释，适合先把结论读一遍再决定；
\Opt{max\_line\_length=} 影响它建议注释的排版；
\Opt{mapping\_file=} 指向一份 \texttt{.imp} 文件，用来告诉 IWYU 那些它无法从语法推断的
项目知识——某个聚合头视为整体对外导出、某个内部生成头不该被直接包含、
某个 \CL\ 头在本项目里的 \CPP\ 替身是谁。

\begin{lstlisting}[style=jsonstyle,caption={.imp 映射文件的大致形态：把约定写成数据},label={lst:ch15-imp}]
[
  { "include": [ "<cstring>", "private",
                 "base/string_util.h", "public" ] },
  { "include": [ "gen/config.h", "public",
                 "generated/config.h", "public" ] },
  { "alt_include": [ "legacy/log.h", "base/log.h" ] }
]
\end{lstlisting}

另一条通道是写在源码里的 pragma 注释，它们只被 IWYU 读取，删掉不影响编译。

\begin{lstlisting}[style=cppstyle,caption={头文件与实现文件里的 IWYU pragma 注释},label={lst:ch15-pragma}]
// widget.h
// IWYU pragma: begin_exports
#include "api_types.h"
#include "geom_types.h"
// IWYU pragma: end_exports

#include <memory>             // IWYU pragma: export
#include <vector>             // IWYU pragma: keep

class Context;                // pointer only: fwd decl

struct Widget {
  std::unique_ptr<Context> ctx_;
};

// widget.cpp
#include "widget.h"           // IWYU pragma: associated
\end{lstlisting}

按清单 \ref{lst:ch15-pragma} 逐条说：\texttt{keep} 表示这行虽被判定多余但必须保留；
\texttt{export} 表示本头文件承诺把该依赖对外提供，聚合头常用；
\texttt{begin\_exports} 与 \texttt{end\_exports} 划定一段范围整体视为导出；
\texttt{provide} 与 \texttt{declare} 配对，指定实体的提供方与声明方；
\texttt{associated} 声明「这个 include 就是当前类的自有头」。
一份长期没人维护 pragma 的仓库，IWYU 的建议会退化成一堆看似荒谬的行，
然后整个工具被团队关掉——这是它最常见的失败方式。

为什么在 \cpp{20} modules 之后 IWYU 的价值曲线会下降？因为「依赖」从隐式的文本包含
移到了显式的 \texttt{import} 声明：模块接口单元编译成 BMI（模块字节码）之后，
消费者不再复制头文本，传递包含的债与「改一行头文件全仓重编」的代价同时消解。
那时要问的不再是「哪些 include 可以删」，而是「模块图是否最小且无环」，
那是第 14 章 \texttt{clang-scan-deps} 的领域。但只要仓库仍以头文件为主要依赖形态
（绝大多数存量项目正是如此），IWYU 仍是编译时间的第一道闸门。

把它和第 14 章的 \Bk include-cleaner\Bk 对照着看更清楚：IWYU 独立于 \CD{}，
需要自己维护 \texttt{.imp} 与 pragma 这套知识，输出是一行行建议；
include-cleaner 与 \CD\ 的后台索引、codeAction 同源，也能走
\texttt{misc-include-cleaner} 这条 \Tidy\ 通道，因此可以直接接入第 14 章的
导出—审阅—应用流水线。想一次性把全仓 include 图洗净，IWYU 更彻底；
想让每个 PR 顺手删掉本文件的多余 include，include-cleaner 更省人力。

\section{clang 内置分析器与云端扫描器}\label{sec:ch15-s4}

\subsection*{路径敏感的符号执行：clang 内置静态分析器}

clang 前端自带一个与 \Tidy\ 机理不同的检查引擎：它沿控制流图做\emph{路径敏感}遍历，
用约束求解维护符号状态，因此能区分「同一个空指针的两条分支」，也能沿调用链跨函数传播状态。
它不需要 \Tidy\ 参与，直接调用编译器即可：

\begin{lstlisting}[style=shellstyle,caption={对若干 TU 运行内置静态分析器},label={lst:ch15-analyze}]
mkdir -p build/analyze
for f in src/*.cpp; do
  clang++ --analyze -o build/analyze "$f" \
    -std=c++20 -Iinclude \
    -Xclang -analyzer-checker=core,cplusplus,security,optin \
    -Xclang -analyzer-output=text
done
\end{lstlisting}

checker 的分组名是真实分类：\texttt{core}（空指针解引用、除零、未初始化值、内存模型）、
\texttt{cplusplus}（对象生命周期、move 之后使用、\Tp{new} 相关）、
\texttt{security}（敏感信息与不安全 API）、\texttt{optin} 及其 \texttt{optin.cpp} 子组
（性能与 API 使用建议，默认关闭，需显式开启），再加上一整套 \texttt{alpha.*} 实验项与
平台组（\texttt{unix}、\texttt{deadcode}、\texttt{valist}）。
要完整的清单，向编译器索取 \texttt{-analyzer-checker-help} 即可。
输出形态由 analyzer-output 选择文本、HTML 或 plist，plist 是并进报表系统时最好用的格式。

代价同样明确：路径爆炸。函数调用深度与分支数量共同决定状态空间，
一个中等规模 TU 吃掉几百兆内存、几十秒时间并不罕见。
这正是它与 \Tidy\ 互补的原因，两者的差别可以逐条列出来：

\begin{longtable}{>{\raggedright\arraybackslash}p{0.16\textwidth}>{\raggedright\arraybackslash}p{0.32\textwidth}>{\raggedright\arraybackslash}p{0.32\textwidth}}
\caption{同一份代码上的两种引擎：模式匹配与符号执行}\label{tab:ch15-vs}\\
\toprule
维度 & \Tidy\ 的自有 check & clang 内置分析器 \\
\midrule
\endfirsthead
\toprule
维度 & \Tidy\ 的自有 check & clang 内置分析器 \\
\midrule
\endhead
机制 & 在 AST 上跑匹配器，命中即报 & 遍历 CFG 加约束求解，维护符号状态 \\
成本 & 与代码行数近似线性，每文件秒级 & 与路径数相关，单文件可达分钟级 \\
误报方向 & 不懂语义：宏、模板、外部约定 & 不懂意图：把理论可行当成现实可能 \\
能否跨函数 & 单个 check 内基本不跨 & 能，可沿调用链传播状态 \\
能否给修复 & 多数族带 fix-it，可导出批量应用 & 通常只报告，不提供修复 \\
CI 预算 & 每个 PR 全量跑 & 限定目录或定时跑，单独计时 \\
\bottomrule
\end{longtable}

两者在配置文件里以 \texttt{clang-analyzer-*} 前缀共处，看上去像一个工具的两组规则，
实为两台引擎；把它们塞进同一个 CI 步骤共享同一个超时，是接入初期最常见的排障噩梦。
构建期扫描的入口这些年在变：过去的常见做法是用 \texttt{scan-build} 那层包装
把分析器插进构建过程，而较新的 LLVM 发行包里这类脚本的随附状态与维护姿态都不稳定，
社区接管了同类功能。因此更稳的写法是像清单 \ref{lst:ch15-analyze} 那样
直接从编译数据库驱动 \texttt{--analyze}、自己收结果，或者把这块需求交给下一节的云端扫描器。

\subsection*{先建模再查询：商业与云端扫描器}

Coverity、Klocwork、PC-lint-plus、Parasoft 与 CodeQL 属于同一种形态：
\textbf{先把代码建成一个可查询的模型，再在模型上跑跨函数、跨文件的查询}。
这给出脚本式分析器给不了的三样东西：留存（能回答「这个问题上周存在吗」）、
团队级趋势曲线（每次扫描对齐到同一份问题标识）、真正的跨 TU 数据流。
它们的接入成本也来自同一个地方——建模阶段必须看到真实构建：
Coverity 要先 capture 一趟构建生成中间数据库，构建被脚本层层拼装时捕获最容易失真；
Klocwork 强调增量索引与 IDE 内视图；PC-lint-plus 以精确与可调著称，
代价是要认真维护它的标准库与项目配置；Parasoft 把静态分析与单元测试、覆盖率、
编码规范合规绑在一个产品里，适合有审计与追溯要求的团队。

GitHub CodeQL 是这一形态在开源生态里的版本：它把代码抽取成关系数据库，
用声明式的 QL 语言查询。一个查询文件的结构是「一段元数据注释加一条
\texttt{from}/\texttt{where}/\texttt{select}」：元数据里的 kind 字段决定它是告警型还是统计型，
\texttt{from} 声明要遍历的实体变量，\texttt{where} 给出约束，\texttt{select} 输出位置、
消息与严重度。写查询的本质是在它提供的 \CPP\ 库（AST、CFG、数据流这些关系的集合）上
做声明式筛选，再配一份同名的查询帮助文档，就能进规则集。

\begin{lstlisting}[style=yamlstyle,caption={Actions 里的 CodeQL 作业：构建由自己跑},label={lst:ch15-codeql}]
name: codeql
on:
  pull_request:
    branches: [ main ]
jobs:
  analyze:
    runs-on: ubuntu-latest
    permissions:
      security-events: write
    steps:
      - uses: actions/checkout@v4
      - uses: github/codeql-action/init@v3
        with:
          languages: cpp
          build-mode: none
      - run: cmake -B build --preset analyze
      - run: cmake --build build
      - uses: github/codeql-action/analyze@v3
\end{lstlisting}

对 \CPP\ 这类编译型语言，抽取阶段必须看到真实编译，所以构建模式的选择很关键：
\texttt{none} 表示 init 不代跑构建，由你在 init 与 analyze 两个步骤之间自己构建
（清单 \ref{lst:ch15-codeql} 就是这个形态，前提是构建命令与第 2 章的
\texttt{compile\_commands.json} 一致）；autobuild 会替你推断构建方式，
对纯 CMake 项目常常成功，但项目里混了自定义脚本或交叉工具链时，
排查「它为什么没抽到文件」比手工指定更费时间。
本地复现与查询调试走命令行：先创建数据库（指定语言与构建命令、源码根目录），
再对数据库跑查询套件，产出 SARIF 供平台展示。
抽取数据库的体积常到数 GB，一套查询跑完也以十分钟计，
这就是「CodeQL 放 CI 而不是放 pre-commit」的直接原因，
与第 16 章按耗时给门禁分层的思路一致。

\begin{warning}[title={「0 findings」几乎从不等于安全}]
任何扫描器返回空结果都可能有三种原因，而三者在报表上长得一模一样。
其一，工具根本没看到代码：\Bk cppcheck\Bk 找不到头文件时会跳过它读不懂的部分，
CodeQL 的抽取没覆盖到某个目标时数据库里就是空的，\Tidy\ 没有编译数据库则整批文件
报同一个错然后被忽略。其二，规则没开：\Opt{enable=} 的类别、\Tidy\ 的
\texttt{Checks} 白名单、CodeQL 的查询套件，默认都只是全量规则的一个子集，
「装了工具」不等于「启用了规则」。其三，被压制淹没：抑制清单与 \texttt{NOLINT}
积累几年之后，报表干净得像新写的代码。所以门禁除了断言零结果，
更该断言「分析到的 TU 数量」与「启用的规则数量」这两个过程指标。
\end{warning}

\section{sanitizer：结论强度的一次跃迁}\label{sec:ch15-s5}

sanitizer 是编译器在每次内存访问与每次可疑运算前后插入的检查代码，外加一份运行时库。
它不做推测：报告即事实。下面表 \ref{tab:ch15-san} 的简称依次对应
Address、UndefinedBehavior、Thread、Memory、Leak、HWAddress 六个 sanitizer。

\begin{longtable}{>{\raggedright\arraybackslash}p{0.11\textwidth}>{\raggedright\arraybackslash}p{0.25\textwidth}>{\raggedright\arraybackslash}p{0.22\textwidth}>{\raggedright\arraybackslash}p{0.23\textwidth}}
\caption{sanitizer 家族对照}\label{tab:ch15-san}\\
\toprule
简称 & 抓什么 & 机制 & 前提与代价 \\
\midrule
\endfirsthead
\toprule
简称 & 抓什么 & 机制 & 前提与代价 \\
\midrule
\endhead
ASan & 堆、栈、全局缓冲区的越界，释放后使用，双重释放，以及泄漏 & 影子内存加红区（redzone） & 约两倍内存与数倍时间，主流平台可用 \\
UBSan & 有符号溢出、空指针、错位访问、越界下标一类的未定义行为 & 在可疑运算点插入断言 & 最便宜；可 trap 也可出诊断 \\
TSan & 数据竞争，即并发访问缺少正确的先后关系 & 向量时钟追踪访问顺序 & 不能与 ASan 同用，开销大 \\
MSan & 未初始化内存的读取 & 影子位跟踪每一位的有效性 & 依赖库必须一并插桩，否则假阳性 \\
LSan & 泄漏的堆对象 & 扫描内存可达性 & 通常由 ASan 默认开启 \\
HWASan & 与 ASan 同类的越界与释放后使用 & 硬件内存标签（AArch64） & 开销远低于 ASan，用于真机 \\
\bottomrule
\end{longtable}

TSan 那一行说的「先后关系」就是 happens-before 关系：只要两次访问之间没有建立同步边，
它就把这判为竞争，而不必等到结果真的错乱。与环境变量的配合决定 sanitizer 能否成为门禁：
\texttt{ASAN\_OPTIONS} 里的 \texttt{detect\_leaks=1} 与 \texttt{halt\_on\_error=1}
让它第一次违规就终止，便于把失败定位到具体用例；
\texttt{UBSAN\_OPTIONS} 里的 \texttt{print\_stacktrace=1} 让报告带上栈回溯，
否则只剩一个文件行号。UBSan 有两种失败姿态：默认打印诊断并可恢复继续执行，
加 \texttt{-fno-sanitize-recover=all} 则命中即停；编译成 trap 指令
（\texttt{-fsanitize-trap}）能去掉运行时库依赖、缩小产物，代价是报告信息很粗，
只剩「这里出过错」。取舍标准是：要栈与可读报告就选诊断，
要嵌进资源受限的测试环境就选 trap。

\begin{lstlisting}[style=shellstyle,caption={CI 里 sanitizer 与单元测试的配对},label={lst:ch15-san-ci}]
# one dedicated instrumented build, never the release one
S="-fsanitize=address,undefined -g -fno-omit-frame-pointer"
cmake -B build-asan -DCMAKE_BUILD_TYPE=Debug \
  -DCMAKE_CXX_FLAGS="$S" -DCMAKE_EXE_LINKER_FLAGS="$S"
cmake --build build-asan --target unit-tests --parallel 8

ASAN_OPTIONS=detect_leaks=1:halt_on_error=1 \
UBSAN_OPTIONS=print_stacktrace=1:halt_on_error=1 \
ctest --test-dir build-asan --output-on-failure

# TSan needs its own build directory, never share one
cmake -B build-tsan -DCMAKE_CXX_FLAGS="-fsanitize=thread -g"
\end{lstlisting}

清单 \ref{lst:ch15-san-ci} 有三个容易被忽略的点。一是必须单开一个插桩构建目录：
sanitizer 与优化、与 \texttt{-flto}、与另一个 sanitizer 的组合都有一堆限制，
复用发布构建只会得到一堆无法解释的行为。二是 \texttt{-fno-omit-frame-pointer} 与 \texttt{-g}
决定报告里的栈能不能读，没有它们就只能靠地址硬猜。三是这里\emph{只跑单元测试}：
sanitizer 的算力预算来自覆盖率，测试越全，它说得越多。

\subsection*{与 \Tidy\ 的重叠处}

两族工具常在同一个问题上碰面，但结论性质不同。
\texttt{bugprone-use-after-move} 在 AST 上看到 move 之后再使用该对象，
却可能被一个语义上安全的自定义 \Tp{swap}、或被分支可达性判断骗过——它报的是「可能」；
ASan 只在被移动对象真的被解引用且内存已失效时才开口，一旦说就是事实，前提是有测试走到那里。
又如 \texttt{cppcoreguidelines-\allowbreak pro-bounds-pointer-\allowbreak arithmetic}
对所有指针算术一律皱眉，噪声很高；把越界的判定交给 ASan 的红区，
再用 \Tidy\ 只保留「显式下标访问」的窄化版本，通常更划算。
同理，\texttt{clang-analyzer-core.NullDereference} 与 UBSan 的空指针检查
是同一件事的两种置信度：前者能覆盖没有测试的分支，后者能确认线上真的会崩。

\begin{compare}[title={静态与动态：四种能力对照}]
\begin{tabular}{@{}>{\raggedright\arraybackslash}p{0.20\textwidth}>{\raggedright\arraybackslash}p{0.27\textwidth}>{\raggedright\arraybackslash}p{0.27\textwidth}@{}}
\toprule
能力 & 无执行静态分析 & sanitizer 运行时检测 \\
\midrule
未执行的路径 & 能报告，这正是它的价值 & 完全没有声音 \\
结论强度 & 「可能」，需人判定后才动手 & 「确实」，可直接终止构建 \\
成本随什么增长 & 代码规模与查询复杂度 & 测试覆盖率与执行时间 \\
失效形态 & 宏、模板、外部库导致假阳性 & 覆盖率不足导致假阴性 \\
\bottomrule
\end{tabular}
\end{compare}

\subsection*{把判定交给模糊测试}

模糊测试不是第五层，而是把第三层放大到输入空间的手段：靶子通常是解析函数，
判定条件仍然是「不崩溃、不违反 sanitizer」。

\begin{lstlisting}[style=shellstyle,caption={fuzz 驱动与 ASan 一起链接},label={lst:ch15-fuzz}]
clang++ src/fuzz_parser.cc src/parser.cc -o fuzz-parser \
  -fsanitize=fuzzer,address -g -O1

./fuzz-parser -max_total_time=600 build/corpus/
\end{lstlisting}

\begin{lstlisting}[style=cppstyle,caption={最小 fuzz 目标：入口只提供字节},label={lst:ch15-fuzztarget}]
#include <cstddef>
#include <cstdint>
#include "parser.h"

extern "C" int LLVMFuzzerTestOneInput(const uint8_t *Data,
                                      size_t Size) {
  Parser p;
  p.Feed(Data, Size);   // crash or sanitizer report
  return 0;
}
\end{lstlisting}

清单 \ref{lst:ch15-fuzztarget} 的入口里没有断言：判定逻辑放在被测库里，而不是放在那个名为
\texttt{LLVMFuzzerTest-\allowbreak OneInput()} 的回调里。
语料（corpus）与最小化后的崩溃用例一起进制品库，
是这类作业可复现的关键，也是它必须与 ASan 或 UBSan 同时构建的原因。
不依赖进程内插桩、需要驱动外部进程或真实文件系统时，
AFL++ 的覆盖率插桩与任务队列更合适；要快速迭代单个目标函数、要进程内并行时，
libFuzzer 更省事。两者的价值上限都由 sanitizer 决定：没有判定器的 fuzz 只是随机执行。

\section{选型：四把尺子与推荐组合}\label{sec:ch15-s6}

评价任何一层只用四把尺子：\textbf{覆盖面}（能触到多少代码与多少路径）、
\textbf{运行成本}（每次消耗的机器时间加上人力维护）、\textbf{误报方向}
（错在报得多还是漏得多）、\textbf{是否需要真实输入}（要不要构建、测试或语料）。
四把尺子量下来，组合就不再是品味问题。

\begin{longtable}{>{\raggedright\arraybackslash}p{0.19\textwidth}>{\raggedright\arraybackslash}p{0.30\textwidth}>{\raggedright\arraybackslash}p{0.31\textwidth}}
\caption{按项目形态推荐的工具组合}\label{tab:ch15-select}\\
\toprule
场景 & 推荐组合 & 取舍理由 \\
\midrule
\endfirsthead
\toprule
场景 & 推荐组合 & 取舍理由 \\
\midrule
\endhead
小仓库、构建不稳定 & 编译警告加 \Bk cppcheck\Bk 手工跑 & 不依赖数据库，误报可当场判读 \\
中型项目、有 CMake & 警告、\Tidy\ 增量、ASan 与 UBSan 跑单测 & 覆盖面与成本最平衡，门禁可长期承受 \\
大型单体、合规压力 & 上述加 Coverity 或 CodeQL 定时全量 & 需要跨函数数据流与趋势留存 \\
汽车与医疗器械 & 上述加 \Bk cppcheck\Bk 的合规 addon 与规则基线 & 每个问题要能追溯到规则编号与抑制记录 \\
嵌入式与 ARM64 交叉 & \Tidy、真机上的 HWASan、少量 UBSan & MSan 的插桩前提做不到，ASan 在真机上太重 \\
攻击面大的解析库 & 上述加 libFuzzer 或 AFL++ 长期作业 & 唯一能自动搜索崩溃输入的一层 \\
\bottomrule
\end{longtable}

其中两行需要额外说明。交叉与嵌入式里 MemorySanitizer 基本出局：它要求参与链接的每个库
都被同样插桩，而厂商提供的二进制往往做不到，缺一片就报出成片假阳性；
ASan 在内存紧张的真机上开销不可接受，于是硬件标签方案 HWASan 成为替代，
在支持内存标签扩展的 ARM64 机器上还能进一步把开销压低。
合规场景的重点不是「找到更多问题」，而是「每个问题可追溯到一条规则编号、
每一次抑制都有记录」，这决定了 addon 配置与抑制清单文件必须进版本控制、必须被 review。

\begin{bestpractice}[title={按结论强度分派 finding，并写清 SLA}]
把不同层的 finding 派给不同的人与流程，否则一定烂尾。
编译警告与 \Tidy\ 的风格族：谁改这段代码谁修，PR 内闭环，门禁直接阻断。
\Tidy\ 的 bugprone 族与内置分析器：模块 owner 判定，允许带 \texttt{NOLINT}
与注释延后处理，但抑制必须写理由，并在下个迭代复审。
sanitizer 报出的崩溃：按缺陷对待，附最小复现用例与语料，当天响应，
因为它给的已经是事实而不是猜测。
静态层的假阳性：进抑制清单并设定周期性复审，清单本身要有 owner。
分派依据是结论强度，不是工具名气：越接近事实，SLA 越短。
\end{bestpractice}

把本章收束成一句可执行的话：\textbf{静态层负责「不要漏」，动态层负责「不要错」，
两者之间用抑制清单与分派流程隔开}。第 16 章会把这几层写进 CMakePresets 与 CI 作业，
让耗时不同的层各归其位；第 17 章的排障手册则从「工具报了但看起来不对」这一类现场问题
反查本章的分界。

% ---- frag_ch16.tex ----
\part{工程化与选型}

\chapter{CMakePresets 与 CI 的一键质量门}\label{ch:ch16}

前十几章把每个工具的机制拆开讲清楚了：编译数据库怎么来（第 2 章）、\CD\ 怎么读它（第 5 章）、\Fmt\ 的字段语义（第 9 章）、\Tidy\ 的 checks 选择（第 12 章）。本章把它们接到两个真实问题上：第一，怎么用一份\texttt{CMakePresets.json} 让任何人 clone 下来都能敲两条命令进入正确状态；第二，怎么把格式检查、静态检查、测试与 sanitizer 收进同一个入口，并且让 CI 与本地跑的是同一条命令。

这里的判断标准只有一个：\textbf{质量门是否可复现}。如果新同事要读三页 README 才能配好环境，或者 CI 里跑的命令在本地无法原样执行，那么这套门禁一定会在半年内退化成一堆各自手敲的别名。

\section{CMakePresets 的分块模型}\label{sec:ch16-s1}

\texttt{CMakePresets.json} 是一份纯 JSON 描述文件，把 configure、build、test、package 四类动作各拆成一个 preset 数组。本章示例统一把顶层的 \texttt{version} 字段写成 3（对应 CMake 3.21 及以上），因为 \texttt{condition}、\texttt{binaryDir} 的宏展开与 \texttt{environment} 到这一版才稳定可用；\texttt{packagePresets} 需要 v5（CMake 3.23），\texttt{workflowPresets} 需要 v8（CMake 3.27）。落笔前先确认 CI 镜像里 CMake 的主版本号支持你写的 schema 版本，否则整份文件会被拒绝解析。较新版本才引入的字段请先核对所装 CMake 的 \texttt{cmake-presets(7)} 手册，本手册只使用 v3 的稳定字段集。

\subsection{configurePresets 的十个字段}\label{sec:ch16-s2}

{\footnotesize
\begin{longtable}{@{}>{\raggedright\arraybackslash}p{0.24\textwidth}
                  >{\raggedright\arraybackslash}p{0.24\textwidth}
                  >{\raggedright\arraybackslash}p{0.33\textwidth}@{}}
\caption{configurePresets 常用字段}\label{tab:ch16-conf}\\
\toprule
字段 & 取值形态 & 作用与注意点 \\
\midrule
\endfirsthead
\toprule
字段 & 取值形态 & 作用与注意点 \\
\midrule
\endhead
\bottomrule
\endlastfoot
\texttt{name} & 字符串 & preset 的唯一标识，命令行 \Opt{preset} 用它；只允许可打印 ASCII。 \\
\texttt{displayName} & 字符串 & 给人看的名字，IDE 与 CMake Tools 面板显示它；不影响命令行。 \\
\texttt{inherits} & 字符串数组 & 继承链，先列出的优先级更低；支持多继承，冲突时后出现者覆盖。 \\
\texttt{hidden} & 布尔 & 只做基块，不允许被 \Opt{preset} 直接点名，通常与 \texttt{inherits} 配套。 \\
\texttt{generator} & 字符串 & 如 \texttt{Ninja}、\texttt{Unix Makefiles}、\texttt{Visual Studio 17 2022}；决定能否导出编译数据库。 \\
\texttt{binaryDir} & 字符串 & 构建目录，习惯写成 \texttt{build/}\-\allowbreak\texttt{\$\{presetName\}}，让每个 preset 独占一个目录。 \\
\texttt{sourceDir} & 字符串 & 默认即 JSON 所在目录；多工程共用一份 presets 时才需要显式改。 \\
\texttt{cacheVariables} & 对象 & 等价于 \texttt{cmake -DV=T}；值里的 \texttt{\$\{presetName\}} 这类宏会被展开，是放质量开关的主位置。 \\
\texttt{toolchainFile} & 字符串 & 交叉编译工具链文件路径，同样支持宏展开；与 \texttt{environment} 里的 \texttt{CXX} 二选一。 \\
\texttt{environment} & 对象 & 只在该次 configure 有效的环境变量，值可写成 \texttt{\$env\{PATH\}} 之类的引用。 \\
\texttt{condition} & 对象 & 平台与环境门控：\texttt{type} 取 \texttt{equals}、\texttt{includes}、\texttt{matches} 等，\texttt{lhs} 常写 \texttt{\$\{hostSystem\allowbreak Name\}}。 \\
\end{longtable}
}

\texttt{condition} 的意义是让一份文件适配所有平台。不满足条件的 preset 在 \texttt{cmake}\Opt{list-presets} 里显示为不可用或直接消失，而不是报错，所以矩阵式 CI 可以共用同一份 presets。

\subsection{buildPresets 与 testPresets}\label{sec:ch16-s3}

build preset 的字段少而关键：\texttt{configurePreset} 指向要构建哪个 configure preset（必填）；\texttt{jobs} 控制并行度，等价于 \texttt{-j}，在 CI 容器里应当显式写死，否则容器按宿主核数起任务容易 OOM；\texttt{targets} 指定目标子集，跑单测时只构建测试目标能省一半时间；\texttt{configuration} 只对多配置生成器有意义；\texttt{verbose} 打印完整命令行，排障时打开；\texttt{nativeToolOptions} 把参数透传给底层构建工具。

这里有一条必须说清的边界：\textbf{\texttt{nativeToolOptions} 传给的是 ninja 或 make，不是编译器}。写由 \texttt{-k} 与 \texttt{0} 组成的数组让 ninja 在报错后继续构建其它目标是合理的；写 \texttt{-fsanitize=address} 则完全无效，因为编译器根本看不到它。sanitizer 属于编译与链接选项，正确落点是 configure preset 的 \texttt{cacheVariables}。

test preset 除了 \texttt{configurePreset} 外，最重要的是 \texttt{output} 块（\texttt{verbosity}、\texttt{outputOnFailure}）、\texttt{execution} 块（把 \texttt{noTestsAction} 设为 \texttt{error} 可以防止「测试根本没跑」被当成绿色）、\texttt{filter} 与 \texttt{exclude} 块（按名称列表或标签排除慢用例），以及 \texttt{coverage} 块——它只在生成器支持覆盖率时有效，会把 \texttt{llvm-cov} 之类的原始数据留在构建目录里供夜间任务汇总。

\section{把质量工具的开关落进 preset}\label{sec:ch16-s4}

对第 5、9、12 章的所有工具来说，它们依赖的只有两样东西：一份\texttt{compile\allowbreak\_\allowbreak commands.json}，和一组正确的编译选项。所以 preset 的职责非常单纯。

\begin{note}[title={编译数据库开关只认生成器}]
\texttt{CMAKE\_EXPORT\_COMPILE\_COMMANDS} 仅对 Makefile Generators 与 Ninja 系列有效。把它写进一个 generator 为 \texttt{Visual Studio 17 2022} 的 preset，CMake 不会报警，但构建目录里不会出现 \texttt{compile\_commands.json}。此时 \CD\ 找不到数据库会退到 fallback command（第 5 章），表现为「能补全标准库，但项目符号大量标红」。这就是 \ref{sec:ch16-s8} 那个 warning 框的由来。
\end{note}

\begin{lstlisting}[style=jsonstyle,caption={一份可直接用于质量门的 CMakePresets.json（节选）},label={lst:ch16-presets}]
{
  "version": 3,
  "cmakeMinimumRequired": {
    "major": 3,
    "minor": 21,
    "patch": 0
  },
  "configurePresets": [
    {
      "name": "base",
      "hidden": true,
      "generator": "Ninja",
      "binaryDir": "${sourceDir}/build/${presetName}",
      "cacheVariables": {
        "CMAKE_EXPORT_COMPILE_COMMANDS": "ON",
        "CMAKE_CXX_STANDARD": "20",
        "CMAKE_CXX_STANDARD_REQUIRED": "ON",
        "CMAKE_CXX_EXTENSIONS": "OFF"
      }
    },
    {
      "name": "dev",
      "displayName": "Local dev",
      "inherits": "base",
      "cacheVariables": {
        "CMAKE_BUILD_TYPE": "Debug",
        "CMAKE_CXX_FLAGS_INIT": "-Wall -Wextra -Wpedantic"
      }
    },
    {
      "name": "asan",
      "displayName": "ASan + UBSan",
      "inherits": "dev",
      "cacheVariables": {
        "CMAKE_CXX_FLAGS_INIT":
          "-fsanitize=address,undefined",
        "CMAKE_EXE_LINKER_FLAGS_INIT":
          "-fsanitize=address,undefined"
      }
    },
    {
      "name": "ci-linux",
      "displayName": "CI Linux",
      "inherits": "dev",
      "environment": {
        "CC": "clang-19",
        "CXX": "clang++-19"
      },
      "condition": {
        "type": "equals",
        "lhs": "${hostSystemName}",
        "rhs": "Linux"
      }
    }
  ],
  "buildPresets": [
    {
      "name": "dev",
      "configurePreset": "dev",
      "jobs": 8,
      "nativeToolOptions": ["-k", "0"]
    },
    {
      "name": "asan",
      "configurePreset": "asan",
      "jobs": 4
    }
  ],
  "testPresets": [
    {
      "name": "asan",
      "configurePreset": "asan",
      "output": {
        "outputOnFailure": true,
        "verbosity": "default"
      },
      "execution": {
        "noTestsAction": "error",
        "stopOnFailure": false
      },
      "exclude": { "tests": ["slow_network"] }
    }
  ]
}
\end{lstlisting}

日常只需两条命令：

\begin{lstlisting}[style=shellstyle,caption={用 preset 完成 configure 与 build}]
cmake --preset dev
cmake --build --preset dev
ctest --preset asan        # requires the asan build
\end{lstlisting}

\subsection{带 INIT 后缀的坑}\label{sec:ch16-s5}

\texttt{CMAKE\_CXX\_FLAGS\_INIT} 只在 \texttt{project()} 首次初始化 \texttt{CMAKE\_CXX\_FLAGS} 时生效。一旦构建目录的缓存里已有 \texttt{CMAKE\_CXX\_FLAGS}，再改 \texttt{*\_INIT} 不会覆盖它，于是出现「preset 里加了新 flag，但 compdb 里仍是旧选项」这类怪现象。三种可靠解法：每个 preset 独占 \texttt{binaryDir}（本章示例的写法）、CI 每次干净 checkout 而不复用构建目录、或者显式清掉缓存变量后重新 configure。跨 TU 的 \Tidy\ 结果依赖编译选项一致，这个坑的代价在第 3 章讨论索引成本时已经付过一次。

\subsection{个人偏好不要污染工程}\label{sec:ch16-s6}

需要固化某个开发者的路径或编译器时，用\texttt{CMakeUserPresets.json}，并把它写进 \texttt{.gitignore}。它通过 \texttt{include} 字段把仓库级 presets 拉进来，再在本地继承与覆写：

\begin{lstlisting}[style=jsonstyle,caption={CMakeUserPresets.json 与 .gitignore 的配合},label={lst:ch16-user}]
// CMakeUserPresets.json
{
  "version": 3,
  "include": ["CMakePresets.json"],
  "configurePresets": [
    {
      "name": "dev-llvm-local",
      "inherits": "dev",
      "environment": {
        "CXX": "/opt/llvm-19/bin/clang++"
      }
    }
  ]
}

// .gitignore
/CMakeUserPresets.json
/build/
\end{lstlisting}

\begin{warning}[title={CMakeUserPresets.json 不能提交}]
这个文件的定位就是「个人机器上的差异」。一旦提交，别人 clone 下来会因为里面写着 \texttt{/home/alice/...} 这类绝对路径而 configure 失败；它还常被用来放含密钥路径的 \texttt{CMAKE\_PREFIX\_PATH}。反过来，若 CI 依赖了一个只在本地存在的同名文件，就会得到「CI 上 preset 不存在」的错误。把 \texttt{cmake}~\Opt{list-presets} 放进 CI 的第一步，可以立刻暴露这类差异。
\end{warning}

\section{唯一入口：verify.sh}\label{sec:ch16-s7}

\begin{bestpractice}[title={一个仓库只允许一个质量门入口脚本}]
把 configure、build、格式化检查、\Tidy、cppcheck、ctest 收进\texttt{scripts/verify.sh}，本地开发者与 CI 都调用同一个脚本、传同一个 preset 名。文档里只写这一条命令，不再罗列零散的工具调用。任何「CI 里额外多做一步」的需求都必须先落到脚本里，否则本地跑绿、CI 跑红的漂移会在两周内让人开始忽略这个门。
\end{bestpractice}

\begin{lstlisting}[style=shellstyle,caption={仓库级质量门脚本 scripts/verify.sh},label={lst:ch16-verify}]
#!/usr/bin/env bash
# The single quality gate entry: same command locally and in CI.
set -uo pipefail
PRESET="${1:-ci-linux}"
BUILD="build/$PRESET"
rc=0
step() { printf "\n=== %s ===\n" "$1"; }

step "configure"
cmake --preset "$PRESET" || exit 1

step "build"
cmake --build --preset "$PRESET" || rc=1

step "format: whole tree, dry-run"
git ls-files "*.c" "*.h" "*.cc" "*.cpp" \
  | xargs -r clang-format --dry-run --Werror || rc=1

step "clang-tidy: changed files only"
base="$(git merge-base HEAD origin/main)"
git diff -U0 "$base" | clang-tidy-diff.py -p1 -path build \
  -checks-file .ci/tidy-ci.json -quiet || rc=1

step "cppcheck"
cppcheck --error-exitcode=2 --enable=warning,style \
  --project="$BUILD/compile_commands.json" \
  --suppressions-list=.ci/cppcheck-suppress.txt \
  --output-file=.ci/cppcheck.log || rc=1

step "tests under sanitizers"
cmake --build --preset asan && ctest --preset asan || rc=1

printf "\n=== quality gate rc=%s ===\n" "$rc"
exit "$rc"
\end{lstlisting}

脚本里有几处刻意的选择值得说明。第一，格式化步骤跑\textbf{全树}的 \Opt{dry-run} 加 \Opt{Werror}，因为它极快且必须绝对干净，任何一处不一致都要在 PR 里暴露；而 \Tidy\ 步骤只跑\textbf{改动行}（第 10 章的增量策略），全量留给夜间任务的 \texttt{run-clang-tidy}。第二，\texttt{clang-tidy-diff.py} 的 \texttt{-path build} 告诉它去\texttt{build/} 下找 compdb，它与 \texttt{-p1} 一起决定 diff 路径能否解析成绝对路径——这两处配错的症状是「一条检查都不报」，详见第 17 章。第三，cppcheck 用 \Opt{error-exitcode=2} 而不是默认的 0，这样脚本能用同一个 \texttt{rc} 变量收集退出码；\Opt{project=} 直接复用 CMake 导出的 compdb，避免第二套 include 路径来源。

\begin{lstlisting}[style=shellstyle,caption={夜间全量与本地最小闭环}]
# Nightly: whole tree, emit a fixes file for the auto PR
run-clang-tidy -p 8 -j 4 -export-fixes nightly.yml \
  > nightly.log 2>&1

# Local minimal loop, must finish in seconds
files=$(git diff --name-only HEAD -- "*.cpp" "*.h")
clang-format -i $files
ctest --preset dev -R 'unit_.*'
\end{lstlisting}

\section{版本对齐：谁约束谁}\label{sec:ch16-s8}

质量门最容易腐坏的地方不是配置，而是版本。四个消费者必须看到同一个 LLVM 大版本：CI 容器镜像、pre-commit hook 的 \texttt{rev}、开发者本机包管理器，以及 \CD\ 的安装方式。约束关系应当是单向的。

{\footnotesize
\begin{longtable}{@{}>{\raggedright\arraybackslash}p{0.19\textwidth}
                  >{\raggedright\arraybackslash}p{0.21\textwidth}
                  >{\raggedright\arraybackslash}p{0.20\textwidth}
                  >{\raggedright\arraybackslash}p{0.21\textwidth}@{}}
\caption{版本由谁定义、由谁消费}\label{tab:ch16-ver}\\
\toprule
约束项 & 由谁定义 & 由谁消费 & 不一致的后果 \\
\midrule
\endfirsthead
\toprule
约束项 & 由谁定义 & 由谁消费 & 不一致的后果 \\
\midrule
\endhead
\bottomrule
\endlastfoot
LLVM 大版本 & CI 镜像，即 \texttt{apt.llvm.\allowbreak org} 的 \texttt{llvm-\allowbreak toolchain-\allowbreak N} & \Fmt、\Tidy、\CD、cppcheck & 同一份代码两种格式化结果；check 名不存在 \\
\texttt{.clang-\allowbreak format} 语义 & 仓库根文件 & \Fmt\ 的全部版本 & 新版本才有的字段被旧版本静默忽略 \\
\texttt{.clang-\allowbreak tidy} 的 check 名 & 仓库文件加 LLVM 版本 & \Tidy & 未知 check 只警告不报错，检查集悄悄缩水 \\
pre-commit 的 \texttt{rev} & \texttt{.pre-commit-\allowbreak config.yaml} & 本地 git hook & 本地放行、CI 阻断 \\
compdb 里的编译器路径 & preset 的 \texttt{environment} 或 \texttt{toolchain\allowbreak File} & \CD、\Tidy、cppcheck & 宏判定不同，条件编译分支解析错 \\
CMake 版本 & preset 的 \texttt{cmake\allowbreak Minimum\allowbreak Required} & \texttt{cmake}~\Opt{preset} & 整份 presets 解析失败 \\
\end{longtable}
}

\begin{lstlisting}[style=yamlstyle,caption={用 apt.llvm.org 固定 LLVM 版本}]
# Install one release once, so the four tools never disagree.
- name: Install LLVM 19
  run: |
    sudo add-apt-repository "$LLVM_APT"
    sudo apt-get update -q
    sudo apt-get install -y --no-install-recommends \
      clang-19 clangd-19 clang-format-19 clang-tidy-19
    sudo update-alternatives --install \
      /usr/bin/clangd clangd /usr/bin/clangd-19 100
    clangd --version
    clang-format --version
    clang-tidy --version
\end{lstlisting}

\begin{warning}[title={VS 生成器下 compdb 缺失导致的静默降级}]
Windows 上团队往往使用 \texttt{Visual Studio 17 2022} 生成器，此时即使 preset 里写了 \texttt{CMAKE\_EXPORT\_COMPILE\_COMMANDS=ON} 也不会生成 compdb。后果分三层：\CD\ 用 fallback command 打开文件，跨文件跳转与「查找引用」大量失效；\Tidy\ 报 \texttt{No compilation database found} 并需要 \texttt{-p} 或手敲 \Opt{compile-commands-dir}；cppcheck 的 \Opt{project} 直接失败。对策是在 Windows preset 里额外提供一个 Ninja 的「分析专用」configure preset，只用于产出 compdb，不参与真实构建。
\end{warning}

\section{GitHub Actions：把诊断变成行内批注}\label{sec:ch16-s9}

\begin{lstlisting}[style=yamlstyle,caption={quality-gate 工作流：矩阵、缓存与行内批注},label={lst:ch16-actions}]
name: quality-gate
on:
  pull_request:
  schedule:
    - cron: "0 3 * * *"

env:
  LLVM_APT: >-
    deb http://apt.llvm.org/noble/
    llvm-toolchain-noble-19 main

jobs:
  gate:
    runs-on: ${{ matrix.os }}
    strategy:
      fail-fast: false
      matrix:
        include:
          - os: ubuntu-24.04
            preset: ci-linux
          - os: windows-latest
            preset: ci-msvc
          - os: macos-14
            preset: ci-macos
    defaults:
      run:
        shell: bash
    steps:
      - uses: actions/checkout@v4
        with:
          fetch-depth: 0        # merge-base needs full history
      - name: Install LLVM
        if: runner.os == "Linux"
        run: |
          sudo apt-get install -y clang-tidy-19 \
            clang-format-19
      - name: Cache compdb and ccache
        uses: actions/cache@v4
        with:
          path: |
            build/${{ matrix.preset }}
            ~/.cache/ccache
          key: >-
            ${{ matrix.preset }}-${{ hashFiles('CMake*') }}
      - name: Quality gate
        run: ./scripts/verify.sh ${{ matrix.preset }}
      - name: Upload tidy fixes
        if: always()
        uses: actions/upload-artifact@v4
        with:
          name: fixes-${{ matrix.preset }}
          path: build/fixes/*.yml
      - name: Inline annotations
        if: always()
        env:
          REVIEWDOG_GITHUB_API_TOKEN: ${{ secrets.GITHUB_TOKEN }}
        run: |
          clang-tidy-diff.py -p1 -quiet -path build \
            | reviewdog -f=clang-tidy -reporter=github-pr-review
\end{lstlisting}

四个要点。其一，\texttt{fetch-depth: 0} 不是可选项：\texttt{clang-tidy-diff.py} 与第 10 章的增量格式化都要算 \texttt{merge-base}，浅克隆会让它们把「整个仓库」当成改动集，从而在 PR 上刷出上千条批注。其二，缓存粒度：\texttt{build/\$\{matrix.preset\}} 目录同时含 compdb 与对象文件，而 \texttt{hashFiles} 只列 \texttt{CMakeLists.txt} 与 \texttt{CMakePresets.json}，于是新增源文件不会使缓存失效、改编译选项则会立刻失效。其三，\Opt{export-fixes} 的产物必须上传，夜间任务的自动修复 PR（第 14 章）消费的就是这些 YAML。其四，行内批注走 reviewdog 而不是把日志贴在 PR 正文里；只有作者不是本人时才贴评论，避免机器人自我评论刷屏。

\begin{lstlisting}[style=yamlstyle,caption={GitLab CI：只在相关文件改动时跑质量门}]
quality-gate:
  image: acme/llvm-19:latest
  rules:
    - changes:
        - "src/**/*.cpp"
        - "include/**/*.h"
        - "CMakePresets.json"
        - ".clang-tidy"
      when: on_success
    - when: never
  script:
    - cmake --preset ci-linux
    - ./scripts/verify.sh ci-linux
  artifacts:
    when: always
    paths: ["build/ci-linux/fixes/"]
\end{lstlisting}

\section{门禁演进与反退化机制}\label{sec:ch16-s10}

门禁上线半年后最常见的结局不是失效，而是被逐个关掉。对抗退化的办法是把「阻断」拆成阶段，并给存量问题一个会自己缩小的出口。

{\footnotesize
\begin{longtable}{@{}>{\raggedright\arraybackslash}p{0.15\textwidth}
                  >{\raggedright\arraybackslash}p{0.28\textwidth}
                  >{\raggedright\arraybackslash}p{0.13\textwidth}
                  >{\raggedright\arraybackslash}p{0.25\textwidth}@{}}
\caption{五个阶段的工具、时长与失败动作}\label{tab:ch16-stage}\\
\toprule
阶段 & 工具与调用 & 允许时长 & 失败动作 \\
\midrule
\endfirsthead
\toprule
阶段 & 工具与调用 & 允许时长 & 失败动作 \\
\midrule
\endhead
\bottomrule
\endlastfoot
编辑器保存 & \CD\ 诊断、\Fmt\ 按文件格式化 & 亚秒 & 仅提示，不阻断 \\
pre-commit & \Fmt\ 的 \Opt{dry-run}、少量 \Tidy\ check & 10 s 以内 & 阻断提交 \\
PR 门 & \ref{lst:ch16-verify}：build、增量 \Tidy、cppcheck、ctest & 5--10 min & 阻断合入 \\
夜间 & \texttt{run-\allowbreak clang-\allowbreak tidy} 全量、sanitizer 长测、覆盖率 & 1--2 h & 开自动修复 PR 或告警 \\
发布前 & 全量 \Tidy\ 基线核对、ASan 与 TSan 回归包 & 30--60 min & 阻断发布 \\
\end{longtable}
}

\subsection{基线：新增零告警，存量只减不增}

不要指望第一天就全绿。可行的做法是让脚本自己生成 suppression：首轮全量跑 \Tidy\ 与 cppcheck，把结果按「文件、check 名、计数」写成基线文件并提交；之后每次门禁把新结果与基线比对，\textbf{超出基线的告警阻断，少于基线的告警自动收紧基线}。cppcheck 有现成的 \Opt{suppress=} 与 \Opt{suppressions-list}；\Tidy\ 侧可以用一个 checks 减法清单达到同样效果。关键点是基线文件的 diff 只会变小，因此「上线前先把检查关掉」这条路会被 review 拦住：关掉一项检查会让基线数字凭空下降，必须由人在 PR 描述里解释。

\subsection{指标与阈值}

三个指标足够描述门禁的健康度：每千行告警数（趋势应当单调下降）、告警修复合入时长的中位数（超过两周说明存量在硬化）、格式化 churn（一次格式化提交触及的行数，突然放大意味着风格文件或 LLVM 版本漂移）。阈值设定原则是「先观测两个月再阻断」：任何新检查项先进入仅报告这一列，稳定后再升到阻断，避免一次全仓扫出的噪声让人把整个门禁静音。

\begin{guideline}[title={门禁的四条不变量}]
\begin{itemize}
  \item 入口唯一：本地与 CI 调同一个脚本，preset 名一致。
  \item 新增零告警：改动行上的新告警必须修，或写成带理由的 \texttt{NOLINT}（第 13 章）。
  \item 版本同源：\Fmt、\Tidy、\CD\ 来自同一个 LLVM release，且镜像 tag 写死。
  \item 存量递减：基线文件每轮只允许变小；变大必须在 PR 里显式说明。
\end{itemize}
\end{guideline}

\section{本章小结}\label{sec:ch16-s11}

CMakePresets 解决「进入正确状态」的成本：一条 \texttt{cmake}~\Opt{preset} 决定生成器、语言标准、警告集、sanitizer 与 compdb 开关，\texttt{condition} 让多平台共用同一份文件。质量门解决「保持这个状态」：一个脚本、五个阶段、基线只减不增。剩下的问题是「出问题怎么办」，那正是第 17 章的手册式清单要回答的。

% ---- frag_ch17.tex ----
\chapter{选型与排障手册}\label{ch:ch17}

本章是全书的「工具页」：上半部分回答「我的项目该配哪几样」，下半部分按症状回答「现在坏了，先敲哪条命令」。它不引入新机制，只把第 1 章到第 16 章的结论压缩成可执行的判据。动手之前先记住下面这条五步顺序，它比任何单条技巧都值钱。

\begin{note}[title={先复现，再改配置}]
排障的第一条纪律是把问题压缩成\textbf{一条能在本机重复的命令}。对 \CD、\Fmt、\Tidy\ 来说，这条命令几乎都是显式的命令行调用，而不是「在编辑器里按快捷键」。编辑器状态、插件缓存、残留的旧进程都会掩盖真相，只有命令行输出才是可比的证据。改配置之前先拿到它的输出，改完之后原样再跑一次，对比才有意义。
\end{note}

\begin{guideline}[title={五步排障顺序}]
\begin{enumerate}
  \item \textbf{compdb}：\texttt{build/} 下有没有 \texttt{compile\_\allowbreak commands.json}，条目里的文件路径与编译器是否命中当前文件。
  \item \textbf{命令行}：工具实际拿到的编译命令是什么。\CD\ 用 \Opt{check} 会原样打印它选中的命令；\Tidy\ 用 \Opt{list-checks} 能确认到底开了哪些检查；\Fmt\ 用 \Opt{dump-config} 打印最终生效的风格。
  \item \textbf{版本}：三个工具的 \Opt{version} 是否同一个 LLVM release，以及它们是否比配置文件里的字段更老。
  \item \textbf{索引}：跨文件能力（引用、实现跳转、全局改名、符号搜索）是否依赖尚未跑完或已被排除的后台索引。
  \item \textbf{查找顺序}：显式的 \Opt{style}/\Opt{config} 参数、上级目录里的同名文件、全局配置、内置默认值，哪一层赢了。
\end{enumerate}
\end{guideline}

\section{选型：四档配置与最小文件清单}\label{sec:ch17-s1}

选型只取决于三件事：有多少人同时改这个仓库、是否开源（风格会不会被外部 PR 打乱）、是否有交叉编译或超大仓。按这三件事分四档，每档给出最小可用文件清单，不要一次配满。

\begin{itemize}
  \item \textbf{档 1，单人小项目}（闭源、两万行以内）：只要 \texttt{.clang-format} 加 \CD\ 的默认行为。落地清单：仓库根一份 \texttt{.clang-format}（选定某个风格族）、\texttt{.gitignore} 里加 \texttt{/build/}、可选 \texttt{.clangd} 写一行 \texttt{CompileFlags} 指向 compdb 目录。目标是「打开就能补全，保存就变形」。
  \item \textbf{档 2，开源库或有外部贡献者}：在档 1 之上加 \texttt{.clang-tidy}（保守检查集，见第 12 章）、\texttt{.pre-commit-config.yaml}（\Fmt\ 的 \Opt{dry-run} 加少量 \Tidy\ 检查）、CI 上的增量门禁（第 10 章）。额外落地：贡献指南里一句「先跑 \texttt{pre-commit install}」。
  \item \textbf{档 3，多人团队}：引入 \texttt{CMake\-\allowbreak Presets.json} 与 \texttt{scripts/verify.sh}（第 16 章），把 preset 名当作团队契约；\texttt{CMakeUserPresets.json} 进 \texttt{.gitignore}。此档的分水岭不是工具数量，而是「只有一个入口脚本」。
  \item \textbf{档 4，超大仓、交叉编译或多工具链}：在档 3 之上增加外部索引与符号服务（第 3 章）、compdb 生成规范（禁止手写，只允许从构建系统导出）、统一容器镜像与 \Opt{query-driver} 白名单、按目录拆分的 \texttt{Index.External} 与 \texttt{HeaderFilterRegex}。
\end{itemize}

{\footnotesize
\begin{longtable}{@{}>{\raggedright\arraybackslash}p{0.20\textwidth}
                  >{\raggedright\arraybackslash}p{0.32\textwidth}
                  >{\raggedright\arraybackslash}p{0.30\textwidth}@{}}
\caption{四档配置的判据与代价}\label{tab:ch17-pick}\\
\toprule
判据 & 该引入什么 & 主要代价 \\
\midrule
\endfirsthead
\toprule
判据 & 该引入什么 & 主要代价 \\
\midrule
\endhead
\bottomrule
\endlastfoot
只有自己在改 & \texttt{.clang-format} 加 \CD & 几乎为零 \\
会被外部 PR 打乱风格 & \texttt{.clang-format}、pre-commit、\Fmt\ 的 \Opt{dry-run} & 首次格式化 churn 大，需独立提交 \\
想让 reviewer 少写重复意见 & \texttt{.clang-tidy} 保守集，只开命名与现代性族 & 要记 check 名；误报需逐条写 \texttt{NOLINT} \\
多人共用构建目录 & \texttt{CMakePresets.json} 与入口脚本 & preset 漂移需要维护 \\
交叉编译或工具链非默认 & \texttt{toolchainFile} 加 query-driver 白名单 & 镜像与本机路径不一致时排查耗时 \\
仓大于百万行 & 外部索引、按目录过滤、夜间全量 & 索引存储与首轮耗时 \\
\end{longtable}
}

\section{\CD\ 无补全或大量标红}\label{sec:ch17-s2}

\textbf{症状}：新打开一个源文件后，标准库有补全、项目自定义类型全部标红；或者 \texttt{goToDefinition} 只在当前文件内有效。

\textbf{最快定位}：

\begin{lstlisting}[style=shellstyle,caption={用 clangd --check 一次看清四层状态}]
clangd --check=src/widget.cpp \
  --compile-commands-dir=build/dev --log=verbose
\end{lstlisting}

输出会依次回答四件事：是否找到 compdb、命中了哪条命令（\texttt{Compile command:} 那一行）、preamble 是否建成、索引是否加载。四种根因对应四种读法。

\begin{itemize}
  \item \textbf{没找到 compdb}：日志出现 \texttt{Failed to find compilation database}。此时 \CD\ 退回 fallback command，只带语言标准与默认搜索路径，项目宏全部缺失。对策是从构建系统导出 compdb（第 2 章），或在 \texttt{.clangd} 的 \texttt{CompileFlags} 段里指定数据库目录。
  \item \textbf{找到了但条目没命中}：文件路径的分隔符或大小写不一致（Windows 高发），或者头文件本身没有条目。\CD\ 会尝试借用同目录源文件的命令，目录约定不成立时就退回 fallback。
  \item \textbf{命令被改坏了}：外部配置用 \texttt{Add}/\texttt{Remove} 动过编译选项，删掉了必需的 \texttt{-I} 或把 \texttt{-std} 改低。把 \Opt{check} 打印的最终命令行与真实构建行逐字对比。
  \item \textbf{编译器不可执行}：开了 \Opt{query-driver} 但路径写错，或驱动不在 \texttt{PATH} 里。后果是系统头与内建宏都取不到，表现为 \texttt{stddef.h not found} 这类基础错误。
\end{itemize}

\textbf{预防}：把 \texttt{clangd}~\Opt{check} 加真实文件路径写成新人手册的第一步，它比任何截图都有说服力。

\section{索引过期、引用结果不全}\label{sec:ch17-s3}

\textbf{症状}：同一份代码，昨天能查到 12 处引用，今天只有 3 处；或者改名后漏掉两个文件。

\textbf{根因}：跨 TU 的能力依赖\textbf{后台索引}，而索引是本机缓存，不进版本控制。缓存目录是与构建目录相邻的 \texttt{.cache/clangd/index}，由 \Opt{background-index}（或配置里的 \texttt{Index} 段）生成。只要索引没跑完，结果就是不全的，这不是缺陷而是设计。

\textbf{对策}：先看日志里 \texttt{Indexed} 行的进度与耗时；确认没有用 \Opt{no-background-index} 关掉它；确认目标目录没有被 \texttt{Index.External} 或黑名单路径排除。缓存被清理脚本带走时，重建成本等于首轮全量索引，第 3 章给过量化。

\begin{warning}[title={后台索引不是数据库}]
\CD\ 的索引是\textbf{尽力而为的缓存}：允许过期、允许被删、允许只覆盖部分 TU。它不能替代 compdb，更不能当作 CI 里的真相来源。任何在 CI 脚本里读取 \texttt{.cache/clangd} 的做法都是错的——那条路径只在编辑器所在机器上存在，内容也不保证完整。CI 里需要跨文件信息时，请直接用 compdb 驱动 \Tidy\ 或第 14 章的重构工具，而不是去猜索引缓存。
\end{warning}

\section{跳到声明而非定义，或跳进第三方头}\label{sec:ch17-s4}

\textbf{症状}：在调用点执行 \texttt{goToDefinition} 落到 \texttt{.h} 里的声明，而不是 \texttt{.cpp} 里的实现；或者一路跳进包管理器下载的第三方头。

\textbf{解释}：声明与实现分离时，「定义」在另一个 TU 里，只看当前 TU 是拿不到它的，必须依赖跨 TU 索引；索引未就绪时 \CD\ 只能给声明。网上流传的 \texttt{declarationAsRHS} 之类键名请谨慎照抄：\CD\ 不认识的键会在日志里记为 unknown key，虽然不会让整份配置失效，但会让人忽略「真正生效的那个键其实写错了」。

\textbf{对策}：先确认 \Opt{background-index} 生效（\ref{sec:ch17-s3}）；声明处可用 \texttt{goToImplementation} 或 code action 里的「转到定义」；第三方头目录要么排除在索引之外，要么反过来用外部索引预生成，避免每次跳进 vendored 代码。

\section{格式化没生效}\label{sec:ch17-s5}

\textbf{症状}：保存后代码纹丝不动，或者只动了缩进没动换行。按下面顺序排除，这四条占九成命中率。

\begin{enumerate}
  \item \textbf{风格文件的查找目录}：\Fmt\ 的 \Opt{style=file} 从\textbf{被格式化文件所在目录}逐级向上找 \texttt{.clang-format}。子目录里若有历史遗留的一份，赢的是那份更近的。用 \Opt{style=file} 配合绝对路径能立刻区分这两种情形。
  \item \textbf{多份风格文件共存}：先列出仓库里所有同名文件。多份共存时「改了没生效」是必然结果。
  \item \textbf{语言分派}：\texttt{Language: Cpp} 只覆盖 \CPP\ 源文件；头文件被当成哪种语言、以及通配段是否被写坏，都会让某一段整体失效。用 \texttt{If} 条件段与 \texttt{Disable} 逐段排查。
  \item \textbf{编辑器把活儿抢走了}：很多环境里格式化由另一个扩展接管（微软系插件或团队自研），它读的是另一套配置。确认「谁执行格式化」是这一节最重要的一步，它解释了「命令行是对的、编辑器里是错的」这种最费时间的组合。
\end{enumerate}

\section{\Tidy\ 一条都不报}\label{sec:ch17-s6}

\textbf{症状}：命令返回码为 0，屏幕干干净净，但你知道代码里有问题。

\begin{itemize}
  \item \textbf{checks 打底写错}：\texttt{Checks} 是按顺序生效的逗号列表，先包含再排除。如果只写了排除项而前面没有包含任何东西，结果就是零检查。拼错的 check 名只会多一行警告，不会失败，很容易看漏。
  \item \textbf{没有 compdb}：\Tidy\ 找不到数据库时打印 \texttt{No compilation database found} 并拒绝分析。用 \texttt{-p build} 或 \Opt{compile-commands-dir} 指定；注意默认是\textbf{从当前工作目录向上}查找，在子目录里跑就会找不到。
  \item \textbf{HeaderFilterRegex 反而排掉了项目头}：默认只显示「被当作 TU 打开的那个文件」里的诊断，头文件中的诊断必须被 \texttt{HeaderFilterRegex} 匹配到才可见。写成通配一切会把第三方头一起报出来，通常应写成仓库内路径的模式。
  \item \textbf{语言标准低于检查要求}：只在 \cpp{17} 及以上成立的现代性检查，在 \texttt{-std=c++11} 下静默不报。
  \item \textbf{观感被抑制干扰}：\Opt{quiet} 只打印诊断本身；\texttt{WarningsAsErrors} 会改变退出码。两者都会让人误判「到底有没有跑起来」。
\end{itemize}

\textbf{定位命令}：先在一个已知有问题的最小文件上单跑一次（附录 B 第一组命令），确认检查集真的生效，再回到项目规模。

\section{\texttt{NOLINT} 不起作用}\label{sec:ch17-s7}

\textbf{规则}：抑制注释必须出现在\textbf{被抑制的那一行}上——诊断在第 40 行，\texttt{// NOLINT} 就得在第 40 行。\texttt{// NOLINTNEXTLINE} 加括号内的 check 名只作用于\textbf{紧接的下一行}，中间夹一个空行就失效。\texttt{NOLINTBEGIN} 与 \texttt{NOLINTEND} 构成区间，必须成对且在同一文件内。带 check 名的抑制只对该 check 有效，写成不存在的名字等于什么都没抑制。

\textbf{宏展开行}：诊断落在宏展开结果上时，源码里可能找不到对应行（展开位置取决于诊断类型，可能指向调用点也可能指向定义处）。此时把注释写在\textbf{宏调用行}，或用区间抑制包住整个调用；写在宏体里往往无效。

\begin{warning}[title={NOLINT 不是文档}]
没有理由说明的 \texttt{NOLINT}，半年后没人记得为什么。请写成 \texttt{// NOLINTNEXTLINE(\allowbreak readability-\allowbreak magic-numbers)} 并在后面补一句原因，同时在门禁里限制单个文件的 \texttt{NOLINT} 数量（第 13 章）。另有一个高频事故：\texttt{NOLINT} 独占一行时会被 \Fmt\ 挪位置，抑制随之失效。
\end{warning}

\section{\texttt{-fix} 改坏代码与修复不幂等}\label{sec:ch17-s8}

\textbf{症状}：加了 \Opt{fix} 之后编译不过，或者同一条命令跑两遍产生两次改动。

\textbf{根因}：其一，\Tidy\ 的 autofix 覆盖率并不完整，许多 check 只报不改；当修复只覆盖表达式的一部分时，结果可能不完整。其二，同一 check 的多个修复落在重叠区间时会被丢弃，要到第二轮才补上，表现就是「不幂等」。其三，不同 check 的修复互相冲突，例如排版型与命名型同时改一行。

\textbf{对策}：只在干净的工作树上跑 \Opt{fix}；先导出 YAML 再落盘：

\begin{lstlisting}[style=shellstyle,caption={导出替换、分目录审阅后再生效}]
clang-tidy src/widget.cpp -p build/dev \
  --checks=-,bugprone-* -export-fixes fix.yml
clang-apply-replacements --dry-run ./fix-dir
\end{lstlisting}

确认无误后由 \texttt{clang-apply-replacements} 真正改写文件（第 14 章）。如果目标是批量重命名，用 \texttt{clang-rename} 而不是靠多个 fix 拼接。

\section{Windows 特有的一组坑}\label{sec:ch17-s9}

\begin{itemize}
  \item \textbf{路径分隔与大小写}：compdb 里的 \texttt{file} 与 \texttt{directory} 若混用反斜杠与正斜杠，或盘符大小写不同，\CD\ 与 \Tidy\ 都可能匹配不到条目。对策是只由 CMake 导出、绝不手写；必要时在 \texttt{.clangd} 里用 \texttt{Add} 补搜索路径。
  \item \textbf{\texttt{/D} 风格参数}：MSVC 的 \texttt{/DFOO=1}、\texttt{/std:c++latest} 是 clang 不认的形式。要么用 \Opt{query-driver} 允许实际执行 \texttt{cl.exe} 去取宏与头路径，要么在 \texttt{.clangd} 里 \texttt{Remove} 掉 \texttt{/*} 这类形式并 \texttt{Add} 等价的 \texttt{-D} 写法。
  \item \textbf{\texttt{cl.exe} 前缀}：compdb 的 \texttt{command} 若写成裸 \texttt{cl}，需要它确实能在 \texttt{PATH} 中找到。CI 上往往要先加载开发者环境脚本，这会让本地与 CI 的 query-driver 结果不同。
  \item \textbf{代码页与中文注释}：文件按 GBK 保存而工具按 UTF-8 解析时，中文注释成为非法序列，症状是诊断列号整体偏移、字符串宽度算错。对策是编辑器强制 UTF-8，并在 \texttt{.editor-\allowbreak config} 里写死。这与第 4 章、第 8 章的 UTF-16 码元偏移是同一类问题的两端：\textbf{偏移基准不一致，永远比偏移本身更危险}。
\end{itemize}

\section{版本错配}\label{sec:ch17-s10}

\textbf{症状}：同一段代码昨天格式化后没变化，今天多出一处改动；或者某个 check 时有时无。三条 \Opt{version} 连查，五分钟定位：

\begin{lstlisting}[style=shellstyle,caption={三个工具必须来自同一个 release}]
clangd --version
clang-format --version
clang-tidy --version
\end{lstlisting}

\CD\ 19 配 \Fmt\ 17 的组合很常见（前者由 IDE 扩展下载，后者由包管理器安装），后果是新版本的字段与换行策略变化会让「保存后又被别人改回去」成为日常。第 16 章的镜像对齐是根治手段；临时缓解是让编辑器扩展使用与 CI 相同的那个二进制。另一侧同样成立：\texttt{.clang-format} 里写了只有新版本认识的字段时，旧版本会忽略而不是报错，所以「本地绿、CI 红」也可能来自这一端。

\section{与既有工程习惯的冲突}\label{sec:ch17-s11}

\begin{itemize}
  \item \textbf{\texttt{-fpermissive} 一类的放宽开关}：真实构建里靠它过编译的老代码，在 \CD\ 下会以错误形式成片暴露。对策是在 \texttt{.clangd} 的 \texttt{Add} 里带上同样的开关，或在 \texttt{Diagnostics} 的 \texttt{Suppress} 列表里点名抑制，而不是去改真实构建。
  \item \textbf{MSVC 扩展语法}：\texttt{\_\_declspec}、\texttt{\_\_try} 之类需要 \texttt{-fms-extensions} 才能被 clang 解析，否则报的是语法错误而不是警告。
  \item \textbf{CUDA 的 \texttt{.cu}}：\texttt{\_\_global\_\_} 与核函数启动语法不在 clang 的 \CPP\ 前端里，直接喂给 \CD\ 会满屏标红。可行的组合是补齐 CUDA 的头搜索路径并接受「不完美」，或者让构建系统导出 compdb 后只在纯 \CPP\ 文件上开启 \Tidy。
  \item \textbf{\texttt{\_\_has\_\_include} 与构建期生成的头}：\texttt{config.h} 这类文件在首次 configure 前不存在，索引与检查都会报缺失。对策：先 configure，再打开编辑器，并把这句写进新人手册。
  \item \textbf{\texttt{\#pragma once} 与 include guard 混用}：同一个头两种防护会让 include 相关检查给出互相矛盾的结论。选定一种，并用检查集强制它。
\end{itemize}

\section{诊断文本怎么读}\label{sec:ch17-s12}

三个工具共享同一套位置前缀，读法固定：

\begin{lstlisting}[style=shellstyle,caption={一条典型诊断的部件拆解}]
src/widget.cpp:42:9: warning: unused parameter
  'scale' [clang-diagnostic-unused-parameter]
   42 |   void Draw(float scale) { }
      |             ^
note: parameter declared here
\end{lstlisting}

\begin{itemize}
  \item \textbf{位置}：\texttt{文件:行:列}，行与列都从 1 开始；列按码元计，第 4 章关于 UTF-16 偏移的结论适用于 LSP 侧的同一位置。
  \item \textbf{等级}：\texttt{error}、\texttt{warning}、\texttt{note} 三级，\texttt{note} 通常承载补充位置而不是独立问题。
  \item \textbf{来源标记}：行尾方括号里是 check 名。\Tidy\ 自己的检查是可读的族名（例如 \texttt{bugprone-\allowbreak use-after-move}）；来自 clang 的编译器诊断被包装成 \texttt{clang-diagnostic-xxx}，对应原始的 \texttt{-Wxxx} 开关。看到 \texttt{clang-diagnostic-} 前缀，就该去调警告开关，而不是在 \texttt{Checks} 里找它。
  \item \textbf{补充信息}：\texttt{note} 或 related information 行指出「根因位置」或另一个相关文件，多数误判来自只读第一行。
  \item \textbf{fix-it 提示}：\texttt{replace 'x' with 'y'} 这类行给出可直接构造的替换，\Tidy\ 的 \Opt{fix} 与 \Opt{export-fixes} 就是把它们变成补丁。
\end{itemize}

\CD\ 的侧信道是日志：\Opt{log=verbose} 会输出诊断、配置加载与索引细节，配合编辑器的输出面板看。级别从 \texttt{error} 到 \texttt{verbose} 递增，排障时 \texttt{verbose} 足够，长期运行别开。

\section{常见误区清单}\label{sec:ch17-s13}

\begin{enumerate}
  \item 把 \Fmt\ 当代码检查器：它只改空白与换行，不理解语义。
  \item 把 compdb 提交进仓库：它含绝对路径与本机细节，正确的要求是「能一键生成」而不是被版本控制。
  \item 认为后台索引等于数据库：索引是缓存，可以缺、可以过期（\ref{sec:ch17-s3}）。
  \item 混用 \Tidy\ 的 \Opt{warnings-as-errors} 与编译器的 \Opt{Werror}：前者作用于 \Tidy\ 诊断，后者作用于 clang 警告，退出码与影响范围都不同。
  \item 把 include-what-you-use 理解成「删 include 的工具」：它的输出需要人工判断间接依赖与导出符号，照单全删容易破坏下游编译。
  \item \Tidy\ 一次性全开：误报成本会让人关掉整个门禁，净收益为负（第 12 章）。
  \item 用 \texttt{NOLINT} 表达「我不想修」：没有原因的抑制等于把技术债隐藏记账。
  \item 让 CI 与本地跑不同的命令：见第 16 章的入口脚本不变量。
  \item 指望 \Opt{fix} 幂等：重叠替换会被丢弃，需要多轮或先导出再应用。
  \item 在多个目录放 \texttt{.clang-format}：查找顺序是就近优先，多份必然漂移。
  \item 认为 \CD\ 报错就等于代码错：先确认它是按真实构建的编译命令在解析。
  \item 把 sanitizer 当静态分析：它们需要真实执行路径，覆盖不到的分支永远报不出来（第 15 章）。
  \item 忽略 \texttt{clang-diagnostic-} 前缀：它说明问题属于编译器警告集，不属于某个 \Tidy\ check。
  \item 指望 code action 修一切：能否自动修取决于插件与 check 是否提供 fix，缺 fix 时它只给建议。
\end{enumerate}

\section{性能与资源类症状}\label{sec:ch17-s16}

\textbf{症状}：编辑器打开一个大文件后风扇狂转；\Tidy\ 在 PR 上跑四十分钟；cppcheck 占满内存。这类问题不该靠「加配置」解决，而该靠\textbf{减少分析量}。

\begin{itemize}
  \item \textbf{\CD\ 首轮与每次改头}：改一个被广泛包含的头会触发大量 preamble 重建，成本随包含链长度而非本次改动行数增长（第 3 章）。缓解手段是把公共头的包含关系收敛、把稳定实现移出被频繁改动的头、必要时用 \Opt{pch-storage} 调整 preamble 的存放位置以降低峰值内存。
  \item \textbf{编辑期跑重检查}：\Tidy\ 集成的路径敏感族（\texttt{clang-analyzer-\allowbreak*}）在每次按键后是不可接受的。用 \texttt{ClangTidy} 段的 \texttt{FastCheckFilter} 让编辑期只跑轻量族，重检查留给 PR 与夜间。
  \item \textbf{\Tidy\ 全量慢}：绝大多数耗时来自三处——过宽的 \texttt{HeaderFilterRegex}（同一头被反复分析）、开启全部 analyzer 族、以及少量巨型函数与深度模板。按第 12 章裁剪检查集后，PR 门只跑改动行（第 10 章），耗时能降一个数量级。
  \item \textbf{cppcheck 慢}：不要用 \texttt{--force} 让它对每个 TU 重复分析同一头文件；用项目文件（compdb）配合 \texttt{-j} 与抑制清单。
  \item \textbf{include-what-you-use 慢}：它要求完整预处理且必须逐 TU 运行，慢是常态；正确用法是只在夜间任务里跑，并把输出当建议而非补丁。
\end{itemize}

\section{排障总表}\label{sec:ch17-s14}

{\footnotesize
\begin{longtable}{@{}>{\raggedright\arraybackslash}p{0.21\textwidth}
                  >{\raggedright\arraybackslash}p{0.21\textwidth}
                  >{\raggedright\arraybackslash}p{0.20\textwidth}
                  >{\raggedright\arraybackslash}p{0.20\textwidth}@{}}
\caption{症状、根因、入口命令与预防}\label{tab:ch17-all}\\
\toprule
症状 & 根因 & 入口命令 & 预防 \\
\midrule
\endfirsthead
\toprule
症状 & 根因 & 入口命令 & 预防 \\
\midrule
\endhead
\bottomrule
\endlastfoot
项目符号全标红而标准库正常 & 没有 compdb，退回 fallback command & \texttt{clangd}~\Opt{check} & 构建后立即导出 compdb \\
成员函数跳不到实现 & 后台索引未完成或该 TU 被排除 & 看 \texttt{Indexed} 日志行 & 首轮等索引跑完 \\
引用结果比同事少 & 本机索引缓存缺失 & 重开 \Opt{background-\allowbreak index} & 别清理 \texttt{.cache/} \\
\texttt{stddef.h} 找不到 & query-driver 指向的驱动不可执行 & \Opt{query-driver} 用绝对路径 & 镜像内固定工具链 \\
改了风格文件没效果 & 就近目录里还有另一份 & 递归查找同名文件 & 只允许仓库根一份 \\
保存后风格又被改回 & 两个扩展抢格式化 & 命令行单独跑一次对比 & 明确谁负责格式化 \\
\Tidy\ 零输出 & checks 打底写错或没有 compdb & \Opt{list-checks} 核对检查集 & 保守集加显式减项 \\
头文件里的告警看不见 & \texttt{HeaderFilterRegex} 不匹配 & \Opt{header-filter=} 试一次 & 写成仓库内路径模式 \\
\texttt{NOLINT} 无效 & 注释行位置或 check 名写错 & 按诊断行列号对齐 & 统一用 \texttt{NOLINT-\allowbreak NEXTLINE} \\
\Opt{fix} 后编译失败 & 修复只覆盖表达式片段 & 先导出 YAML 审阅 & 干净工作树才跑 fix \\
同一命令两遍结果不同 & 重叠替换被丢弃 & 分两轮或按族分次跑 & 修复独立成一个提交 \\
Windows 上条目匹配不到 & 分隔符与盘符大小写混用 & \Opt{check} 看命中的命令行 & 只由 CMake 导出 \\
中文注释处列号怪异 & 文件不是 UTF-8 & 查文件编码 & 配置里写死 UTF-8 \\
某检查时有时无 & 语言标准低于其要求 & 看 \texttt{-std} 实参 & preset 固化标准 \\
本地绿而 CI 红 & 三方版本不同源 & 三条 \Opt{version} 连查 & 镜像、rev、本机同源 \\
诊断刷屏且都在第三方 & 没有排除第三方头 & 配置里 \texttt{Suppress} & 索引与检查都限定仓内 \\
\end{longtable}
}

\section{本章小结}\label{sec:ch17-s15}

选型的答案是「按人数与开源程度分四档，每档只补一样东西」。排障的答案是「按 compdb、命令行、版本、索引、查找顺序五步走，每一步都有一条命令可以作证」。把这两句话贴在团队 wiki 上，多数问题能在十分钟内定位。附录 A、B、C 提供逐字段与逐命令的速查。

% ---- frag_app.tex ----
\appendix

\chapter{配置字段速查}\label{ch:appA}

本附录把三类配置文件的主要字段集中列出，取值一律以对应工具的官方文档为准：同一字段在不同 LLVM 版本里可能新增枚举值，字段名也可能被拆分或重命名。表中的「典型取值」指实际会写进文件的东西，不是完整取值域。所有表都按字母序排列，方便直接查找。

\section{\texttt{.clang-format} 字段速查}\label{sec:appA-s1}

字段作用于由 \texttt{Language} 段决定的语言；带 \texttt{If} 的条件段可以覆写下面任何一项。

{\footnotesize
\begin{longtable}{@{}>{\raggedright\arraybackslash}p{0.30\textwidth}
                  >{\raggedright\arraybackslash}p{0.34\textwidth}
                  >{\raggedright\arraybackslash}p{0.18\textwidth}@{}}
\caption{\texttt{.clang-format} 常用字段}\label{tab:appA-fmt}\\
\toprule
字段 & 作用 & 典型取值 \\
\midrule
\endfirsthead
\toprule
字段 & 作用 & 典型取值 \\
\midrule
\endhead
\bottomrule
\endlastfoot
\texttt{BasedOnStyle} & 继承一个基础风格族，其余字段在其上覆写 & \texttt{LLVM}、\texttt{Google}、\texttt{Mozilla} \\
\texttt{Language} & 本段适用的语言；\texttt{None} 表示无条件段 & \texttt{Cpp}、\texttt{None} \\
\texttt{Disable} & 在 \texttt{If} 段里整体关闭该段的格式化 & \texttt{true} \\
\texttt{AccessModifierOffset} & 访问修饰符相对类体的缩进量 & \texttt{-4}、\texttt{0} \\
\texttt{AlignAfterOpenBracket} & 左括号后换行时的对齐方式 & \texttt{Align}、\texttt{DontAlign} \\
\texttt{AlignArrayOfStructures} & 结构体数组的成员列对齐 & \texttt{Left}、\texttt{Right} \\
\texttt{AlignConsecutive-\allowbreak Assignments} & 连续赋值语句按等号对齐 & \texttt{None}、\texttt{Left} \\
\texttt{AlignConsecutive-\allowbreak Macros} & 连续宏定义按值对齐 & \texttt{None}、\texttt{Left} \\
\texttt{AlignConsecutive-\allowbreak Declarations} & 连续声明按名字对齐 & \texttt{None}、\texttt{Left} \\
\texttt{AlignEscapedNewlines} & 反斜杠续行的位置 & \texttt{Left}、\texttt{Right} \\
\texttt{AlignOperands} & 二元运算符左右操作数对齐 & \texttt{true}、\texttt{false} \\
\texttt{AlignTrailingComments} & 行尾注释列对齐 & \texttt{true} \\
\texttt{AllowAllParametersOf-\allowbreak DeclarationOnNextLine} & 允许参数全部放下一行 & \texttt{true} \\
\texttt{AllowShortBlocks\allowbreak OnASingleLine} & 允许短块压成一行 & \texttt{Never}、\texttt{Empty} \\
\texttt{AllowShortCaseLabels-\allowbreak OnASingleLine} & 允许短 case 与语句同行 & \texttt{false}、\texttt{true} \\
\texttt{AllowShortEnums\allowbreak OnASingleLine} & 允许短枚举压成一行 & \texttt{true} \\
\texttt{AllowShortFunctions-\allowbreak OnASingleLine} & 允许短函数体同行 & \texttt{All}、\texttt{Empty} \\
\texttt{AllowShortIfStatements-\allowbreak OnASingleLine} & 允许短 if 与语句同行 & \texttt{Never}、\texttt{All} \\
\texttt{AllowShortLambdas-\allowbreak OnASingleLine} & 允许短 lambda 压成一行 & \texttt{None}、\texttt{All} \\
\texttt{AllowShortLoops\allowbreak OnASingleLine} & 允许短循环体同行 & \texttt{false} \\
\texttt{AlwaysBreakBefore-\allowbreak MultilineStrings} & 多行字符串前的换行策略 & \texttt{true} \\
\texttt{BinPackArguments} & 实参列表是否允许紧凑填充 & \texttt{true}、\texttt{false} \\
\texttt{BinPackParameters} & 形参列表是否允许紧凑填充 & \texttt{true}、\texttt{false} \\
\texttt{BreakBeforeBraces} & 大括号换行策略的总开关 & \texttt{Attach}、\texttt{Allman} \\
\texttt{BraceWrapping} & 逐场景覆写 \texttt{Custom} 下的括号位置 & 子对象，见下 \\
\texttt{BreakAfterAttributes} & 属性与声明之间的换行 & \texttt{Never}、\texttt{Always} \\
\texttt{BreakBefore\allowbreak BinaryOperators} & 二元运算符前置换行 & \texttt{None}、\texttt{All} \\
\texttt{BreakBeforeConcept-\allowbreak Body} & concept 定义的 \texttt{\{} 前是否换行 & \texttt{true} \\
\texttt{BreakBefore\allowbreak InheritanceComma}? & 已废弃，改用 \texttt{BreakInheritanceList} & —— \\
\texttt{BreakBefore\allowbreak TernaryOperators} & 三目运算符前置换行 & \texttt{true} \\
\texttt{BreakConstructor\allowbreak Initializers} & 构造初始化列表的断行位置 & \texttt{BeforeColon}、\texttt{AfterColon} \\
\texttt{BreakInheritanceList} & 继承列表的断行位置 & \texttt{BeforeColon} \\
\texttt{BreakStringLiterals} & 是否拆分超长字符串字面量 & \texttt{true}、\texttt{false} \\
\texttt{ColumnLimit} & 软列宽上限，0 表示不限 & \texttt{80}、\texttt{100}、\texttt{0} \\
\texttt{CommentPragmas} & 被视为普通文本保留的注释前缀正则 & \texttt{doxygen} 相关写法 \\
\texttt{ConstructorInitializer-\allowbreak AllOnOneLineOrOnePerLine} & 初始化列表一行放不下就每行一项 & \texttt{true} \\
\texttt{Cpp11BracedListStyle} & 花括号列表按函数调用式排版 & \texttt{true} \\
\texttt{DeriveLineEnding} & 从多数行推断行尾符（CRLF 或 LF） & \texttt{true} \\
\texttt{DerivePointerAlignment} & 从文件里推断指针星号位置 & \texttt{false} 更可控 \\
\texttt{EmptyLineBefore-\allowbreak AccessModifier} & 访问修饰符前保留的空行 & \texttt{Never}、\texttt{Always} \\
\texttt{FixNamespaceComments} & 自动补全闭合格子的注释 & \texttt{true} \\
\texttt{IndentCaseLabels} & case 标签相对 switch 再缩进 & \texttt{true}、\texttt{false} \\
\texttt{IndentAccessModifiers} & 访问修饰符是否额外缩进 & \texttt{false} \\
\texttt{IndentPPDirectives} & 预处理指令的缩进策略 & \texttt{None}、\texttt{AfterHash} \\
\texttt{IndentRequires} & concept 的 \texttt{requires} 是否缩进 & \texttt{true} \\
\texttt{IndentWidth} & 一级缩进的空格数 & \texttt{2}、\texttt{4} \\
\texttt{KeepEmptyLinesAtThe-\allowbreak StartOfBlocks} & 块首空行是否保留 & \texttt{false} \\
\texttt{LambdaBodyExpressions}? & 较新版本才有的 lambda 排版控制 & 见对应版本文档 \\
\texttt{MaxEmptyLinesToKeep} & 连续空行的保留上限 & \texttt{1}、\texttt{2} \\
\texttt{NamespaceIndentation} & namespace 体的缩进策略 & \texttt{None}、\texttt{All} \\
\texttt{PointerAlignment} & 星号与 \texttt{\&} 靠谁 & \texttt{Left}、\texttt{Right} \\
\texttt{RawStringDelimiters-\allowbreak Patterns} & 识别原生字符串定界符的正则列表 & 正则列表 \\
\texttt{ReflowQuotedArrays}? & 较新版本才有的引号数组重排 & 见对应版本文档 \\
\texttt{ReferenceAlignment}? & 较新版本才有的引用独立对齐 & 见对应版本文档 \\
\texttt{RequiresClausePosition} & \texttt{requires} 子句放在哪一行的位置 & \texttt{OwnLine}、\texttt{SameLine} \\
\texttt{SeparateDefinitionBlocks} & 定义块之间插入空行 & \texttt{Leave}、\texttt{Insert} \\
\texttt{SortIncludes} & 是否排序 include 块 & \texttt{CaseSensitive} \\
\texttt{SortUsingDeclarations} & 是否排序 using 声明 & \texttt{true} \\
\texttt{SpaceAfterCStyleCast} & C 风格转换后是否加空格 & \texttt{false} \\
\texttt{SpaceAfterLogicalNot} & 逻辑非后是否加空格 & \texttt{false} \\
\texttt{SpaceBeforeCaseColon}? & 较新版本才有的 case 冒号前空格 & 见对应版本文档 \\
\texttt{SpaceBeforeCpp11-\allowbreak BracedList} & 花括号列表前是否加空格 & \texttt{true} \\
\texttt{SpaceBeforeParens} & 关键字与括号之间是否加空格 & \texttt{Control\allowbreak Statements} \\
\texttt{SpaceInEmptyBlock} & 空块内是否留空格 & \texttt{false} \\
\texttt{SpacesBeforeTrailing-\allowbreak Comments} & 行尾注释前的空格数 & \texttt{1}、\texttt{2} \\
\texttt{SpacesInAngles} & 模板尖括号内侧是否加空格 & \texttt{Never}、\texttt{Always} \\
\texttt{SpacesInParentheses} & 圆括号内侧是否加空格 & \texttt{false} \\
\texttt{TabWidth} & 制表符的显示宽度，用于列计算 & \texttt{4} \\
\texttt{UseTab} & 缩进使用制表符还是空格 & \texttt{Never}、\texttt{ForIndentation} \\
\end{longtable}
}

带问号的字段只在较新版本存在，写进旧版本会被忽略而不是报错（第 17 章的版本错配一节）。若需要逐字段语义，第 9 章按「空白、换行、对齐、include」四族做了完整展开。

\subsection{\texttt{BraceWrapping} 的子字段}\label{sec:appA-s2}

{\footnotesize
\begin{longtable}{@{}>{\raggedright\arraybackslash}p{0.34\textwidth}
                  >{\raggedright\arraybackslash}p{0.24\textwidth}
                  >{\raggedright\arraybackslash}p{0.24\textwidth}@{}}
\caption{\texttt{BreakBeforeBraces: Custom} 时的子开关}\label{tab:appA-brace}\\
\toprule
子字段 & 控制对象 & 常见取向 \\
\midrule
\endfirsthead
\toprule
子字段 & 控制对象 & 常见取向 \\
\midrule
\endhead
\bottomrule
\endlastfoot
\texttt{AfterCaseLabel} & case 标签后的大括号 & 与 \texttt{Attach} 一致 \\
\texttt{AfterClass} & 类与结构体 & \texttt{false} 或 \texttt{true} \\
\texttt{AfterControlStatement} & \texttt{if}/\texttt{for}/\texttt{while} & \texttt{MultiLine} 折中 \\
\texttt{AfterEnum}、\texttt{AfterStruct} & 枚举与结构体 & 常取 \texttt{true} \\
\texttt{AfterFunction} & 函数定义与声明可分别设置 & 定义 \texttt{true} \\
\texttt{AfterNamespace} & 命名空间 & \texttt{false} 居多 \\
\texttt{AfterObjCDeclaration} & Objective-C 声明 & 与类一致 \\
\texttt{AfterUnion}、\texttt{AfterExternBlock} & 联合体与 extern 块 & 常取 \texttt{true} \\
\texttt{BeforeCatch}、\texttt{BeforeElse} & \texttt{catch} 与 \texttt{else} 前换行 & Allman 取 \texttt{true} \\
\texttt{IndentBraces} & 大括号本身是否缩进 & 多为 \texttt{false} \\
\texttt{SplitEmptyFunction} & 空函数体是否拆分 & 多为 \texttt{false} \\
\texttt{SplitEmptyRecord} & 空类体是否拆分 & 多为 \texttt{false} \\
\texttt{SplitEmptyNamespace} & 空命名空间是否拆分 & 多为 \texttt{false} \\
\end{longtable}
}

\section{\texttt{.clang-tidy} 配置速查}\label{sec:appA-s3}

顶层键只有七个，其余细节全在检查选项里。较新版本支持目录级嵌套配置文件，子目录中的文件会与上级合并。

{\footnotesize
\begin{longtable}{@{}>{\raggedright\arraybackslash}p{0.26\textwidth}
                  >{\raggedright\arraybackslash}p{0.30\textwidth}
                  >{\raggedright\arraybackslash}p{0.26\textwidth}@{}}
\caption{\texttt{.clang-tidy} 的顶层键}\label{tab:appA-top}\\
\toprule
键 & 作用 & 注意点 \\
\midrule
\endfirsthead
\toprule
键 & 作用 & 注意点 \\
\midrule
\endhead
\bottomrule
\endlastfoot
\texttt{Checks} & 逗号分隔的检查集，按顺序做包含与排除 & 拼错的名字只警告不失败 \\
\texttt{CheckOptions} & 逐 check 的参数表，新版也接受 \texttt{ChecksOptions} & 名字要写全限定 check 名 \\
\texttt{HeaderFilterRegex} & 决定哪些头的诊断可见 & 是正则而不是通配符 \\
\texttt{ExcludeHeader\allowbreak FilterRegex} & 从上述集合里再减去一部分 & 常用于排除 vendored 头 \\
\texttt{User} & 记入报告的用户标识 & CI 里写死便于归因 \\
\texttt{WarningsAsErrors} & 把指定诊断升格为 error & 影响退出码与门禁 \\
\texttt{FormatStyle} & 让排版型检查顺带改空白 & 默认不格式化 \\
\end{longtable}
}

{\footnotesize
\begin{longtable}{@{}>{\raggedright\arraybackslash}p{0.36\textwidth}
                  >{\raggedright\arraybackslash}p{0.30\textwidth}
                  >{\raggedright\arraybackslash}p{0.16\textwidth}@{}}
\caption{\texttt{CheckOptions} 中最常改动的项}\label{tab:appA-opt}\\
\toprule
选项全名 & 作用 & 建议取值 \\
\midrule
\endfirsthead
\toprule
选项全名 & 作用 & 建议取值 \\
\midrule
\endhead
\bottomrule
\endlastfoot
\texttt{readability-function-\allowbreak cognitive-complexity.\allowbreak Threshold} & 认知复杂度上限，超过即报 & \texttt{15}、\texttt{25} \\
\texttt{readability-function-size.\allowbreak LineThreshold} & 函数行数上限 & \texttt{80}、\texttt{150} \\
\texttt{readability-function-size.\allowbreak StatementThreshold} & 函数语句数上限 & \texttt{120} \\
\texttt{readability-function-size.\allowbreak ParameterThreshold} & 形参个数上限 & \texttt{5}、\texttt{7} \\
\texttt{readability-magic-numbers.\allowbreak IgnoredIntegerValues} & 整数魔数白名单，分号分隔 & \texttt{0;1;2;-1;100} \\
\texttt{readability-magic-numbers.\allowbreak IgnoredFloatingPointValues} & 浮点魔数白名单 & \texttt{0.0;1.0} \\
\texttt{readability-identifier-\allowbreak length.MinimumNameLength} & 标识符最短长度 & \texttt{2}、\texttt{3} \\
\texttt{readability-identifier-\allowbreak length.IgnoredExceptions?} & 旧版曾有的豁免列表，先核对版本 & —— \\
\texttt{readability-implicit-bool-\allowbreak conversion.\allowbreak AllowIntegerConditions} & 允许整数隐式转 bool & \texttt{true} \\
\texttt{readability-implicit-bool-\allowbreak conversion.\allowbreak AllowPointerConditions} & 允许指针隐式转 bool & \texttt{true} \\
\texttt{readability-simplify-boolean-\allowbreak expr.ChainedConditionalReturn} & 是否化简链式条件 return & \texttt{false} \\
\texttt{readability-uniqueptr-delete-\allowbreak release.PreferResetCall} & 修复用 \texttt{reset()} 而非重新赋值 & \texttt{true} \\
\texttt{readability-named-parameter.\allowbreak Replacement} & 未命名参数的占位写法 & 见第 13 章 \\
\texttt{readability-\allowbreak identifier-naming.\allowbreak ClassCase} & 类型命名风格 & \texttt{CamelCase} \\
\texttt{readability-\allowbreak identifier-naming.\allowbreak FunctionCase} & 函数命名风格 & \texttt{camelBack} \\
\texttt{readability-\allowbreak identifier-naming.\allowbreak VariableCase} & 变量命名风格 & \texttt{camelBack} \\
\texttt{readability-\allowbreak identifier-naming.\allowbreak ConstantCase} & 常量命名风格 & \texttt{UPPER\_CASE} \\
\texttt{readability-\allowbreak identifier-naming.\allowbreak MacroCase} & 宏命名风格 & \texttt{UPPER\_CASE} \\
\texttt{readability-\allowbreak identifier-naming.\allowbreak NamespaceCase} & 命名空间风格 & \texttt{lower\_case} \\
\texttt{modernize-loop-convert.\allowbreak MaxCopySize} & 允许多大的按值拷贝以转范围 for & \texttt{16} \\
\texttt{modernize-loop-convert.\allowbreak MakeRange} & 是否建议改用 range 版本 & \texttt{true} \\
\texttt{modernize-use-auto.\allowbreak MinTypeNameLength} & 类型名长度阈值，低于则不建议 auto & \texttt{5}、\texttt{8} \\
\texttt{modernize-use-nullptr.\allowbreak NullMacros} & 哪些旧宏应替换为 \texttt{nullptr} & \texttt{NULL;0} \\
\texttt{modernize-use-override.\allowbreak UseNoexcept} & 补 \texttt{override} 时是否带 noexcept & \texttt{true} \\
\texttt{modernize-avoid-c-arrays.\allowbreak AllowStringConstants} & 允许字符串常量用数组形式 & \texttt{true} \\
\texttt{modernize-use-equals-default.\allowbreak IgnoreMacros} & 宏构造的成员函数不提示默认化 & \texttt{true} \\
\texttt{performance-unnecessary-\allowbreak value-param.AllowedTypes} & 允许按值传参的类型列表 & 正则列表 \\
\texttt{performance-move-const-arg.\allowbreak CheckTriviallyCopyableMove} & 平凡可拷贝类型是否也报 move & \texttt{false} \\
\texttt{bugprone-argument-comment.\allowbreak StrictMode} & 要求所有字面量参数都写注释 & \texttt{false} \\
\texttt{bugprone-assert-side-effect.\allowbreak AssertFunctions} & 被视为断言的函数名列表 & 自定义断言名 \\
\texttt{bugprone-reserved-identifier.\allowbreak AllowedIdentifiers} & 允许的保留标识符例外 & 双下划线接口名 \\
\texttt{cppcoreguidelines-pro-type-\allowbreak member-init.IgnoreArrays} & 数组成员不强制初始化 & \texttt{true} \\
\texttt{misc-include-cleaner.\allowbreak IgnoreHeaders} & 不计入「多余 include」的头列表 & 系统头、生成头 \\
\texttt{misc-include-cleaner.\allowbreak AllowedHeaders} & 允许始终保留的头列表 & 预编译头 \\
\end{longtable}
}

\section{\texttt{.clangd} 配置速查}\label{sec:appA-s4}

\CD\ 的配置分全局与项目两层：全局文件在用户配置目录的 \texttt{clangd/\allowbreak config.yaml}，项目文件是仓库（或任意上级目录）里的 \texttt{.clangd}。两者字段同名，项目文件优先级更高；同一目录还可放 \texttt{.clangd} 目录存放多份条件配置。下表的「作用域」列指该键是否可被 \texttt{If} 条件限定。

{\footnotesize
\begin{longtable}{@{}>{\raggedright\arraybackslash}p{0.28\textwidth}
                  >{\raggedright\arraybackslash}p{0.24\textwidth}
                  >{\raggedright\arraybackslash}p{0.12\textwidth}
                  >{\raggedright\arraybackslash}p{0.18\textwidth}@{}}
\caption{\CD\ 的配置键}\label{tab:appA-clangd}\\
\toprule
键 & 取值与默认 & 能否条件化 & 作用 \\
\midrule
\endfirsthead
\toprule
键 & 取值与默认 & 能否条件化 & 作用 \\
\midrule
\endhead
\bottomrule
\endlastfoot
\texttt{If} & \texttt{PathPrefix}、\texttt{FilenameRegex}、\texttt{Directory}、\texttt{Index} 及 \texttt{AnyOf}/\texttt{AllOf}/\texttt{Not} & —— & 给整段配置设定生效条件 \\
\texttt{CompileFlags.\allowbreak Compilation-\allowbreak Database} & 目录路径，默认为工作目录向上查找 & 是 & 指定 compdb 位置 \\
\texttt{CompileFlags.Add} & 字符串列表，追加到编译命令尾部 & 是 & 补标准、宏、搜索路径 \\
\texttt{CompileFlags.Remove} & 字符串列表，支持尾随 \texttt{*} 通配 & 是 & 删掉 clang 不认的参数 \\
\texttt{CompileFlags.Compiler} & 单个可执行名或路径 & 是 & 覆写驱动，配合 query-driver \\
\texttt{FallbackStyle} & 风格族名，默认 \texttt{LLVM} & 是 & 找不到风格文件时用 \\
\texttt{Diagnostics.\allowbreak Suppress} & 诊断名列表 & 是 & 屏蔽指定错误与警告 \\
\texttt{Diagnostics.\allowbreak UnusedIncludes} & \texttt{Strict}、\texttt{None} & 是 & 多余 include 提示强度 \\
\texttt{Diagnostics.\allowbreak MissingIncludes} & \texttt{Strict}、\texttt{None} & 是 & 缺失 include 提示强度 \\
\texttt{Diagnostics.\allowbreak ClangTidy.Checks} & 与 \Tidy\ 同形的检查集 & 是 & 编辑期跑哪些检查 \\
\texttt{Diagnostics.\allowbreak ClangTidy.Header-\allowbreak Filter} & 正则 & 是 & 头文件里的 \Tidy\ 诊断可见性 \\
\texttt{Diagnostics.\allowbreak ClangTidy.FastCheck-\allowbreak Filter} & \texttt{None}、\texttt{Sane}、\texttt{Strict} & 是 & 编辑期跳过昂贵检查 \\
\texttt{Index.Background} & \texttt{Build}、\texttt{Skip}、\texttt{Types} & 否（全局） & 是否写本地索引缓存 \\
\texttt{Index.\allowbreak External.\allowbreak Blacklist-\allowbreak Paths} & 路径前缀列表 & 否 & 排除不必索引的目录 \\
\texttt{InlayHints.\allowbreak Enabled} & 布尔 & 是 & 总开关 \\
\texttt{InlayHints.\allowbreak ParameterNames} & 布尔 & 是 & 实参位置名提示 \\
\texttt{InlayHints.\allowbreak DeducedTypes} & 布尔 & 是 & \texttt{auto} 推导类型提示 \\
\texttt{InlayHints.\allowbreak Designators} & 布尔 & 是 & 聚合初始化字段名提示 \\
\texttt{InlayHints.\allowbreak ConstantValues} & 布尔 & 是 & 常量取值提示 \\
\texttt{InlayHints.\allowbreak DefaultArguments} & 布尔 & 是 & 缺省实参提示 \\
\texttt{InlayHints.\allowbreak TypeNameLimit} & 整数，默认 24 & 是 & 超过该长度不提示类型 \\
\texttt{SemanticTokens.\allowbreak Disabled} & 布尔 & 否 & 关闭细粒度语义着色 \\
\texttt{Style.\allowbreak FullyQualifiedName} & 限定名规则串 & 是 & 补全里是否显示全限定名 \\
\texttt{Completion.\allowbreak AllScopes} & 布尔 & 是 & 补全是否跨作用域 \\
\texttt{Completion.\allowbreak Snippets} & \texttt{LCPO}、\texttt{Disabled} & 是 & 是否用 snippet 占位 \\
\end{longtable}
}

\section{三类配置的查找顺序与冲突}\label{sec:appA-s5}

{\footnotesize
\begin{longtable}{@{}>{\raggedright\arraybackslash}p{0.16\textwidth}
                  >{\raggedright\arraybackslash}p{0.34\textwidth}
                  >{\raggedright\arraybackslash}p{0.32\textwidth}@{}}
\caption{谁先被读到}\label{tab:appA-order}\\
\toprule
工具 & 查找顺序（由先到后） & 冲突时的赢家 \\
\midrule
\endfirsthead
\toprule
工具 & 查找顺序（由先到后） & 冲突时的赢家 \\
\midrule
\endhead
\bottomrule
\endlastfoot
\Fmt & 命令行的 \Opt{style} 或 \Opt{config}、被格式化文件所在目录向上第一个 \texttt{.clang-format}、名为 \texttt{\_clang-format} 的文件、\texttt{.clang-format} 加文件扩展名后缀的形式、用户级配置、\Opt{fallback-style} & 显式参数赢；文件级就近赢；段内 \texttt{If} 覆写主段 \\
\Tidy & 命令行的 \Opt{checks} 等参数、当前文件所在目录向上的 \texttt{.clang-tidy}、用户级配置 & 命令行参数赢；文件级就近赢并与上级合并 \\
\CD & 命令行参数、目录层级里的 \texttt{.clangd}（越近越优先）、全局 \texttt{config.yaml} & 命令行赢，其后就近项目文件赢 \\
\end{longtable}
}

两条常见误读值得单独强调。第一，\Opt{style=file} 的「file」指的是\textbf{相对被格式化文件}的查找，而不是相对当前工作目录；在仓库根跑格式化命令时，子目录里遗留的风格文件依然会赢。第二，\CD\ 的 \texttt{.clangd} 与 \Tidy\ 的 \texttt{.clang-tidy} 是两套互不相干的查找链：前者只影响编辑期诊断，后者只影响独立运行的 \Tidy；\CD\ 内嵌 \Tidy\ 时读的是 \texttt{Diagnostics.\allowbreak ClangTidy} 段，而不是项目里的 \texttt{.clang-tidy}——这正是第 17 章里「本地与命令行结论不一致」的高频来源。

\chapter{命令行与诊断速查}\label{ch:appB}

本附录按工具给出常用调用。所有示例假定工具来自同一个 LLVM release（第 16 章），并且 compdb 位于 \texttt{build/dev}。每张表只列真正会改行为的选项，完整选项请用 \Opt{help} 或 \Opt{help-hidden} 查看。

\section{\Fmt\ 的常用调用}\label{sec:appB-s1}

\begin{lstlisting}[style=shellstyle,caption={clang-format 的典型用法}]
# check only, never touch the file
clang-format --dry-run --Werror $(git ls-files '*.cpp')
# rewrite in place
clang-format -i src/widget.cpp src/widget.h
# only some line ranges, useful for incremental gates
clang-format --lines=10:24,60:80 -i src/widget.cpp
# format a fragment as if it were a header
echo 'int f( int a ){}' | clang-format --assume-filename=x.h
# machine-readable output for editors and tools
clang-format --output-replacements-xml src/widget.cpp
# print the resolved style, then start a config file from it
clang-format --dump-config > .clang-format
\end{lstlisting}

{\footnotesize
\begin{longtable}{@{}>{\raggedright\arraybackslash}p{0.34\textwidth}
                  >{\raggedright\arraybackslash}p{0.48\textwidth}@{}}
\caption{clang-format 选项速查}\label{tab:appB-fmt}\\
\toprule
选项 & 含义与使用场景 \\
\midrule
\endfirsthead
\toprule
选项 & 含义与使用场景 \\
\midrule
\endhead
\bottomrule
\endlastfoot
\Opt{style=file} & 按就近查找规则取风格文件；找不到时用 fallback \\
\Opt{fallback-style} & 上一步的兜底族名；\texttt{none} 可防止意外改动 \\
\Opt{config} & 直接在命令行内联一段风格描述，适合一次性验证 \\
\Opt{dry-run} & 只报告会有哪些改动而不写文件 \\
\Opt{Werror} & 把 \Opt{dry-run} 的报告变成非零退出码，门禁常用 \\
\Opt{assume-filename} & 从 stdin 读时假装文件名，用于语言分派 \\
\Opt{output-replacements-xml} & 输出替换列表，供 diff 工具与编辑器消费 \\
\Opt{dump-config} & 打印当前生效的全部字段，可作配置起点 \\
\Opt{sort-includes} & 单独控制是否排序 include（新版可能已移除） \\
\Opt{style} & 也可写 \texttt{Chromium}、\texttt{Mozilla} 等内置族 \\
\end{longtable}
}

\section{\Tidy\ 的常用调用}\label{sec:appB-s2}

\begin{lstlisting}[style=shellstyle,caption={clang-tidy 的典型用法}]
# single file without a compilation database
clang-tidy src/widget.cpp -- -std=c++20
# with a database, and show diagnostics from project headers
clang-tidy src/widget.cpp -p build/dev \
  --header-filter='include/.*'
# verify which checks are really enabled
clang-tidy --list-checks --checks='-*,readability-*'
# export instead of applying, then review
clang-tidy src/widget.cpp -p build/dev -export-fixes fix.yml
# apply a previously exported fix file
clang-tidy --fixes=fix.yml src/widget.cpp -p build/dev
# batch entry points
run-clang-tidy -p build/dev -j 4 -export-fixes all.yml
clang-tidy-diff.py -p1 -path build/dev -quiet < pr.diff
\end{lstlisting}

{\footnotesize
\begin{longtable}{@{}>{\raggedright\arraybackslash}p{0.34\textwidth}
                  >{\raggedright\arraybackslash}p{0.48\textwidth}@{}}
\caption{clang-tidy 选项速查}\label{tab:appB-tidy}\\
\toprule
选项 & 含义与使用场景 \\
\midrule
\endfirsthead
\toprule
选项 & 含义与使用场景 \\
\midrule
\endhead
\bottomrule
\endlastfoot
\Opt{checks} & 逗号分隔的包含与排除序列，后写者生效 \\
\Opt{list-checks} & 打印实际启用的检查名，排障第一步 \\
\texttt{-p} & 指定 compdb 所在目录，等价于 \Opt{compile-commands-dir} \\
\Opt{header-filter} & 决定头的诊断是否显示，是正则 \\
\Opt{line-filter} & 只分析指定行区间，增量门禁的核心 \\
\Opt{export-fixes} & 把可自动修的部分写成 YAML \\
\Opt{fixes} & 从 YAML 读取并实际应用替换 \\
\Opt{fix} & 直接改写源文件，务必在干净工作树上跑 \\
\Opt{warnings-as-errors} & 指定哪些警告升级为 error，影响退出码 \\
\Opt{quiet} & 只输出诊断，去掉进度与统计信息 \\
\Opt{system-headers} & 连系统头一起分析，通常只在排障时开 \\
\Opt{extra-arg} & 向编译命令追加参数，无需改 compdb \\
\end{longtable}
}

\section{\CD\ 的常用调用}\label{sec:appB-s3}

\CD\ 一般由编辑器启动，但下面这些命令行参数在排障时同样有效，尤其是 \Opt{check}。

\begin{lstlisting}[style=shellstyle,caption={clangd 的命令行排障用法}]
# the single most useful debugging entry
clangd --check=src/widget.cpp \
  --compile-commands-dir=build/dev --log=verbose
# only a line range, much faster on large files
clangd --check=src/widget.cpp --check-lines=10-40
# let clangd run the real compiler to collect macros
clangd --query-driver='/opt/tc/*/bin/clang++'
# globs are allowed, but every path must be executable
clangd --query-driver=/usr/bin/clang++,/usr/bin/gcc
# measure the cost without writing an index
clangd --check=src/widget.cpp --no-background-index
\end{lstlisting}

{\footnotesize
\begin{longtable}{@{}>{\raggedright\arraybackslash}p{0.34\textwidth}
                  >{\raggedright\arraybackslash}p{0.48\textwidth}@{}}
\caption{clangd 选项速查}\label{tab:appB-clangd}\\
\toprule
选项 & 含义与使用场景 \\
\midrule
\endfirsthead
\toprule
选项 & 含义与使用场景 \\
\midrule
\endhead
\bottomrule
\endlastfoot
\Opt{check} & 单文件自检模式，打印命令、preamble 与诊断 \\
\Opt{check-lines} & 配合 \Opt{check} 限定行区间 \\
\Opt{compile-commands-dir} & 指定 compdb 目录，命令行优先于配置文件 \\
\Opt{query-driver} & 允许执行编译器以获取宏与系统头，逗号分隔的匹配串 \\
\Opt{background-index} & 建立本地索引，跨文件能力的来源 \\
\Opt{no-background-index} & 关闭索引，用于对比「缺索引」的表现 \\
\Opt{header-insertion} & include 插入策略，\texttt{never} 或 \texttt{iwyu} \\
\Opt{clang-tidy} & 编辑期启用内嵌 \Tidy\ 检查 \\
\Opt{pch-storage} & preamble 存内存或磁盘，影响峰值占用 \\
\Opt{limit-references} & 限制引用结果条数，超大仓保护 \\
\Opt{enable-config} & 是否读取全局与项目配置文件 \\
\Opt{log} & 日志级别，从 \texttt{error} 到 \texttt{verbose} \\
\end{longtable}
}

\section{重构与查询工具}\label{sec:appB-s4}

{\footnotesize
\begin{longtable}{@{}>{\raggedright\arraybackslash}p{0.32\textwidth}
                  >{\raggedright\arraybackslash}p{0.50\textwidth}@{}}
\caption{apply-replacements、rename、query、include-fixer}\label{tab:appB-refac}\\
\toprule
调用 & 用途与注意 \\
\midrule
\endfirsthead
\toprule
调用 & 用途与注意 \\
\midrule
\endhead
\bottomrule
\endlastfoot
\texttt{clang-apply-replacements \Opt{dry-run} ./fixdir} & 汇总一批 YAML 并按目录应用；先 dry-run 看冲突 \\
\texttt{clang-apply-replacements \Opt{style=file} ./fixdir} & 应用后立刻按风格文件排版，少一次往返 \\
\texttt{clang-rename -p build -offset=4212 \Opt{new-name=Widget}} & 按字节偏移重命名，跨 TU 依赖索引完整性 \\
\texttt{clang-rename -line=37 -column=9 \Opt{new-name=Widget} -i} & 按行列定位，加 \texttt{-i} 才真正写文件 \\
\texttt{clang-query build/compile\_\allowbreak commands.json} & 交互式匹配 AST，进入后可用 \texttt{match}、\texttt{p}、\texttt{set} \\
\texttt{match cxxRecordDecl()} & 会话里列出所有类，再用 \texttt{p} 或 \texttt{dump} 看细节 \\
\texttt{clang-query -c 'match \ldots' <dir>} & 非交互模式，一条命令跑完一个匹配 \\
\texttt{clang-include-fixer -p build \Opt{symbol=Widget}} & 为未定义符号插入最合适的 include \\
\texttt{clang-include-fixer -p build \Opt{extra-arg}} & 需要更多搜索路径时用，不改动 compdb \\
\texttt{clang-find-\allowbreak all-symbols} & 扫描 compdb 覆盖的文件生成符号表，供 include-fixer 使用 \\
\end{longtable}
}

\section{cppcheck}\label{sec:appB-s5}

\begin{lstlisting}[style=shellstyle,caption={cppcheck 的两种输入方式与抑制}]
# driven by the compilation database (preferred)
cppcheck --project=build/dev/compile_commands.json \
  --error-exitcode=2 -j 8 --quiet
# plain directory scan with explicit suppression file
cppcheck --enable=warning,style,performance \
  --suppressions-list=.ci/suppr.txt --inline-suppr src/
# machine output for CI annotations
cppcheck --xml --xml-version=3 src/ 2> cppcheck.xml
\end{lstlisting}

{\footnotesize
\begin{longtable}{@{}>{\raggedright\arraybackslash}p{0.34\textwidth}
                  >{\raggedright\arraybackslash}p{0.48\textwidth}@{}}
\caption{cppcheck 选项速查}\label{tab:appB-cppcheck}\\
\toprule
选项 & 含义与使用场景 \\
\midrule
\endfirsthead
\toprule
选项 & 含义与使用场景 \\
\midrule
\endhead
\bottomrule
\endlastfoot
\Opt{project} & 用 compdb 或 cppcheck 自有工程文件驱动分析 \\
\Opt{error-exitcode} & 出错时返回指定码，脚本收集结果的前提 \\
\Opt{enable} & 打开 \texttt{warning}、\texttt{style}、\texttt{performance} 等类别 \\
\Opt{suppress} & 单条抑制，形如 \texttt{id:path} \\
\Opt{suppressions-list} & 抑制清单文件，基线机制的载体 \\
\Opt{inline-suppr} & 允许源码里的内联抑制注释 \\
\Opt{template} & 自定义输出格式，可对齐诊断前缀 \\
\Opt{std} & 指定语言标准，避免误判新语法 \\
\Opt{force} & 强制分析所有配置组合，慢，慎用 \\
\end{longtable}
}

\section{include-what-you-use}\label{sec:appB-s6}

\begin{lstlisting}[style=shellstyle,caption={iwyu 的批量与单文件用法}]
# batch mode through the helper script
iwyu_tool.py -p build/dev -- -Xiwyu --verbose=3
# one file, with a mapping file to keep chosen headers
iwyu -Xiwyu --mapping_file=acme.imp -std=c++20 \
  -Iinclude src/widget.cpp
\end{lstlisting}

它的输出是\textbf{建议}：哪些 include 可以换成前置声明、哪些符号需要自己包含，都需要人工确认间接依赖与导出符号（第 15 章）。映射文件与注释风格开关要提交进仓库，否则每次运行结果不可复现。

\section{\texttt{clang} 直译器侧的分析与 sanitizer}\label{sec:appB-s7}

{\footnotesize
\begin{longtable}{@{}>{\raggedright\arraybackslash}p{0.38\textwidth}
                  >{\raggedright\arraybackslash}p{0.44\textwidth}@{}}
\caption{编译期分析与运行期 sanitizer}\label{tab:appB-san}\\
\toprule
调用 & 说明 \\
\midrule
\endfirsthead
\toprule
调用 & 说明 \\
\midrule
\endhead
\bottomrule
\endlastfoot
\texttt{clang -fsyntax-only -Wall -Wextra a.c} & 最快的一遍解析，只出编译诊断 \\
\texttt{-Xanalyzer} \texttt{-analyzer-\allowbreak checker=core,security} & clang 静态分析器，不产出目标文件 \\
\texttt{clang++ -fsanitize=address,undefined} \texttt{-g} & ASan 与 UBSan 可以同时使用 \\
\texttt{clang++ -fsanitize=thread a.cpp} & TSan，\textbf{不能与 ASan 同开} \\
\texttt{clang++ -fsanitize=leak a.cpp} & LSan，通常由 ASan 顺带启用 \\
\texttt{clang++ -fsanitize=fuzzer-no-link ...} & 覆盖率引导测试的插桩形式 \\
\texttt{ASAN\_OPTIONS=detect\_leaks=1} & 运行期开关，与编译期插桩配套 \\
\texttt{UBSAN\_OPTIONS=\allowbreak print\_stacktrace=1} & 让未定义行为带上栈回溯 \\
\texttt{-fsanitize-recover=address} & 报错后继续运行，适合批量收集 \\
\end{longtable}
}

\section{编辑器侧动作对照}\label{sec:appB-s8}

{\footnotesize
\begin{longtable}{@{}>{\raggedright\arraybackslash}p{0.17\textwidth}
                  >{\raggedright\arraybackslash}p{0.23\textwidth}
                  >{\raggedright\arraybackslash}p{0.21\textwidth}
                  >{\raggedright\arraybackslash}p{0.22\textwidth}@{}}
\caption{同一件事在三种编辑器里的入口}\label{tab:appB-editor}\\
\toprule
动作 & VSCode 与扩展 & Neovim 内置 LSP & Emacs 与 LSP 模式 \\
\midrule
\endfirsthead
\toprule
动作 & VSCode 与扩展 & Neovim 内置 LSP & Emacs 与 LSP 模式 \\
\midrule
\endhead
\bottomrule
\endlastfoot
格式化缓冲区 & 命令 \texttt{Format Document}，开关 \texttt{editor.\allowbreak formatOnSave} & \texttt{vim.lsp.buf.\allowbreak format()} & \texttt{lsp-format-\allowbreak buffer} \\
格式化选定区域 & \texttt{Format Selection} & \texttt{format()} 传 \texttt{range} & \texttt{lsp-format-\allowbreak region} \\
organize includes & clangd 的 Source Action & code action 里选同名项 & \texttt{lsp-execute-\allowbreak code-action} \\
快速修复 & \texttt{Quick Fix}，快捷键 \texttt{Ctrl+.} & \texttt{code\_\allowbreak action()} & \texttt{lsp-execute-\allowbreak code-action} \\
重命名符号 & \texttt{Rename Symbol}，\texttt{F2} & \texttt{vim.lsp.buf.\allowbreak rename()} & \texttt{lsp-rename} \\
跳转到实现 & \texttt{Go to Implementation} & \texttt{implementation()} & \texttt{lsp-find-\allowbreak implementation} \\
引用列表 & \texttt{Find All References} & \texttt{lsp.buf.\allowbreak references()} & \texttt{lsp-find-\allowbreak references} \\
内联提示开关 & \texttt{editor.inlay\-\allowbreak Hints.enabled} & \texttt{inlay\_hint.\allowbreak enable(true)} & 由独立插件提供 \\
语义高亮 & 扩展默认开启 & \texttt{lsp.semantic\_\allowbreak tokens.start()} & \texttt{lsp-semantic-\allowbreak tokens} \\
\CD\ 二进制路径 & \texttt{clangd.path} & \texttt{cmd} 配置项 & \texttt{lsp-clients} 变量 \\
额外启动参数 & \texttt{clangd.\allowbreak arguments} 数组 & \texttt{before\_init} 回调 & \texttt{lsp-clangd} 函数 \\
日志面板 & 输出面板选 clangd & \texttt{LspLog} 窗口 & \texttt{*lsp-log*} 缓冲区 \\
\end{longtable}
}

\begin{note}[title={编辑器只负责「发起」，不负责「正确」}]
上表的入口都可能被第三方扩展接管（尤其是格式化与 organize includes）。判断责任归属的办法只有一个：在命令行跑出正确结果，再回到编辑器里比对。两侧不一致时，问题一定在编辑器的分派链上，而不是配置文件的语义上。
\end{note}

\section{诊断文本结构与 grep 技巧}\label{sec:appB-s9}

诊断的统一结构是\textbf{位置、等级、消息、来源标记}四段，冒号分隔位置：

\begin{itemize}
  \item \texttt{file:line:col: warning: msg [check-name]} 是最常见的一行形式。
  \item \texttt{error} 与 \texttt{warning} 之外，\texttt{note} 行承载补充位置，不是独立问题。
  \item 来源标记可能是 \Tidy\ 的族名（\texttt{readability-\allowbreak magic-numbers}），可能是被包装的编译诊断（\texttt{clang-diagnostic-\allowbreak unused-variable}），也可能是 cppcheck 的短 id（\texttt{uninitvar}）。
\end{itemize}

按来源统计与筛除噪声的几条命令：

\begin{lstlisting}[style=shellstyle,caption={用 grep 与 awk 把诊断变成可统计的数据}]
# count diagnostics per check name
grep -Eo '\[[a-zA-Z_.-]+\]$' tidy.log | sort | uniq -c | sort -rn
# keep only one family, then drop duplicates
grep -F '[readability-' tidy.log | sort -u > readability.log
# top offending files: cut the path and count occurrences
grep -o '^[^:]*' tidy.log | sort | uniq -c | sort -rn | head -5
# collapse diagnostics to file:line for a suppression list
awk -F: '/warning:|error:/ {print $1" "$2}' tidy.log | sort -u
# how many diagnostics, and against what size baseline
grep -c 'warning:' tidy.log ;  wc -l src/*.cpp | tail -1
\end{lstlisting}

\Opt{export-fixes} 生成的 YAML 是另一个高价值数据源，结构与字段如下：

\begin{itemize}
  \item 顶层 \texttt{MainSourceFile}：这批替换属于哪个翻译单元。
  \item 顶层 \texttt{Diagnostics}：诊断数组，每项含 \texttt{ErrorMessage}（消息文本）、\texttt{Diagnostic\allowbreak Name}（check 名）、\texttt{Filesystem\allowbreak RelativePath}（相对路径）、\texttt{LineNumber} 与 \texttt{ColumnNumber}。
  \item 每项内的 \texttt{Fixes} 列表：逐条替换，字段为 \texttt{Offset}、\texttt{Length}、\texttt{ReplacementText} 与说明用的 \texttt{Message}。
  \item 顶层 \texttt{SharedFixNotes}：把出现多次的同一组替换抽出去重，应用时再展开。
\end{itemize}

把这份 YAML 按 \texttt{DiagnosticName} 聚合计数，就得到第 16 章基线文件的原始形态；按 \texttt{FilesystemRelativePath} 聚合，则得到「哪个目录最脏」的清单。

\chapter{术语表}\label{ch:appC}

术语按主题排列而非字母序：先协议与位置计算，再解析与索引，然后格式化与检查，最后是流程类词汇。英文列给出书写形式，方便检索官方文档。

{\footnotesize
\begin{longtable}{@{}>{\raggedright\arraybackslash}p{0.19\textwidth}
                  >{\raggedright\arraybackslash}p{0.25\textwidth}
                  >{\raggedright\arraybackslash}p{0.37\textwidth}@{}}
\caption{中文术语、原文与一句话定义}\label{tab:appC-gloss}\\
\toprule
术语 & 原文 & 一句话定义 \\
\midrule
\endfirsthead
\toprule
术语 & 原文 & 一句话定义 \\
\midrule
\endhead
\bottomrule
\endlastfoot
语言服务器协议 & LSP, Language Server Protocol & 编辑器与语言分析器之间的 JSON-RPC 协议 \\
帧 & LSP frame, Content-Length & 报文由头字段与体组成，靠 \texttt{Content-Length} 定界 \\
能力协商 & capability negotiation & 初始化时双方声明各自支持的功能集合 \\
请求与响应 & request, response & 有 id、需要回答的调用；失败走 error 字段 \\
通知 & notification & 无 id、不期望回答的单向消息 \\
取消 & cancellation & 用已完成的请求 id 通知服务端停止后续工作 \\
工作区文件夹 & workspace folder & 服务端共享根目录与索引范围的项目集合 \\
偏移 & offset, byte vs UTF-16 & 文档内位置；\textbf{不同基准之间不可直接换算} \\
码元 & UTF-16 code unit & LSP 侧计列宽时使用的单位 \\
翻译单元 & translation unit & 一个源文件递归展开所有 include 后的整体 \\
编译数据库 & compilation database & 记录每个文件真实编译命令的 JSON 列表 \\
回退命令 & fallback command & 找不到数据库条目时临时拼出的编译命令 \\
前导块 & preamble & 头文件部分的可复用预编译结果 \\
前导块复用 & preamble reuse & 改动只在文件尾部时跳过重复解析 \\
后台索引 & background index & 在服务端进程内渐进构建的跨文件符号索引 \\
外部索引 & external index & 由构建产物或符号服务器提供的只读索引 \\
符号索引 & symbol index & 供全局搜索与跨 TU 跳转使用的名字表 \\
语义 token & semantic tokens & 服务端返回的逐 token 分类，用于精确着色 \\
内联提示 & inlay hints & 编辑器文本之间插入的参数名与推导类型 \\
代码动作 & code action & 光标处可执行的修复或重构候选 \\
代码透镜 & code lens & 行上方显示的可点击元信息 \\
fix-it 提示 & fix-it hint & 编译器诊断里「把 A 替换为 B」的部分 \\
引用 & references & 某符号在整个项目中的使用位置 \\
调用层次 & call hierarchy & 谁调用它、它调用谁的双向视图 \\
AST 匹配器 & AST matcher & 以组合子表达式描述语法树查询目标 \\
数据流分析 & dataflow analysis & 沿控制流传播抽象状态以判断变量取值 \\
路径敏感 & path-sensitive & 区分不同分支条件组合的分析方式 \\
误报 & false positive & 报告了并不存在的问题 \\
漏报 & false negative & 没能报告真实存在的问题 \\
抑制 & suppression, NOLINT & 用注释或清单让某条诊断不再出现 \\
头过滤正则 & HeaderFilterRegex & 决定哪些头的诊断会被显示出来 \\
检查族 & check family & 以 \texttt{bugprone-}、\texttt{readability-} 等前缀分组的检查集合 \\
幂等格式化 & idempotent formatting & 再跑一次不产生任何改动 \\
增量格式化 & incremental formatting & 只处理改动行或改动文件 \\
替换列表 & replacements & 工具输出的「偏移加长度加新文本」的修改描述 \\
自动修复 & autofix & 检查自带的、可直接改写源码的部分 \\
调试适配器协议 & DAP & 调试器与编辑器之间的独立协议 \\
问题匹配器 & problem matcher & 把构建输出正则化后填进诊断列表的机制 \\
清洁构建缓存 & ccache & 复用编译结果，但会跳过真实编译命令 \\
静态分析门禁 & quality gate & 在合入前强制执行的一组检查与阻断规则 \\
基线告警 & baseline, existing debt & 已提交的问题清单，只允许缩小 \\
存量与新增 & pre-existing, introduced & 门禁区分两类告警的处理方式 \\
清理抖动 & formatting churn & 风格变更导致的大量无关改动 \\
消毒器 & sanitizer & 编译期插桩加运行期检查的一类动态检测器 \\
影子内存 & shadow memory & sanitizer 用来记录原始内存元数据的平行空间 \\
数据竞争 & data race & 两个线程无同步地访问同一位置且至少一方写 \\
未定义行为 & UB, undefined behavior & 标准不规定结果的行为，运行期可随时暴露 \\
模糊测试插桩 & fuzzing instrumentation & 为覆盖率引导而插入的计数器与回调 \\
容器镜像对齐 & image pinning & 让工具版本固定在同一个发行版镜像里 \\
浅克隆 & shallow clone & 缺少完整历史，会使增量分析退化成全量 \\
\end{longtable}
}

\begin{guideline}[title={查术语的三个习惯}]
同一概念在文档、代码与聊天里常有三种写法：\textbf{先写出英文原名}，再补中文；涉及位置计算时把「字节」「码元」「字符」三者区分清楚；讨论检查项时始终带族名前缀，避免「那个空指针检查」这类无法复现的指代。
\end{guideline}
\end{document}
